\begin{longtable}{p{0.03\textwidth}|p{0.22\textwidth}|p{0.22\textwidth}|p{0.22\textwidth}|p{0.22\textwidth}}
\hline
\textbf{id} & \textbf{Hugtök} & \textbf{Skilgreining} & \textbf{Skýring} & \textbf{Athugasemdir} \\
\hline
0 & MISSING & ... &  &  \\
\hline
1 & -und, -d &  &  &  \\
\hline
2 & \(\phi\)-fall Eulers &  &  &  \\
\hline
3 & \(A\)-stöðugleiki &  &  &  \\
\hline
4 & \(B\)-splæsifall &  &  &  \\
\hline
5 & \(K\)-fræði &  &  &  \\
\hline
6 & \(n\)-und, \(n\)-d, röðuð \(n\)-und, röðuð \hbox{\(n\)-d} & $n$-\textit{und}, þar sem $n$ er <ref idTerm=3475>náttúruleg tala</ref>, samanstendur af $n$ <ref idTerm=2123>hlutum</ref> í tiltekinni röð. $n$-undin sem inniheldur hlutina $a_1, a_2, \ldots, a_n$ í þessari röð er táknuð með $(a_1, a_2, \ldots, a_n)$. &  &  \\
\hline
7 & \(p\)-leg heiltala, heil \(p\)-leg tala &  &  &  \\
\hline
8 & \(p\)-leg tala &  &  &  \\
\hline
9 & \(p\)-leg virðing &  &  &  \\
\hline
10 & \(p\)-legur &  &  &  \\
\hline
11 & \(t\)-dreifing &  &  &  \\
\hline
12 & \(t\)-próf &  &  &  \\
\hline
13 & á köflum &  &  &  \\
\hline
14 & Abel-, abelskur &  &  &  \\
\hline
15 & Abel-deildi af fyrstu gerð, fágað deildi, Abel-diffur af fyrstu gerð, fágað diffur &  &  &  \\
\hline
16 & Abel-deildi, Abel-diffur &  &  &  \\
\hline
17 & Abel-grúpa, víxlgrúpa, abelsk grúpa, víxlin grúpa &  &  &  \\
\hline
18 & Abel-útvíkkun, Abel-víkkun, abelsk útvíkkun &  &  &  \\
\hline
19 & Abel-víðátta, Abel-víðerni &  &  &  \\
\hline
20 & aðalhornalína, meginhornalína &  &  &  \\
\hline
21 & aðdráttarmengi &  &  &  \\
\hline
22 & aðdráttarpunktur &  &  &  \\
\hline
23 & aðfaldaður, hrópmerktur &  &  &  \\
\hline
24 & aðfeldaröð &  &  &  \\
\hline
25 & aðfella &  &  &  \\
\hline
26 & aðfelld hálflína &  &  &  \\
\hline
27 & aðfelld samsíða lína &  &  &  \\
\hline
28 & aðfelldur þríhyrningur &  &  &  \\
\hline
29 & aðfelldur, aðfellu- &  &  &  \\
\hline
30 & aðfellu- &  &  &  \\
\hline
31 & aðfelluhringur &  &  &  \\
\hline
32 & aðfellujafn &  &  &  \\
\hline
33 & aðfelluliðun, aðfelld liðun &  &  &  \\
\hline
34 & aðfellunormlegur &  &  &  \\
\hline
35 & aðfelluréttur &  &  &  \\
\hline
36 & aðfelluröð &  &  &  \\
\hline
37 & aðfelluskilvirkni &  &  &  \\
\hline
38 & aðfelluskilvirkur &  &  &  \\
\hline
39 & aðfellustefna, aðfelluátt &  &  &  \\
\hline
40 & aðfellustöðugur &  &  &  \\
\hline
41 & aðferð &  &  &  \\
\hline
42 & aðferð breytilegra stuðla &  &  &  \\
\hline
43 & aðferð minnstu ferninga, aðferð minnstu fervika &  &  &  \\
\hline
44 & aðgengilegur jaðarpunktur &  &  &  \\
\hline
45 & aðgerð &  &  &  \\
\hline
46 & aðgerð, reikniaðgerð & <ref idTerm=3282>Vörpun</ref> af gerðinni $X \times \cdots \times X \to X$ þar sem $X$ er <ref idTerm=3270>mengi</ref> kallast \textit{reikniaðgerð} á $X$ &  & Skilgreiningin í Hugtakasafninu er ekki nógu góð, hún krefst þess ekki að $X_k=X$ og eitt dæmið sem er gefið er meira að segja $(x,n) \mapsto x^n : \mathbb{R} \times \mathbb{N} \to \mathbb{R}$ sem er ekki reikniaðgerð á $\mathbb{R}$ \\
\hline
47 & aðgerðagreining &  &  &  \\
\hline
48 & aðgerðartákn &  &  &  \\
\hline
49 & aðgerðartákn &  &  &  \\
\hline
50 & aðgreina &  &  &  \\
\hline
52 & aðgreiningareiginleiki Hausdorffs, frumsenda Hausdorffs, frumsetning Hausdorffs &  &  &  \\
\hline
53 & aðhæf heildun, aðhæf tegrun &  &  &  \\
\hline
54 & aðhæft ítrekunarreiknirit &  &  &  \\
\hline
55 & aðhæfur &  &  &  \\
\hline
56 & aðhlutur, bakhlutur &  &  &  \\
\hline
57 & aðhvarf &  &  &  \\
\hline
58 & aðhvarfsfall &  &  &  \\
\hline
59 & aðhvarfsferill &  &  &  \\
\hline
60 & aðhvarfsfrávik, leif &  &  &  \\
\hline
61 & aðhvarfsgreining &  &  &  \\
\hline
62 & aðhvarfsjafna &  &  &  \\
\hline
63 & aðhvarfslína &  &  &  \\
\hline
64 & aðhvarfsstuðull &  &  &  \\
\hline
65 & aðlæg hlið & Tvær <ref idTerm=2048>hliðar</ref> <ref idTerm=3161>marghyrnings</ref> sem skerast í nákvæmlega einum <ref idTerm=4001>punkti</ref> (í tveim víddum) eða tvær <ref idTerm=1325>hliðar</ref> <ref idTerm=3133>margflötungs</ref> sem skerast í einum eða fleiri punktum sem liggja allir á sömu <ref idTerm=2911>línu</ref> (í þrem víddum) kallast \textit{aðlægar hliðar} &  & \\
\hline
66 & aðlæg horn, samlæg horn, samliggjandi horn &  &  &  \\
\hline
67 & aðlægur &  &  &  \\
\hline
68 & aðlægur punktur &  &  &  \\
\hline
69 & aðlægur, samlægur, samliggjandi, grannstæður &  &  &  \\
\hline
70 & aðlaga &  &  &  \\
\hline
71 & aðlögun &  &  &  \\
\hline
72 & aðlögunarpróf, mátunarpróf &  &  &  \\
\hline
74 & aðmengi, ítak, bakmengi & Sjá <ref idTerm=3282>vörpun</ref> &  &  \\
\hline
75 & aðoka &  &  &  \\
\hline
76 & aðoka, aðoka frá vinstri &  &  &  \\
\hline
77 & aðskilja, skilja að, greina á milli, aðgreina &  &  &  \\
\hline
78 & aðskiljanleg jafna &  &  &  \\
\hline
79 & aðskiljanleg lokun, aðskiljanlegur hjúpur &  &  &  \\
\hline
80 & aðskiljanleg margliða &  &  &  \\
\hline
81 & aðskiljanleg útvíkkun, aðskiljanleg víkkun &  &  &  \\
\hline
82 & aðskiljanlega lokaður &  &  &  \\
\hline
83 & aðskiljanlegt net &  &  &  \\
\hline
84 & aðskiljanlegur &  &  &  \\
\hline
85 & aðskiljanleikastig &  &  &  \\
\hline
86 & aðskiljanleiki &  &  &  \\
\hline
87 & aðskilnaðarfrumsenda &  &  &  \\
\hline
88 & aðskilnaðargreining &  &  &  \\
\hline
89 & aðskilnaðarregla, aðskilnaðarháttur &  &  &  \\
\hline
90 & aðskilnaðarstærð &  &  &  \\
\hline
91 & aðskilnaðarstyrkur &  &  &  \\
\hline
92 & aðskilnaður &  &  &  \\
\hline
93 & aðskilnaður breytistærða &  &  &  \\
\hline
94 & aðstoðar-, hjálpar-, viðbótar- &  &  &  \\
\hline
95 & aðstreymi &  &  &  \\
\hline
96 & æðri stærðfræði &  &  &  \\
\hline
97 & æstæður, sístæður, í jafnvægi, jafnvægis- &  &  &  \\
\hline
98 & ætt &  &  &  \\
\hline
99 & af handahófi &  &  &  \\
\hline
100 & afgangsfervik &  &  &  \\
\hline
101 & afgangsliður, leif &  &  &  \\
\hline
102 & afgangsróf &  &  &  \\
\hline
103 & afgangssetning &  &  &  \\
\hline
104 & afgangsskekkja &  &  &  \\
\hline
105 & afgangur, minnsta leif &  &  &  \\
\hline
106 & afkóta, dulráða, aflykla &  &  &  \\
\hline
107 & afkótun, dulráðning, aflyklun &  &  &  \\
\hline
108 & afleiða, afleitt fall &  &  &  \\
\hline
109 & afleiðanlegur &  &  &  \\
\hline
110 & afleiðanleiki &  &  &  \\
\hline
111 & afleiddur &  &  &  \\
\hline
112 & afleiddur varpi &  &  &  \\
\hline
113 & afleiðing &  &  &  \\
\hline
114 & afleiðingur, afleiðuvirki &  &  &  \\
\hline
115 & afleiðsla &  &  &  \\
\hline
116 & afleiðsla, ályktun &  &  &  \\
\hline
117 & afleiðslu- &  &  &  \\
\hline
118 & afleiðsluaðferð &  &  &  \\
\hline
119 & afleiðslusönnun &  &  &  \\
\hline
120 & afleitt mengi &  &  &  \\
\hline
121 & afleitt ríki &  &  &  \\
\hline
122 & aflfræði &  &  &  \\
\hline
123 & aflfræðigreining &  &  &  \\
\hline
124 & afmáanleg ósamfelldni &  &  &  \\
\hline
125 & afmáanlegur sérstöðupunktur &  &  &  \\
\hline
126 & afsanna, hrekja &  &  &  \\
\hline
127 & afsannanlegur, hrekjanlegur &  &  &  \\
\hline
128 & afsannanleiki &  &  &  \\
\hline
129 & afsannanleiki, hrekjanleiki &  &  &  \\
\hline
130 & afsönnun, hrakning &  &  &  \\
\hline
131 & afstæð samkvæmni &  &  &  \\
\hline
132 & afstæð svipgrúpa &  &  &  \\
\hline
133 & afstæði &  &  &  \\
\hline
134 & aftasti tölustafur &  &  &  \\
\hline
135 & aftasti tölustafur &  &  &  \\
\hline
136 & afturfærsla &  &  &  \\
\hline
137 & afturinnsetning, endurinnsetning &  &  &  \\
\hline
138 & afturmunur &  &  &  \\
\hline
139 & afturvirk skekkjugreining &  &  &  \\
\hline
140 & afvinding &  &  &  \\
\hline
141 & Agnesarhaddur &  &  &  \\
\hline
142 & áhætta &  &  &  \\
\hline
143 & áhætta neytandans &  &  &  \\
\hline
144 & áhættufall &  &  &  \\
\hline
145 & áhættufræði &  &  &  \\
\hline
146 & áhættuleikur &  &  &  \\
\hline
147 & ákvarðaður &  &  &  \\
\hline
148 & ákvarðanafræði &  &  &  \\
\hline
149 & ákvarðanlegur, úrskurðanlegur &  &  &  \\
\hline
150 & ákveða &  &  &  \\
\hline
151 & ákveðið heildi, ákveðið tegur &  &  &  \\
\hline
152 & ákveðnipunktur &  &  &  \\
\hline
153 & ákveðu- &  &  &  \\
\hline
154 & ákveðujafna &  &  &  \\
\hline
155 & ákvörðun &  &  &  \\
\hline
156 & ákvörðunaraðferð &  &  &  \\
\hline
157 & ákvörðunarfall &  &  &  \\
\hline
158 & ákvörðunarferli &  &  &  \\
\hline
159 & ákvörðunarverkefni &  &  &  \\
\hline
160 & álægur &  &  &  \\
\hline
161 & alef &  &  &  \\
\hline
162 & Aleksandrov-þjöppun, einspunktsþjöppun &  &  &  \\
\hline
163 & alflokkur, alheimsflokkur &  &  &  \\
\hline
164 & algebra með einingarstaki &  &  &  \\
\hline
165 & algebra, algefra &  &  &  \\
\hline
166 & algebruferill, algebrulegur ferill &  &  &  \\
\hline
167 & algebruflötur, algebrulegur flötur &  &  &  \\
\hline
168 & algebrugrannfræði, algebruleg grannfræði &  &  &  \\
\hline
169 & algebrujafna, algebruleg jafna, margliðujafna &  &  &  \\
\hline
170 & algebruleg heiltala, heil algebruleg tala &  &  &  \\
\hline
171 & algebruleg lokun &  &  &  \\
\hline
172 & algebruleg margfeldni &  &  &  \\
\hline
173 & algebruleg stæða, stæða &  &  &  \\
\hline
174 & algebruleg tala, algebrutala &  &  &  \\
\hline
175 & algebruleg útvíkkun, algebruleg víkkun &  &  &  \\
\hline
176 & algebruleg víðátta &  &  &  \\
\hline
177 & algebrulega háður &  &  &  \\
\hline
178 & algebrulega lokaður &  &  &  \\
\hline
179 & algebrulega óháður &  &  &  \\
\hline
180 & algebrulega samoka stak, samoka stak &  &  &  \\
\hline
181 & algebrulegt fall &  &  &  \\
\hline
182 & algebrulegt mengi &  &  &  \\
\hline
183 & algebrulegt rúm &  &  &  \\
\hline
184 & algebrulegt talnasvið, talnasvið, talnakroppur &  &  &  \\
\hline
185 & algebrulegt varpvíðerni &  &  &  \\
\hline
186 & algebrulegt víðerni &  &  &  \\
\hline
187 & algebrulegur sléttuferill &  &  &  \\
\hline
188 & algebrulegur, algebru- &  &  &  \\
\hline
189 & algebrumótun &  &  &  \\
\hline
190 & algebrumynstur, algebrulegt mynstur &  &  &  \\
\hline
191 & algebrurúmfræði, algebruleg rúmfræði &  &  &  \\
\hline
192 & algebrutalnafræði, algebruleg talnafræði &  &  &  \\
\hline
193 & algengasta gildi, líklegasta gildi &  &  &  \\
\hline
194 & algildi, lengd, tölugildi &  &  &  \\
\hline
195 & algildisflötur &  &  &  \\
\hline
196 & algildismerki, tölugildismerki &  &  &  \\
\hline
197 & algildisminnsta leif &  &  &  \\
\hline
198 & algildisvægi, alvægi &  &  &  \\
\hline
199 & algjörlega ósamanhangandi, algjörlega ósamhangandi &  &  &  \\
\hline
200 & algjörlega ósamfelldur &  &  &  \\
\hline
201 & algjörlega takmarkaður, algjörlega skorðaður &  &  &  \\
\hline
202 & alhæfa &  &  &  \\
\hline
209 & alheildanlegur, altegranlegur &  &  &  \\
\hline
210 & alheimur &  &  &  \\
\hline
211 & aljafna &  &  &  \\
\hline
212 & alls staðar samleitinn, alstaðar samleitinn &  &  &  \\
\hline
213 & alls staðar þéttur, alstaðar þéttur &  &  &  \\
\hline
214 & allsherjar- &  &  &  \\
\hline
215 & allsherjaralgebra &  &  &  \\
\hline
216 & allsherjareiginleiki &  &  &  \\
\hline
217 & allsherjarflokkur &  &  &  \\
\hline
218 & allsherjarhlutur &  &  &  \\
\hline
219 & allsherjarlokun &  &  &  \\
\hline
220 & allsherjarmagnari, allsherjarvirki &  &  &  \\
\hline
221 & allsherjarsvið &  &  &  \\
\hline
222 & allsherjartákn &  &  &  \\
\hline
223 & allsherjarþekjuflötur &  &  &  \\
\hline
224 & allsherjarþekjurúm &  &  &  \\
\hline
225 & allt eða ekkert &  &  &  \\
\hline
226 & almengi &  &  &  \\
\hline
227 & almenn algebrujafna, almenn margliðujafna &  &  &  \\
\hline
228 & almenn grannfræði, mengjagrannfræði &  &  &  \\
\hline
229 & almenn línuleg grúpa &  &  &  \\
\hline
230 & almenn línuleg varpgrúpa &  &  &  \\
\hline
231 & almenn samfellutilgáta, alhæfð samfellutilgáta &  &  &  \\
\hline
232 & almenn staða &  &  &  \\
\hline
233 & almenna afstæðiskenningin &  &  &  \\
\hline
234 & almennt brot & <ref idTerm=2072>Hlutfall</ref> tveggja <ref idTerm=1896>heiltalna</ref> $a$ og $b$, táknað með $\frac{a}{b}$. Talan $a$ kallast \textit{teljari} og talan $b$ \textit{nefnari}  &  &  \\
\hline
235 & almennt tilfelli, almennt tilvik &  &  &  \\
\hline
236 & almennt, að jafnaði &  &  &  \\
\hline
237 & almennur &  &  &  \\
\hline
239 & almennur ferill &  &  &  \\
\hline
240 & almennur liður &  &  &  \\
\hline
241 & almennur punktur &  &  &  \\
\hline
242 & alsamfelldni &  &  &  \\
\hline
243 & alsamfelldur &  &  &  \\
\hline
244 & alsamleggjanlegur &  &  &  \\
\hline
245 & alsamleitinn &  &  &  \\
\hline
246 & alsamleitinn í jöfnum mæli &  &  &  \\
\hline
247 & alsamleitni &  &  &  \\
\hline
248 & alsamleitni í jöfnum mæli &  &  &  \\
\hline
249 & alsamleitnihálfslétta &  &  &  \\
\hline
250 & alsamleitnimark, alsamleitniláhnit &  &  &  \\
\hline
251 & altæk fullyrðing, alhæf fullyrðing &  &  &  \\
\hline
252 & ályktunartölfræði &  &  &  \\
\hline
253 & ámóta &  &  &  \\
\hline
254 & ámóta &  &  &  \\
\hline
255 & ámóta fylki &  &  &  \\
\hline
256 & ámótun &  &  &  \\
\hline
257 & ámyndun &  &  &  \\
\hline
258 & án formerkis &  &  &  \\
\hline
259 & án skerðingar á víðgildi &  &  &  \\
\hline
260 & andbrigður þinur, mótbrigður þinur &  &  &  \\
\hline
261 & andbrigður, mótbrigður &  &  &  \\
\hline
262 & andeinsmóta &  &  &  \\
\hline
263 & andeinsmótun &  &  &  \\
\hline
264 & andfætispunktar, gagnstæðir punktar &  &  &  \\
\hline
265 & andfætisvörpun &  &  &  \\
\hline
266 & andfágaður &  &  &  \\
\hline
267 & andhermískt form &  &  &  \\
\hline
268 & andhermískur &  &  &  \\
\hline
269 & andhverf vensl &  &  &  \\
\hline
270 & andhverfa &  &  &  \\
\hline
271 & andhverfa (falls), andhverft fall &  &  &  \\
\hline
272 & andhverfa (fylkis), umhverfa (fylkis), andhverft fylki, umhverft fylki &  &  &  \\
\hline
273 & andhverfa (vörpunar), andhverf vörpun & \textit{Andhverfa} <ref idTerm=3282>vörpunnar</ref> $f: X \to Y$, ef hún er til, er vörpun $g: Y \to X$ þannig að <ref idTerm=4518>samskeyting</ref> varpananna $f$ og $g$ annars vegar og $g$ og $f$ hins vegar sé viðeigandi <ref idTerm=4498>samsemdarvörpun</ref>, þ.e. \[ g \circ f = \mathrm{id}_X \quad \text{og} \quad f \circ g = \mathrm{id}_Y. \] Andhverfa $f$ er þá yfirleitt táknuð með $f^{-1}$ &  &  \\
\hline
274 & andhverfanlegur & <ref idTerm=3282>Vörpun</ref> kallast \textit{andhverfanleg} ef hún á sér <ref idTerm=273>andhverfu</ref> &  & \\
\hline
275 & andhverfður, umhverfður &  &  &  \\
\hline
276 & andhverfing &  &  &  \\
\hline
277 & andhverft breiðbogafall &  &  &  \\
\hline
278 & andhverfuformúla, andhverfuformkorn &  &  &  \\
\hline
279 & andhverfur &  &  &  \\
\hline
280 & andhverfur breiðbogakósínus &  &  &  \\
\hline
281 & andhverfur breiðbogakótangens &  &  &  \\
\hline
282 & andhverfur breiðbogasínus &  &  &  \\
\hline
283 & andhverfur breiðbogatangens &  &  &  \\
\hline
284 & andhverfuregla Fouriers &  &  &  \\
\hline
285 & andhverfusetning &  &  &  \\
\hline
286 & andinnanmótun &  &  &  \\
\hline
287 & andlínulegur &  &  &  \\
\hline
288 & andlogri &  &  &  \\
\hline
289 & andrárástand &  &  &  \\
\hline
290 & andsælis, rangsælis &  &  &  \\
\hline
293 & andsamhverft fall, oddhverft fall &  &  &  \\
\hline
294 & andsamhverft fylki &  &  &  \\
\hline
295 & andsamhverft tvílínulegt form, oddhverft tvílínulegt form &  &  &  \\
\hline
296 & andsamhverfur &  &  &  \\
\hline
297 & andsamhverfur, oddhverfur &  &  &  \\
\hline
298 & andsammótun, andmótun &  &  &  \\
\hline
299 & andsamsíða &  &  &  \\
\hline
300 & andsamsíðungur, ferhyrningur með engar hliðar samsíða &  &  &  \\
\hline
301 & andsjálfhverfur &  &  &  \\
\hline
302 & andsjálfmótun &  &  &  \\
\hline
303 & andvíxill &  &  &  \\
\hline
304 & andvíxlanlegur &  &  &  \\
\hline
305 & andvíxlast &  &  &  \\
\hline
306 & andvíxlinn &  &  &  \\
\hline
307 & andvíxlunarjafna &  &  &  \\
\hline
308 & annað meginform, annað undirstöðuform &  &  &  \\
\hline
309 & annarrar raðar &  &  &  \\
\hline
310 & annarrar stéttar &  &  &  \\
\hline
311 & annars stigs &  &  &  \\
\hline
312 & annars stigs algebruflötur &  &  &  \\
\hline
313 & annars stigs ferill &  &  &  \\
\hline
314 & annars stigs ferill, ferningsferill &  &  &  \\
\hline
315 & annars stigs háflötur, ferningsháflötur &  &  &  \\
\hline
316 & annars stigs margliða, ferningsmargliða & <ref idTerm=1948>Margliða</ref> með stig 2 &  & Minnast á lausnarformúlu í skýringu \\
\hline
317 & annars stigs nálgun &  &  &  \\
\hline
318 & annars stigs sviðsútvíkkun, ferningssviðsútvíkkun, annars stigs sviðsvíkkun, ferningssviðsvíkkun &  &  &  \\
\hline
319 & annars stigs, fernings- &  &  &  \\
\hline
320 & áreiðanleiki &  &  &  \\
\hline
321 & áreiðanleiki &  &  &  \\
\hline
322 & arfgengur eiginleiki &  &  &  \\
\hline
323 & arftekið grannmynstur &  &  &  \\
\hline
324 & arftekið hlutnet, einskorðað hlutnet &  &  &  \\
\hline
325 & arftekinn, einskorðaður &  &  &  \\
\hline
326 & arkarfall, bogafall &  &  &  \\
\hline
327 & arkarkósínus &  &  &  \\
\hline
328 & arkarkótangens &  &  &  \\
\hline
329 & arkarsínus &  &  &  \\
\hline
330 & arkartangens &  &  &  \\
\hline
331 & arkimedískt algildi, arkimedísk virðing &  &  &  \\
\hline
332 & arkimedískt raðsvið, arkimedískt raðað svið &  &  &  \\
\hline
333 & armur &  &  &  \\
\hline
334 & armur & Sjá <ref idTerm=2198>horn</ref> &  &  \\
\hline
335 & ás &  &  &  \\
\hline
336 & ás, möndull &  &  &  \\
\hline
337 & asnabrú &  &  &  \\
\hline
338 & ássamhvarf, ássamhverfni &  &  &  \\
\hline
339 & ássamhverfing, ássamhverfa &  &  &  \\
\hline
340 & ássamhverfur &  &  &  \\
\hline
341 & ássamlínumótun &  &  &  \\
\hline
342 & ássnið &  &  &  \\
\hline
343 & ástand &  &  &  \\
\hline
344 & ástand &  &  &  \\
\hline
345 & ástand, stand &  &  &  \\
\hline
346 & ástandsrúm &  &  &  \\
\hline
347 & ástandstala &  &  &  \\
\hline
348 & ástandsvigur &  &  &  \\
\hline
349 & ásvigur, snúningsvigur &  &  &  \\
\hline
350 & átæk mótun, átæk sammótun &  &  &  \\
\hline
351 & átæk vörpun & <ref idTerm=3282>Vörpun</ref> sem hefur allt <ref idTerm=74>bakmengið</ref> sem <ref idTerm=3417>myndmengi</ref>, þ.e. vörpun $f:X \to Y$ sem uppfyllir $f(X)=Y$ &  &  \\
\hline
352 & átækur &  &  &  \\
\hline
353 & atburðaalgebra &  &  &  \\
\hline
354 & atburður (í líkindafræði), tilviljanakenndur atburður &  &  &  \\
\hline
355 & atburður, atvik, tilvik &  &  &  \\
\hline
356 & átekt, átaka &  &  &  \\
\hline
357 & athugun &  &  &  \\
\hline
358 & atlas, kortasafn &  &  &  \\
\hline
359 & atóm &  &  &  \\
\hline
360 & atómlaust mál &  &  &  \\
\hline
361 & áttað horn, stefnt horn, snúningshorn &  &  &  \\
\hline
362 & áttað rúm &  &  &  \\
\hline
363 & áttaður &  &  &  \\
\hline
364 & áttaður hringur &  &  &  \\
\hline
365 & áttaður, stefndur, stefnubundinn &  &  &  \\
\hline
366 & áttanlegur &  &  &  \\
\hline
367 & áttanleiki &  &  &  \\
\hline
368 & áttflötungur &  &  &  \\
\hline
369 & átthyrningstala, átthyrnd tala &  &  &  \\
\hline
370 & átthyrningur &  &  &  \\
\hline
371 & áttrækin einsfirðun &  &  &  \\
\hline
372 & áttrækinn &  &  &  \\
\hline
373 & áttskiptin einsfirðun &  &  &  \\
\hline
374 & áttskiptinn &  &  &  \\
\hline
375 & áttuð slétta &  &  &  \\
\hline
376 & áttuð víðátta &  &  &  \\
\hline
377 & áttun &  &  &  \\
\hline
378 & áttunda- &  &  &  \\
\hline
379 & áttundakerfi &  &  &  \\
\hline
380 & áttundareikningur &  &  &  \\
\hline
381 & áttungur &  &  &  \\
\hline
382 & auðkenning, greinimark &  &  &  \\
\hline
383 & auðugt vigurbundin &  &  &  \\
\hline
384 & auðugur &  &  &  \\
\hline
385 & augljós &  &  &  \\
\hline
386 & augljós lausn &  &  &  \\
\hline
387 & augljósleiki &  &  &  \\
\hline
388 & augljóst bundin &  &  &  \\
\hline
389 & augljóst grannmynstur, fábrotið grannmynstur &  &  &  \\
\hline
390 & aukahending &  &  &  \\
\hline
391 & aukahornalína &  &  &  \\
\hline
392 & aukahringur, aukahvirfilhringur &  &  &  \\
\hline
393 & aukahvirfilpunktur &  &  &  \\
\hline
394 & aukaskilyrði, hliðarskilyrði &  &  &  \\
\hline
395 & aukið fylki &  &  &  \\
\hline
396 & aukinn &  &  &  \\
\hline
397 & autt sæti &  &  &  \\
\hline
398 & ávöntun &  &  &  \\
\hline
399 & bæti, tölvustafur &  &  &  \\
\hline
400 & bagall &  &  &  \\
\hline
401 & Baire-fall &  &  &  \\
\hline
402 & Baire-setning, setning Baires &  &  &  \\
\hline
403 & bakliður &  &  &  \\
\hline
404 & bakpokaverkefni &  &  &  \\
\hline
405 & Banach-algebra &  &  &  \\
\hline
406 & Banach-rúm &  &  &  \\
\hline
407 & bandalag &  &  &  \\
\hline
408 & bandbreidd &  &  &  \\
\hline
409 & bandfylki &  &  &  \\
\hline
410 & bann &  &  &  \\
\hline
411 & bann- &  &  &  \\
\hline
412 & bannjafnvægi &  &  &  \\
\hline
413 & bannlíkindi &  &  &  \\
\hline
414 & Bartlett-próf &  &  &  \\
\hline
415 & baugað rúm &  &  &  \\
\hline
416 & baugamótun &  &  &  \\
\hline
417 & baugmynstur &  &  &  \\
\hline
418 & baugsútvíkkun, baugsvíkkun &  &  &  \\
\hline
419 & baugur &  &  &  \\
\hline
420 & Bayes-ályktunaraðferð &  &  &  \\
\hline
421 & Bayes-setning, setning Bayes &  &  &  \\
\hline
422 & Bayes-tölfræði &  &  &  \\
\hline
423 & bein mynd &  &  &  \\
\hline
424 & bein skekkja &  &  &  \\
\hline
425 & bein sönnun &  &  &  \\
\hline
426 & bein summa &  &  &  \\
\hline
427 & bein þríliða &  &  &  \\
\hline
428 & beinlínu- &  &  &  \\
\hline
429 & beinlínumynd &  &  &  \\
\hline
430 & beinn &  &  &  \\
\hline
431 & beinn eftirkomandi, næsti eftirkomandi &  &  &  \\
\hline
432 & beinn liður &  &  &  \\
\hline
433 & beinn stöðugleiki &  &  &  \\
\hline
434 & beinn undanfari, næsti undanfari &  &  &  \\
\hline
435 & beinn þáttur &  &  &  \\
\hline
436 & beint horn & <ref idTerm=2198>Horn</ref> kallast \textit{beint horn} ef <ref idTerm=334>armar</ref> þess eru <ref idTerm=3400>gagnstæðir</ref>. Slíkt horn <ref idTerm=2223>mælist</ref> 180 <ref idTerm=1667>gráður</ref> eða $\pi$ <ref idTerm=478>radíana</ref> &  &  \\
\hline
437 & beint kerfi &  &  &  \\
\hline
438 & beint margfeldi &  &  &  \\
\hline
439 & beint markgildi &  &  &  \\
\hline
440 & belti &  &  &  \\
\hline
441 & Bernoulli-dreifing, tvíliðudreifing, tvíkostadreifing &  &  &  \\
\hline
442 & Bernoulli-hending &  &  &  \\
\hline
443 & Bernoulli-tala &  &  &  \\
\hline
444 & Bernoulli-tilraun &  &  &  \\
\hline
445 & Bessel-fall, sívalningsfall &  &  &  \\
\hline
446 & besta leikáætlun, besta áætlun &  &  &  \\
\hline
447 & besta nálgun &  &  &  \\
\hline
448 & bestunarverkefni &  &  &  \\
\hline
449 & bestur &  &  &  \\
\hline
450 & betadreifing &  &  &  \\
\hline
451 & betafall &  &  &  \\
\hline
452 & beyging, sveiging &  &  &  \\
\hline
453 & beygjuskilakrosspunktur &  &  &  \\
\hline
454 & biðraðafræði &  &  &  \\
\hline
457 & bil & <ref idTerm=2098>Hlutmengi</ref> $I$ í <ref idTerm=3270>mengi</ref> <ref idTerm=4093>rauntalna</ref> kallast \textit{bil} ef fyrir sérhver $x, y \in I$ og $z \in \mathbb{R}$ þannig að $x < z < y$ gildi að $z \in I$. Ef mengi rauntalna $a$ þannig að $a \leq x$ fyrir öll $x \in I$ á sér efra mark $L$ kallast $L$ \textit{vinstri endapunktur} bilsins. Sömuleiðis ef mengi rauntalna $b$ þannig að $x \leq b$ fyrir öll $x \in I$ á sér neðra mark $U$ þá kallast $U$ \textit{hægri endapunktur} bilsins &  &  \\
\hline
462 & bjagaður, vilhallur, skakkur, hlutdrægur &  &  &  \\
\hline
463 & bjögun &  &  &  \\
\hline
464 & bjögun, togun &  &  &  \\
\hline
465 & blað &  &  &  \\
\hline
466 & blað, þynni &  &  &  \\
\hline
467 & blandað hliðarskilyrði, blönduð þvingun &  &  &  \\
\hline
468 & blandað jaðargildisverkefni &  &  &  \\
\hline
469 & blandað úrtak &  &  &  \\
\hline
470 & blandið lotutugabrot, blandið umferðartugabrot &  &  &  \\
\hline
471 & blandin hlutafleiða, blönduð hlutafleiða &  &  &  \\
\hline
472 & blandin tala &  &  &  \\
\hline
473 & blandinn liður, blandaður liður &  &  &  \\
\hline
474 & blandinn þinur &  &  &  \\
\hline
475 & blendið slembiferli &  &  &  \\
\hline
476 & blöndun &  &  &  \\
\hline
477 & boðleið, rás &  &  &  \\
\hline
478 & radíani, bogaeining, bogamálseining & <ref idTerm=3072>Mælieining</ref> á <ref idTerm=5102>stærð</ref> <ref idTerm=2198>horna</ref>. Einn \textit{radíani} er skilgreindur sem stærð hornsins með <ref idTerm=3661>oddpunkt</ref> í miðju <ref idTerm=2295>hrings</ref> sem fæst þegar að <ref idTerm=479>bogalengd</ref> hringsins er jöfn geisla hans &  & \\
\hline
479 & bogalengd & <ref idTerm=2860>Lengd</ref> <ref idTerm=2276>hringboga</ref> &  &  \\
\hline
480 & bogalengdarfall &  &  &  \\
\hline
481 & bogalengdarfrymi &  &  &  \\
\hline
482 & bogalengdarfrymi &  &  &  \\
\hline
483 & bogamál &  &  &  \\
\hline
484 & bogamínúta &  &  &  \\
\hline
485 & bogasekúnda &  &  &  \\
\hline
486 & bogi, ferilbogi &  &  &  \\
\hline
487 & bogin lína, boglína, sveigð lína, sveiglína &  &  &  \\
\hline
488 & boginn, sveigður &  &  &  \\
\hline
489 & Boole-aðgerð &  &  &  \\
\hline
490 & Boole-algebra &  &  &  \\
\hline
491 & Boole-baugur &  &  &  \\
\hline
492 & Boole-breyta &  &  &  \\
\hline
493 & Boole-fall &  &  &  \\
\hline
494 & Boole-grind &  &  &  \\
\hline
495 & Boole-virki &  &  &  \\
\hline
496 & Borel-algebra &  &  &  \\
\hline
497 & Borel-atburðaalgebra &  &  &  \\
\hline
498 & Borel-fall &  &  &  \\
\hline
499 & Borel-mengi &  &  &  \\
\hline
500 & braut &  &  &  \\
\hline
501 & braut, ferill &  &  &  \\
\hline
502 & brautarúm &  &  &  \\
\hline
503 & breiðboga kótangens &  &  &  \\
\hline
504 & breiðbogafall, gleiðbogafall &  &  &  \\
\hline
505 & breiðbogaflötur &  &  &  \\
\hline
506 & breiðbogakósekans &  &  &  \\
\hline
507 & breiðbogakósínus &  &  &  \\
\hline
508 & breiðbogariti &  &  &  \\
\hline
509 & breiðbogasekans &  &  &  \\
\hline
510 & breiðbogasínus &  &  &  \\
\hline
511 & breiðbogasnúðflötur &  &  &  \\
\hline
512 & breiðbogatangens &  &  &  \\
\hline
513 & breiðbogi, gleiðbogi &  &  &  \\
\hline
514 & breidd &  &  &  \\
\hline
515 & breiðger &  &  &  \\
\hline
516 & breiðger fleygbogaflötur, söðulflötur &  &  &  \\
\hline
517 & breiðger hlutafleiðujafna &  &  &  \\
\hline
518 & breiðger punktur &  &  &  \\
\hline
519 & breiðger Riemann-flötur &  &  &  \\
\hline
520 & breiðger rúmfræði &  &  &  \\
\hline
521 & breiðger samsíðufrumsenda &  &  &  \\
\hline
522 & breiðger sjálfmótun, breiðger færsla &  &  &  \\
\hline
523 & breiðger slétta &  &  &  \\
\hline
524 & breiðger vefja &  &  &  \\
\hline
525 & breiðgerð &  &  &  \\
\hline
526 & breiðgerð &  &  &  \\
\hline
527 & breiðgert rúm &  &  &  \\
\hline
528 & brenniás &  &  &  \\
\hline
529 & brennidepill, brennipunktur &  &  &  \\
\hline
530 & brenniflötur &  &  &  \\
\hline
531 & brennigeisli &  &  &  \\
\hline
532 & brennilengd &  &  &  \\
\hline
533 & brennilína &  &  &  \\
\hline
534 & brennilína &  &  &  \\
\hline
535 & brennimengi &  &  &  \\
\hline
536 & brennimiðstrengur &  &  &  \\
\hline
537 & brennistrengur &  &  &  \\
\hline
538 & breyta, breytistærð, breytileg stærð &  &  &  \\
\hline
539 & breytast &  &  &  \\
\hline
540 & breytast með samfelldum hætti &  &  &  \\
\hline
541 & breytilegur &  &  &  \\
\hline
542 & breytingarhraði &  &  &  \\
\hline
543 & breytistuðull, fráviksstuðull &  &  &  \\
\hline
544 & breytulaus &  &  &  \\
\hline
545 & breytuskipti, innsetning &  &  &  \\
\hline
546 & brigðgengur &  &  &  \\
\hline
547 & brjóstvitsregla, þumalfingursregla &  &  &  \\
\hline
548 & broddur &  &  &  \\
\hline
549 & brot &  &  &  \\
\hline
550 & brotabaugur &  &  &  \\
\hline
551 & brotabrot &  &  &  \\
\hline
552 & brotaregla, hlutfallsregla &  &  &  \\
\hline
553 & brotareikningur &  &  &  \\
\hline
554 & brotareikningur &  &  &  \\
\hline
555 & brotastrik, lárétt brotastrik &  &  &  \\
\hline
556 & brotastrik, skástrik, brotaskástrik &  &  &  \\
\hline
557 & brotasvið &  &  &  \\
\hline
558 & brothluti &  &  &  \\
\hline
559 & broti &  &  &  \\
\hline
560 & brotið íðal &  &  &  \\
\hline
561 & brotið línustrik &  &  &  \\
\hline
562 & brotin línuleg færsla, Möbíusarfærsla &  &  &  \\
\hline
563 & brotin línuleg grúpa &  &  &  \\
\hline
564 & brotin tala, brot &  &  &  \\
\hline
565 & brotinn línulegur &  &  &  \\
\hline
566 & brotinn veldisvísir &  &  &  \\
\hline
567 & brotinn, brota- &  &  &  \\
\hline
568 & brotmark, hlutfallsmark &  &  &  \\
\hline
569 & brotpunktur &  &  &  \\
\hline
570 & brottskurður &  &  &  \\
\hline
571 & Brown-hreyfing, brownsk hreyfing &  &  &  \\
\hline
572 & brú &  &  &  \\
\hline
573 & brúa, inngiska, innreikna &  &  &  \\
\hline
574 & brugðningafræði &  &  &  \\
\hline
575 & brugðningagrúpa &  &  &  \\
\hline
576 & brugðningur &  &  &  \\
\hline
577 & brún &  &  &  \\
\hline
578 & brún & <ref idTerm=2958>Línustrik</ref> þar sem tvær <ref idTerm=1325>hliðar</ref> <ref idTerm=6454>þrívíðs</ref> <ref idTerm=3133>margflötungs</ref> <ref idTerm=4771>skerast</ref> &  & Ræða almenna margflötunga í skýringu \\
\hline
579 & brúnahorn &  &  &  \\
\hline
580 & brúnahorn &  &  &  \\
\hline
581 & brúnalengd &  &  &  \\
\hline
582 & brúun, innreiknun, inngiskun &  &  &  \\
\hline
583 & brúunaraðferð &  &  &  \\
\hline
584 & brúunarfall &  &  &  \\
\hline
585 & brúunarskekkja &  &  &  \\
\hline
586 & bryddað fylki &  &  &  \\
\hline
587 & brydding &  &  &  \\
\hline
588 & bryddingaraðferð &  &  &  \\
\hline
589 & bryggja, útgizka, útreikna &  &  &  \\
\hline
590 & bryggjun, útgiskun, útreiknun &  &  &  \\
\hline
591 & bryggjunarfall &  &  &  \\
\hline
592 & bryggjunarskekkja &  &  &  \\
\hline
593 & bundin &  &  &  \\
\hline
594 & bundin &  &  &  \\
\hline
595 & bundin af hvelum, hvelabundin &  &  &  \\
\hline
596 & bundin af línum, línubundin &  &  &  \\
\hline
597 & bundin af samsíða línum, samsíðubundin &  &  &  \\
\hline
598 & bundin af sléttum, sléttubundin &  &  &  \\
\hline
599 & bundin breyta, bundin breytistærð, leppbreyta &  &  &  \\
\hline
600 & bundinn Riemann-flötur &  &  &  \\
\hline
601 & bundinn vigur &  &  &  \\
\hline
602 & burðarmengi &  &  &  \\
\hline
603 & burður, flutningur &  &  &  \\
\hline
604 & burtséð frá &  &  &  \\
\hline
605 & bútaaðferð, bútunaraðferð &  &  &  \\
\hline
606 & bútur &  &  &  \\
\hline
607 & bútur, bogi &  &  &  \\
\hline
608 & byggja, byggja upp, búa til, smíða &  &  &  \\
\hline
609 & bylgja, alda &  &  &  \\
\hline
610 & bylgjujafna, öldujafna &  &  &  \\
\hline
611 & bylta &  &  &  \\
\hline
612 & bylta (fylkis), bylt fylki &  &  &  \\
\hline
613 & bylting &  &  &  \\
\hline
614 & byrjunarlausn &  &  &  \\
\hline
615 & Cantor-mengi, Cantorsmengi, talnamengi Cantors, þríundamengi Cantors &  &  &  \\
\hline
616 & Cauchy-dreifing &  &  &  \\
\hline
617 & Cauchy-formúla &  &  &  \\
\hline
618 & Cauchy-margfeldi &  &  &  \\
\hline
619 & Cauchy-margföldun &  &  &  \\
\hline
620 & Cauchy-próf, samleitnipróf Cauchys &  &  &  \\
\hline
621 & Cauchy-runa &  &  &  \\
\hline
622 & Cauchy-setningin &  &  &  \\
\hline
623 & Cayley-tala, áttundatala &  &  &  \\
\hline
624 & Ceva-strik &  &  &  \\
\hline
625 & Ceva-þríhyrningur &  &  &  \\
\hline
626 & Cholesky-þáttun &  &  &  \\
\hline
627 & Cramer-regla, regla Cramers &  &  &  \\
\hline
628 & d'Alembert-virki &  &  &  \\
\hline
629 & dæmi, sýnidæmi &  &  &  \\
\hline
630 & dæmi, þraut, reikningsdæmi, æfingardæmi &  &  &  \\
\hline
631 & dæmigerður &  &  &  \\
\hline
632 & dæmigert stak &  &  &  \\
\hline
633 & dæmigert úrtak &  &  &  \\
\hline
634 & dálkaðgerð &  &  &  \\
\hline
635 & dálkauppstokkun, dálkaumröðun &  &  &  \\
\hline
636 & dálkfylki, eindálkafylki, súlufylki &  &  &  \\
\hline
637 & dálkmeðaltal &  &  &  \\
\hline
638 & dálkstétt &  &  &  \\
\hline
639 & dálksumma, dálksamtala &  &  &  \\
\hline
640 & dálkur, súla &  &  &  \\
\hline
641 & dálkvigur, súluvigur &  &  &  \\
\hline
642 & dálkvísir &  &  &  \\
\hline
643 & dánartíðni &  &  &  \\
\hline
644 & Dedekind-baugur &  &  &  \\
\hline
645 & Dedekind-snið &  &  &  \\
\hline
646 & deila &  &  &  \\
\hline
647 & deilaflokkur &  &  &  \\
\hline
648 & deilanlegur &  &  &  \\
\hline
649 & deilanleiki &  &  &  \\
\hline
650 & deilasumma &  &  &  \\
\hline
651 & deilda- og heildareikningur, diffur- og tegurreikningur, örsmæðareikningur, reiknivísi &  &  &  \\
\hline
652 & deilda, diffra &  &  &  \\
\hline
653 & deildaalgebra &  &  &  \\
\hline
654 & deildabaugur &  &  &  \\
\hline
655 & deildaform, diffurform &  &  &  \\
\hline
656 & deildagrannmynstur &  &  &  \\
\hline
657 & deildagrúpa &  &  &  \\
\hline
658 & deildajafna, afleiðujafna, diffurjafna &  &  &  \\
\hline
659 & deildamengi &  &  &  \\
\hline
660 & deildamóta, diffurmóta &  &  &  \\
\hline
661 & deildamótull &  &  &  \\
\hline
662 & deildamótun, diffurmótun &  &  &  \\
\hline
663 & deildamyndun &  &  &  \\
\hline
664 & deildamynstur &  &  &  \\
\hline
665 & deildanleg víðátta, diffranleg víðátta &  &  &  \\
\hline
666 & deildanlegur á köflum, diffranlegur á köflum &  &  &  \\
\hline
667 & deildanlegur frá hægri, diffranlegur frá hægri &  &  &  \\
\hline
668 & deildanlegur frá vinstri, diffranlegur frá vinstri &  &  &  \\
\hline
669 & deildanlegur í punkti, diffranlegur í punkti &  &  &  \\
\hline
670 & deildanlegur, diffranlegur &  &  &  \\
\hline
671 & deildanleiki, diffranleiki &  &  &  \\
\hline
672 & deildareikningur, diffurreikningur &  &  &  \\
\hline
673 & deildarúm &  &  &  \\
\hline
675 & deildarúmfræði Riemanns, Riemann-deildarúmfræði &  &  &  \\
\hline
676 & deildarúmfræði, diffurrúmfræði &  &  &  \\
\hline
677 & deildastuðull, deildahlutfall, diffurkvóti, afleiða í punkti &  &  &  \\
\hline
678 & deildavensl &  &  &  \\
\hline
679 & deildavirki, diffurvirki, afleiðuvirki &  &  &  \\
\hline
680 & deildi, diffur &  &  &  \\
\hline
681 & deildun, diffrun &  &  &  \\
\hline
682 & deildunarvirki, diffrunarvirki &  &  &  \\
\hline
683 & deilibaugur &  &  &  \\
\hline
684 & deiling &  &  &  \\
\hline
685 & deiling án afgangs & <ref idTerm=686>Deiling</ref> <ref idTerm=3475>náttúrulegrar tölu</ref> með annarri náttúrulegri tölu þar sem afgangurinn er 0 &  &  \\
\hline
686 & deiling með afgangi & Ef $n$ og $d$ eru <ref idTerm=1896>heilar tölur</ref> þannig að $d \neq 0$ þá er til heil tala $q$ og <ref idTerm=3475>náttúruleg tala</ref> $r$ sem uppfylla \[n = q \cdot d + r \quad \text{og} \quad 0 \leq r < |d|.\] Að rita $n$ á þessu formi nefnist að \textit{deila} <ref idTerm=5397>tölunni</ref> $d$ upp í $n$ og kallast þá talan $q$ \textit{kvóti} deilingarinnar og talan $r$ \textit{afgangur} hennar &  & Óþarfa óþjál skilgreining kannski \\
\hline
687 & deilingaraðferð, deilingarreiknirit &  &  &  \\
\hline
688 & deilingarmerki &  &  &  \\
\hline
689 & deilir &  &  &  \\
\hline
690 & deilistofn &  &  &  \\
\hline
691 & deltafelli, deltafall, deltafelli Diracs, deltafall Diracs &  &  &  \\
\hline
692 & deltaferill &  &  &  \\
\hline
693 & depilfeldi, innfeldi, innra margfeldi &  &  &  \\
\hline
694 & Desargues-, argeskur &  &  &  \\
\hline
695 & Desargues-setningin meiri, stóra Desargues-setningin &  &  &  \\
\hline
696 & Desargues-setningin minni, litla Desargues-setningin &  &  &  \\
\hline
697 & Desargues-slétta, argesk slétta &  &  &  \\
\hline
698 & deyfð sveifla &  &  &  \\
\hline
699 & Díofantosargreining &  &  &  \\
\hline
700 & Díofantosarjafna &  &  &  \\
\hline
701 & Dirichlet-dreifing &  &  &  \\
\hline
702 & Dirichlet-röð &  &  &  \\
\hline
703 & draga &  &  &  \\
\hline
704 & draga (rót) &  &  &  \\
\hline
705 & draga frá &  &  &  \\
\hline
706 & draga saman &  &  &  \\
\hline
707 & dráttarferill &  &  &  \\
\hline
708 & dráttur án skila &  &  &  \\
\hline
709 & dráttur með skilum &  &  &  \\
\hline
710 & draugalest &  &  &  \\
\hline
711 & dreif, dreififelli, dreififall &  &  &  \\
\hline
712 & dreifafræði, dreififellafræði &  &  &  \\
\hline
713 & dreifarlausn, dreififellislausn &  &  &  \\
\hline
714 & dreifður &  &  &  \\
\hline
715 & dreifing &  &  &  \\
\hline
716 & dreifing, líkindadreifing &  &  &  \\
\hline
717 & dreifingafjölskylda &  &  &  \\
\hline
718 & dreifingarfall, dreififall, líkindadreifingarfall &  &  &  \\
\hline
719 & dreifingarferill &  &  &  \\
\hline
720 & dreifingargerð &  &  &  \\
\hline
721 & dreifingaróháður &  &  &  \\
\hline
722 & dreifinn & <ref idTerm=46>Reikniaðgerð</ref> $\odot$ er sögð vera \textit{dreifin} yfir $\oplus$ ef $\odot$ og $\oplus$ fullnægja <ref idTerm=723>dreifireglunni</ref> &  &  \\
\hline
723 & dreifiregla & Ef $\odot$ og $\oplus$ eru <ref idTerm=5702>tvístæðar reikniaðgerðir</ref> á <ref idTerm=>mengi</ref> $X$ er sagt að $\odot$ og $\oplus$ fullnægi \textit{dreifireglunni} ef fyrir sérhver stök $x, y, z \in X$ gildir að \[ x \odot (y \oplus z) = (x \odot y) \oplus (x \odot z) \] og \[ (y \oplus z) \odot x = (y \odot x) \oplus (z \odot x) \] &  & Setja í skýringu: Sagt að $\odot$ og $\oplus$ fullnægi dreifireglunni frá vinstri ef fyrir sérhver stök $x, y, z \in X$ gildir að \[ x \odot (y \oplus z) = (x \odot y) \oplus (x \odot z). \] Á sama hátt er sagt að $\odot$ og $\oplus$ fullnægi dreifireglunni frá hægri ef fyrir sérhver stök $x, y, z \in X$ gildir að \[ (y \oplus z) \odot x = (y \odot x) \oplus (z \odot x). \] Því er sagt að $\odot$ og $\oplus$ fullnægi dreifireglunni ef $\odot$ og $\oplus$ fullnægi dreifireglunni bæði frá vinstri og hægri. Minnast á tengsl við víxlni $\odot$ \\
\hline
724 & dreifirit, punktaský, punktarit &  &  &  \\
\hline
725 & dreifni &  &  &  \\
\hline
726 & dreifnihlutfall &  &  &  \\
\hline
727 & dreift hlutmengi &  &  &  \\
\hline
728 & dreifvirðing &  &  &  \\
\hline
729 & dreifvirðingarbaugur &  &  &  \\
\hline
730 & drómasetning &  &  &  \\
\hline
731 & drómi &  &  &  \\
\hline
732 & Dynkin-rit &  &  &  \\
\hline
733 & Dynkin-safn &  &  &  \\
\hline
734 & dýpt &  &  &  \\
\hline
735 & eða &  &  &  \\
\hline
736 & eða-tákn &  &  &  \\
\hline
737 & eðun & Fyrir sérhverjar <ref idTerm=1515>yrðingar</ref> $p$ og $q$ er $p \vee q$ (lesið: „$p$ eða $q$“) sú yrðing sem segir að að minnsta kosti önnur yrðinganna $p$ og $q$ sé <ref idTerm=4565>sönn</ref>. Yrðingin $p \vee q$ er þess vegna <ref idTerm=3890>ósönn</ref> þegar yrðingarnar $p$ og $q$ eru báðar ósannar, en annars er hún sönn &  & Setja í skýringu: Þessari staðreynd má lýsa með töflunni að neðan, sem sýnir sanngildi $p \vee q$ eftir ólíkum sanngildum $p$ annars vegar og $q$ hins vegar. (ásamt mynd) \\
\hline
738 & eðunarhluti, eðunarliður &  &  &  \\
\hline
739 & ef og aðeins ef, þá og því aðeins að &  &  &  \\
\hline
740 & efra flokksmark &  &  &  \\
\hline
741 & efra flokksmark &  &  &  \\
\hline
742 & efra heildi, efra tegur &  &  &  \\
\hline
743 & efra mark, hámark &  &  &  \\
\hline
744 & efra markgildi, stærsti þéttipunktur, efsti þéttipunktur &  &  &  \\
\hline
745 & efra öryggisbil &  &  &  \\
\hline
746 & efra stýrimark &  &  &  \\
\hline
747 & efra stýrimark &  &  &  \\
\hline
748 & efra þríhyrningsfylki &  &  &  \\
\hline
749 & efri miðjuruna, efri miðjuturn &  &  &  \\
\hline
750 & efri stallagerð, efri þrepagerð, efri þrepaskipan &  &  &  \\
\hline
751 & efri stéttar rökfræði &  &  &  \\
\hline
752 & efri þríhyrningsgrúpa &  &  &  \\
\hline
753 & efsta gildi, hæsta gildi, stærsta gildi, stærsta hágildi, víðfeðmt hágildi &  &  &  \\
\hline
754 & eftir á &  &  &  \\
\hline
755 & eftirádreifing &  &  &  \\
\hline
756 & eftirálíkindi &  &  &  \\
\hline
757 & eftirámat &  &  &  \\
\hline
758 & eftirfararaðtala &  &  &  \\
\hline
759 & eftirfari &  &  &  \\
\hline
760 & eftirlit &  &  &  \\
\hline
761 & eftirlitsrit &  &  &  \\
\hline
762 & eftirtáknun &  &  &  \\
\hline
763 & egg-, egglaga &  &  &  \\
\hline
764 & eggferill, eggbaugur &  &  &  \\
\hline
765 & eggflötur &  &  &  \\
\hline
766 & eggin á fleygnum &  &  &  \\
\hline
767 & eiði, slitmengi &  &  &  \\
\hline
768 & eiginaðhvarf, sjálfsaðhvarf &  &  &  \\
\hline
769 & eiginaðhverfur, sjálfsaðhverfur &  &  &  \\
\hline
770 & eigindlegur &  &  &  \\
\hline
771 & eigindlegur eiginleiki &  &  &  \\
\hline
772 & eiginfall &  &  &  \\
\hline
773 & eiginfylgni, sjálffylgni &  &  &  \\
\hline
774 & eiginfylgnifall, sjálffylgnifall &  &  &  \\
\hline
775 & eiginfylgnirit, sjálffylgnirit &  &  &  \\
\hline
776 & eiginfylgnistuðull, sjálffylgnistuðull &  &  &  \\
\hline
777 & eigingildi &  &  &  \\
\hline
778 & eigingildisverkefni &  &  &  \\
\hline
779 & eiginlausn &  &  &  \\
\hline
780 & eiginleg frumsenda &  &  &  \\
\hline
781 & eiginleg hlutgrúpa &  &  &  \\
\hline
782 & eiginleg hlutvera &  &  &  \\
\hline
783 & eiginleg ósamleitni &  &  &  \\
\hline
784 & eiginleg tvinntala, þverkynja tala, þverkyns tala &  &  &  \\
\hline
785 & eiginleg útvíkkun, eiginleg víkkun &  &  &  \\
\hline
786 & eiginleg vörpun &  &  &  \\
\hline
787 & eiginleg þverstaðalgrúpa, snúningagrúpa, snúningsgrúpa &  &  &  \\
\hline
788 & eiginleg þverstaðalvörpun, snúningur &  &  &  \\
\hline
789 & eiginlega ósamleitinn &  &  &  \\
\hline
790 & eiginlega samsniða &  &  &  \\
\hline
791 & eiginlegt brot &  &  &  \\
\hline
792 & eiginlegt hlutmengi &  &  &  \\
\hline
793 & eiginlegt hlutrúm &  &  &  \\
\hline
794 & eiginlegt hlutsvið &  &  &  \\
\hline
795 & eiginlegt rætt fall &  &  &  \\
\hline
796 & eiginlegt þverstaðlað fylki, snúningsfylki &  &  &  \\
\hline
797 & eiginlegur &  &  &  \\
\hline
798 & eiginlegur deilir &  &  &  \\
\hline
799 & eiginlegur flokkur &  &  &  \\
\hline
800 & eiginlegur flutningur, eiginleg hreyfing &  &  &  \\
\hline
801 & eiginlegur óendanleiki &  &  &  \\
\hline
802 & eiginlegur punktur &  &  &  \\
\hline
803 & eiginleiki &  &  &  \\
\hline
804 & eiginleikinn um endanleg snið &  &  &  \\
\hline
805 & eiginrófsþéttleiki &  &  &  \\
\hline
806 & eiginrúm &  &  &  \\
\hline
807 & eiginslétta &  &  &  \\
\hline
808 & eiginvídd, rúmfræðileg margfeldni &  &  &  \\
\hline
809 & eiginvigrafylki &  &  &  \\
\hline
810 & eiginvigur, eiginvektor &  &  &  \\
\hline
811 & eileif &  &  &  \\
\hline
812 & einangraður &  &  &  \\
\hline
813 & einangraður hluti &  &  &  \\
\hline
814 & einangraður punktur &  &  &  \\
\hline
815 & einangraður sérstöðupunktur &  &  &  \\
\hline
816 & einblaða breiðbogaflötur, einskála breiðbogaflötur &  &  &  \\
\hline
817 & eind &  &  &  \\
\hline
818 & eindagrúpa, margföldunargrúpa &  &  &  \\
\hline
819 & eindregið frumsendukerfi, eindregið frumsetningakerfi &  &  &  \\
\hline
820 & eindreginn &  &  &  \\
\hline
821 & einfalda &  &  &  \\
\hline
822 & einfalda &  &  &  \\
\hline
823 & einfaldlega gegnvirkur &  &  &  \\
\hline
824 & einfaldlega lotubundinn &  &  &  \\
\hline
825 & einfaldlega samanhangandi, einfaldlega samhangandi &  &  &  \\
\hline
826 & einfaldur &  &  &  \\
\hline
827 & einfaldur atburður, frumatburður &  &  &  \\
\hline
828 & einfaldur baugur &  &  &  \\
\hline
829 & einfaldur bogi, einfaldur ferill &  &  &  \\
\hline
830 & einfaldur hnútur, smárahnútur &  &  &  \\
\hline
831 & einfaldur leggur &  &  &  \\
\hline
832 & einfaldur marghyrningur &  &  &  \\
\hline
833 & einfaldur punktur &  &  &  \\
\hline
834 & einfaldur vegur, braut &  &  &  \\
\hline
835 & einfaldur, grunnstæður &  &  &  \\
\hline
836 & einfalt eigingildi &  &  &  \\
\hline
837 & einfalt fall &  &  &  \\
\hline
838 & einfalt fylki, frumfylki &  &  &  \\
\hline
839 & einfalt heildi, einvítt heildi, einfalt tegur, einvítt tegur &  &  &  \\
\hline
840 & einfalt net &  &  &  \\
\hline
841 & einfalt róf &  &  &  \\
\hline
842 & einfalt samhverft fall &  &  &  \\
\hline
843 & einfalt skaut, einfaldur póll &  &  &  \\
\hline
844 & einfalt úrtak, hreint úrtak &  &  &  \\
\hline
845 & einfari &  &  &  \\
\hline
846 & einföld aðgerð &  &  &  \\
\hline
847 & einföld hending &  &  &  \\
\hline
848 & einföld núllstöð &  &  &  \\
\hline
849 & einföld örvarás &  &  &  \\
\hline
850 & einföld rás &  &  &  \\
\hline
851 & einföld rót &  &  &  \\
\hline
852 & einföld rúmfræði, grunnstæð rúmfræði &  &  &  \\
\hline
853 & einföld talnafræði, grunnstæð talnafræði &  &  &  \\
\hline
854 & einföld úrtaka, hrein úrtaka &  &  &  \\
\hline
855 & einföld útvíkkun, einföld víkkun &  &  &  \\
\hline
856 & einföld þríliða &  &  &  \\
\hline
857 & einföldun &  &  &  \\
\hline
858 & einföldun &  &  &  \\
\hline
859 & einföldunarregla &  &  &  \\
\hline
860 & einföldunarregla (fyrir hornaföll) &  &  &  \\
\hline
861 & eingild vörpun &  &  &  \\
\hline
862 & eingildur &  &  &  \\
\hline
863 & eingilt fall &  &  &  \\
\hline
864 & einhalla &  &  &  \\
\hline
865 & einhalla á köflum &  &  &  \\
\hline
866 & einhalla fall & <ref idTerm=1078>Fall</ref> $f: X \to Y$ þar sem $X$ og $Y$ eru <ref idTerm=2098>hlutmengi</ref> í <ref idTerm=3270>mengi</ref> <ref idTerm=4093>rauntalna</ref> er sagt vera \textit{einhalla} ef það er <ref idTerm=0>vaxandi</ref> eða <ref idTerm=0>minnkandi</ref> &  &  \\
\hline
867 & einhalla runa &  &  &  \\
\hline
868 & einhalla samleitni &  &  &  \\
\hline
869 & einhalli &  &  &  \\
\hline
870 & einhliða &  &  &  \\
\hline
871 & einhliða afleiða &  &  &  \\
\hline
872 & einhliða deildanleiki, einhliða diffranleiki &  &  &  \\
\hline
873 & einhliða flötur &  &  &  \\
\hline
874 & einhliða íðal &  &  &  \\
\hline
875 & einhliða markgildi &  &  &  \\
\hline
876 & einhliða markgildi &  &  &  \\
\hline
877 & einhliða öryggisbil &  &  &  \\
\hline
878 & einhliða próf &  &  &  \\
\hline
879 & einhliða samfelldni, samfelldni frá annarri hlið &  &  &  \\
\hline
880 & einhluta &  &  &  \\
\hline
881 & einhluta keilusnið &  &  &  \\
\hline
882 & einingarbil &  &  &  \\
\hline
883 & einingarfall &  &  &  \\
\hline
884 & einingarferningur &  &  &  \\
\hline
885 & einingarfylki &  &  &  \\
\hline
886 & einingarhringskífa, einingarskífa &  &  &  \\
\hline
887 & einingarhringur &  &  &  \\
\hline
888 & einingarhvel &  &  &  \\
\hline
889 & einingarkúla &  &  &  \\
\hline
890 & einingarrót &  &  &  \\
\hline
891 & einingarskipting, einingarhlutun, hlutun á einum &  &  &  \\
\hline
892 & einingarvigur, einingarvektor &  &  &  \\
\hline
893 & einkenna &  &  &  \\
\hline
895 & einkenni, eigind &  &  &  \\
\hline
896 & einkryppudreifing &  &  &  \\
\hline
897 & einliða & <ref idTerm=1948>Margliða</ref> sem hefur aðeins einn lið &  &  \\
\hline
898 & einliða jafna &  &  &  \\
\hline
899 & einmargur &  &  &  \\
\hline
900 & einmáta fylki &  &  &  \\
\hline
901 & einmáta grúpa &  &  &  \\
\hline
902 & einmenningsspil, einmenningur &  &  &  \\
\hline
903 & einmörg samsvörun &  &  &  \\
\hline
904 & einmótull &  &  &  \\
\hline
905 & einmótun, ímótun &  &  &  \\
\hline
906 & einn &  &  &  \\
\hline
907 & einn og aðeins einn, nákvæmlega einn &  &  &  \\
\hline
908 & einoka fylki &  &  &  \\
\hline
909 & einoka grúpa &  &  &  \\
\hline
910 & einoka grúpa yfir rauntölurnar, raunþverstaðalgrúpa, þverstaðalgrúpa yfir rauntölurnar &  &  &  \\
\hline
911 & einoka grúpa yfir tvinntölurnar &  &  &  \\
\hline
912 & einoka jafngildur &  &  &  \\
\hline
913 & einoka virki &  &  &  \\
\hline
914 & einpunktungur, einspunktsrúm &  &  &  \\
\hline
915 & einrás, \(1\)-rás &  &  &  \\
\hline
916 & einsátta &  &  &  \\
\hline
917 & einsátta rúm &  &  &  \\
\hline
918 & einsetupunktur &  &  &  \\
\hline
919 & einsfirða &  &  &  \\
\hline
920 & einsfirðanagrúpa &  &  &  \\
\hline
921 & einsfirðun, firðrækin vörpun &  &  &  \\
\hline
922 & einshyrndur & Tveir <ref idTerm=6438>þríhyrningar</ref> kallast \textit{einshyrndir} ef hægt er að para <ref idTerm=2198>horn</ref> annars þríhyrningsins við horn hins þríhyrningsins þannig að hornin í sama pari eru jafn stór &  &  \\
\hline
923 & einskorða &  &  &  \\
\hline
924 & einskorðað grannmynstur, hlutrúmsgrannmynstur &  &  &  \\
\hline
925 & einskorðarður &  &  &  \\
\hline
926 & einskorðuð lokun &  &  &  \\
\hline
927 & einskorðun & Fyrir <ref idTerm=3282>vörpun</ref> $f: X \to Y$ kallast vörpunin \[ f|_A : A \to Y; \quad f|_A(x) = f(x) \] þar sem $A$ er <ref idTerm=2098>hlutmengi</ref> í $X$ \textit{einskorðun} vörpunarinnar $f$ við mengið $A$ &  &  \\
\hline
928 & einskorðunarmótun &  &  &  \\
\hline
929 & einskrefa ítrekun &  &  &  \\
\hline
930 & einslæg horn &  &  &  \\
\hline
931 & einslaga þríhyrningur, einshyrndur þríhyrningur &  &  &  \\
\hline
932 & einslaga, einslagaður & Tvær <ref idTerm=6532>rúmmyndir</ref> $A$ og $B$ í <ref idTerm=1034>evklíðsku rúmi</ref> $\mathbb{E}$ kallast \textit{einslaga} ef til er <ref idTerm=5779>færsla</ref> $\phi$ á $\mathbb{E}$ þannig að $\phi(A)$ og $B$ séu <ref idTerm=0>eins</ref> rúmmyndir &  &  \\
\hline
933 & einsleit deildajafna &  &  &  \\
\hline
934 & einsleitt fall &  &  &  \\
\hline
935 & einsleitt íðal &  &  &  \\
\hline
936 & einsleitt Markov-ferli &  &  &  \\
\hline
937 & einsleitt rúm &  &  &  \\
\hline
938 & einsleitt stak &  &  &  \\
\hline
939 & einsleitur &  &  &  \\
\hline
940 & einsleitur &  &  &  \\
\hline
941 & einsleitur þáttur &  &  &  \\
\hline
942 & einslögun &  &  &  \\
\hline
943 & einslögunarfærsla & <ref idTerm=5779>Færsla</ref> $\phi$ á tiltekinni <ref idTerm=4884>sléttu</ref> kallast \textit{einslögunarfærsla} ef <ref idTerm=2072>hlutfallið</ref> milli <ref idTerm=2860>lengda</ref> <ref idTerm=2958>strikanna</ref> $\phi(AB)$ og $AB$ er það sama fyrir öll strik $AB$, sem kallast þá \textit{kvarði} einslögunarfærslunnar &  &  \\
\hline
944 & einslögunarmiðja, stríkkunarmiðja &  &  &  \\
\hline
945 & einsmóta &  &  &  \\
\hline
946 & einsmótun &  &  &  \\
\hline
947 & einsmótun &  &  &  \\
\hline
948 & einsmótun &  &  &  \\
\hline
949 & einsmótunarflokkur &  &  &  \\
\hline
950 & einsmótunarsetning &  &  &  \\
\hline
951 & einstæð aðgerð & <ref idTerm=46>Aðgerð</ref> á <ref idTerm=3270>mengi</ref> $X$ af gerðinni $X \to X$ kallast \textit{einstæð aðgerð} &  &  \\
\hline
952 & einstæð umsögn &  &  &  \\
\hline
953 & einstæð vensl &  &  &  \\
\hline
954 & einstæður &  &  &  \\
\hline
955 & einstafa tala &  &  &  \\
\hline
956 & einstaklingur &  &  &  \\
\hline
957 & einstika fjölskylda &  &  &  \\
\hline
958 & einstika grúpa &  &  &  \\
\hline
959 & einstikun &  &  &  \\
\hline
960 & einstikunarsetning &  &  &  \\
\hline
961 & einstoga &  &  &  \\
\hline
962 & einstogun &  &  &  \\
\hline
963 & einstökungur, einsstaksmengi & <ref idTerm=3270>Mengi</ref> með nákvæmlega eitt <ref idTerm=5128>stak</ref> &  &  \\
\hline
964 & einsþrepsaðferð &  &  &  \\
\hline
965 & eintæk vörpun & <ref idTerm=3282>Vörpun</ref> sem varpar ólíkum <ref idTerm=5128>stökum</ref> úr <ref idTerm=1343>formenginu</ref> í ólík stök úr <ref idTerm=74>bakmenginu</ref> &  &  \\
\hline
966 & eintækt fall &  &  &  \\
\hline
967 & eintækur &  &  &  \\
\hline
968 & eintekt, eintaka &  &  &  \\
\hline
969 & einungamótun &  &  &  \\
\hline
970 & einungur, hálfgrúpa með hlutleysu &  &  &  \\
\hline
971 & einvalda &  &  &  \\
\hline
972 & einvíð dreifing &  &  &  \\
\hline
973 & einvíð stærð, tölustærð, stigastærð &  &  &  \\
\hline
974 & einvíður &  &  &  \\
\hline
975 & einvítt hvel, einhvel &  &  &  \\
\hline
976 & einþríhyrningsfylki &  &  &  \\
\hline
977 & ekki augljós &  &  &  \\
\hline
978 & ekki fastur &  &  &  \\
\hline
979 & ekki í jöfnum mæli &  &  &  \\
\hline
980 & ekki jákvæður, frekar neikvæður, óviðlægur &  &  &  \\
\hline
981 & ekki neikvæður, frekar jákvæður, ófrádrægur &  &  &  \\
\hline
982 & ekki núll &  &  &  \\
\hline
983 & ekki núlldeilir, reglulegt stak &  &  &  \\
\hline
984 & ekki ótvíræður, ekki einræður &  &  &  \\
\hline
985 & ekki tómur &  &  &  \\
\hline
986 & eltingarverkefni &  &  &  \\
\hline
987 & endahnútur &  &  &  \\
\hline
988 & endahnútur, endapunktur &  &  &  \\
\hline
989 & endanleg aðferð &  &  &  \\
\hline
990 & endanleg fjöldatala &  &  &  \\
\hline
991 & endanleg framsetning &  &  &  \\
\hline
992 & endanleg röð &  &  &  \\
\hline
993 & endanleg runa &  &  &  \\
\hline
994 & endanleg slétta &  &  &  \\
\hline
995 & endanleg tala &  &  &  \\
\hline
996 & endanleg útvíkkun, endanleg víkkun &  &  &  \\
\hline
997 & endanleg þrepun &  &  &  \\
\hline
998 & endanlega framsettur &  &  &  \\
\hline
999 & endanlega frumsendanlegur &  &  &  \\
\hline
1000 & endanlega margir &  &  &  \\
\hline
1001 & endanlega samanhangandi, endanlega samhangandi &  &  &  \\
\hline
1002 & endanlega spannaður, endanlega framleiddur &  &  &  \\
\hline
1003 & endanlega viðbætinn, endanlega samleggjandi &  &  &  \\
\hline
1004 & endanlega víður &  &  &  \\
\hline
1005 & endanlega víður &  &  &  \\
\hline
1006 & endanlegt mengi &  &  &  \\
\hline
1007 & endanlegt svið, Galois-svið &  &  &  \\
\hline
1008 & endanlegt tugabrot &  &  &  \\
\hline
1009 & endanlegur & <ref idTerm=3270>Mengi</ref> kallast \textit{endanlegt} ef <ref idTerm=1211>fjöldatala</ref> þess er <ref idTerm=3475>náttúrleg tala</ref> &  &  \\
\hline
1010 & endanlegur fjöldi &  &  &  \\
\hline
1011 & endanlegur munur &  &  &  \\
\hline
1012 & endanleikahyggja &  &  &  \\
\hline
1013 & endanleiki &  &  &  \\
\hline
1014 & endapunktur &  &  &  \\
\hline
1015 & endi &  &  &  \\
\hline
1016 & endi &  &  &  \\
\hline
1017 & endurkoma, afturkoma &  &  &  \\
\hline
1018 & endurkomutími, endurkomustund &  &  &  \\
\hline
1019 & endurkvæmt ástand, afturkvæmt ástand &  &  &  \\
\hline
1020 & endurkvæmt ferli, afturkvæmt ferli &  &  &  \\
\hline
1021 & endurkvæmur, afturkvæmur &  &  &  \\
\hline
1022 & endurnýjunarferli &  &  &  \\
\hline
1023 & endurnýjunarpunktferli &  &  &  \\
\hline
1024 & endurtölusetning, tölusetning á nýjan leik &  &  &  \\
\hline
1025 & endurúrtaka &  &  &  \\
\hline
1026 & Euler- &  &  &  \\
\hline
1027 & Euler-aðferð &  &  &  \\
\hline
1028 & Euler-net &  &  &  \\
\hline
1029 & evklíðsk firð &  &  &  \\
\hline
1030 & evklíðsk rúmfræði & <ref idTerm=4271>Rúmfræði</ref> sem byggir á <ref idTerm=1452>frumsendum</ref> forngríska <ref idTerm=5117>stærðfræðingsins</ref> Evklíðs &  & Setja fram frumsendurnar í skýringu \\
\hline
1031 & evklíðsk slétta & <ref idTerm=5729>Tvívítt</ref> <ref idTerm=4266>rúm</ref> sem uppfyllir frumsendur <ref idTerm=1030>Evklíðskar rúmfræði</ref> &  &  \\
\hline
1032 & evklíðskt innfeldi, innfeldi &  &  &  \\
\hline
1033 & evklíðskt innfeldisrúm, evklíðskt rúm, innfeldisrúm &  &  &  \\
\hline
1034 & evklíðskt rúm & evklíðsk slétta & <ref idTerm=4266>Rúm</ref> sem uppfyllir frumsendur <ref idTerm=1030>Evklíðskar rúmfræði</ref>  &  &  \\
\hline
1035 & evklíðskur baugur &  &  &  \\
\hline
1036 & evklíðskur flutningur &  &  &  \\
\hline
1037 & evklíðskur staðall &  &  &  \\
\hline
1038 & evklíðskur, Evklíðs- &  &  &  \\
\hline
1039 & eyða &  &  &  \\
\hline
1040 & eyðingaraðferð &  &  &  \\
\hline
1041 & eyðingarrúm &  &  &  \\
\hline
1043 & fábrotinn &  &  &  \\
\hline
1044 & fábrotinn hnútur &  &  &  \\
\hline
1045 & fábrotleiki &  &  &  \\
\hline
1046 & fæðingaferli &  &  &  \\
\hline
1047 & fæðingar- og dánarferli &  &  &  \\
\hline
1048 & fæðingatíðni &  &  &  \\
\hline
1049 & færsla &  &  &  \\
\hline
1050 & færslufylki &  &  &  \\
\hline
1051 & færslugrúpa &  &  &  \\
\hline
1052 & fágað fall, tvinnfágað fall &  &  &  \\
\hline
1053 & fágað línubundin &  &  &  \\
\hline
1054 & fágað mengi &  &  &  \\
\hline
1055 & fágað mynstur, tvinnfágað mynstur &  &  &  \\
\hline
1056 & fágað rúm, tvinnfágað rúm &  &  &  \\
\hline
1057 & fágað víðerni &  &  &  \\
\hline
1058 & fágaður ferill &  &  &  \\
\hline
1059 & fágaður, tvinnfágaður &  &  &  \\
\hline
1060 & fáguð framlenging &  &  &  \\
\hline
1061 & fáguð rúmfræði &  &  &  \\
\hline
1062 & fáguð talnafræði, talnafræðigreining &  &  &  \\
\hline
1063 & fáguð víðátta, tvinnfáguð víðátta, tvinnvíðátta &  &  &  \\
\hline
1064 & fáguð vörpun &  &  &  \\
\hline
1065 & fágun &  &  &  \\
\hline
1066 & fágunaraðgreinanlegur &  &  &  \\
\hline
1067 & fágunarfullkominn &  &  &  \\
\hline
1068 & fágunarhjúpur &  &  &  \\
\hline
1069 & fágunarkúptur &  &  &  \\
\hline
1070 & fágunarsvæði &  &  &  \\
\hline
1071 & fágunarsvæði &  &  &  \\
\hline
1072 & faldfirð, firðamargfeldi &  &  &  \\
\hline
1073 & faldheildi, faldtegur &  &  &  \\
\hline
1074 & faldmál, margfeldismál &  &  &  \\
\hline
1075 & faldmengi, mengjamargfeldi, margfeldismengi & &  & \\
\hline
1076 & faldnet, netamargfeldi &  &  &  \\
\hline
1077 & faldrúm, margfeldisrúm &  &  &  \\
\hline
1078 & fall & <ref idTerm=3282>Vörpun</ref> sem hefur $\mathbb{R}$ eða $\mathbb{C}$ sem <ref idTerm=74>bakmengi</ref> &  &  \\
\hline
1079 & falla saman &  &  &  \\
\hline
1080 & fallaalgebra &  &  &  \\
\hline
1081 & fallajafna &  &  &  \\
\hline
1082 & fallandi keðja, minnkandi keðja &  &  &  \\
\hline
1083 & fallandi, minnkandi &  &  &  \\
\hline
1084 & fallaröð &  &  &  \\
\hline
1085 & fallarúm &  &  &  \\
\hline
1086 & fallaruna &  &  &  \\
\hline
1087 & fallasvið &  &  &  \\
\hline
1088 & fallfasti &  &  &  \\
\hline
1089 & fallfrymi &  &  &  \\
\hline
1090 & fallgildi & Fyrir <ref idTerm=1078>fall</ref> $f : X \to Y$ og <ref idTerm=5128>stak</ref> $x$ úr $X$ kallast stakið úr $Y$ sem $f$ úthlutar $x$ \textit{gildi fallsins} $f$ í $x$ og táknað með $f(x)$ &  &  \\
\hline
1091 & fallgildi &  &  &  \\
\hline
1092 & falljafna Eulers &  &  &  \\
\hline
1093 & fallörvanet &  &  &  \\
\hline
1094 & falltákn &  &  &  \\
\hline
1095 & fallvensl &  &  &  \\
\hline
1096 & fánavíðátta, flaggvíðátta &  &  &  \\
\hline
1097 & fáni, flagg &  &  &  \\
\hline
1098 & fánýti &  &  &  \\
\hline
1099 & farandsalaverkefni &  &  &  \\
\hline
1100 & fasaferill &  &  &  \\
\hline
1101 & fasameðaltal &  &  &  \\
\hline
1102 & fasamynd &  &  &  \\
\hline
1103 & fasarúm &  &  &  \\
\hline
1104 & fasi &  &  &  \\
\hline
1105 & fastafall, fasti, fast fall &  &  &  \\
\hline
1106 & fastakommureikningur &  &  &  \\
\hline
1107 & fastaliður &  &  &  \\
\hline
1108 & fastasvið &  &  &  \\
\hline
1109 & fastheldnihætta &  &  &  \\
\hline
1110 & fastheldnimistök &  &  &  \\
\hline
1111 & fasti Eulers &  &  &  \\
\hline
1112 & fasti, fastastærð &  &  &  \\
\hline
1113 & fastur &  &  &  \\
\hline
1114 & fátækleg tala &  &  &  \\
\hline
1115 & fella-, falla- &  &  &  \\
\hline
1116 & fellagreining &  &  &  \\
\hline
1117 & felli &  &  &  \\
\hline
1118 & felliafleiða &  &  &  \\
\hline
1119 & fellideildun, fellidiffrun &  &  &  \\
\hline
1120 & felliheildi, fellitegur &  &  &  \\
\hline
1121 & felliheildun, fellitegrun &  &  &  \\
\hline
1122 & felujafna &  &  &  \\
\hline
1123 & ferð &  &  &  \\
\hline
1124 & ferfaldur, fjórfaldur &  &  &  \\
\hline
1125 & ferflötungur, fjórflötungur, þrístrend strýta, þrístrendur píramíti &  &  &  \\
\hline
1126 & ferhliðungur, fjórhliðungur &  &  &  \\
\hline
1127 & ferhyrningur & <ref idTerm=3161>Marghyrningur</ref> með fjórar <ref idTerm=2048>hliðar</ref> &  & Er verið að moka of mikilli ábyrgð inn í skilgreininguna á marghyrningi? Er ferhyrningur nógu mikilvægt sértilvik að það ætti að koma bein skilgreining hér? \\
\hline
1128 & ferilaðlögun &  &  &  \\
\hline
1129 & ferilheildi, feriltegur &  &  &  \\
\hline
1130 & ferilhorn &  &  &  \\
\hline
1131 & ferill, boglína, sveiglína & <ref idTerm=3417>Myndmengi</ref> <ref idTerm=4359>samfellds falls</ref> á <ref idTerm=3010>lokuðu bili</ref> &  & Nota evklíðska skilgreiningu frekar en grannfræðilega? \\
\hline
1132 & ferill, vegur, leið &  &  &  \\
\hline
1133 & ferillyfting &  &  &  \\
\hline
1134 & ferilris, reisn &  &  &  \\
\hline
1135 & ferilsamanhangandi, vegsamanhangandi, ferilsamhangandi, vegsamhangandi &  &  &  \\
\hline
1136 & ferilsamhengisþáttur, vegsamhengisþáttur &  &  &  \\
\hline
1137 & ferlarúm, vegarúm, leiðarúm &  &  &  \\
\hline
1138 & ferli &  &  &  \\
\hline
1139 & ferlínulegur, fjórlínulegur &  &  &  \\
\hline
1140 & Fermat-frumtala &  &  &  \\
\hline
1141 & Fermat-vefja, fleygger vefja &  &  &  \\
\hline
1142 & fernd &  &  &  \\
\hline
1143 & ferninga &  &  &  \\
\hline
1144 & ferningsbestun &  &  &  \\
\hline
1145 & ferningseileif &  &  &  \\
\hline
1146 & ferningsflötur &  &  &  \\
\hline
1147 & ferningsform &  &  &  \\
\hline
1148 & ferningsfylki, ferningslaga fylki &  &  &  \\
\hline
1149 & ferningsgagnkvæmni &  &  &  \\
\hline
1150 & ferningsgagnkvæmnissetning, setning um ferningsgagnkvæmni &  &  &  \\
\hline
1151 & ferningsheildanlegur, ferningstegranlegur &  &  &  \\
\hline
1152 & ferningslaga &  &  &  \\
\hline
1153 & ferningslaus &  &  &  \\
\hline
1154 & ferningsleif &  &  &  \\
\hline
1155 & ferningsmeðaltal &  &  &  \\
\hline
1156 & ferningsmeðaltalsrót &  &  &  \\
\hline
1157 & ferningsrót, önnur rót & Ef $a$ er <ref idTerm=4093>rauntala</ref> þannig að $a \geq 0$ þá hefur <ref idTerm=2478>jafnan</ref> $x^2 = a$ nákvæmlega eina lausn $t \geq 0$ og þessi lausn er kölluð \textit{kvaðratrótin} af $a$ og táknuð með $\sqrt{a}$ &  & Á að bæta við vísun í $n$-tu rót? \\
\hline
1158 & ferningssumma frávika, ferningsfrávik &  &  &  \\
\hline
1159 & ferningssvið &  &  &  \\
\hline
1160 & ferningstala, ferhyrnd tala &  &  &  \\
\hline
1161 & ferningun &  &  &  \\
\hline
1162 & ferningun hrings &  &  &  \\
\hline
1163 & ferningur & <ref idTerm=4175>Rétthyrningur</ref> með allar <ref idTerm=2048>hliðar</ref> jafn <ref idTerm=2860>langar</ref> &  &  \\
\hline
1164 & ferstrendingur &  &  &  \\
\hline
1165 & fertala &  &  &  \\
\hline
1166 & ferverpill, fjórvíður teningur &  &  &  \\
\hline
1167 & fervik &  &  &  \\
\hline
1168 & fervikagreining &  &  &  \\
\hline
1169 & Fibonacci-runa &  &  &  \\
\hline
1170 & Fibonacci-tala &  &  &  \\
\hline
1171 & fimmfaldur &  &  &  \\
\hline
1172 & fimmflötungur &  &  &  \\
\hline
1173 & fimmhyrndur &  &  &  \\
\hline
1174 & fimmhyrningstala, fimmhyrnd tala &  &  &  \\
\hline
1175 & fimmhyrningur &  &  &  \\
\hline
1176 & fimmsetning &  &  &  \\
\hline
1177 & fimmstirningur, reglulegur fimmstirningur &  &  &  \\
\hline
1178 & fimmta stigs &  &  &  \\
\hline
1179 & fimmta stigs ferill &  &  &  \\
\hline
1180 & fimmta stigs flötur &  &  &  \\
\hline
1181 & fimmta stigs jafna &  &  &  \\
\hline
1182 & fimmtánhyrningur &  &  &  \\
\hline
1183 & fimmtungsmark &  &  &  \\
\hline
1184 & fimmund &  &  &  \\
\hline
1185 & fín upplausn &  &  &  \\
\hline
1186 & fíngerving &  &  &  \\
\hline
1187 & fingrareikningur &  &  &  \\
\hline
1188 & fínleiki &  &  &  \\
\hline
1189 & fínn &  &  &  \\
\hline
1190 & finna gildi (einhvers), reikna út &  &  &  \\
\hline
1191 & fínna grannmynstur, sterkara grannmynstur &  &  &  \\
\hline
1192 & fínni þakning &  &  &  \\
\hline
1193 & fínt knippi &  &  &  \\
\hline
1194 & firð & Firð á <ref idTerm=3270>mengi</ref> $M$ er <ref idTerm=1078>fall</ref> $d:M \times M \to \mathbb{R}_+$, sem uppfyllir 

\begin{itemize}
    \item $d(x,x)=0$ fyrir öll $x \in M$
    \item $d(x,y)>0$ fyrir öll $x \neq y \in M$
    \item $d(x,y)=d(y,x)$ fyrir öll $x,y \in M$
    \item $d(x,y) \leq d(x,z) + d(z,y)$ fyrir öll $x,y,z \in M$
\end{itemize} <ref idTerm=5102>Stærðin</ref> $d(x,y)$ er gjarnan kölluð <ref idTerm=1204>fjarlægðin</ref> milli $x$ og $y$ &  &  \\
\hline
1195 & firð- &  &  &  \\
\hline
1196 & firðanlegur &  &  &  \\
\hline
1197 & firðanleiki &  &  &  \\
\hline
1198 & firðargrannmynstur, firðgrannmynstur &  &  &  \\
\hline
1199 & firðrækinn &  &  &  \\
\hline
1200 & firðrúm & <ref idTerm=5627>Tvennd</ref> $(M,d)$ þar sem $M$ er <ref idTerm=3270>mengi</ref> og $d:M \times M \to \mathbb{R}_+$ er <ref idTerm=1194>firð</ref> &  &  \\
\hline
1201 & firðun &  &  &  \\
\hline
1202 & firðunarsetning &  &  &  \\
\hline
1203 & firðþinur &  &  &  \\
\hline
1204 & fjarlægð &  &  &  \\
\hline
1205 & fjarlægt innra horn & Þau <ref idTerm=6295>innri horn</ref> <ref idTerm=3161>marghyrnings</ref> sem ekki eru <ref idTerm=1677>grannhorn</ref> tiltekins ytra horns &  &  \\
\hline
1206 & fjarstæða &  &  &  \\
\hline
1207 & fjögraferningasetning, fjögurraferningasetning &  &  &  \\
\hline
1208 & fjögralitasetning, fjögurralitasetning &  &  &  \\
\hline
1209 & fjölbreytugreining &  &  &  \\
\hline
1210 & fjöldatala &  &  &  \\
\hline
1211 & fjöldatala & <ref idTerm=2489>Jafngildisflokkur</ref> <ref idTerm=3270>mengja</ref> þar sem <ref idTerm=3332>milli</ref> sérhverra tveggja <ref idTerm=3270>mengja</ref> er til <ref idTerm=1586>gagntæk vörpun</ref> &  &  \\
\hline
1212 & fjöldatalnareikningur, reikningur með fjöldatölum &  &  &  \\
\hline
1213 & fjöldi &  &  &  \\
\hline
1214 & fjöldi, fjöldatala, stétt &  &  &  \\
\hline
1215 & fjölfylgni &  &  &  \\
\hline
1216 & fjölfylgnistuðull &  &  &  \\
\hline
1217 & fjölgeisli &  &  &  \\
\hline
1218 & fjölhlutanet &  &  &  \\
\hline
1219 & fjölhnattaverkefni, fjölagnaverkefni &  &  &  \\
\hline
1220 & fjölkryppudreifing, fjöltindadreifing &  &  &  \\
\hline
1221 & fjölliða &  &  &  \\
\hline
1222 & fjölliðudreifing, fjölkostadreifing &  &  &  \\
\hline
1223 & fjölliðuregla &  &  &  \\
\hline
1224 & fjöllotubundinn, margfaldlega lotubundinn &  &  &  \\
\hline
1225 & fjölmættisfræði &  &  &  \\
\hline
1226 & fjölnet, net &  &  &  \\
\hline
1227 & fjölskífa &  &  &  \\
\hline
1228 & fjölskipting &  &  &  \\
\hline
1229 & fjölskylda &  &  &  \\
\hline
1230 & fjölstafa tala &  &  &  \\
\hline
1231 & fjölstig &  &  &  \\
\hline
1232 & fjölundirþýður &  &  &  \\
\hline
1233 & fjölvísir &  &  &  \\
\hline
1234 & fjölyfirþýður &  &  &  \\
\hline
1235 & fjölþýður &  &  &  \\
\hline
1236 & fjórða rót &  &  &  \\
\hline
1237 & fjórða stigs &  &  &  \\
\hline
1238 & fjórða stigs ferill &  &  &  \\
\hline
1239 & fjórða stigs flötur &  &  &  \\
\hline
1240 & fjórða stigs jafna &  &  &  \\
\hline
1241 & fjórðungsbogi &  &  &  \\
\hline
1242 & fjórðungsgeiri &  &  &  \\
\hline
1243 & fjórðungsmark &  &  &  \\
\hline
1244 & fjórðungur &  &  &  \\
\hline
1245 & fjórflötungsgrúpa &  &  &  \\
\hline
1246 & fjórflötungshnit &  &  &  \\
\hline
1247 & fjórflötungstala &  &  &  \\
\hline
1248 & fjórgrúpa Kleins &  &  &  \\
\hline
1249 & fjórhliða rúmhorn &  &  &  \\
\hline
1250 & fjórlaufaferill, fjögralaufaferill, fjögralaufasmári &  &  &  \\
\hline
1251 & fjórsetning &  &  &  \\
\hline
1252 & fjórskipting, fjórhlutun &  &  &  \\
\hline
1253 & fjórveldi, fjórða veldi &  &  &  \\
\hline
1254 & fjórveldi, fjórða veldi (af heilli tölu) &  &  &  \\
\hline
1255 & fjórvigur, fjórvíður vigur &  &  &  \\
\hline
1256 & flæði &  &  &  \\
\hline
1257 & flæðihraði &  &  &  \\
\hline
1258 & flæðinet &  &  &  \\
\hline
1259 & flæðirit &  &  &  \\
\hline
1260 & flæðistærð &  &  &  \\
\hline
1261 & flæðisþéttleiki &  &  &  \\
\hline
1262 & flækja, flækjustig &  &  &  \\
\hline
1263 & flatarheildi, flatartegur &  &  &  \\
\hline
1264 & flatarmælir &  &  &  \\
\hline
1265 & flatarmál &  &  &  \\
\hline
1266 & flatarmál, yfirborðsmál, yfirborð &  &  &  \\
\hline
1267 & flatarmálsform &  &  &  \\
\hline
1268 & flatarmálsfrymi &  &  &  \\
\hline
1269 & flatarrækinn &  &  &  \\
\hline
1270 & flatarrækni &  &  &  \\
\hline
1271 & flatarskiki &  &  &  \\
\hline
1272 & flatarskiki &  &  &  \\
\hline
1273 & flatneskja &  &  &  \\
\hline
1274 & flattur sporvölusnúðflötur &  &  &  \\
\hline
1275 & flattur sporvölusnúður &  &  &  \\
\hline
1276 & flatur &  &  &  \\
\hline
1277 & flatur punktur, flatpunktur &  &  &  \\
\hline
1278 & flaumrænn &  &  &  \\
\hline
1279 & flaumtölva, flaumræn tölva &  &  &  \\
\hline
1280 & fléttingur &  &  &  \\
\hline
1281 & fléttufræði, samtakagreining, samtektagreining &  &  &  \\
\hline
1282 & fleygaður &  &  &  \\
\hline
1283 & fleygaður frá hægri, nákvæmur frá hægri &  &  &  \\
\hline
1284 & fleygaður frá vinstri, nákvæmur frá vinstri &  &  &  \\
\hline
1285 & fleygaður varpi &  &  &  \\
\hline
1286 & fleygbogaflötur &  &  &  \\
\hline
1287 & fleygbogariti &  &  &  \\
\hline
1288 & fleygbogasneið &  &  &  \\
\hline
1289 & fleygbogasnúðflötur &  &  &  \\
\hline
1290 & fleygbogi &  &  &  \\
\hline
1291 & fleygger &  &  &  \\
\hline
1292 & fleygger hlutafleiðujafna &  &  &  \\
\hline
1293 & fleygger hnit &  &  &  \\
\hline
1294 & fleygger punktur &  &  &  \\
\hline
1295 & fleygger Riemannflötur &  &  &  \\
\hline
1296 & fleygger rúmfræði &  &  &  \\
\hline
1297 & fleygger sívalningur &  &  &  \\
\hline
1298 & fleygger sjálfmótun, fleygger færsla &  &  &  \\
\hline
1299 & fleygger slétta &  &  &  \\
\hline
1300 & fleyggerð &  &  &  \\
\hline
1301 & fleyggerð &  &  &  \\
\hline
1302 & fleyggert aðhvarf &  &  &  \\
\hline
1303 & fleyguð lest &  &  &  \\
\hline
1304 & flísalagning &  &  &  \\
\hline
1305 & flokkafjöldi &  &  &  \\
\hline
1306 & flokkatala &  &  &  \\
\hline
1307 & flokksbreidd, flokkslengd &  &  &  \\
\hline
1308 & flokksmark &  &  &  \\
\hline
1309 & flokksmeðaltal &  &  &  \\
\hline
1310 & flokksmiðja, miðpunktur flokkunarbils &  &  &  \\
\hline
1311 & flokksnúmer &  &  &  \\
\hline
1312 & flokkuð gögn &  &  &  \\
\hline
1313 & flokkun &  &  &  \\
\hline
1314 & flokkun &  &  &  \\
\hline
1315 & flokkunarbil &  &  &  \\
\hline
1316 & flokkunarbreyta &  &  &  \\
\hline
1317 & flokkunargögn &  &  &  \\
\hline
1318 & flokkunarhending &  &  &  \\
\hline
1319 & flokkunarrúm &  &  &  \\
\hline
1320 & flokkunarsetning &  &  &  \\
\hline
1321 & flökkupunktur, hlaupapunktur, breytilegur punktur &  &  &  \\
\hline
1322 & flokkur &  &  &  \\
\hline
1323 & flokkur &  &  &  \\
\hline
1324 & flötur &  &  &  \\
\hline
1325 & flötur, hliðarflötur, hlið & <ref idTerm=3161>Marghyrningur</ref> sem myndar <ref idTerm=2089>hluta</ref> <ref idTerm=6256>yfirborðs</ref> <ref idTerm=3133>margflötungs</ref> &  &  \\
\hline
1326 & flutningagrúpa, hreyfingagrúpa &  &  &  \\
\hline
1327 & flutningsíðal &  &  &  \\
\hline
1328 & flutningur, hreyfing & <ref idTerm=5779>Færsla</ref> á tiltekinni <ref idTerm=4884>sléttu</ref> sem <ref idTerm=5963>varðveitir</ref> <ref idTerm=1204>fjarlægðir</ref> &  & Skilgreiningin í hugtakabankanum er svoldið ókjarnyrt: "Flutningur er færsla á tiltekinni sléttu sem hefur þann eiginleika að þegar strik er fært, þá er færða strikið jafn langt og upphaflega strikið. Flutningar breyta því ekki lengdum strika þegar þeir færa þau til." \\
\hline
1329 & flytja &  &  &  \\
\hline
1330 & földun &  &  &  \\
\hline
1331 & földunarjafna &  &  &  \\
\hline
1332 & fólgið fall &  &  &  \\
\hline
1333 & fólginn &  &  &  \\
\hline
1334 & for-Hilbert-rúm &  &  &  \\
\hline
1335 & forgrunnur &  &  &  \\
\hline
1336 & forhlutur, fráhlutur, óðal &  &  &  \\
\hline
1337 & forknippi &  &  &  \\
\hline
1338 & forkönnun &  &  &  \\
\hline
1339 & forliður &  &  &  \\
\hline
1340 & form &  &  &  \\
\hline
1341 & form, jafnþætt margliða, einsleit margliða &  &  &  \\
\hline
1343 & formengi, frámengi, skilgreiningarmengi, óðal & Sjá <ref idTerm=3282>vörpun</ref> &  &  \\
\hline
1344 & formerki &  &  &  \\
\hline
1346 & formerkisfall &  &  &  \\
\hline
1347 & formerkispróf &  &  &  \\
\hline
1348 & formerkjaregla &  &  &  \\
\hline
1349 & formerkjaskipti &  &  &  \\
\hline
1350 & formerkjatala &  &  &  \\
\hline
1351 & formgera &  &  &  \\
\hline
1352 & formhyggja &  &  &  \\
\hline
1353 & formleg afleiða &  &  &  \\
\hline
1354 & formleg margliða &  &  &  \\
\hline
1355 & formleg rökfræði &  &  &  \\
\hline
1356 & formleg sönnun &  &  &  \\
\hline
1357 & formleg veldaröð &  &  &  \\
\hline
1358 & formlega samkvæmur &  &  &  \\
\hline
1359 & formlegt kerfi, formkerfi &  &  &  \\
\hline
1360 & formlegt mál &  &  &  \\
\hline
1361 & formlegt raunsvið &  &  &  \\
\hline
1362 & formlegur &  &  &  \\
\hline
1363 & formúla, formkorn, formáli &  &  &  \\
\hline
1364 & formúlusafn &  &  &  \\
\hline
1365 & formynd &  &  &  \\
\hline
1366 & formynd, frummynd, öfug mynd & Fyrir <ref idTerm=3282>vörpun</ref> $f: X \to Y$ kallast <ref idTerm=3270>mengi</ref> þeirra <ref idTerm=5128>staka</ref> úr $X$ sem $f$ varpar í stak úr hlutmengi $B$ í $Y$ \textit{frummynd} vörpunarinnar $f$ af menginu $B$. Hana má rita á forminu \[ f^{-1}(B) = \{x \in X \mid f(x) \in B\}. \] &  &  \\
\hline
1367 & forraðað mengi &  &  &  \\
\hline
1368 & forrit &  &  &  \\
\hline
1369 & forrit, tölvuforrit &  &  &  \\
\hline
1370 & forritun &  &  &  \\
\hline
1371 & forröðun &  &  &  \\
\hline
1372 & forsagnar- og leiðréttingaraðferð &  &  &  \\
\hline
1373 & forsegjandi, spáfall &  &  &  \\
\hline
1374 & forsenda &  &  &  \\
\hline
1375 & forsenda &  &  &  \\
\hline
1376 & forsenda til einföldunar &  &  &  \\
\hline
1377 & forsenda, krafa &  &  &  \\
\hline
1378 & forsía &  &  &  \\
\hline
1379 & forskema &  &  &  \\
\hline
1380 & forskeyti &  &  &  \\
\hline
1381 & fortáknun, pólskur ritháttur &  &  &  \\
\hline
1382 & fortilraun &  &  &  \\
\hline
1383 & forystuliður &  &  &  \\
\hline
1384 & forystustuðull &  &  &  \\
\hline
1385 & forþjappaður &  &  &  \\
\hline
1386 & fossaferli, skriðuferli &  &  &  \\
\hline
1387 & föst færslulíkindi, stöðug færslulíkindi &  &  &  \\
\hline
1388 & fótpunktaþríhyrningur &  &  &  \\
\hline
1389 & fótpunktur &  &  &  \\
\hline
1390 & Fourier-, sveiflu- &  &  &  \\
\hline
1391 & Fourier-færsla, Fourier-myndun, Fourier-ummyndun &  &  &  \\
\hline
1392 & Fourier-greining, sveiflugreining &  &  &  \\
\hline
1393 & Fourier-heildi, Fourier-tegur &  &  &  \\
\hline
1394 & Fourier-liðun, sveifluliðun &  &  &  \\
\hline
1395 & Fourier-mynd, Fourier-færsla &  &  &  \\
\hline
1396 & Fourier-röð &  &  &  \\
\hline
1397 & Fourier-stuðull &  &  &  \\
\hline
1398 & frádragi &  &  &  \\
\hline
1399 & frádráttarstofn &  &  &  \\
\hline
1400 & frádráttur &  &  &  \\
\hline
1401 & fræði &  &  &  \\
\hline
1402 & fræðin um keilusnið &  &  &  \\
\hline
1403 & frændhorn & Tvö <ref idTerm=2198>horn</ref> sem samanlagt <ref idTerm=2223>mælast</ref> 180 <ref idTerm=1667>gráður</ref> &  & Skýrari skilgreining en sú á stæ.is?: "Ef hægt er að flytja horn þannig að það verði grannhorn annars horns, þá er sagt að hornin séu frændhorn." \\
\hline
1404 & framandlegt hvel &  &  &  \\
\hline
1405 & framfærsla &  &  &  \\
\hline
1406 & framinnsetning, forinnsetning &  &  &  \\
\hline
1407 & framleiða &  &  &  \\
\hline
1408 & framleiðandi fall &  &  &  \\
\hline
1409 & framleiðandi lína, framleiðandi &  &  &  \\
\hline
1410 & framleiðandi veldaröð &  &  &  \\
\hline
1411 & framlenging, útvíkkun &  &  &  \\
\hline
1412 & frammunur &  &  &  \\
\hline
1413 & framsetjandi &  &  &  \\
\hline
1414 & framsetjanlegur &  &  &  \\
\hline
1415 & framsetjanleiki &  &  &  \\
\hline
1416 & framsetning &  &  &  \\
\hline
1417 & framsetning &  &  &  \\
\hline
1418 & framsetningarregla, framsetningarsetning &  &  &  \\
\hline
1419 & framvinda, tilhneiging &  &  &  \\
\hline
1420 & framvindill, útbreiðslufall &  &  &  \\
\hline
1421 & framvindu- &  &  &  \\
\hline
1422 & framvindujafna &  &  &  \\
\hline
1423 & frávik &  &  &  \\
\hline
1424 & frávikshlutfall &  &  &  \\
\hline
1425 & Fréchet-deildanlegur, Fréchet-diffranlegur &  &  &  \\
\hline
1426 & Fréchet-rúm &  &  &  \\
\hline
1427 & Fredholm-skilyrði &  &  &  \\
\hline
1428 & frekar jákvæður deilir &  &  &  \\
\hline
1429 & fremsti tölustafur &  &  &  \\
\hline
1430 & fremsti tölustafur &  &  &  \\
\hline
1431 & frjáls &  &  &  \\
\hline
1432 & frjáls breyta, frjáls breytistærð &  &  &  \\
\hline
1433 & frjáls grunnur &  &  &  \\
\hline
1434 & frjáls grúpa &  &  &  \\
\hline
1435 & frjáls Riemann-flötur &  &  &  \\
\hline
1436 & frjáls spönnuður &  &  &  \\
\hline
1437 & frjáls útvíkkun, frjáls víkkun &  &  &  \\
\hline
1438 & frjáls vigur &  &  &  \\
\hline
1439 & frjáls vísir &  &  &  \\
\hline
1440 & frjáls þáttur &  &  &  \\
\hline
1441 & frjálst margfeldi &  &  &  \\
\hline
1442 & frjálst stak &  &  &  \\
\hline
1443 & frum- &  &  &  \\
\hline
1444 & fruma &  &  &  \\
\hline
1445 & frumbreyta, óháð breyta, óháð breytistærð, frumstærð &  &  &  \\
\hline
1446 & frumbútur (í rúmfræði) &  &  &  \\
\hline
1447 & frumdeilir &  &  &  \\
\hline
1448 & frumhugtak, grunnhugtak, undirstöðuhugtak &  &  &  \\
\hline
1449 & frumíðal &  &  &  \\
\hline
1450 & frummargliða &  &  &  \\
\hline
1451 & frumsenda um lítil mengi, frumsetning um lítil mengi &  &  &  \\
\hline
1452 & frumsenda, frumsetning &  &  &  \\
\hline
1453 & frumsenduaðferð, frumsetningaraðferð &  &  &  \\
\hline
1454 & frumsendugrip, frumsetningargrip, frumsenduskema, frumsetningaskema &  &  &  \\
\hline
1455 & frumsendukerfi, frumsetningakerfi &  &  &  \\
\hline
1456 & frumsenduleg mengjafræði, frumsendumengjafræði &  &  &  \\
\hline
1457 & frumsendulegur, frumsendu-, frumsetninga- &  &  &  \\
\hline
1458 & frumstæð margliða &  &  &  \\
\hline
1459 & frumstæð rót &  &  &  \\
\hline
1460 & frumstæð sviðsútvíkkun, frumstæð sviðsvíkkun &  &  &  \\
\hline
1461 & frumstæður &  &  &  \\
\hline
1462 & frumstætt íðal &  &  &  \\
\hline
1463 & frumstætt rakið fall &  &  &  \\
\hline
1464 & frumstætt stak &  &  &  \\
\hline
1465 & frumstak &  &  &  \\
\hline
1467 & frumsvið, frumkroppur &  &  &  \\
\hline
1468 & frumtala, prímtala & <ref idTerm=3475>Náttúruleg tala</ref> $p > 1$ kallast \textit{frumtala} ef einu náttúrulegu tölurnar sem <ref idTerm=1599>ganga upp í</ref> henni eru <ref idTerm=5397>talan</ref> $1$ og talan $p$ &  &  \\
\hline
1469 & frumtalnasetning &  &  &  \\
\hline
1470 & frumtalnatafla &  &  &  \\
\hline
1471 & frumtalnatvíburar &  &  &  \\
\hline
1472 & frumu- &  &  &  \\
\hline
1473 & frumuklasi &  &  &  \\
\hline
1474 & frumþáttun &  &  &  \\
\hline
1475 & frumþáttur & Sjá <ref idTerm=5824>undirstöðusetningu reikningslistarinnar</ref> &  &  \\
\hline
1476 & frymi &  &  &  \\
\hline
1477 & frymið íðal &  &  &  \\
\hline
1478 & frymin framsetning &  &  &  \\
\hline
1479 & frymin þáttun &  &  &  \\
\hline
1480 & fryminn þáttur &  &  &  \\
\hline
1481 & fullgerð lokun &  &  &  \\
\hline
1482 & fullgerð útvíkkun, heilsteypt útvíkkun, fullgerð víkkun, heilsteypt víkkun &  &  &  \\
\hline
1483 & fullgert svið, heilsteypt svið &  &  &  \\
\hline
1484 & fullheildanlegur, fulltegranlegur &  &  &  \\
\hline
1485 & fullkomið firðrúm &  &  &  \\
\hline
1486 & fullkomið frumsendukerfi, fullkomið frumsetningakerfi &  &  &  \\
\hline
1487 & fullkomið leifamengi, fullkomið leifakerfi &  &  &  \\
\hline
1488 & fullkomið mál &  &  &  \\
\hline
1489 & fullkomið raðsvið, fullkomið raðað svið &  &  &  \\
\hline
1490 & fullkomin kenning &  &  &  \\
\hline
1491 & fullkomin spyrðing &  &  &  \\
\hline
1492 & fullkomin tala &  &  &  \\
\hline
1493 & fullkominn &  &  &  \\
\hline
1494 & fullkominn ferhliðungur &  &  &  \\
\hline
1495 & fullkominn ferhyrningur &  &  &  \\
\hline
1496 & fullkomlega samfelldur &  &  &  \\
\hline
1497 & fullkomlega viðbætinn, fullkomlega samleggjandi &  &  &  \\
\hline
1498 & fullkomleikafrumsenda, fullkomleikafrumsetning &  &  &  \\
\hline
1499 & fullkomleikafrumsenda, fullkomleikafrumsetning, frumsenda um fullkomleika, frumsetning um fullkomleika &  &  &  \\
\hline
1500 & fullkomleikasetning &  &  &  \\
\hline
1501 & fullkomleikasetning Gödels &  &  &  \\
\hline
1502 & fullkomleiki &  &  &  \\
\hline
1503 & fullkomnun &  &  &  \\
\hline
1504 & fullnægja, uppfylla &  &  &  \\
\hline
1505 & fullnægjanlegur, uppfyllanlegur &  &  &  \\
\hline
1506 & fullskipað hlutnet &  &  &  \\
\hline
1507 & fullstytt brot &  &  &  \\
\hline
1508 & fullt hlutríki &  &  &  \\
\hline
1509 & fullt horn, \(360^\circ\) horn &  &  &  \\
\hline
1510 & fullt rúmhorn &  &  &  \\
\hline
1511 & fulltengt net &  &  &  \\
\hline
1512 & fulltrúa- &  &  &  \\
\hline
1513 & fulltrúi &  &  &  \\
\hline
1514 & fullur &  &  &  \\
\hline
1515 & fullyrðing, yrðing & <ref idTerm=5072>Staðhæfing</ref> sem er annaðhvort <ref idTerm=4565>sönn</ref> eða <ref idTerm=3890>ósönn</ref> &  &  \\
\hline
1516 & fylgibreyta &  &  &  \\
\hline
1517 & fylgibreyta, fylgistærð, háð breyta, háð breytistærð &  &  &  \\
\hline
1518 & fylgibreyta, fylgistærð, háð breyta, háð breytistærð &  &  &  \\
\hline
1519 & fylginn &  &  &  \\
\hline
1520 & fylgisetning &  &  &  \\
\hline
1521 & fylgiþáttur &  &  &  \\
\hline
1522 & fylgnar breytur &  &  &  \\
\hline
1523 & fylgni &  &  &  \\
\hline
1524 & fylgnifall &  &  &  \\
\hline
1525 & fylgnifylki &  &  &  \\
\hline
1526 & fylgnigreining &  &  &  \\
\hline
1527 & fylgnihlutfall &  &  &  \\
\hline
1528 & fylgnireikningur &  &  &  \\
\hline
1529 & fylgnirit &  &  &  \\
\hline
1530 & fylgnisporbaugur &  &  &  \\
\hline
1531 & fylgnisporvala &  &  &  \\
\hline
1532 & fylgnistuðull &  &  &  \\
\hline
1533 & fylgnitafla &  &  &  \\
\hline
1534 & fylki &  &  &  \\
\hline
1535 & fylkja- &  &  &  \\
\hline
1536 & fylkjaalgebra &  &  &  \\
\hline
1537 & fylkjafall &  &  &  \\
\hline
1538 & fylkjafræði &  &  &  \\
\hline
1539 & fylkjagildur &  &  &  \\
\hline
1540 & fylkjagrúpa &  &  &  \\
\hline
1541 & fylkjajafna &  &  &  \\
\hline
1542 & fylkjamargfeldi &  &  &  \\
\hline
1543 & fylkjareikningur &  &  &  \\
\hline
1544 & fylkjaritháttur &  &  &  \\
\hline
1545 & fylkjastaðall &  &  &  \\
\hline
1546 & fylkjaumhverfing, fylkjaandhverfing &  &  &  \\
\hline
1547 & fylkjaútsetning &  &  &  \\
\hline
1548 & fylkjungur &  &  &  \\
\hline
1549 & fylliaðgerð &  &  &  \\
\hline
1550 & fyllilega óbreyttur &  &  &  \\
\hline
1551 & fyllilega tryggur &  &  &  \\
\hline
1552 & fyllilíkindi &  &  &  \\
\hline
1553 & fyllimengi & <ref idTerm=3280>Mengjamunur</ref> á forminu $G \setminus A$ þar sem $A \subseteq G$ kallast \textit{fyllimengi} $A$ í $G$. Ef enginn vafi leikur á hvert mengið $G$ sé og við höfum ekki áhuga á því sem er utan $G$ (sem kallast þá \textit{grunnmengi}) köllum við þennan mengjamun einfaldlega fyllimengi $A$ og taknum hann með $A^c$ &  & Koma fyrir í skýringu: Oft er það svo að öll mengi, sem verið er að skoða í einhverju sambandi, eru hlutmengi af einu og sama menginu $G$. Þá getur verið gagnlegt að líta svo á að mengið $G$ innihaldi allt sem við höfum áhuga á að fjalla um og að það sem sé utan $G$ varði okkur engu. Þegar svo er gert er mengið $G$ stundum kallað grunnmengi.

Við þessar aðstæður þarf oft að fjalla um mengjamun af taginu $G \setminus A$, þar sem $A$ er hlutmengi í $G$. Ef enginn vafi leikur á hvert mengið $G$ sé, þá er þessi mengjamunur kallaður fyllimengi mengisins $A$ og hann er táknaður með $A^c$. Ýmis önnur tákn þekkjast fyrir fyllimengi $A$; til dæmis $\overline{A}$, $A’$ og $\mathbf{C}A$. Við höfum sem sagt: \[ A^c = G \setminus A. \] Óformlega má segja að fyllimengið $A^c$ sé sá hluti grunnmengisins sem liggur utan mengisins $A$. Á Venn-myndum er venja að tákna grunnmengið með rétthyrningi og á myndinni að neðan er fyllimengið $A^c$ litað rautt. (ásamt mynd)

Um fyllimengi sam- og sniðmengja gilda reglur De Morgans, sem segja að fyrir sérhver tvö mengi $A$ og $B$ gildi: \[ (A \cup B)^c = A^c \cap B^c \quad \text{og} \quad (A \cap B)^c = A^c \cup B^c. \] Jafnframt gilda eftirfarandi reglur fyrir mengin $A$ og $B$:

$A \cap A^c = \varnothing$.
$A \cup A^c = G$.
$(A^c)^c = A$.
$A \subset B$ ef og aðeins ef $B^c \subset A^c$. - Minnast líkar á aðrar leiðir til að tákna þetta \\
\hline
1554 & fylling &  &  &  \\
\hline
1555 & fyllingarskilyrði &  &  &  \\
\hline
1556 & fyllirúm &  &  &  \\
\hline
1557 & fyllitala, tölufylling &  &  &  \\
\hline
1558 & fyrir fram &  &  &  \\
\hline
1559 & fyrirframdreifing &  &  &  \\
\hline
1560 & fyrirframlíkindi &  &  &  \\
\hline
1561 & fyrirframmat &  &  &  \\
\hline
1562 & fyrri teljanleikafrumsendan &  &  &  \\
\hline
1563 & fyrsta meginform, fyrsta undirstöðuform &  &  &  \\
\hline
1564 & fyrsta stigs &  &  &  \\
\hline
1565 & fyrstu raðar &  &  &  \\
\hline
1566 & fyrstu stéttar &  &  &  \\
\hline
1567 & gæðaeftirlit &  &  &  \\
\hline
1568 & gagnabanki &  &  &  \\
\hline
1569 & gagnasafn, gagnagrunnur &  &  &  \\
\hline
1570 & gagnasafnsfræði &  &  &  \\
\hline
1571 & gagnavinnsla &  &  &  \\
\hline
1572 & gagnfágaður &  &  &  \\
\hline
1573 & gagnkvæmnissetning &  &  &  \\
\hline
1574 & gagnkvæmur &  &  &  \\
\hline
1575 & gagnræður &  &  &  \\
\hline
1576 & gagnsnið &  &  &  \\
\hline
1577 & gagnsniða &  &  &  \\
\hline
1578 & gagnsníða &  &  &  \\
\hline
1579 & gagnsniðill &  &  &  \\
\hline
1580 & gagnsniðsskilyrði &  &  &  \\
\hline
1581 & gagnstæð stefna, gagnstæð átt &  &  &  \\
\hline
1582 & gagnstæð tala &  &  &  \\
\hline
1583 & gagnstæður &  &  &  \\
\hline
1584 & gagnstætt formerki &  &  &  \\
\hline
1585 & gagntæk samsvörun, gagnkvæm samsvörun &  &  &  \\
\hline
1586 & gagntæk vörpun & <ref idTerm=3282>Vörpun</ref> sem er bæði <ref idTerm=965>eintæk</ref> og <ref idTerm=351>átæk</ref> &  &  \\
\hline
1587 & gagntækur &  &  &  \\
\hline
1588 & gagntekt, gagntaka &  &  &  \\
\hline
1589 & gagntilgáta &  &  &  \\
\hline
1590 & gagnvega- &  &  &  \\
\hline
1591 & gagnvegafjarlægð &  &  &  \\
\hline
1592 & gagnvegafullkominn &  &  &  \\
\hline
1593 & gagnvegahnit &  &  &  \\
\hline
1594 & gagnvegastreymi &  &  &  \\
\hline
1598 & gammafall &  &  &  \\
\hline
1599 & ganga upp í & Sagt er að <ref idTerm=3475>náttúruleg tala</ref> $n$ \textit{gengur upp í} náttúrulegri tölu $m$ ef hægt er að <ref idTerm=685>deila $m$ með $n$ án afgangs</ref> &  &  \\
\hline
1600 & gangbestun, kvik bestun, hreyfibestun, hreyfin bestun &  &  &  \\
\hline
1601 & gangkerfi, hreyfikerfi &  &  &  \\
\hline
1602 & gataður &  &  &  \\
\hline
1603 & Gauss-dreifing, normleg dreifing &  &  &  \\
\hline
1604 & Gauss-dreifingarfall, normlegt dreifingarfall, normlegt dreififall &  &  &  \\
\hline
1605 & Gauss-ferill, normlegur dreififerill &  &  &  \\
\hline
1606 & Gauss-ferli, normlegt ferli &  &  &  \\
\hline
1607 & Gauss-krappi, Gauss-sveigja, heildarkrappi, heildarsveigja &  &  &  \\
\hline
1608 & Gauss-Markov-metill &  &  &  \\
\hline
1609 & Gauss-tala, Gauss-heiltala &  &  &  \\
\hline
1610 & gegnvirk grúpa &  &  &  \\
\hline
1611 & gegnvirkni &  &  &  \\
\hline
1612 & gegnvirkniregla, gegnvirkniháttur &  &  &  \\
\hline
1613 & gegnvirkniregla, gegnvirknilögmál &  &  &  \\
\hline
1614 & gegnvirkur &  &  &  \\
\hline
1615 & geirahorn &  &  &  \\
\hline
1616 & geiri, hringgeiri &  &  &  \\
\hline
1617 & geislahnit &  &  &  \\
\hline
1618 & geislamarkgildi &  &  &  \\
\hline
1619 & geisli, hringgeisli & Sjá <ref idTerm=2295>hringur</ref> &  &  \\
\hline
1620 & geisli, kúlugeisli &  &  &  \\
\hline
1621 & gera ráð fyrir, setja sem svo, ganga út frá, gefa sér &  &  &  \\
\hline
1622 & gerð fylkis &  &  &  \\
\hline
1623 & gerð, tag &  &  &  \\
\hline
1624 & gerðarfræði &  &  &  \\
\hline
1625 & gerðarsetning &  &  &  \\
\hline
1626 & gerðarstöðugleiki &  &  &  \\
\hline
1627 & gervi- &  &  &  \\
\hline
1628 & gervihvel &  &  &  \\
\hline
1629 & gervikúptur &  &  &  \\
\hline
1630 & gervimál &  &  &  \\
\hline
1631 & gervislembitölur &  &  &  \\
\hline
1632 & geyma &  &  &  \\
\hline
1633 & geyma &  &  &  \\
\hline
1634 & geymd &  &  &  \\
\hline
1635 & geymd, varðveisla &  &  &  \\
\hline
1636 & geymdur tölustafur &  &  &  \\
\hline
1637 & giftingarsetning &  &  &  \\
\hline
1638 & giftingarverkefni &  &  &  \\
\hline
1639 & gildamengi &  &  &  \\
\hline
1640 & gildi, mynd (staks) &  &  &  \\
\hline
1641 & gildisdreifing &  &  &  \\
\hline
1642 & gildisdreifingarfræði &  &  &  \\
\hline
1643 & gildleiki &  &  &  \\
\hline
1644 & gjaldgeng lausn &  &  &  \\
\hline
1645 & gjaldþrot, hrun &  &  &  \\
\hline
1646 & gleiðhyrndur, sljóhyrndur & <ref idTerm=6438>Þríhyrningur</ref> með eitt <ref idTerm=1647>gleitt</ref> <ref idTerm=2198>horn</ref> &  &  \\
\hline
1647 & gleiður, sljór & <ref idTerm=2198>Horn</ref> kallast \textit{gleitt} ef það <ref idTerm=2223>mælist</ref> stærra en <ref idTerm=4167>rétt horn</ref> &  &  \\
\hline
1648 & gleitt horn, gleiðhorn, sljótt horn, sljóhorn &  &  &  \\
\hline
1649 & gleymskuvarpi &  &  &  \\
\hline
1650 & gleypið ástand &  &  &  \\
\hline
1651 & gleypið mengi &  &  &  \\
\hline
1652 & gleypilíkindi &  &  &  \\
\hline
1653 & gleyping &  &  &  \\
\hline
1654 & gleypinn &  &  &  \\
\hline
1655 & gleypinn tálmi &  &  &  \\
\hline
1656 & gleypiregla &  &  &  \\
\hline
1657 & gloppa &  &  &  \\
\hline
1658 & gloppa, eyða &  &  &  \\
\hline
1659 & gloppuröð, eyðuröð &  &  &  \\
\hline
1660 & gloppusetning, eyðusetning &  &  &  \\
\hline
1661 & gná &  &  &  \\
\hline
1662 & gnæfur &  &  &  \\
\hline
1663 & Gödel-tölusetning &  &  &  \\
\hline
1664 & gögn &  &  &  \\
\hline
1665 & götuð hringskífa &  &  &  \\
\hline
1666 & götuð tvinntalnaslétta &  &  &  \\
\hline
1667 & gráða, horngráða, bogagráða & <ref idTerm=3072>Mælieining</ref> sem oft er notuð til að mæla stærð <ref idTerm=2276>hringboga</ref> og <ref idTerm=2198>horna</ref> sem er fengin með því að skipta boga hrings í 360 eins boga sem kallast \textit{gráður} og eru táknaðar með $1^\circ$ &  &  \\
\hline
1668 & gráðubogi &  &  &  \\
\hline
1669 & grannalgebra &  &  &  \\
\hline
1670 & grannbaugur &  &  &  \\
\hline
1671 & grannfræði &  &  &  \\
\hline
1672 & grannfræðileg óbreyta &  &  &  \\
\hline
1673 & grannfræðilegur, grann- &  &  &  \\
\hline
1674 & grannfræðingur &  &  &  \\
\hline
1675 & granngrúpa &  &  &  \\
\hline
1676 & grannhnútur, grannstæður hnútur &  &  &  \\
\hline
1677 & grannhorn, aðlæg frændhorn & Ef tvö <ref idTerm=2198>horn</ref> hafa einn sameiginlegan <ref idTerm=334>arm</ref> og hinir armarnir eru <ref idTerm=3400>gagnstæðar</ref> <reg idTerm=1840>hálflínur</ref>, þá er sagt að hornin séu \textit{grannhorn} &  &  \\
\hline
1678 & grannhvel &  &  &  \\
\hline
1679 & grannleggur, grannstæður leggur &  &  &  \\
\hline
1680 & grannmóta &  &  &  \\
\hline
1681 & grannmótun &  &  &  \\
\hline
1682 & grannmynstur &  &  &  \\
\hline
1683 & grannnet &  &  &  \\
\hline
1684 & grannör, grannstæð ör &  &  &  \\
\hline
1685 & grannrúm &  &  &  \\
\hline
1686 & grannstæður hornpunktur &  &  &  \\
\hline
1687 & grannsvið &  &  &  \\
\hline
1688 & grannvíðátta &  &  &  \\
\hline
1689 & Grassmann-margfeldi, Grassmann-feldi, ytra margfeldi, útfeldi &  &  &  \\
\hline
1690 & Grassmann-víðátta &  &  &  \\
\hline
1691 & grein breiðboga, grein gleiðboga &  &  &  \\
\hline
1692 & greinargott mengi, Súslín-mengi &  &  &  \\
\hline
1693 & greining, stærðfræðigreining, stærðfræðileg greining &  &  &  \\
\hline
1694 & grennd &  &  &  \\
\hline
1695 & grenndagrunnur, grunnur af grenndum &  &  &  \\
\hline
1696 & grenndakerfi &  &  &  \\
\hline
1697 & grenndarinndragi &  &  &  \\
\hline
1698 & grenndasía &  &  &  \\
\hline
1699 & grennsl &  &  &  \\
\hline
1700 & grennslafylki, hnútafylki &  &  &  \\
\hline
1701 & greyping &  &  &  \\
\hline
1702 & greypt frymið íðal &  &  &  \\
\hline
1703 & greyptur &  &  &  \\
\hline
1704 & grind &  &  &  \\
\hline
1705 & grind &  &  &  \\
\hline
1706 & grindafræði &  &  &  \\
\hline
1707 & grindamótun &  &  &  \\
\hline
1708 & grindargrúpa &  &  &  \\
\hline
1709 & gróf nálgun &  &  &  \\
\hline
1710 & gróf tala, snyrt tala &  &  &  \\
\hline
1711 & grófara grannmynstur, veikara grannmynstur &  &  &  \\
\hline
1712 & gróft mat &  &  &  \\
\hline
1714 & grunnbaugur &  &  &  \\
\hline
1715 & grunneining &  &  &  \\
\hline
1716 & grunnflötur &  &  &  \\
\hline
1717 & grunnform &  &  &  \\
\hline
1718 & grunnfylki &  &  &  \\
\hline
1719 & grunnhorn &  &  &  \\
\hline
1720 & grunnhringur &  &  &  \\
\hline
1721 & grunnlausn &  &  &  \\
\hline
1722 & grunnlausn, stofnlausn &  &  &  \\
\hline
1723 & grunnlína &  &  &  \\
\hline
1724 & grunnlína &  &  &  \\
\hline
1725 & grunnlota &  &  &  \\
\hline
1726 & grunnlota, stofnlota, minnsta lota &  &  &  \\
\hline
1727 & grunnlotukerfi &  &  &  \\
\hline
1728 & grunnpunktur &  &  &  \\
\hline
1729 & grunnpunktur &  &  &  \\
\hline
1730 & grunnpunktur, miðja &  &  &  \\
\hline
1731 & grunnreikningsaðgerð &  &  &  \\
\hline
1732 & grunnrúm &  &  &  \\
\hline
1733 & grunnstæð sönnun &  &  &  \\
\hline
1734 & grunnstætt fall &  &  &  \\
\hline
1735 & grunnstak, stak í grunni &  &  &  \\
\hline
1736 & grunnsvið &  &  &  \\
\hline
1737 & grunntala náttúrlegs logra, stofntala náttúrlegs logra &  &  &  \\
\hline
1738 & grunntala, stofntala &  &  &  \\
\hline
1740 & grunnur, opinn grunnur &  &  &  \\
\hline
1741 & grunnvigur &  &  &  \\
\hline
1742 & grúpa &  &  &  \\
\hline
1743 & grúpa með virkjum &  &  &  \\
\hline
1744 & grúpuaðgerð &  &  &  \\
\hline
1745 & grúpueinsmótun &  &  &  \\
\hline
1746 & grúpufræði &  &  &  \\
\hline
1747 & grúpumiðja, miðja &  &  &  \\
\hline
1748 & grúpumótun, grúpusammótun &  &  &  \\
\hline
1749 & grúpumynstur &  &  &  \\
\hline
1750 & grúpuríki &  &  &  \\
\hline
1751 & grúputafla &  &  &  \\
\hline
1752 & grúpuútvíkkun, grúpuvíkkun &  &  &  \\
\hline
1753 & grúpuverkun, verkun grúpu &  &  &  \\
\hline
1755 & grýpi &  &  &  \\
\hline
1756 & gullinn rétthyrningur &  &  &  \\
\hline
1757 & gullinsnið &  &  &  \\
\hline
1758 & há-, óstækkanlegur &  &  &  \\
\hline
1759 & Haar-mál &  &  &  \\
\hline
1760 & háðar hendingar &  &  &  \\
\hline
1761 & háðir atburðir &  &  &  \\
\hline
1762 & háður &  &  &  \\
\hline
1763 & hæð &  &  &  \\
\hline
1764 & hæðafótpunktaþríhyrningur &  &  &  \\
\hline
1765 & hæðarflötur &  &  &  \\
\hline
1767 & hæðarlína, hæðarferill &  &  &  \\
\hline
1768 & hæðarsetning &  &  &  \\
\hline
1769 & hæðarslétta &  &  &  \\
\hline
1770 & hæðasetning &  &  &  \\
\hline
1771 & hæðaskurðpunktur, hæðapunktur &  &  &  \\
\hline
1772 & hæði &  &  &  \\
\hline
1773 & hægra einingarstak &  &  &  \\
\hline
1774 & hægra hjámengi &  &  &  \\
\hline
1775 & hægra íðal &  &  &  \\
\hline
1776 & hægri aðoki &  &  &  \\
\hline
1777 & hægri afleiða, afleiða frá hægri &  &  &  \\
\hline
1778 & hægri andhverfa, andhverfa frá hægri &  &  &  \\
\hline
1779 & hægri deildanleiki, deildanleiki frá hægri, hægri diffranleiki, diffranleiki frá hægri &  &  &  \\
\hline
1780 & hægri dreifiregla & Sjá <ref idTerm=723>dreifiregla</ref> &  &  \\
\hline
1781 & hægri endapunktur &  &  &  \\
\hline
1782 & hægri grennd &  &  &  \\
\hline
1783 & hægri handar hnitakerfi &  &  &  \\
\hline
1784 & hægri hávísir &  &  &  \\
\hline
1785 & hægri hlið &  &  &  \\
\hline
1786 & hægri hornklofi &  &  &  \\
\hline
1787 & hægri lágvísir &  &  &  \\
\hline
1788 & hægri mótull &  &  &  \\
\hline
1789 & hægri oddklofi &  &  &  \\
\hline
1790 & hægri oddsvigi &  &  &  \\
\hline
1791 & hægri samsemdarmótun &  &  &  \\
\hline
1792 & hægri svigi &  &  &  \\
\hline
1793 & hægri umhverfa, umhverfa frá hægri, hægri andhverfa, andhverfa frá hægri &  &  &  \\
\hline
1794 & hægrisnúin skrúflína &  &  &  \\
\hline
1795 & hægrisnúinn &  &  &  \\
\hline
1796 & hægrisnúinn ferill &  &  &  \\
\hline
1797 & hægvaxandi, hægstígandi &  &  &  \\
\hline
1798 & hækka upp &  &  &  \\
\hline
1799 & hækkun &  &  &  \\
\hline
1800 & hækkunarskekkja &  &  &  \\
\hline
1801 & hærri afleiða &  &  &  \\
\hline
1802 & hærri liður, liður af hærri röð, liður af hærra stigi &  &  &  \\
\hline
1803 & hæsta stig, hástig &  &  &  \\
\hline
1804 & hafa (lausn) &  &  &  \\
\hline
1805 & hafa í för með sér, leiða af sér, leiða til &  &  &  \\
\hline
1806 & háflötur &  &  &  \\
\hline
1807 & hafna, vísa frá &  &  &  \\
\hline
1808 & hágildispunktur &  &  &  \\
\hline
1810 & hágildissetning, hágildislögmál, háalgildissetning, háalgildislögmál &  &  &  \\
\hline
1811 & hagstæður atburður &  &  &  \\
\hline
1812 & hagstætt tilfelli &  &  &  \\
\hline
1813 & hagtölfræði &  &  &  \\
\hline
1814 & háíðal, óstækkanlegt íðal &  &  &  \\
\hline
1815 & háíðalarúm &  &  &  \\
\hline
1816 & hákeðja, óstækkanleg keðja &  &  &  \\
\hline
1817 & hákeðjulögmál Hausdorffs &  &  &  \\
\hline
1818 & hálágmark &  &  &  \\
\hline
1819 & hálágmarkslausn &  &  &  \\
\hline
1820 & hálágmarksleikáætlun, hálágmarksáætlun &  &  &  \\
\hline
1821 & hálágmarkslögmál, fjalldalaregla &  &  &  \\
\hline
1822 & hálágmarkssetning &  &  &  \\
\hline
1823 & haldagripur &  &  &  \\
\hline
1824 & haldahvel, kúluhvel með handföngum &  &  &  \\
\hline
1825 & hálf lota, hálflota &  &  &  \\
\hline
1826 & hálf-, næstum &  &  &  \\
\hline
1827 & hálfákveðinn &  &  &  \\
\hline
1828 & hálfás &  &  &  \\
\hline
1829 & hálfbaugur &  &  &  \\
\hline
1830 & hálfeinfaldur &  &  &  \\
\hline
1831 & hálffirð &  &  &  \\
\hline
1832 & hálffleygaður varpi &  &  &  \\
\hline
1833 & hálfgildingsfirð, gervifirð &  &  &  \\
\hline
1834 & hálfgrúpa &  &  &  \\
\hline
1835 & hálfgrúpumótun &  &  &  \\
\hline
1836 & hálfhliðarregla &  &  &  \\
\hline
1837 & hálfhornrækinn &  &  &  \\
\hline
1838 & hálfhringur &  &  &  \\
\hline
1839 & hálfhvel &  &  &  \\
\hline
1840 & hálflína & <ref idTerm=457>Bil</ref> sem hefur nákvæmlega einn endapunkt &  &  \\
\hline
1841 & hálflínulegur &  &  &  \\
\hline
1842 & hálflograpappír &  &  &  \\
\hline
1843 & hálflokað bil & Samheiti fyrir <ref idTerm=1844>hálfopið bil</ref> &  &  \\
\hline
1844 & hálfopið bil & <ref idTerm=457>Bil</ref> sem inniheldur nákvæmlega einn endapunkta sinna &  &  \\
\hline
1845 & hálfrúm &  &  &  \\
\hline
1846 & hálfsamfelldni &  &  &  \\
\hline
1847 & hálfsamfelldur &  &  &  \\
\hline
1848 & hálfsamfelldur að neðan &  &  &  \\
\hline
1849 & hálfsamfelldur að ofan &  &  &  \\
\hline
1850 & hálfslétta &  &  &  \\
\hline
1851 & hálfstaðall, hálfgildingsstaðall &  &  &  \\
\hline
1852 & hálfstig &  &  &  \\
\hline
1853 & hálftvílínulegur &  &  &  \\
\hline
1854 & hálfur langás &  &  &  \\
\hline
1855 & hálfur meginás &  &  &  \\
\hline
1856 & hálfur miðstrengur &  &  &  \\
\hline
1857 & hálfur skammás &  &  &  \\
\hline
1858 & hálfur þverás &  &  &  \\
\hline
1859 & hali &  &  &  \\
\hline
1860 & hallahorn &  &  &  \\
\hline
1861 & halli &  &  &  \\
\hline
1862 & halli &  &  &  \\
\hline
1863 & halli, hallatala &  &  &  \\
\hline
1864 & hámarka &  &  &  \\
\hline
1865 & hámarkalágmörkun &  &  &  \\
\hline
1866 & hámarksleikáætlun, hámarksáætlun &  &  &  \\
\hline
1867 & hámarksskekkja &  &  &  \\
\hline
1868 & hámarksstaðall &  &  &  \\
\hline
1869 & hamfarafræði &  &  &  \\
\hline
1870 & hamför &  &  &  \\
\hline
1871 & Hamilton-fertala &  &  &  \\
\hline
1872 & Hamilton-rás &  &  &  \\
\hline
1873 & Hamilton-vegur &  &  &  \\
\hline
1874 & Hamilton-virki &  &  &  \\
\hline
1875 & hámörkun &  &  &  \\
\hline
1876 & handfang &  &  &  \\
\hline
1877 & happ &  &  &  \\
\hline
1878 & harðger &  &  &  \\
\hline
1879 & háreisn &  &  &  \\
\hline
1880 & háreistur &  &  &  \\
\hline
1881 & Hartogs-svæði &  &  &  \\
\hline
1882 & háslétta &  &  &  \\
\hline
1883 & hástak, óstækkanlegt stak &  &  &  \\
\hline
1884 & háttargrind &  &  &  \\
\hline
1885 & háttarlögmál &  &  &  \\
\hline
1886 & háttarökfræði &  &  &  \\
\hline
1887 & háttur &  &  &  \\
\hline
1888 & haus eða sporður &  &  &  \\
\hline
1889 & Hausdorff-grannmynstur &  &  &  \\
\hline
1890 & Hausdorff-rúm &  &  &  \\
\hline
1891 & Hausdorff-vídd &  &  &  \\
\hline
1892 & hávísir, brjóstvísir &  &  &  \\
\hline
1893 & hefja (í veldi) &  &  &  \\
\hline
1894 & hefja í annað veldi &  &  &  \\
\hline
1895 & heil línuleg færsla &  &  &  \\
\hline
1896 & heil tala, heiltala & <ref idTerm=3270>Mengið</ref> $\mathbb{Z}=\mathbb{N} \cup (-\mathbb{N})$ kallast mengi \textit{heillra talna} &  & Þessari skilgreiningu er mjög ábótavant en salta hana í bili \\
\hline
1897 & heil veldaröð &  &  &  \\
\hline
1898 & heilbaugur &  &  &  \\
\hline
1899 & heild &  &  &  \\
\hline
1900 & heilda-, heildunar-, heildis-, tegur- &  &  &  \\
\hline
1901 & heilda, tegra &  &  &  \\
\hline
1902 & heildajafna, tegurjafna &  &  &  \\
\hline
1903 & heildanlegur, tegranlegur &  &  &  \\
\hline
1904 & heildanleikaskilyrði, tegranleikaskilyrði &  &  &  \\
\hline
1905 & heildanleiki, tegranleiki &  &  &  \\
\hline
1906 & heildarafleiða &  &  &  \\
\hline
1907 & heildareikningur, tegurreikningur &  &  &  \\
\hline
1908 & heildarfylgni &  &  &  \\
\hline
1909 & heildari, tegrari &  &  &  \\
\hline
1910 & heildarlausn, almenn lausn, fullkomin lausn &  &  &  \\
\hline
1911 & heildarskekkja &  &  &  \\
\hline
1912 & heildarsumma &  &  &  \\
\hline
1913 & heildarsumma, lokasumma &  &  &  \\
\hline
1914 & heildarvik &  &  &  \\
\hline
1915 & heildi, tegur &  &  &  \\
\hline
1916 & heildisfærsla, heildismyndun, heildisummyndun, tegurfærsla, tegurmyndun, tegurummyndun &  &  &  \\
\hline
1917 & heildisformúla, tegurformúla &  &  &  \\
\hline
1918 & heildismælir, tegurmælir &  &  &  \\
\hline
1919 & heildispróf, tegurpróf &  &  &  \\
\hline
1920 & heildissetning, tegursetning &  &  &  \\
\hline
1921 & heildisstofn, tegurstofn &  &  &  \\
\hline
1922 & heildisteiknari, tegurteiknari &  &  &  \\
\hline
1923 & heildisvíðátta &  &  &  \\
\hline
1924 & heildisvirki, heildunarvirki, tegurvirki, tegrunarvirki &  &  &  \\
\hline
1925 & heildun með innsetningu, tegrun með innsetningu, innsetningaraðferð &  &  &  \\
\hline
1926 & heildun, tegrun &  &  &  \\
\hline
1927 & heildunarbil, tegrunarbil &  &  &  \\
\hline
1928 & heildunarbreyta, tegrunarbreyta &  &  &  \\
\hline
1929 & heildunarfasti, tegrunarfasti &  &  &  \\
\hline
1930 & heildunarferill, heildunarvegur, tegrunarferill, tegrunarvegur &  &  &  \\
\hline
1931 & heildunarfræði, tegurfræði &  &  &  \\
\hline
1932 & heildunarsvæði, tegrunarsvæði &  &  &  \\
\hline
1933 & heildunarvirki, tegrunarvirki &  &  &  \\
\hline
1934 & heildunarþáttur, tegrunarþáttur &  &  &  \\
\hline
1935 & heill &  &  &  \\
\hline
1936 & heill snúningur &  &  &  \\
\hline
1937 & heilleg lokun &  &  &  \\
\hline
1938 & heilleg útvíkkun, heilleg baugsútvíkkun, heilleg víkkun, heilleg baugsvíkkun &  &  &  \\
\hline
1939 & heillega lokaður &  &  &  \\
\hline
1940 & heillegur &  &  &  \\
\hline
1941 & heilsteypt grúpa &  &  &  \\
\hline
1942 & heilsteypt mengi &  &  &  \\
\hline
1943 & heilsteypt net &  &  &  \\
\hline
1944 & heilsteyptur töfraferningur &  &  &  \\
\hline
1945 & heilsteyptur, fullgerður &  &  &  \\
\hline
1946 & heilt fall &  &  &  \\
\hline
1947 & heilt margfeldi, margfeldi &  &  &  \\
\hline
1948 & heilt rætt fall, margliða & <ref idTerm=5097>Stæða</ref> af gerðinni \[a_{n}x^n + a_{n-1}x^{n-1} + \ldots + a_{1}x + a_{0},\] þar sem $a_{0},\ldots,a_{n}$ eru <ref idTerm=4093>rauntölur</ref> og $a_{n} \neq 0$, kallast \textit{margliða}. 

Tölurnar $a_{0}, \ldots, a_{n}$ kallast \textit{stuðlar} margliðunnar og kallast þá $a_{n}$ \textit{forystustuðull} hennar og $a_{0}$ \textit{fastastuðull} hennar. <ref idTerm=5102>Stærðirnar</ref> $a_{i}x^i$ kallast \textit{liðir} margliðunnar og kallast þá $a_{n}x^n$ \textit{forystuliður} hennar og $a_{0}$ \textit{fastaliður} hennar.

Talan $n$ kallast \textit{stig} margliðunnar og er það táknað með $\text{deg}(P)$ &  & Minnast á ekki raungildar margliður í skýringu. Vantar færslu fyrir fastamargliðu: " Margliða sem hefur enga liði aðra en fastaliðinn kallast \textit{fastamargliða}" - Hægt að gera skilgreininguna ennþá kjarnyrtari? \\
\hline
1949 & heilt stak, heilstak &  &  &  \\
\hline
1950 & heiltölubestun &  &  &  \\
\hline
1951 & heiltölufylki &  &  &  \\
\hline
1952 & heiltölugildur &  &  &  \\
\hline
1953 & heiltöluhjúpur &  &  &  \\
\hline
1954 & heiltöluhlutafall, gólffall &  &  &  \\
\hline
1955 & heiltöluhluti &  &  &  \\
\hline
1956 & heiltölulausn &  &  &  \\
\hline
1957 & heimfærð stærðfræði &  &  &  \\
\hline
1958 & heimfærður &  &  &  \\
\hline
1959 & heiti &  &  &  \\
\hline
1960 & helminga &  &  &  \\
\hline
1961 & helmingalína &  &  &  \\
\hline
1962 & helmingalína horns & <ref idTerm=2911>Lína</ref> sem fer í gegnum <ref idTerm=3661>oddpunkt</ref> <ref idTerm=2198>horns</ref> og skiptir horninu í tvö eins horn kallast \textit{helmingalína hornsins} &  &  \\
\hline
1963 & helmingaslétta &  &  &  \\
\hline
1964 & helmingun &  &  &  \\
\hline
1965 & hending, slembibreyta, slembistærð, slemba &  &  &  \\
\hline
1966 & herma, meðsveiflun &  &  &  \\
\hline
1967 & herming &  &  &  \\
\hline
1968 & hermísk firð, Hermite-firð &  &  &  \\
\hline
1969 & hermískt form, Hermite-form &  &  &  \\
\hline
1970 & hermískt innfeldi, innfeldi &  &  &  \\
\hline
1971 & hermískt innfeldisrúm, hermískt rúm, innfeldisrúm &  &  &  \\
\hline
1972 & hermískur &  &  &  \\
\hline
1973 & hermískur virki &  &  &  \\
\hline
1974 & Heronsregla & Ef <ref idTerm=6438>þríhyrningur</ref> hefur <ref idTerm=0>hliðarlengdir</ref> $a, b$ og $c$ og $s=\frac{1}{2}(a+b+c)$ er hálft <ref idTerm=5774>ummálið</ref>, þá segir \textit{regla Herons} að <ref idTerm=1265>flatarmál</ref> þríhyrningsins sé $$ F = \sqrt{s (s - a)(s - b)(s - c)}. $$ &  & Vantar færslu fyrir hliðarlengd \\
\hline
1975 & herping, samdráttur &  &  &  \\
\hline
1976 & Hesse-ákveða &  &  &  \\
\hline
1977 & Hesse-fylki &  &  &  \\
\hline
1978 & Hilbert-rúm &  &  &  \\
\hline
1979 & Hilbert-teningur &  &  &  \\
\hline
1980 & himinaflfræði, himinhreyfingafræði, hnatthreyfingafræði &  &  &  \\
\hline
1981 & himinhvel, himinhvolf, himinhvelfing, himinkúla &  &  &  \\
\hline
1982 & hindranafræði &  &  &  \\
\hline
1983 & hindrun &  &  &  \\
\hline
1984 & hitajafna, varmaleiðnijafna &  &  &  \\
\hline
1985 & hjáaðoka, aðoka frá hægri &  &  &  \\
\hline
1986 & hjáalgebra &  &  &  \\
\hline
1987 & hjáendanlegur &  &  &  \\
\hline
1988 & hjágrúpa &  &  &  \\
\hline
1989 & hjáhæð, dýpt &  &  &  \\
\hline
1990 & hjájaðar &  &  &  \\
\hline
1991 & hjájaðarvirki &  &  &  \\
\hline
1992 & hjájaðra &  &  &  \\
\hline
1993 & hjájöðrun &  &  &  \\
\hline
1994 & hjákarteskur ferningur &  &  &  \\
\hline
1995 & hjákeðja &  &  &  \\
\hline
1996 & hjákeðjufléttingur &  &  &  \\
\hline
1997 & hjákjarni &  &  &  \\
\hline
1998 & hjálokaður &  &  &  \\
\hline
1999 & hjálparlína &  &  &  \\
\hline
2000 & hjálparsetning &  &  &  \\
\hline
2001 & hjálparsetning Zorns &  &  &  \\
\hline
2002 & hjámagurt mengi &  &  &  \\
\hline
2003 & hjámargfeldi, summa &  &  &  \\
\hline
2004 & hjámengi &  &  &  \\
\hline
2005 & hjámynd &  &  &  \\
\hline
2006 & hjárás &  &  &  \\
\hline
2007 & hjartaferill &  &  &  \\
\hline
2008 & hjásamtoga &  &  &  \\
\hline
2009 & hjásamtogun &  &  &  \\
\hline
2010 & hjásvipfræði &  &  &  \\
\hline
2011 & hjásvipfræði &  &  &  \\
\hline
2012 & hjásvipfræðibaugur &  &  &  \\
\hline
2013 & hjásvipfræðiflokkur &  &  &  \\
\hline
2014 & hjásvipfræðigrúpa &  &  &  \\
\hline
2015 & hjátrefjasumma, trefjasumma &  &  &  \\
\hline
2016 & hjátrefjun &  &  &  \\
\hline
2017 & hjávarpi, andbrigður varpi, mótbrigður varpi &  &  &  \\
\hline
2018 & hjávídd, fyllivídd &  &  &  \\
\hline
2019 & hjávigur, hjávektor &  &  &  \\
\hline
2020 & hjáþáttafylki, fylgiþáttafylki &  &  &  \\
\hline
2021 & hjáþáttur, fylgiþáttur &  &  &  \\
\hline
2022 & hjól &  &  &  \\
\hline
2023 & hjól &  &  &  \\
\hline
2024 & hjólferill &  &  &  \\
\hline
2025 & hjólflatarhnútur &  &  &  \\
\hline
2026 & hjólflötungur, hjólþykkvi &  &  &  \\
\hline
2027 & hjólflötur, hringfeldi &  &  &  \\
\hline
2028 & hjúfrun &  &  &  \\
\hline
2029 & hjúfur- &  &  &  \\
\hline
2030 & hjúfurhringur, krappahringur, sveigjuhringur &  &  &  \\
\hline
2031 & hjúfurhvel, krappahvel, sveigjuhvel, hjúfurkúla, krappakúla, sveigjukúla &  &  &  \\
\hline
2032 & hjúfurslétta &  &  &  \\
\hline
2033 & hjúfursnertill &  &  &  \\
\hline
2034 & hjúpferill &  &  &  \\
\hline
2035 & hjúpflötur &  &  &  \\
\hline
2036 & hjúpur &  &  &  \\
\hline
2037 & hjúpur &  &  &  \\
\hline
2038 & hlaupakommureikningur &  &  &  \\
\hline
2039 & hlaupandi meðaltal &  &  &  \\
\hline
2040 & hlaupandi, breytilegur &  &  &  \\
\hline
2041 & hlaupavísir, breytilegur vísir &  &  &  \\
\hline
2042 & hleðsla &  &  &  \\
\hline
2043 & hlekking, hlekkjafesti, festi &  &  &  \\
\hline
2044 & hlekkjaður, samhlekkjaður, samslunginn &  &  &  \\
\hline
2045 & hlekkjunartala &  &  &  \\
\hline
2046 & hlekkur &  &  &  \\
\hline
2048 & hlið, hliðarlína, hliðarstrik & Sjá <ref idTerm=3161>marghyrningur</ref> &  &  \\
\hline
2049 & hliðarhorn &  &  &  \\
\hline
2050 & hliðarsetning, samantekt &  &  &  \\
\hline
2051 & hliðarvirki &  &  &  \\
\hline
2052 & hliðhjólferill &  &  &  \\
\hline
2053 & hliðraður &  &  &  \\
\hline
2054 & hliðranagrúpa &  &  &  \\
\hline
2055 & hliðruð línuleg deildajafna, hliðruð línuleg diffurjafna &  &  &  \\
\hline
2056 & hliðrun & <ref idTerm=5779>Færsla</ref> $\tau$ á tiltekinni <ref idTerm=4884>sléttu</ref> kallast \textit{hliðrun} ef hún hefur eftirfarandi tvo eiginleika:

\begin{enumerate}
    \item <ref idTerm=1204>fjarlægðin</ref> milli $A$ og $\tau(A)$ er sú sama fyrir alla <ref idTerm=4001>punkta</ref> $A$,
    \item ef $A \neq \tau(A)$ gildir fyrir alla punkta $A$ að <ref idTerm=1840>hálflínurnar</ref> frá $A$ í gegnum $\tau(A)$ tilgreini sömu <ref idTerm=5135>stefnuna</ref>
\end{enumerate} &  &  \\
\hline
2057 & hliðrun til hægri, hægri hliðrun &  &  &  \\
\hline
2058 & hliðrun til vinstri, vinstri hliðrun &  &  &  \\
\hline
2059 & hlutafleiða &  &  &  \\
\hline
2060 & hlutafleiða (í punkti) &  &  &  \\
\hline
2061 & hlutafleiðujafna &  &  &  \\
\hline
2062 & hlutafleiðuvirki &  &  &  \\
\hline
2063 & hlutákveða &  &  &  \\
\hline
2064 & hlutalgebra &  &  &  \\
\hline
2065 & hlutbaugur &  &  &  \\
\hline
2066 & hlutbil &  &  &  \\
\hline
2067 & hlutbrot &  &  &  \\
\hline
2068 & hlutdeildanlegur, hlutdiffranlegur &  &  &  \\
\hline
2069 & hlutdeildun, hlutdiffrun &  &  &  \\
\hline
2070 & hluteinsfirðun &  &  &  \\
\hline
2071 & hluteinungur &  &  &  \\
\hline
2072 & hlutfall &  &  &  \\
\hline
2073 & hlutfallajafna &  &  &  \\
\hline
2074 & hlutfallsleg skekkja &  &  &  \\
\hline
2075 & hlutfallsleg tíðni &  &  &  \\
\hline
2076 & hlutfallslegt tíðnifall &  &  &  \\
\hline
2077 & hlutfallslegur vöxtur &  &  &  \\
\hline
2078 & hlutfallspróf &  &  &  \\
\hline
2079 & hlutfallsstuðull, hlutfallsþáttur &  &  &  \\
\hline
2080 & hlutfallstala, hlutfallsþáttur &  &  &  \\
\hline
2081 & hlutfjölskylda &  &  &  \\
\hline
2082 & hlutflokkur &  &  &  \\
\hline
2083 & hlutfylgni &  &  &  \\
\hline
2084 & hlutfylgnistuðull &  &  &  \\
\hline
2085 & hlutfylki &  &  &  \\
\hline
2086 & hlutgrind &  &  &  \\
\hline
2087 & hlutgrúpa &  &  &  \\
\hline
2088 & hlutheildun, hluttegrun &  &  &  \\
\hline
2089 & hluti, partur &  &  &  \\
\hline
2090 & hlutíðal &  &  &  \\
\hline
2091 & hlutkenning &  &  &  \\
\hline
2092 & hlutkerfi &  &  &  \\
\hline
2093 & hlutlaus &  &  &  \\
\hline
2094 & hlutlaus rúmfræði &  &  &  \\
\hline
2095 & hlutleysuþáttur &  &  &  \\
\hline
2096 & hlutlíkan &  &  &  \\
\hline
2097 & hlutmargfeldi, margfeldisstúfur &  &  &  \\
\hline
2098 & hlutmengi & <ref idTerm=3270>Mengi</ref> $A$ er sagt vera \textit{hlutmengi} í mengi $B$ ef sérhvert <ref idTerm=5128>stak</ref> í $A$ er líka stak í $B$. Þetta er táknað með $A \subset B$ &  & Nota $\subseteq$ í staðinn? \\
\hline
2099 & hlutmismunajafna &  &  &  \\
\hline
2100 & hlutmismunur &  &  &  \\
\hline
2101 & hlutmótull &  &  &  \\
\hline
2102 & hlutmynstur &  &  &  \\
\hline
2103 & hlutnefnari &  &  &  \\
\hline
2104 & hlutnet &  &  &  \\
\hline
2105 & hlutnet &  &  &  \\
\hline
2106 & hlutniðurstaða, hlutútkoma &  &  &  \\
\hline
2107 & hlutorð &  &  &  \\
\hline
2108 & hlutrakinn &  &  &  \\
\hline
2109 & hlutreiknanlegur &  &  &  \\
\hline
2110 & hlutríki &  &  &  \\
\hline
2111 & hlutröð &  &  &  \\
\hline
2112 & hlutrúm &  &  &  \\
\hline
2113 & hlutruna &  &  &  \\
\hline
2114 & hlutskilgreind aðgerð, hlutskilgreind reikniaðgerð &  &  &  \\
\hline
2115 & hlutskilgreind vörpun &  &  &  \\
\hline
2116 & hlutskilgreint fall &  &  &  \\
\hline
2117 & hlutskilgreint rakið fall, hlutrakið fall &  &  &  \\
\hline
2118 & hlutsumma, summustúfur &  &  &  \\
\hline
2119 & hlutsvæði &  &  &  \\
\hline
2120 & hlutsvið &  &  &  \\
\hline
2121 & hlutteljari &  &  &  \\
\hline
2122 & hluttré &  &  &  \\
\hline
2123 & hlutur, gripur &  &  &  \\
\hline
2124 & hlutúrtak &  &  &  \\
\hline
2125 & hlutvera, hlutmengisvensl &  &  &  \\
\hline
2126 & hlutverutákn &  &  &  \\
\hline
2127 & hlutvíðátta &  &  &  \\
\hline
2128 & hlutvíðerni &  &  &  \\
\hline
2129 & hlutvigurrúm, línulegt hlutrúm &  &  &  \\
\hline
2130 & hlutþakning &  &  &  \\
\hline
2131 & hnappelda &  &  &  \\
\hline
2132 & hneppishraði, hnapphraði &  &  &  \\
\hline
2133 & hnig &  &  &  \\
\hline
2134 & hnigaðferð &  &  &  \\
\hline
2135 & hnigaðferð mesta bratta &  &  &  \\
\hline
2136 & hnigskref &  &  &  \\
\hline
2137 & hnik &  &  &  \\
\hline
2138 & hnikaafleiða &  &  &  \\
\hline
2139 & hnikajafna &  &  &  \\
\hline
2140 & hnikalögmál &  &  &  \\
\hline
2141 & hnikareikningur &  &  &  \\
\hline
2142 & hnikaverkefni &  &  &  \\
\hline
2143 & hnit &  &  &  \\
\hline
2144 & hnit, stak, stuðull &  &  &  \\
\hline
2145 & hnitakerfi & <ref idTerm=4541>Samsvörun</ref> milli <ref idTerm=4001>punkta</ref> tiltekins <ref idTerm=4266>rúms</ref> og <ref idTerm=5397>talna</ref> sem kallast \textit{hnit} sem gerir manni kleift að staðsetja punkta í rúminu &  &  \\
\hline
2146 & hnitapappír &  &  &  \\
\hline
2147 & hnitarúmfræði, hnitafræði &  &  &  \\
\hline
2148 & hnitás, hnitaás &  &  &  \\
\hline
2149 & hnitaskiptafylki, grunnskiptafylki &  &  &  \\
\hline
2150 & hnitaskiptavörpun &  &  &  \\
\hline
2151 & hnitaskipti, hnitafærsla &  &  &  \\
\hline
2152 & hnitfall, hnitafall &  &  &  \\
\hline
2153 & hnitgrennd, hnitagrennd &  &  &  \\
\hline
2154 & hnitrúm, hnitarúm &  &  &  \\
\hline
2155 & hnitslétta, hnitaslétta &  &  &  \\
\hline
2156 & hnöttóttur, kúlulaga &  &  &  \\
\hline
2157 & hnútafræði &  &  &  \\
\hline
2158 & hnútalitun &  &  &  \\
\hline
2159 & hnútamengi &  &  &  \\
\hline
2160 & hnútflötur &  &  &  \\
\hline
2161 & hnútgrúpa &  &  &  \\
\hline
2162 & hnútsnet, örnet &  &  &  \\
\hline
2163 & hnútur &  &  &  \\
\hline
2164 & hnútur, gálma &  &  &  \\
\hline
2165 & hnútur, punktur &  &  &  \\
\hline
2166 & hnýttur &  &  &  \\
\hline
2167 & höfnun, frávísun &  &  &  \\
\hline
2168 & höfnunarfall, styrkleikafall, styrkfall &  &  &  \\
\hline
2169 & höfnunarhætta &  &  &  \\
\hline
2170 & höfnunarlíkindi, frávísunarlíkindi &  &  &  \\
\hline
2171 & höfnunarlíkindi, frávísunarlíkindi &  &  &  \\
\hline
2172 & höfnunarlína, frávísunarlína &  &  &  \\
\hline
2173 & höfnunarmark, frávísunarmark &  &  &  \\
\hline
2174 & höfnunarmistök &  &  &  \\
\hline
2175 & höfnunarsvæði, frávísunarsvæði &  &  &  \\
\hline
2176 & höfuðásaaðferð, meginásaaðferð &  &  &  \\
\hline
2177 & höfuðdeilir &  &  &  \\
\hline
2178 & höfuðgildi, megingildi &  &  &  \\
\hline
2179 & höfuðgrein, megingrein &  &  &  \\
\hline
2180 & höfuðhlutákveða &  &  &  \\
\hline
2181 & höfuðhluti &  &  &  \\
\hline
2182 & höfuðíðal &  &  &  \\
\hline
2183 & höfuðíðalbaugur &  &  &  \\
\hline
2184 & höfuðíðalsetning &  &  &  \\
\hline
2185 & höfuðkrappageisli, meginkrappageisli, höfuðsveigjugeisli, meginsveigjugeisli &  &  &  \\
\hline
2186 & höfuðkrappamiðja, krappamiðja, sveigjumiðja &  &  &  \\
\hline
2187 & höfuðkrappi, meginkrappi, höfuðsveigja, meginsveigja &  &  &  \\
\hline
2188 & höfuðlausn &  &  &  \\
\hline
2189 & höfuðsnertill, meginsnertill &  &  &  \\
\hline
2190 & höfuðstóll &  &  &  \\
\hline
2191 & höfuðturn &  &  &  \\
\hline
2192 & hófur &  &  &  \\
\hline
2193 & höggbylgja &  &  &  \\
\hline
2194 & högun &  &  &  \\
\hline
2195 & hóptíðni &  &  &  \\
\hline
2196 & hópur &  &  &  \\
\hline
2197 & hópþáttur, flokksþáttur &  &  &  \\
\hline
2198 & horn & Tvær <ref idTerm=1840>hálflínur</ref> með sama <ref idTerm=5850>upphafspunkt</ref> eru sagðar mynda \textit{horn}. Upphafspunktur hálflínanna kallast þá \textit{oddpunktur} hornsins og hálflínurnar kallast \textit{armar} þess &  & Setja í skýringu: Sérhver tvö <ref idTerm=2958>línustrik</ref> með sameiginlegan endapunkt sem eru innihaldin í horninu mynda jafnframt sama horn \\
\hline
2199 & horn, stefnuhorn, áttarhorn, horngildi &  &  &  \\
\hline
2200 & horn, stefnuhorn, áttarhorn, hornhnit &  &  &  \\
\hline
2201 & horna-, hornafræði- &  &  &  \\
\hline
2202 & hornafall &  &  &  \\
\hline
2203 & hornafallamargliða &  &  &  \\
\hline
2204 & hornafallaröð &  &  &  \\
\hline
2205 & hornafleiða &  &  &  \\
\hline
2206 & hornafræði &  &  &  \\
\hline
2207 & hornalína &  &  &  \\
\hline
2208 & hornalínu-reitafylki &  &  &  \\
\hline
2209 & hornalínuaðferð &  &  &  \\
\hline
2210 & hornalínuaðferð Cantors &  &  &  \\
\hline
2211 & hornalínuaðhvarf &  &  &  \\
\hline
2212 & hornalínufylki &  &  &  \\
\hline
2213 & hornalínugera, hornalínufæra &  &  &  \\
\hline
2214 & hornalínugerð &  &  &  \\
\hline
2215 & hornalínugerlegur &  &  &  \\
\hline
2216 & hornalínugerningur &  &  &  \\
\hline
2217 & hornalínugná &  &  &  \\
\hline
2218 & hornalínugnæfur &  &  &  \\
\hline
2219 & hornalínureitur &  &  &  \\
\hline
2220 & hornalínustak, hornalínustuðull &  &  &  \\
\hline
2221 & hornalínuvörpun &  &  &  \\
\hline
2222 & hornalínuþríhyrningur &  &  &  \\
\hline
2223 & hornamál &  &  &  \\
\hline
2224 & hornamál, hornmál &  &  &  \\
\hline
2225 & hornasumma & <ref idTerm=2820>Samanlagt</ref> <ref idTerm=2223>gráðumál</ref> <ref idTerm=2198>horna</ref> <ref idTerm=3161>marghyrnings</ref> &  &  \\
\hline
2226 & hornaukasetning &  &  &  \\
\hline
2227 & hornfjarlægð &  &  &  \\
\hline
2228 & hornhraði &  &  &  \\
\hline
2229 & hornhröðun &  &  &  \\
\hline
2231 & hornklofamargfeldi, Lie-margfeldi &  &  &  \\
\hline
2232 & hornklofamargföldun, Lie-margföldun &  &  &  \\
\hline
2233 & hornklofi &  &  &  \\
\hline
2234 & hornmínúta &  &  &  \\
\hline
2235 & hornpunktur, oddpunktur & <ref idTerm=4001>Punktur</ref> þar sem tvær <ref idTerm=2048>hliðar</ref> <ref idTerm=3161>marghyrnings</ref> <ref idTerm=4771>skerast</ref> (í tveimur <ref idTerm=6087>víddum</ref>) eða punktur þar sem þrjár eða fleiri <ref idTerm=1325>hliðar</ref> <ref idTerm=3133>margflötungs</ref> skerast (í þrem víddum) &  & Ræða almenna margflötunga í skýringu \\
\hline
2236 & hornrækin vörpun &  &  &  \\
\hline
2237 & hornrækinn &  &  &  \\
\hline
2238 & hornrækni &  &  &  \\
\hline
2239 & hornræknirúmfræði &  &  &  \\
\hline
2240 & hornrétt liðun, þverstæð liðun, þverliðun &  &  &  \\
\hline
2241 & hornrétt ofanvarp &  &  &  \\
\hline
2242 & hornréttur &  &  &  \\
\hline
2243 & hornsekúnda &  &  &  \\
\hline
2244 & hraði &  &  &  \\
\hline
2245 & hraði &  &  &  \\
\hline
2246 & hraðminnkandi, hraðfallandi &  &  &  \\
\hline
2247 & hraður &  &  &  \\
\hline
2248 & hraðvaxandi, hraðstækkandi &  &  &  \\
\hline
2249 & hreiðraður &  &  &  \\
\hline
2250 & hreiðruð bil &  &  &  \\
\hline
2251 & hreiðruð margföldun &  &  &  \\
\hline
2252 & hreiðruð tvískipting &  &  &  \\
\hline
2253 & hreiður &  &  &  \\
\hline
2254 & hrein algebra &  &  &  \\
\hline
2255 & hrein rúmfræði &  &  &  \\
\hline
2256 & hrein stærðfræði &  &  &  \\
\hline
2257 & hrein torræð útvíkkun, hrein torræð víkkun &  &  &  \\
\hline
2258 & hreinlega óaðskiljanleg útvíkkun, óaðskiljanleg útvíkkun, hreinlega óaðskiljanleg víkkun, óaðskiljanleg víkkun &  &  &  \\
\hline
2259 & hreint lotutugabrot, hreint umferðartugabrot &  &  &  \\
\hline
2260 & hreintóna deildajafna &  &  &  \\
\hline
2261 & hreintóna hreyfing &  &  &  \\
\hline
2262 & hreintóna sveifill &  &  &  \\
\hline
2263 & hreintóna sveifla, hrein sveifla &  &  &  \\
\hline
2264 & hreinvíður &  &  &  \\
\hline
2265 & hreinvíður &  &  &  \\
\hline
2266 & hreyfanleg sérstaða, hreyfanlegur sérstöðupunktur &  &  &  \\
\hline
2267 & hreyfifræði &  &  &  \\
\hline
2268 & hreyfihnitakerfi &  &  &  \\
\hline
2269 & hreyfilýsing, gangfræði &  &  &  \\
\hline
2270 & hreyfing &  &  &  \\
\hline
2271 & hreyfingarfasti &  &  &  \\
\hline
2272 & hreyfirammi, hreyfigrunnur &  &  &  \\
\hline
2273 & hringabundin, bundin af hringjum &  &  &  \\
\hline
2274 & hringaðferð &  &  &  \\
\hline
2275 & hringbogaþríhyrningur &  &  &  \\
\hline
2276 & hringbogi, bogi, örk & Sá hluti <ref idTerm=2295>hrings</ref> sem liggur milli tveggja <ref idTerm=4001>punkta</ref> hans kallast \textit{hringbogi} &  & Setja í skýringu: Þegar talað er um bogann $AB$ sem liggur milli tveggja punkta hringsins $A$ og $B$ er jafnan átt við þann minni \\
\hline
2277 & hringfari &  &  &  \\
\hline
2278 & hringfeldisgrúpa &  &  &  \\
\hline
2279 & hringgrúpa &  &  &  \\
\hline
2280 & hringhreyfing &  &  &  \\
\hline
2281 & hringhverfing punkts, spegilmynd punkts í hring &  &  &  \\
\hline
2282 & hringhverfing, hringspeglun, speglun í hring &  &  &  \\
\hline
2283 & hringjaðar, hringferill &  &  &  \\
\hline
2284 & hringkragi, kragi &  &  &  \\
\hline
2285 & hringlaga, kringlóttur, hring- &  &  &  \\
\hline
2286 & hringsamhverfni &  &  &  \\
\hline
2287 & hringskiptijafna &  &  &  \\
\hline
2288 & hringskiptimargliða &  &  &  \\
\hline
2289 & hringskipting &  &  &  \\
\hline
2290 & hringskiptisvið &  &  &  \\
\hline
2291 & hringsneið, sneið &  &  &  \\
\hline
2292 & hringsvæði &  &  &  \\
\hline
2294 & hringur &  &  &  \\
\hline
2295 & hringur, hringferill, hringlína & <ref idTerm=3270>Mengi</ref> þeirra <ref idTerm=4001>punkta</ref> í tiltekinni <ref idTerm=4884>sléttu</ref> sem eru í sömu <ref idTerm=1204>fjarlægð</ref> frá punkti sem liggur á sléttunni. Sá punktur kallast þá \textit{miðja} hringsins og fjarlægðin frá honum kallast \textit{geisli} hringsins &  & \\
\hline
2296 & hringvik &  &  &  \\
\hline
2297 & hringvika &  &  &  \\
\hline
2298 & hringvikshorn &  &  &  \\
\hline
2299 & hrísla &  &  &  \\
\hline
2300 & hröð samleitni &  &  &  \\
\hline
2301 & hröðun &  &  &  \\
\hline
2302 & hugarreikningur &  &  &  \\
\hline
2303 & hugartilraun &  &  &  \\
\hline
2304 & hugbúnaður &  &  &  \\
\hline
2305 & hugsaður &  &  &  \\
\hline
2306 & hugtak &  &  &  \\
\hline
2307 & huldureikningur &  &  &  \\
\hline
2308 & hulduvísir &  &  &  \\
\hline
2309 & hundraðshlutamark, hundraðsmark &  &  &  \\
\hline
2310 & hvarf &  &  &  \\
\hline
2311 & hvarf &  &  &  \\
\hline
2312 & hvarfflötur &  &  &  \\
\hline
2313 & hvarfgildi &  &  &  \\
\hline
2314 & hvarfgildi, stöðugildi &  &  &  \\
\hline
2315 & hvarfpunktur, hvarf, jafnvægispunktur, stöðupunktur &  &  &  \\
\hline
2316 & hvarfsetning &  &  &  \\
\hline
2317 & hvass & <ref idTerm=2198>Horn</ref> kallast \textit{hvasst} ef það <ref idTerm=2223>mælist</ref> minna en <ref idTerm=4167>rétt horn</ref> &  &  \\
\hline
2318 & hvasshyrndur & <ref idTerm=6438>Þríhyrningur</ref> með eitt <ref idTerm=2317>hvasst</ref> <ref idTerm=2198>horn</ref>. &  &  \\
\hline
2319 & hvel, kúluhvel, kúluflötur, kúluyfirborð & <ref idTerm=3270>Mengi</ref> þeirra <ref idTerm=4001>punkta</ref> sem eru í sömu <ref idTerm=1204>fjarlægð</ref> frá tilteknum punkti &  & Taka fram að punktarnir þurfi að vera í firðrúmi í skýringu og minnast á miðju kúluhvelsins \\
\hline
2321 & hvelfdur &  &  &  \\
\hline
2322 & hvelft fall &  &  &  \\
\hline
2323 & hvelft horn &  &  &  \\
\hline
2324 & hverfa alls staðar, vera alls staðar núll &  &  &  \\
\hline
2325 & hverfa, verða núll &  &  &  \\
\hline
2326 & hverfandi, sem hverfur &  &  &  \\
\hline
2327 & hverfipunktur, beygjuskil &  &  &  \\
\hline
2328 & hverfisnertill, beygjuskilasnertill &  &  &  \\
\hline
2329 & hverfitregða, tregðuvægi &  &  &  \\
\hline
2330 & hverfull, svipull &  &  &  \\
\hline
2331 & hverfulleiki &  &  &  \\
\hline
2332 & hverfult ástand &  &  &  \\
\hline
2333 & hvergi deildanlegur, hvergi diffranlegur &  &  &  \\
\hline
2334 & hvergi núll &  &  &  \\
\hline
2335 & hvergi samleitinn &  &  &  \\
\hline
2336 & hvergi þéttur &  &  &  \\
\hline
2337 & hversdagsleg mengjafræði &  &  &  \\
\hline
2338 & hversu lítill sem vera skal &  &  &  \\
\hline
2339 & hversu stór sem vera skal &  &  &  \\
\hline
2340 & hvirfilhringur &  &  &  \\
\hline
2341 & hvirfilpunktur &  &  &  \\
\hline
2342 & hvirfilpunktur &  &  &  \\
\hline
2343 & hvirfilpunktur, miðja &  &  &  \\
\hline
2344 & hyrna &  &  &  \\
\hline
2345 & hyrnt hvel &  &  &  \\
\hline
2346 & hyrnuaðferð &  &  &  \\
\hline
2347 & hyrnubákn, hyrnuþyrping &  &  &  \\
\hline
2348 & hyrnufruma &  &  &  \\
\hline
2349 & hyrnukeðja &  &  &  \\
\hline
2350 & hyrnunálgun &  &  &  \\
\hline
2351 & hyrnurit &  &  &  \\
\hline
2352 & hyrnusvipfræði &  &  &  \\
\hline
2353 & hyrnuvörpun &  &  &  \\
\hline
2354 & í hið óendanlega &  &  &  \\
\hline
2355 & í jöfnum mæli &  &  &  \\
\hline
2356 & í öfugu hlutfalli við &  &  &  \\
\hline
2357 & í punkti, í hverjum punkti &  &  &  \\
\hline
2358 & í samræmi við röðun, raðvís &  &  &  \\
\hline
2359 & iða &  &  &  \\
\hline
2360 & íðal &  &  &  \\
\hline
2361 & íðal- &  &  &  \\
\hline
2362 & íðalflokkagrúpa &  &  &  \\
\hline
2363 & íðalflokkur &  &  &  \\
\hline
2364 & íðalgrúpa &  &  &  \\
\hline
2365 & íðalkeðja &  &  &  \\
\hline
2366 & iðurfeldi &  &  &  \\
\hline
2367 & íholur, íbjúgur &  &  &  \\
\hline
2368 & ílag &  &  &  \\
\hline
2369 & ílags-frálagsgreining &  &  &  \\
\hline
2370 & ílagsgögn &  &  &  \\
\hline
2371 & ílangur sporvölusnúðflötur &  &  &  \\
\hline
2372 & ílangur sporvölusnúðflötur &  &  &  \\
\hline
2373 & ílangur sporvölusnúður &  &  &  \\
\hline
2374 & ílangur sporvölusnúður &  &  &  \\
\hline
2375 & illa framsett verkefni &  &  &  \\
\hline
2376 & ilpunktur &  &  &  \\
\hline
2377 & ímyndaður hringur &  &  &  \\
\hline
2378 & innanmótun &  &  &  \\
\hline
2379 & innanverð víxlhorn, innri víxlhorn &  &  &  \\
\hline
2380 & innbyrðis óháðir atburðir &  &  &  \\
\hline
2381 & innbyrðis óháður &  &  &  \\
\hline
2382 & inndragi &  &  &  \\
\hline
2383 & inndráttur &  &  &  \\
\hline
2386 & inngeisli, geisli innritaðs hrings, geisli innritaðrar kúlu &  &  &  \\
\hline
2387 & inngeisli, innri geisli, skammgeisli &  &  &  \\
\hline
2388 & innhjólferill &  &  &  \\
\hline
2389 & innhorn &  &  &  \\
\hline
2390 & innhyrndur &  &  &  \\
\hline
2391 & innhyrndur marghyrningur &  &  &  \\
\hline
2392 & innifalinn, umlukinn &  &  &  \\
\hline
2393 & innihald, Jordan-innihald &  &  &  \\
\hline
2394 & innihalda &  &  &  \\
\hline
2395 & innihaldslaus &  &  &  \\
\hline
2396 & innihaldslaus staðhæfing &  &  &  \\
\hline
2397 & innliðavíxl &  &  &  \\
\hline
2398 & innmengi, iður, innri kjarni, opinn kjarni &  &  &  \\
\hline
2399 & innmiðja, miðja innritaðs hrings, miðja innritaðrar kúlu &  &  &  \\
\hline
2400 & innmyndun, staðbundin greyping &  &  &  \\
\hline
2401 & innör &  &  &  \\
\hline
2402 & innra innihald, innra Jordan-innihald &  &  &  \\
\hline
2403 & innra mál, innanmál &  &  &  \\
\hline
2404 & innra margfeldi &  &  &  \\
\hline
2405 & innra meðalfervik, innri dreifni &  &  &  \\
\hline
2406 & innri aðgerð, innri reikniaðgerð &  &  &  \\
\hline
2407 & innri eiginleiki &  &  &  \\
\hline
2408 & innri firð &  &  &  \\
\hline
2409 & innri fylgni &  &  &  \\
\hline
2410 & innri helmingalína &  &  &  \\
\hline
2411 & innri jafna &  &  &  \\
\hline
2412 & innri margföldun &  &  &  \\
\hline
2413 & innri punktur &  &  &  \\
\hline
2414 & innri rúmfræði &  &  &  \\
\hline
2415 & innri sjálfmótun &  &  &  \\
\hline
2416 & innrita &  &  &  \\
\hline
2417 & innritaður &  &  &  \\
\hline
2418 & innritaður hringur, innanverður snertihringur & <ref idTerm=2295>Hringur</ref> er \textit{innritaður} í <ref idTerm=3161>marghyrning</ref> ef allar <ref idTerm=2048>hliðar</ref> marghyrningsins eru <ref idTerm=4912>snertlar</ref> við hringinn &  &  \\
\hline
2419 & innrituð kúla, innanverð snertikúla &  &  &  \\
\hline
2420 & innsæi, skjótsæi, snarskyggni &  &  &  \\
\hline
2421 & innsæiskenning &  &  &  \\
\hline
2422 & innsæisrökfræði &  &  &  \\
\hline
2423 & innsetjanlegur &  &  &  \\
\hline
2424 & innsetning &  &  &  \\
\hline
2425 & innsetning &  &  &  \\
\hline
2426 & innsetningaraðferð &  &  &  \\
\hline
2427 & innsetningarfrumsenda, innsetningarfrumsetning, frumsenda um innsetningu, frumsetning um innsetningu &  &  &  \\
\hline
2428 & innsetningarregla &  &  &  \\
\hline
2429 & innsetningarregla &  &  &  \\
\hline
2430 & innstæður þvervigur, innstæður þverill, innri þvervigur, innri þverill &  &  &  \\
\hline
2431 & innstæður, innri &  &  &  \\
\hline
2432 & innstæður, innri &  &  &  \\
\hline
2433 & innstætt horn &  &  &  \\
\hline
2434 & innstig, inngráða &  &  &  \\
\hline
2435 & innsveifarferill &  &  &  \\
\hline
2436 & ítreka &  &  &  \\
\hline
2437 & ítrekað heildi, ítrekað tegur, endurtekið heildi, endurtekið tegur &  &  &  \\
\hline
2438 & ítrekun &  &  &  \\
\hline
2439 & ítrekun (vörpunar), ítrekuð vörpun &  &  &  \\
\hline
2440 & ítrekunaraðferð &  &  &  \\
\hline
2441 & ítrekunarfræði &  &  &  \\
\hline
2442 & ítrekunarrás, ítrekunarhringur, ítrekunarlykkja &  &  &  \\
\hline
2443 & ítrekunarreikningur &  &  &  \\
\hline
2444 & ítrekunarskref &  &  &  \\
\hline
2445 & ívera &  &  &  \\
\hline
2446 & íverutákn &  &  &  \\
\hline
2447 & íveruvensl &  &  &  \\
\hline
2448 & Jacobi-ákveða &  &  &  \\
\hline
2449 & Jacobi-fylki &  &  &  \\
\hline
2450 & Jacobi-tákn &  &  &  \\
\hline
2451 & Jacobi-víðerni &  &  &  \\
\hline
2452 & Jacobson-rót &  &  &  \\
\hline
2453 & jaðar &  &  &  \\
\hline
2454 & jaðar &  &  &  \\
\hline
2455 & jaðar, jaðarferill &  &  &  \\
\hline
2456 & jaðarbogi &  &  &  \\
\hline
2457 & jaðarbútur &  &  &  \\
\hline
2458 & jaðardreifing &  &  &  \\
\hline
2459 & jaðarferill &  &  &  \\
\hline
2460 & jaðargildi &  &  &  \\
\hline
2461 & jaðargildisverkefni &  &  &  \\
\hline
2462 & jaðarhegðun &  &  &  \\
\hline
2463 & jaðarhringur &  &  &  \\
\hline
2464 & jaðarkúla &  &  &  \\
\hline
2465 & jaðarlag &  &  &  \\
\hline
2466 & jaðarpunktur &  &  &  \\
\hline
2467 & jaðarskilyrði &  &  &  \\
\hline
2468 & jaðarvíðátta &  &  &  \\
\hline
2469 & jaðarvirki &  &  &  \\
\hline
2470 & jaðarvörpun &  &  &  \\
\hline
2471 & jaðarþéttleiki &  &  &  \\
\hline
2472 & jaðraður &  &  &  \\
\hline
2473 & jafn að ummáli, ummálsjafn &  &  &  \\
\hline
2474 & jafn, jafnstór &  &  &  \\
\hline
2475 & jafn, sléttur &  &  &  \\
\hline
2476 & jafna Pells &  &  &  \\
\hline
2477 & jafna út &  &  &  \\
\hline
2478 & jafna, líking & <ref idTerm=1515>Yrðing</ref> sem lýsir því að tvær <ref idTerm=5102>stærðir</ref> séu jafnar &  &  \\
\hline
2479 & jafnaðarmerki, jöfnunarmerki, samasemmerki, samsemdarmerki &  &  &  \\
\hline
2480 & jafnarma & <ref idTerm=6438>Þríhyrningur</ref> kallast \textit{jafnarma} ef hann hefur tvær eins <ref idTerm=2048>hliðar</ref> &  &  \\
\hline
2481 & jafnása breiðbogi, jafnhliða breiðbogi, réttstæður breiðbogi &  &  &  \\
\hline
2482 & jafnbreiður ferill &  &  &  \\
\hline
2483 & jafnbreytinn &  &  &  \\
\hline
2484 & jafndreifður &  &  &  \\
\hline
2485 & jafnfirðarferill &  &  &  \\
\hline
2486 & jafnfirðarferill &  &  &  \\
\hline
2487 & jafnfjarlægur &  &  &  \\
\hline
2488 & jafngildi &  &  &  \\
\hline
2489 & jafngildisflokkur & <ref idTerm=2098>Hlutmengi</ref> $A$ í <ref idTerm=3270>mengi</ref> $X$ sem uppfyllir fyrir öll $x,y \in A$ að $x$ og $y$ séu <ref idTerm=0>vensluð saman</ref> með tilliti til tiltekinna <ref idTerm=6592>jafngildisvensla</ref> og að $x \in A$ sé ekki venslað saman við neitt <ref idTerm=5128>stak</ref> í $X \setminus A$ &  & \\
\hline
2490 & jafngildur &  &  &  \\
\hline
2491 & jafngildur í jöfnum mæli &  &  &  \\
\hline
2492 & jafngilt fylki &  &  &  \\
\hline
2493 & jafnhallaferill &  &  &  \\
\hline
2494 & jafnhallaferill &  &  &  \\
\hline
2495 & jafnhliða &  &  &  \\
\hline
2496 & jafnhliða þríhyrningur & <ref idTerm=6438>Þríhyrningur</ref> með allar <ref idTerm=2048>hliðar</ref> jafn <ref idTerm=2860>langar</ref> &  &  \\
\hline
2497 & jafnhlutfallaröð, kvótaröð, rúmkynja röð &  &  &  \\
\hline
2498 & jafnhlutfallaruna, kvótaruna &  &  &  \\
\hline
2499 & jafnhorna, sem mynda sama horn &  &  &  \\
\hline
2500 & jafnhornaferill &  &  &  \\
\hline
2501 & jafnhyrndur, með öll horn jöfn &  &  &  \\
\hline
2502 & jafnhyrningur, jafnhyrndur marghyrningur &  &  &  \\
\hline
2503 & jafnlíklegur &  &  &  \\
\hline
2504 & jafnmælisgrannmynstur &  &  &  \\
\hline
2505 & jafnmælismynstur &  &  &  \\
\hline
2506 & jafnmælisrúm &  &  &  \\
\hline
2507 & jafnmælisskorða &  &  &  \\
\hline
2508 & jafnmættisflötur &  &  &  \\
\hline
2509 & jafnmættislína &  &  &  \\
\hline
2510 & jafnmunaröð &  &  &  \\
\hline
2511 & jafnmunaruna &  &  &  \\
\hline
2512 & jafnsamfelldni &  &  &  \\
\hline
2513 & jafnsamfelldur &  &  &  \\
\hline
2514 & jafnsamleitinn &  &  &  \\
\hline
2515 & jafnsamleitni &  &  &  \\
\hline
2516 & jafnstæðni, jafnleiki, sléttleiki &  &  &  \\
\hline
2517 & jafnstæðugrúpa &  &  &  \\
\hline
2518 & jafnstæður & <ref idTerm=1078>Fall</ref> $f: X \to Y$ þar sem $X$ og $Y$ eru <ref idTerm=2098>hlutmengi</ref> í <ref idTerm=3270>mengi</ref> <ref idTerm=4093>rauntalna</ref> kallast \textit{jafnstætt} ef fyrir öll $x \in X$ gildi að \[ f(-x) = f(x)\] &  &  \\
\hline
2519 & jafnstæður &  &  &  \\
\hline
2520 & jafntímaferill &  &  &  \\
\hline
2521 & jafnvæg skekkja &  &  &  \\
\hline
2522 & jafnvægi &  &  &  \\
\hline
2523 & jafnvægt mengi &  &  &  \\
\hline
2524 & jafnvægur &  &  &  \\
\hline
2525 & jafnvíður &  &  &  \\
\hline
2526 & jafnvígur hnútur &  &  &  \\
\hline
2527 & jafnþætt hnit, einsleit hnit &  &  &  \\
\hline
2528 & jafnþættur, einsleitur &  &  &  \\
\hline
2529 & jafnþátta &  &  &  \\
\hline
2530 & jafnþáttuð margliða &  &  &  \\
\hline
2531 & jákvæð fylgni &  &  &  \\
\hline
2532 & jákvæð stríkkun &  &  &  \\
\hline
2533 & jákvæður snúningur, snúningur rangsælis &  &  &  \\
\hline
2534 & jákvæður snúningur, snúningur rangsælis &  &  &  \\
\hline
2535 & jákvæður, viðlægur & <ref idTerm=4093>Rauntala</ref> kallast \textit{jákvæð} ef hún er stærri en 0 &  & Vantar færslu fyrir samanburð á rauntölum (samanburður er til en e.t.v. of vítt hugtak)  \\
\hline
2536 & jákvætt ákveðinn &  &  &  \\
\hline
2537 & jákvætt áttaður &  &  &  \\
\hline
2538 & jákvætt endurkvæmur, jákvætt afturkvæmur &  &  &  \\
\hline
2539 & jákvætt hálfákveðinn &  &  &  \\
\hline
2540 & japanskur baugur &  &  &  \\
\hline
2541 & japönsk kúlnagrind, japönsk talnagrind, japönsk reiknigrind &  &  &  \\
\hline
2542 & jöðrun &  &  &  \\
\hline
2543 & jöðrunarfræði &  &  &  \\
\hline
2544 & jöfn dreifing &  &  &  \\
\hline
2545 & jöfn hröðun &  &  &  \\
\hline
2546 & jöfn tala, slétt tala & <ref idTerm=1896>Heil tala</ref> sem $2$ <ref idTerm=1599>gengur upp í</ref>, þ.e. heil tala $n$ sem rita má á forminu $n = 2 \cdot k$ þar sem $k$ er heil tala &  &  \\
\hline
2547 & jöfnuður &  &  &  \\
\hline
2548 & jöfnuhneppi &  &  &  \\
\hline
2549 & Jordan-ferill, einfaldur lokaður ferill &  &  &  \\
\hline
2550 & kápubrún, möttulbrún &  &  &  \\
\hline
2551 & kápuflatarmál, möttulflatarmál &  &  &  \\
\hline
2552 & kápuflötur, kápa, möttull, möttulflötur &  &  &  \\
\hline
2553 & kápuhlið, kápuhliðarflötur, möttulhlið, möttulhliðarflötur &  &  &  \\
\hline
2554 & kartesk hnit, Kartesarhnit & <ref idTerm=2145>Hnitakerfi</ref> á <ref idTerm=1031>evklíðska sléttu</ref> sem fæst með því að setja upp <ref idTerm=6469>þverstæðar</ref> <ref idTerm=5416>talnalínur</ref> sem kallast \textit{hnitaásar} og skerast í núllpunktum sínum. Þessi <ref idTerm=4845>skurðpunktur</ref> kallast þá \textit{upphafspunktur} hnitakerfisins sem er yfirleitt táknaður með $O$ og venja er að kalla aðra talnalínuna $x$-ás og hina $y$-ás &  & Er þetta nógu nákvæm skilgreining? \\
\hline
2556 & karteskt rúm, Kartesarrúm &  &  &  \\
\hline
2557 & karteskur ferningur &  &  &  \\
\hline
2559 & kassi, kubbur, rétthyrndur samhliðungur & <ref idTerm=4625>Sexflötungur</ref> þar sem <ref idTerm=65>aðlægar hliðar</ref> mynda <ref idTerm=4167>rétt horn</ref> &  &  \\
\hline
2560 & kast &  &  &  \\
\hline
2561 & keðja &  &  &  \\
\hline
2562 & keðjuályktun &  &  &  \\
\hline
2563 & keðjubrot &  &  &  \\
\hline
2564 & keðjubrotsliðun, liðun í keðjubrot &  &  &  \\
\hline
2565 & keðjuferill &  &  &  \\
\hline
2566 & keðjufléttingur &  &  &  \\
\hline
2567 & keðjuflötur &  &  &  \\
\hline
2568 & keðjugrúpa &  &  &  \\
\hline
2569 & keðjuherping &  &  &  \\
\hline
2570 & keðjujaðar &  &  &  \\
\hline
2571 & keðjujafngildi &  &  &  \\
\hline
2572 & keðjuregla &  &  &  \\
\hline
2573 & keðjusamtoga &  &  &  \\
\hline
2574 & keðjusamtogun &  &  &  \\
\hline
2575 & keðjuskilyrði &  &  &  \\
\hline
2576 & keila & ... &  & Almenn keila þar sem punktar flöts með sveigðum jaðri eru tengdir með línustriki í einn punkt. Líka hægt að segja að formið afmarkist af skífu og fletinum (sem kallast víst líka möttull eins og í sívalningi?) sem fæst með að tengja punkta hringsins við einn annan punkt (má punkturinn ekki liggja í grunnfletinum?) \\
\hline
2577 & keilu- &  &  &  \\
\hline
2578 & keilu- &  &  &  \\
\hline
2579 & keiluflötur &  &  &  \\
\hline
2580 & keilulaga &  &  &  \\
\hline
2581 & keilulaga &  &  &  \\
\hline
2582 & keilusnið með miðju, miðpunktskeilusnið &  &  &  \\
\hline
2583 & keilustúfur &  &  &  \\
\hline
2584 & keilustúfur, stýfð keila &  &  &  \\
\hline
2585 & kennifall &  &  &  \\
\hline
2587 & kenniferill &  &  &  \\
\hline
2588 & kenniferill &  &  &  \\
\hline
2589 & kenniflokkur &  &  &  \\
\hline
2590 & kenniflötur &  &  &  \\
\hline
2591 & kenniflötur &  &  &  \\
\hline
2592 & kennigrúpa, nykurgrúpa &  &  &  \\
\hline
2593 & kennijafna &  &  &  \\
\hline
2594 & kennileg hlutgrúpa &  &  &  \\
\hline
2595 & kenning &  &  &  \\
\hline
2596 & kenning um mikil frávik &  &  &  \\
\hline
2597 & kennir, grúpukennir, kennimótun &  &  &  \\
\hline
2598 & kennir, kennimótun &  &  &  \\
\hline
2599 & kennirót &  &  &  \\
\hline
2600 & kennitala &  &  &  \\
\hline
2601 & kennitala &  &  &  \\
\hline
2602 & kennitala sviðs &  &  &  \\
\hline
2603 & keppnisnet &  &  &  \\
\hline
2604 & kerfi &  &  &  \\
\hline
2605 & kerfisbundin skekkja, bjagi, vilhalli, hliðrunarskekkja, hlutdrægni &  &  &  \\
\hline
2606 & kerfisbundin úrtakskönnun &  &  &  \\
\hline
2607 & kerfisgreining &  &  &  \\
\hline
2608 & kí-dreifing, \(\chi\)-dreifing &  &  &  \\
\hline
2609 & kí-í-öðru-dreifing, \(\chi^2\)-dreifing &  &  &  \\
\hline
2610 & kí-í-öðru-próf, \(\chi^2\)-próf &  &  &  \\
\hline
2611 & kíkisröð &  &  &  \\
\hline
2612 & kím &  &  &  \\
\hline
2613 & kínversk talnagrind, kínversk kúlnagrind, kínversk reiknigrind &  &  &  \\
\hline
2614 & kínverska leifasetningin &  &  &  \\
\hline
2615 & kjálki &  &  &  \\
\hline
2616 & kjarnafall &  &  &  \\
\hline
2617 & kjarngóður virki &  &  &  \\
\hline
2618 & kjarngott rúm &  &  &  \\
\hline
2619 & kjarni &  &  &  \\
\hline
2620 & kjarni &  &  &  \\
\hline
2621 & kjörgildi, besta gildi &  &  &  \\
\hline
2622 & kjörgildisskilyrði, útgildisskilyrði &  &  &  \\
\hline
2623 & kjörlausn, besta lausn &  &  &  \\
\hline
2624 & kjörpunktur, besti punktur &  &  &  \\
\hline
2625 & klasaúrtak, hreiðrað úrtak &  &  &  \\
\hline
2626 & klasaúrtaka &  &  &  \\
\hline
2627 & klasi &  &  &  \\
\hline
2628 & klasi &  &  &  \\
\hline
2629 & Klein-flaska &  &  &  \\
\hline
2630 & Klein-grúpa &  &  &  \\
\hline
2631 & klifurferill &  &  &  \\
\hline
2632 & klofin fleyguð lest &  &  &  \\
\hline
2633 & klofin útvíkkun, klofin víkkun &  &  &  \\
\hline
2634 & klofinn &  &  &  \\
\hline
2635 & klofna, þættast &  &  &  \\
\hline
2636 & klofningssvið &  &  &  \\
\hline
2637 & klofningur, klofnun &  &  &  \\
\hline
2638 & knappasti tæmandi metill &  &  &  \\
\hline
2639 & knippafræði &  &  &  \\
\hline
2640 & knippamótun &  &  &  \\
\hline
2641 & knippi &  &  &  \\
\hline
2642 & knippisrúm &  &  &  \\
\hline
2643 & koma á eftir &  &  &  \\
\hline
2644 & koma á undan &  &  &  \\
\hline
2645 & komma &  &  &  \\
\hline
2646 & komma, tugakomma &  &  &  \\
\hline
2647 & kompáslína &  &  &  \\
\hline
2648 & kompássjálfmótun, kompásfærsla &  &  &  \\
\hline
2649 & könnun &  &  &  \\
\hline
2650 & kórbundin &  &  &  \\
\hline
2651 & kórfærsla, kórhnitaskipti &  &  &  \\
\hline
2652 & kórgerð &  &  &  \\
\hline
2653 & kórgerð Jordans, Jordan-kórgerð, Jordan-gerð &  &  &  \\
\hline
2654 & kórsamoka &  &  &  \\
\hline
2655 & kortalitun &  &  &  \\
\hline
2656 & kósekans &  &  &  \\
\hline
2657 & kósínus &  &  &  \\
\hline
2658 & kósínusregla & Um <ref idTerm=6438>þríhyrning</ref> $\Delta A B C$ með <ref idTerm=0>hliðarlengdir</ref> $a$, $b$ og $c$ gildir \textit{kósínusreglan} $$ c^2 = a^2 + b^2 - 2 a b\, \cos(C). $$ Þegar <ref idTerm=2198>hornið</ref> $C$ er <ref idTerm=4167>rétt</ref>, þá er síðasti liðurinn $0$ og jafnan gefur <ref idTerm=4002>reglu Pýþagórasar</ref> &  &  \\
\hline
2659 & kostnaðarfall, viðfangsfall &  &  &  \\
\hline
2660 & kóta, dulrita, lykla &  &  &  \\
\hline
2661 & kótangens &  &  &  \\
\hline
2662 & kóti, lykill, dulmálslykill, lykilmál &  &  &  \\
\hline
2663 & kótun, dulritun, lyklun &  &  &  \\
\hline
2664 & kótunarfræði, lyklunarfræði, dulmálsfræði &  &  &  \\
\hline
2665 & krafa &  &  &  \\
\hline
2666 & krappageisli, sveigjugeisli &  &  &  \\
\hline
2667 & krappamiðja, sveigjumiðja &  &  &  \\
\hline
2668 & krappamiðjuferill &  &  &  \\
\hline
2669 & krappaþinur &  &  &  \\
\hline
2670 & krappi, sveigja, dás ({\it hk.}) &  &  &  \\
\hline
2671 & krefjast, gefa sem forsendu &  &  &  \\
\hline
2672 & kristalgrúpa &  &  &  \\
\hline
2673 & Kronecker-tákn &  &  &  \\
\hline
2674 & krosshlutfall &  &  &  \\
\hline
2675 & krosshúfa &  &  &  \\
\hline
2676 & krossmargfeldi, vigurmargfeldi, vektormargfeldi, krossfeldi, vigurfeldi &  &  &  \\
\hline
2677 & krosspunktur &  &  &  \\
\hline
2678 & kúffeldi &  &  &  \\
\hline
2679 & kúfur &  &  &  \\
\hline
2680 & kúla, hnöttur, knöttur & <ref idTerm=6532>Rúmmynd</ref> sem afmarkast af <ref idTerm=2319>kúluhveli</ref> &  &  \\
\hline
2681 & kúla, reiknikúla &  &  &  \\
\hline
2682 & kúlnarúmfræði &  &  &  \\
\hline
2683 & kúlnatroðsla &  &  &  \\
\hline
2684 & kúlu-, hnatt-, hvel- &  &  &  \\
\hline
2685 & kúlubelti &  &  &  \\
\hline
2686 & kúlufall, hvelfall &  &  &  \\
\hline
2687 & kúlufleygur &  &  &  \\
\hline
2688 & kúluflutningur, kúluhreyfing &  &  &  \\
\hline
2689 & kúlugeiri &  &  &  \\
\hline
2690 & kúlugrennd &  &  &  \\
\hline
2691 & kúluhattur, kúluhetta &  &  &  \\
\hline
2692 & kúluhelmingur, hálfkúla &  &  &  \\
\hline
2693 & kúluhnit &  &  &  \\
\hline
2694 & kúluhorn, hnatthorn &  &  &  \\
\hline
2695 & kúluhornafræði &  &  &  \\
\hline
2696 & kúluhringgeisli &  &  &  \\
\hline
2697 & kúluhringmiðja &  &  &  \\
\hline
2698 & kúluhverfing punkts, spegilmynd punkts í kúlu &  &  &  \\
\hline
2699 & kúluhverfing, kúluspeglun, speglun í kúlu &  &  &  \\
\hline
2700 & kúlulíki &  &  &  \\
\hline
2701 & kúlumarghyrningur, kúluhyrningur &  &  &  \\
\hline
2702 & kúluofanvarp &  &  &  \\
\hline
2703 & kúlurúmfræði &  &  &  \\
\hline
2704 & kúlusamhverfni &  &  &  \\
\hline
2705 & kúluskel &  &  &  \\
\hline
2706 & kúluskífa &  &  &  \\
\hline
2707 & kúlustrýta, kúlupíramíti &  &  &  \\
\hline
2708 & kúluþríhyrningur &  &  &  \\
\hline
2709 & kúpni &  &  &  \\
\hline
2710 & kúpniskilyrði &  &  &  \\
\hline
2711 & kúpt fall &  &  &  \\
\hline
2712 & kúpt línuleg samantekt &  &  &  \\
\hline
2713 & kúpt mengi &  &  &  \\
\hline
2714 & kúptur &  &  &  \\
\hline
2715 & kúptur ferill &  &  &  \\
\hline
2716 & kúptur hjúpur &  &  &  \\
\hline
2717 & kúptur marghyrningur, úthyrndur marghyrningur &  &  &  \\
\hline
2718 & kúptur þykkvi &  &  &  \\
\hline
2719 & kúvending &  &  &  \\
\hline
2720 & kvarðagrúpa, kvörðunargrúpa &  &  &  \\
\hline
2721 & kvarðahlutfall &  &  &  \\
\hline
2722 & kvarðaskipting &  &  &  \\
\hline
2723 & kvarði &  &  &  \\
\hline
2724 & kvarði, mælikvarði, stigi &  &  &  \\
\hline
2725 & kvísl, grein &  &  &  \\
\hline
2726 & kvíslað fall, marggilt fall &  &  &  \\
\hline
2727 & kvíslað þekjurúm &  &  &  \\
\hline
2728 & kvíslaður &  &  &  \\
\hline
2729 & kvíslideilir &  &  &  \\
\hline
2730 & kvísliferli, kvíslunarferli &  &  &  \\
\hline
2731 & kvíslipunktur, kvíslunarpunktur &  &  &  \\
\hline
2732 & kvíslisnið &  &  &  \\
\hline
2733 & kvíslivísir, kvíslunarvísir &  &  &  \\
\hline
2734 & kvísluð lausn &  &  &  \\
\hline
2736 & kvíslun, forkun, tvíkvíslun &  &  &  \\
\hline
2737 & kvíslunarfræði, forkunarfræði &  &  &  \\
\hline
2738 & kvörðun &  &  &  \\
\hline
2739 & kvörðunaraðferð &  &  &  \\
\hline
2740 & kvóti, hlutfall, hlutatala &  &  &  \\
\hline
2741 & kyrr, kyrra-, kyrrðar-, fastur, fasta- &  &  &  \\
\hline
2742 & kyrragrúpa &  &  &  \\
\hline
2743 & kyrralína &  &  &  \\
\hline
2744 & kyrrapunktssetning Banachs, fastapunktssetning Banachs, herpingarsetning Banachs &  &  &  \\
\hline
2745 & kyrrapunktssetning, fastapunktssetning &  &  &  \\
\hline
2746 & kyrrapunktur, kyrr punktur, fastapunktur &  &  &  \\
\hline
2747 & kyrrastak, kyrrt stak, fastastak &  &  &  \\
\hline
2748 & kyrrasvið, kyrrapunktasvið, fastapunktasvið &  &  &  \\
\hline
2749 & l'Hôpital-regla &  &  &  \\
\hline
2750 & lægri miðjuruna, lægri miðjuturn &  &  &  \\
\hline
2751 & lægsta stak, neðsta stak, fyrsta stak, minnsta stak &  &  &  \\
\hline
2752 & lækka niður &  &  &  \\
\hline
2753 & lækkun &  &  &  \\
\hline
2754 & lækkunarskekkja &  &  &  \\
\hline
2755 & lag &  &  &  \\
\hline
2756 & lág-, óminnkanlegur, ósmækkanlegur &  &  &  \\
\hline
2757 & lágflötur &  &  &  \\
\hline
2758 & lágfrumíðal, óminnkanlegt frumíðal &  &  &  \\
\hline
2760 & lággildispunktur &  &  &  \\
\hline
2762 & lággildissetning, lággildislögmál, lágalgildissetning, lágalgildislögmál &  &  &  \\
\hline
2763 & lággrunnur, óminnkanlegur grunnur, ósmækkanlegur grunnur &  &  &  \\
\hline
2764 & lágíðal, óminnkanlegt íðal, ósmækkanlegt íðal &  &  &  \\
\hline
2765 & lágjafna &  &  &  \\
\hline
2766 & lágmargliða &  &  &  \\
\hline
2767 & lágmarka &  &  &  \\
\hline
2768 & lágmarksleikáætlun, lágmarksáætlun &  &  &  \\
\hline
2769 & lágmörkun &  &  &  \\
\hline
2770 & lagnet &  &  &  \\
\hline
2771 & lágnet &  &  &  \\
\hline
2772 & lágreisn &  &  &  \\
\hline
2773 & lágreistur &  &  &  \\
\hline
2774 & lagsfylki &  &  &  \\
\hline
2775 & lagshorn &  &  &  \\
\hline
2776 & lagshorn & Tvö <ref idTerm=2198>horn</ref> sem samanlagt <ref idTerm=2223>mælast</ref> 90 <ref idTerm=1667>gráður</ref> &  & Hugtakabankinn segir að hornin þurfa að vera aðlæg en Wikipedia gerir enga slíka kröfu \\
\hline
2777 & lagskipt úrtak &  &  &  \\
\hline
2778 & lagskipting &  &  &  \\
\hline
2779 & lágstak, óminnkanlegt stak, ósmækkanlegt stak &  &  &  \\
\hline
2780 & lágvísir, knévísir &  &  &  \\
\hline
2781 & láhnit &  &  &  \\
\hline
2782 & landabréf, kort &  &  &  \\
\hline
2783 & Landau-tákn &  &  &  \\
\hline
2784 & landmæling, landmælingafræði &  &  &  \\
\hline
2785 & langadeiling &  &  &  \\
\hline
2786 & langalína &  &  &  \\
\hline
2787 & langamargföldun &  &  &  \\
\hline
2788 & langás &  &  &  \\
\hline
2789 & langhlið &  &  &  \\
\hline
2790 & langsveifarferill &  &  &  \\
\hline
2791 & Laplace-færsla, Laplace-myndun, Laplace-ummyndun &  &  &  \\
\hline
2792 & Laplace-liðun &  &  &  \\
\hline
2793 & Laplace-mynd, Laplace-færsla &  &  &  \\
\hline
2794 & Laplace-virki &  &  &  \\
\hline
2795 & láréttur &  &  &  \\
\hline
2796 & lásnertill &  &  &  \\
\hline
2797 & láta kyrrt &  &  &  \\
\hline
2798 & latneskur ferningur, rómverskur ferningur &  &  &  \\
\hline
2799 & lauf &  &  &  \\
\hline
2800 & laufaður &  &  &  \\
\hline
2801 & laufferill Descartes &  &  &  \\
\hline
2802 & laufun &  &  &  \\
\hline
2803 & Laurent-röð &  &  &  \\
\hline
2804 & lausn &  &  &  \\
\hline
2805 & lausn með rótardrætti &  &  &  \\
\hline
2806 & lausn, úrlausn &  &  &  \\
\hline
2807 & lausnagrunnur &  &  &  \\
\hline
2808 & lausnamengi &  &  &  \\
\hline
2809 & lausnarferill, heildisferill &  &  &  \\
\hline
2810 & lausnarflötur, heildisflötur &  &  &  \\
\hline
2811 & lausnarformúla &  &  &  \\
\hline
2812 & lausnartilgáta &  &  &  \\
\hline
2813 & lausnarúm &  &  &  \\
\hline
2814 & láþverill &  &  &  \\
\hline
2815 & Lebesgue-heildanlegur, Lebesgue-tegranlegur &  &  &  \\
\hline
2816 & Lebesgue-heildi, Lebesgue-tegur &  &  &  \\
\hline
2817 & Lebesgue-mál &  &  &  \\
\hline
2818 & leg &  &  &  \\
\hline
2819 & lega &  &  &  \\
\hline
2820 & leggja saman &  &  &  \\
\hline
2821 & leggja við, bæta við &  &  &  \\
\hline
2822 & leggjalitun &  &  &  \\
\hline
2823 & leggjamengi &  &  &  \\
\hline
2824 & leggjaruna, vegur &  &  &  \\
\hline
2825 & leggur &  &  &  \\
\hline
2826 & legufrumsenda, legufrumsetning &  &  &  \\
\hline
2827 & legufylki &  &  &  \\
\hline
2828 & legumynstur &  &  &  \\
\hline
2829 & legurúmfræði &  &  &  \\
\hline
2830 & leguvensl &  &  &  \\
\hline
2831 & Leibniz-próf, víxlmerkjapróf Leibniz &  &  &  \\
\hline
2832 & Leibniz-röð &  &  &  \\
\hline
2833 & leið &  &  &  \\
\hline
2834 & leiða af &  &  &  \\
\hline
2835 & leiðiliður &  &  &  \\
\hline
2836 & leiðing &  &  &  \\
\hline
2837 & leiðingarmerki &  &  &  \\
\hline
2838 & leiðingarregla &  &  &  \\
\hline
2839 & leiðréttandi &  &  &  \\
\hline
2840 & leiðrétting &  &  &  \\
\hline
2841 & leiðréttingarliður &  &  &  \\
\hline
2842 & leiðréttingarþáttur &  &  &  \\
\hline
2843 & leif &  &  &  \\
\hline
2844 & leif &  &  &  \\
\hline
2845 & leifaflokkasvið &  &  &  \\
\hline
2846 & leifaflokkur &  &  &  \\
\hline
2847 & leifajafna, samleifing &  &  &  \\
\hline
2848 & leifakerfi &  &  &  \\
\hline
2849 & leifareikningur &  &  &  \\
\hline
2850 & leifasetning &  &  &  \\
\hline
2851 & leifastofn &  &  &  \\
\hline
2852 & leikáætlun, áætlun &  &  &  \\
\hline
2853 & leikmaður, keppandi &  &  &  \\
\hline
2854 & leikregla &  &  &  \\
\hline
2855 & leiktré &  &  &  \\
\hline
2856 & leikur &  &  &  \\
\hline
2857 & leit &  &  &  \\
\hline
2858 & leitaraðferð &  &  &  \\
\hline
2859 & leitni, hneigð &  &  &  \\
\hline
2860 & lengd &  &  &  \\
\hline
2861 & lengd, algildi, tölugildi &  &  &  \\
\hline
2862 & lengd, staðall, tölugildi &  &  &  \\
\hline
2863 & lengdarbaugur &  &  &  \\
\hline
2864 & lengdarfrymi &  &  &  \\
\hline
2865 & lengdarsnið &  &  &  \\
\hline
2866 & leppvísir &  &  &  \\
\hline
2867 & lest &  &  &  \\
\hline
2868 & leyfilegur &  &  &  \\
\hline
2869 & leyfilegur &  &  &  \\
\hline
2870 & leysa &  &  &  \\
\hline
2871 & leysanleg grúpa &  &  &  \\
\hline
2872 & leysanleg jafna &  &  &  \\
\hline
2873 & leysanlegur &  &  &  \\
\hline
2874 & leysanleiki &  &  &  \\
\hline
2875 & leysifylki &  &  &  \\
\hline
2876 & leysijafna &  &  &  \\
\hline
2877 & leysimengi &  &  &  \\
\hline
2878 & leysivirki &  &  &  \\
\hline
2879 & lið fyrir lið &  &  &  \\
\hline
2880 & liða &  &  &  \\
\hline
2881 & liðagreining &  &  &  \\
\hline
2882 & liðhnútur &  &  &  \\
\hline
2883 & liðun &  &  &  \\
\hline
2884 & liðun &  &  &  \\
\hline
2885 & liðun í beina summu &  &  &  \\
\hline
2886 & liðun í Dirichlet-röð &  &  &  \\
\hline
2887 & liðun í rásir &  &  &  \\
\hline
2888 & liðun í röð, raðarliðun &  &  &  \\
\hline
2889 & liður &  &  &  \\
\hline
2890 & liður, þáttur &  &  &  \\
\hline
2891 & Lie-algebra &  &  &  \\
\hline
2892 & Lie-grúpa &  &  &  \\
\hline
2893 & Lie-hornklofi, Lie-margfeldi &  &  &  \\
\hline
2894 & líftölfræði &  &  &  \\
\hline
2895 & líkan &  &  &  \\
\hline
2896 & líkanafræði &  &  &  \\
\hline
2897 & líkanagerð &  &  &  \\
\hline
2898 & líkindafall &  &  &  \\
\hline
2899 & líkindafræði &  &  &  \\
\hline
2900 & líkindafræðingur &  &  &  \\
\hline
2901 & líkindaframleiðandi fall &  &  &  \\
\hline
2902 & líkindafylki &  &  &  \\
\hline
2903 & líkindamál &  &  &  \\
\hline
2904 & líkindareikningur &  &  &  \\
\hline
2905 & líkindarúm &  &  &  \\
\hline
2906 & líkindavísir &  &  &  \\
\hline
2907 & líkindi, líkur &  &  &  \\
\hline
2908 & líklegt frávik &  &  &  \\
\hline
2909 & líming &  &  &  \\
\hline
2910 & lína &  &  &  \\
\hline
2911 & lína, bein lína &  &  &  \\
\hline
2912 & lína, ferill &  &  &  \\
\hline
2913 & línuaðgerð &  &  &  \\
\hline
2914 & línubundin &  &  &  \\
\hline
2915 & línuflötur &  &  &  \\
\hline
2916 & línufylki, einlínufylki &  &  &  \\
\hline
2917 & línugera &  &  &  \\
\hline
2918 & línugerningur &  &  &  \\
\hline
2919 & línuleg algebra &  &  &  \\
\hline
2920 & línuleg bestun &  &  &  \\
\hline
2921 & línuleg deildajafna, línuleg diffurjafna &  &  &  \\
\hline
2922 & línuleg fellagreining, línuleg fallagreining &  &  &  \\
\hline
2923 & línuleg grúpa &  &  &  \\
\hline
2924 & línuleg hlutafleiðujafna &  &  &  \\
\hline
2925 & línuleg jafna &  &  &  \\
\hline
2926 & línuleg lokun &  &  &  \\
\hline
2927 & línuleg röðun, fullröðun &  &  &  \\
\hline
2928 & línuleg samantekt &  &  &  \\
\hline
2929 & línuleg spönn, línulegur hjúpur &  &  &  \\
\hline
2930 & línuleg útsetning &  &  &  \\
\hline
2931 & línuleg vörpun, línuleg færsla &  &  &  \\
\hline
2932 & línulega háður, línulega venslaður &  &  &  \\
\hline
2933 & línulega óháður, línulega óvenslaður &  &  &  \\
\hline
2934 & línulega raðað mengi, fullraðað mengi, keðja &  &  &  \\
\hline
2935 & línulegt aðhvarf &  &  &  \\
\hline
2936 & línulegt fall &  &  &  \\
\hline
2937 & línulegt fall, línulegt form &  &  &  \\
\hline
2938 & línulegt felli &  &  &  \\
\hline
2939 & línulegt grannmynstur &  &  &  \\
\hline
2940 & línulegt grannrúm &  &  &  \\
\hline
2941 & línulegt hæði &  &  &  \\
\hline
2942 & línulegt hlutrúm &  &  &  \\
\hline
2943 & línulegt jöfnuhneppi &  &  &  \\
\hline
2944 & línulegt kerfi &  &  &  \\
\hline
2945 & línulegt óhæði &  &  &  \\
\hline
2946 & línulegt rúm, vigurrúm, vektorrúm &  &  &  \\
\hline
2947 & línulegur &  &  &  \\
\hline
2948 & línulegur á köflum &  &  &  \\
\hline
2949 & línulegur grunnur, vigurgrunnur, Hamel-grunnur, grunnur &  &  &  \\
\hline
2950 & línulegur kvarði &  &  &  \\
\hline
2951 & línulegur virki &  &  &  \\
\hline
2952 & línuleiki &  &  &  \\
\hline
2953 & línumeðaltal &  &  &  \\
\hline
2954 & linur &  &  &  \\
\hline
2955 & graf, línurit, fallrit, varprit, ferill (falls) & Fyrir <ref idTerm=3282>vörpun</ref> $f: X \to Y$ kallast <ref idTerm=3270>mengi</ref> allra <ref idTerm=5627>tvennda</ref> af gerðinni $(x,f(x))$, þar sem $x \in X$ \textit{graf} vörpunarinnar $f$ og er táknað með \[ G_f = \{(x,f(x)) \mid x \in X\} \] &  & Vantar þýðinguna "graf" Tala um gröf falla frekar? Láta fylgja myndir af gröfum algengra falla í skýringu \\
\hline
2956 & línustallagerð, línuþrepagerð, línuþrepaskipan &  &  &  \\
\hline
2957 & línustétt &  &  &  \\
\hline
2958 & línustrik, strik &  &  & Skilgreina endapunkt hér \\
\hline
2959 & línusumma, línusamtala &  &  &  \\
\hline
2960 & línuuppstokkun, línuumröðun &  &  &  \\
\hline
2961 & línuvigur &  &  &  \\
\hline
2962 & línuvísir &  &  &  \\
\hline
2963 & lit-, lita- &  &  &  \\
\hline
2964 & litamargliða &  &  &  \\
\hline
2965 & litatala &  &  &  \\
\hline
2966 & lítið ríki, smáríki &  &  &  \\
\hline
2967 & lítið úrtak &  &  &  \\
\hline
2968 & lítill &  &  &  \\
\hline
2969 & litla Fermat-setningin &  &  &  \\
\hline
2970 & litun &  &  &  \\
\hline
2971 & litvísir, litaflokkur &  &  &  \\
\hline
2972 & ljósgeislafræði &  &  &  \\
\hline
2973 & lóðhnit &  &  &  \\
\hline
2974 & loðið mengi &  &  &  \\
\hline
2975 & lóðlína, þverlína, þverill &  &  &  \\
\hline
2976 & lóðmengi &  &  &  \\
\hline
2977 & lóðrétt ofanvarp &  &  &  \\
\hline
2978 & lóðréttur &  &  &  \\
\hline
2979 & lóðréttur &  &  &  \\
\hline
2980 & loftfall &  &  &  \\
\hline
2981 & lögbundinn, orsaka-, orsakar- &  &  &  \\
\hline
2982 & löggeng bestun &  &  &  \\
\hline
2983 & löggengur &  &  &  \\
\hline
2984 & lögmál Bortkiewicz, markgildissetning Poissons, lögmál lítils fjölda, lögmál sjaldgæfra atburða &  &  &  \\
\hline
2985 & lögmál mikils fjölda, meðalgildislögmál &  &  &  \\
\hline
2986 & lögmál minnstu verkunar &  &  &  \\
\hline
2987 & lögmál um annað tveggja &  &  &  \\
\hline
2988 & lögmál um ofurendanlega þrepun &  &  &  \\
\hline
2989 & lögmál um stærðfræðilega þrepun, þrepunarlögmál &  &  &  \\
\hline
2990 & lögmál um takmörkun í jöfnum mæli, lögmál um skorðun í jöfnum mæli, jafnmælisskorðunarlögmál &  &  &  \\
\hline
2991 & lögmál, regla &  &  &  \\
\hline
2992 & logra- &  &  &  \\
\hline
2993 & lograafleiða &  &  &  \\
\hline
2994 & lografall &  &  &  \\
\hline
2995 & lograheildi, logrategur &  &  &  \\
\hline
2996 & lograkúptur &  &  &  \\
\hline
2997 & lograkvarði &  &  &  \\
\hline
2998 & lograkvíslipunktur, lograkvíslunarpunktur &  &  &  \\
\hline
2999 & lograpappír &  &  &  \\
\hline
3000 & lograpróf &  &  &  \\
\hline
3001 & lograstuðull &  &  &  \\
\hline
3002 & logratafla &  &  &  \\
\hline
3003 & logravefja, lograkuðungur, lograkuðungslína &  &  &  \\
\hline
3004 & logri &  &  &  \\
\hline
3005 & logri með grunntölu, logri með stofntölu &  &  &  \\
\hline
3006 & lögun, snið, gerð, mót &  &  &  \\
\hline
3007 & loka- &  &  &  \\
\hline
3008 & lokaákvörðun &  &  &  \\
\hline
3009 & lokaástand &  &  &  \\
\hline
3010 & lokað bil & <ref idTerm=457>Bil</ref> sem inniheldur endapunkta sína &  &  \\
\hline
3011 & lokað deildaform, lokað diffurform &  &  &  \\
\hline
3012 & lokað hálfrúm &  &  &  \\
\hline
3013 & lokað hlutrúm &  &  &  \\
\hline
3014 & lokað mengi, lokað hlutmengi &  &  &  \\
\hline
3015 & lokaður &  &  &  \\
\hline
3016 & lokaður ferill, lokaður vegur, lykkja & <ref idTerm=1131>Ferill</ref> þar sem báðir endapunktar ferilsins eru sami punkturinn &  &  \\
\hline
3017 & lokaður flötur, þjappaður flötur &  &  &  \\
\hline
3018 & lokaður gagnvegur &  &  &  \\
\hline
3019 & lokaður hjúpur, lokun &  &  &  \\
\hline
3020 & lokaður kúptur hjúpur, kúpt lokun &  &  &  \\
\hline
3021 & lokaður línulegur virki &  &  &  \\
\hline
3022 & lokaður með tilliti til margföldunar, margföldunarlokaður &  &  &  \\
\hline
3023 & lokaður Riemann-flötur, þjappaður Riemann-flötur &  &  &  \\
\hline
3024 & lokaður straumur &  &  &  \\
\hline
3025 & lokagrannmynstur &  &  &  \\
\hline
3026 & lokahlutur &  &  &  \\
\hline
3027 & lokahnútur &  &  &  \\
\hline
3028 & lokajafna &  &  &  \\
\hline
3029 & lokanlegur &  &  &  \\
\hline
3030 & lokapunktur &  &  &  \\
\hline
3031 & lokastaða &  &  &  \\
\hline
3032 & lokastak &  &  &  \\
\hline
3033 & lokuð fullyrðing &  &  &  \\
\hline
3034 & lokuð hringskífa &  &  &  \\
\hline
3035 & lokuð kúla &  &  &  \\
\hline
3036 & lokuð segð, lokuð fullyrðing, lokuð yrðing & <ref idTerm=1515>Yrðing</ref> sem inniheldur engar óþekktar <ref idTerm=5102>stærðir</ref> &  &  \\
\hline
3037 & lokuð vörpun &  &  &  \\
\hline
3038 & lokuð þakning &  &  &  \\
\hline
3039 & lokunarvirki &  &  &  \\
\hline
3040 & lota, raðstig &  &  &  \\
\hline
3041 & lota, umferð &  &  &  \\
\hline
3042 & lotubinding &  &  &  \\
\hline
3043 & lotubundið flökt &  &  &  \\
\hline
3044 & lotubundið keðjubrot &  &  &  \\
\hline
3045 & lotubundin framlenging &  &  &  \\
\hline
3046 & lotubundin Markov-keðja, rásuð Markov-keðja &  &  &  \\
\hline
3047 & lotubundin úrtakskönnun &  &  &  \\
\hline
3048 & lotubundinn, umferðar- &  &  &  \\
\hline
3049 & lotufall, lotubundið fall & <ref idTerm=1078>Fall</ref> $f: \mathbb{R} \to \mathbb{R}$ er sagt vera \textit{lotubundið} ef til er <ref idTerm=4093>rauntala</ref> $p > 0$ þannig að fyrir öll $x \in \mathbb{R}$ gildi: \[ f(x + p) = f(x). \] Talan $p$ kallast þá \textit{lota} fallsins &  & Þarf skilgreiningarmengið að vera allt $\mathbb{R}$? \\
\hline
3050 & lotugrind &  &  &  \\
\hline
3051 & lotugrunnur &  &  &  \\
\hline
3052 & lotugrúpa &  &  &  \\
\hline
3053 & loturæma &  &  &  \\
\hline
3054 & lotusamhliðungur &  &  &  \\
\hline
3055 & lotusamsíðungur &  &  &  \\
\hline
3056 & lotusamsíðungur, lotusamhliðungur &  &  &  \\
\hline
3057 & lotusvæði &  &  &  \\
\hline
3058 & lotutugabrot, umferðartugabrot &  &  &  \\
\hline
3059 & lotutvennd &  &  &  \\
\hline
3060 & lúpa &  &  &  \\
\hline
3061 & lyfta &  &  &  \\
\hline
3062 & lygaraþversögn &  &  &  \\
\hline
3063 & lykkjuheildi, lykkjutegur, heildi yfir lokaðan feril, tegur yfir lokaðan feril &  &  &  \\
\hline
3064 & lykkjurúm &  &  &  \\
\hline
3065 & lýsandi tölfræði &  &  &  \\
\hline
3066 & mælanleg vörpun &  &  &  \\
\hline
3067 & mælanlegt fall &  &  &  \\
\hline
3068 & mælanlegt mengi &  &  &  \\
\hline
3069 & mælanlegt rúm &  &  &  \\
\hline
3070 & mælanlegur &  &  &  \\
\hline
3071 & mælanleiki &  &  &  \\
\hline
3072 & mælieining & \textit{Mælieining} er ákveðið <ref idTerm=3081>magn</ref> tiltekinnar <ref idTerm=5102>stærðar</ref> sem er notað sem viðmið við mælingu slíkra stærða &  & Þarf að bæta við að allar mælingar eru þá margfeldi einingarinnar eða setja það í skýringu eða segir það sig kannski sjálft? \\
\hline
3073 & mælingarskekkja &  &  &  \\
\hline
3074 & mælirúmfræði &  &  &  \\
\hline
3075 & mæliskekkja, athugunarskekkja &  &  &  \\
\hline
3076 & mælistika &  &  &  \\
\hline
3077 & mætti &  &  &  \\
\hline
3078 & mættisfall &  &  &  \\
\hline
3079 & mættisfræði &  &  &  \\
\hline
3080 & magn &  &  &  \\
\hline
3081 & magn &  &  &  \\
\hline
3082 & magnaralaus &  &  &  \\
\hline
3083 & magnari, magnvirki, kvantari, kvöðull &  &  &  \\
\hline
3084 & magur &  &  &  \\
\hline
3085 & magurt mengi, magurt hlutmengi &  &  &  \\
\hline
3086 & mál &  &  &  \\
\hline
3087 & mál fyrstu stéttar &  &  &  \\
\hline
3088 & mál, tungumál &  &  &  \\
\hline
3089 & málalgebra &  &  &  \\
\hline
3090 & málfræði, mælifræði &  &  &  \\
\hline
3091 & málrækinn, málheldinn &  &  &  \\
\hline
3092 & málrúm &  &  &  \\
\hline
3093 & málsamleitni, samleitni í máli &  &  &  \\
\hline
3094 & málskipan &  &  &  \\
\hline
3095 & málskipanarbreyta &  &  &  \\
\hline
3096 & málskipanarfræði &  &  &  \\
\hline
3097 & málskipti &  &  &  \\
\hline
3098 & margblaða flötur &  &  &  \\
\hline
3099 & margbreytuaðhvarf, fjölbreytuaðhvarf &  &  &  \\
\hline
3100 & margein samsvörun &  &  &  \\
\hline
3101 & margfalda &  &  &  \\
\hline
3102 & margfalda í kross &  &  &  \\
\hline
3103 & margfalda í kross &  &  &  \\
\hline
3104 & margfalda upp úr sviga &  &  &  \\
\hline
3105 & margfaldanlegur &  &  &  \\
\hline
3106 & margfaldari &  &  &  \\
\hline
3107 & margfaldlega &  &  &  \\
\hline
3108 & margfaldlega samanhangandi, margfaldlega samhangandi &  &  &  \\
\hline
3109 & margfaldur &  &  &  \\
\hline
3110 & margfaldur &  &  &  \\
\hline
3111 & margfaldur leggur &  &  &  \\
\hline
3112 & margfaldur punktur &  &  &  \\
\hline
3113 & margfaldur samanburður &  &  &  \\
\hline
3114 & margfaldur snertill &  &  &  \\
\hline
3115 & margfalt heildi, margfalt tegur &  &  &  \\
\hline
3116 & margfalt próf &  &  &  \\
\hline
3117 & margfeldi &  &  &  \\
\hline
3118 & margfeldi &  &  &  \\
\hline
3119 & margfeldi, feldi, fald &  &  &  \\
\hline
3120 & margfeldið fall &  &  &  \\
\hline
3121 & margfeldinn &  &  &  \\
\hline
3122 & margfeldisfrumsenda, margfeldisfrumsetning &  &  &  \\
\hline
3123 & margfeldisgrannmynstur, faldgrannmynstur &  &  &  \\
\hline
3124 & margfeldisregla &  &  &  \\
\hline
3125 & margfeldistákn, feldistákn, faldtákn &  &  &  \\
\hline
3126 & margfeldni &  &  &  \\
\hline
3127 & margfeldni &  &  &  \\
\hline
3128 & margfeldni &  &  &  \\
\hline
3129 & margflötungsgrúpa &  &  &  \\
\hline
3130 & margflötungshorn &  &  &  \\
\hline
3131 & margflötungsnet &  &  &  \\
\hline
3132 & margflötungstala &  &  &  \\
\hline
3133 & margflötungur, flötungur & <ref idTerm=6532>Rúmmynd</ref> sem afmarkast af <ref idTerm=1009>endanlega</ref> mörgum <ref idTerm=3161>marghyrningum</ref> &  &  \\
\hline
3141 & margföldun &  &  &  \\
\hline
3148 & margföldunargrúpa &  &  &  \\
\hline
3149 & margföldunarmerki &  &  &  \\
\hline
3150 & margföldunarsetning &  &  &  \\
\hline
3151 & margföldunarstofn &  &  &  \\
\hline
3152 & margföldunartafla &  &  &  \\
\hline
3154 & marggild rökfræði &  &  &  \\
\hline
3155 & marggild vörpun &  &  &  \\
\hline
3156 & marggildni &  &  &  \\
\hline
3157 & marggildur &  &  &  \\
\hline
3158 & marghliðungur &  &  &  \\
\hline
3159 & marghyrninganet &  &  &  \\
\hline
3160 & marghyrningstala &  &  &  \\
\hline
3161 & marghyrningur, hyrningur & <ref idTerm=4891>Flatarmynd</ref> með <ref idTerm=2453>jaðri</ref> sem samanstendur af <ref idTerm=1009>endanlega</ref> mörgum <ref idTerm=2958>línustrikum</ref> sem kallast \textit{hliðar} marghyrningsins &  & Hvað er hægt að segja annað en jaðar í tveimur víddum? \\
\hline
3162 & margkynjuð tala &  &  &  \\
\hline
3163 & margliða með mörgum breytum &  &  &  \\
\hline
3164 & margliðuaðhvarf &  &  &  \\
\hline
3165 & margliðubaugur &  &  &  \\
\hline
3166 & margliðufall &  &  &  \\
\hline
3167 & margliðuíðal &  &  &  \\
\hline
3168 & margliðujafna &  &  &  \\
\hline
3169 & margliðujafna &  &  &  \\
\hline
3170 & margliðuliðun &  &  &  \\
\hline
3171 & margliðunálgun &  &  &  \\
\hline
3172 & margliðuvöxtur &  &  &  \\
\hline
3173 & margliðuþáttun &  &  &  \\
\hline
3174 & marglínuleg algebra, fjöllínuleg algebra &  &  &  \\
\hline
3175 & marglínulegur, fjöllínulegur &  &  &  \\
\hline
3176 & margra manna spil &  &  &  \\
\hline
3177 & margræð jafna &  &  &  \\
\hline
3178 & margræðni &  &  &  \\
\hline
3179 & margrætt jöfnukerfi, margrætt jöfnuhneppi &  &  &  \\
\hline
3180 & margstæð umsögn, umsögn með mörgum breytum &  &  &  \\
\hline
3181 & margtækur &  &  &  \\
\hline
3182 & margvíð dreifing &  &  &  \\
\hline
3183 & margvíð dreifing &  &  &  \\
\hline
3184 & margvíð fervikagreining, fjölbreytufervikagreining &  &  &  \\
\hline
3185 & margvíð tölfræði &  &  &  \\
\hline
3186 & margviðbætinn, margsamleggjandi &  &  &  \\
\hline
3187 & margvíður &  &  &  \\
\hline
3188 & margvíður &  &  &  \\
\hline
3189 & margvítt rúm &  &  &  \\
\hline
3190 & margþætt tala &  &  &  \\
\hline
3191 & margþrepa aðferð &  &  &  \\
\hline
3192 & margþrepa bestun, margþrepabestun &  &  &  \\
\hline
3193 & margþrepa leikur &  &  &  \\
\hline
3194 & markfall, markgildisfall &  &  &  \\
\hline
3195 & markgildi &  &  &  \\
\hline
3196 & markgildi frá hægri, hægra markgildi &  &  &  \\
\hline
3197 & markgildi frá vinstri, vinstra markgildi &  &  &  \\
\hline
3198 & markgildi í jöfnum mæli &  &  &  \\
\hline
3199 & markgildispunktur &  &  &  \\
\hline
3200 & markgildissetning &  &  &  \\
\hline
3201 & markgildissetning Bernoullis, Bernoulli-lögmál mikils fjölda &  &  &  \\
\hline
3202 & markgildisstak &  &  &  \\
\hline
3203 & markgildistaka &  &  &  \\
\hline
3204 & markmeðaltal &  &  &  \\
\hline
3205 & Markov-ferli &  &  &  \\
\hline
3206 & Markov-keðja &  &  &  \\
\hline
3207 & Markov-reiknanlegur &  &  &  \\
\hline
3208 & markraðtala &  &  &  \\
\hline
3209 & markstaða &  &  &  \\
\hline
3210 & marktækur &  &  &  \\
\hline
3211 & marktækur tölustafur &  &  &  \\
\hline
3212 & marktekt &  &  &  \\
\hline
3213 & marktektarstig, prófstig &  &  &  \\
\hline
3214 & marktilfelli &  &  &  \\
\hline
3215 & marktökulíkur, prófstig &  &  &  \\
\hline
3216 & markverður &  &  &  \\
\hline
3217 & massi &  &  &  \\
\hline
3218 & mat &  &  &  \\
\hline
3219 & mat &  &  &  \\
\hline
3220 & mát &  &  &  \\
\hline
3221 & mat að ofan, yfirmat &  &  &  \\
\hline
3222 & mat, stikabundið mat &  &  &  \\
\hline
3223 & mátað við &  &  &  \\
\hline
3224 & mátfall &  &  &  \\
\hline
3225 & mátform &  &  &  \\
\hline
3226 & mátgæði &  &  &  \\
\hline
3227 & mátgrúpa &  &  &  \\
\hline
3228 & mátrúm &  &  &  \\
\hline
3229 & matsfræði &  &  &  \\
\hline
3230 & mátstikar &  &  &  \\
\hline
3231 & mátverkefni &  &  &  \\
\hline
3232 & með endanleg gildi &  &  &  \\
\hline
3233 & með sömu brennipunkta &  &  &  \\
\hline
3234 & meðal-, meðaltals-, að meðaltali &  &  &  \\
\hline
3235 & meðalferskekkja &  &  &  \\
\hline
3236 & meðalfervik úrtaks, úrtaksdreifni &  &  &  \\
\hline
3237 & meðalfervik, dreifni, brigði &  &  &  \\
\hline
3238 & meðalframleiðslugæði &  &  &  \\
\hline
3239 & meðalfrávik &  &  &  \\
\hline
3240 & meðalgildi &  &  &  \\
\hline
3241 & meðalgildiseiginleiki &  &  &  \\
\hline
3242 & meðalgildissetning &  &  &  \\
\hline
3243 & meðalkrappi, meðalsveigja &  &  &  \\
\hline
3244 & meðalmunur &  &  &  \\
\hline
3245 & meðalreisn &  &  &  \\
\hline
3246 & meðalreistur &  &  &  \\
\hline
3247 & meðalsamleitni, meðaltalssamleitni, samleitni í \(L_1\)-staðli &  &  &  \\
\hline
3248 & meðalskekkja &  &  &  \\
\hline
3249 & meðalstærð úrtaks &  &  &  \\
\hline
3251 & meðaltal, hreint meðaltal, venjulegt meðaltal &  &  &  \\
\hline
3252 & meðaltap &  &  &  \\
\hline
3253 & meðalþéttleiki &  &  &  \\
\hline
3254 & megin-, höfuð-, aðal- &  &  &  \\
\hline
3255 & meginákveða, höfuðákveða &  &  &  \\
\hline
3256 & meginás &  &  &  \\
\hline
3257 & meginás, höfuðás &  &  &  \\
\hline
3258 & meginásafærsla, höfuðásafærsla &  &  &  \\
\hline
3259 & meginásasetning &  &  &  \\
\hline
3260 & megináslausn &  &  &  \\
\hline
3261 & megingrein, höfuðgrein &  &  &  \\
\hline
3262 & meginhorn &  &  &  \\
\hline
3263 & meginhringur &  &  &  \\
\hline
3264 & meginhringur, meginhvirfilhringur, hjálparhringur &  &  &  \\
\hline
3265 & meginkrappastefna, höfuðkrappastefna, meginsveigjustefna, höfuðsveigjustefna &  &  &  \\
\hline
3266 & meginmarkgildissetning tölfræðinnar, höfuðsetning tölfræðinnar &  &  &  \\
\hline
3267 & meginþáttur &  &  &  \\
\hline
3268 & meginþverill, höfuðþverill &  &  &  \\
\hline
3269 & meirihlutaákvörðun &  &  &  \\
\hline
3270 & mengi & <ref idTerm=4308>Safn</ref> <ref idTerm=5114>hluta</ref> &  & Taka fram að safnið (flokkurinn) þarf að uppfylla ZFC frumsemdurnar í skýringu. Setja hlekk beint á flokk frekar en safn? \\
\hline
3271 & mengjaaðgerð &  &  &  \\
\hline
3272 & mengjaalgebra &  &  &  \\
\hline
3273 & mengjaalgebra &  &  &  \\
\hline
3274 & mengjafall &  &  &  \\
\hline
3275 & mengjafræði &  &  &  \\
\hline
3276 & mengjafræðilegur &  &  &  \\
\hline
3277 & mengjagildur &  &  &  \\
\hline
3278 & mengjahlutvera &  &  &  \\
\hline
3279 & mengjamargfeldi, faldmengi, Kartesarmargfeldi, karteskt margfeldi & \textit{Faldmengi} tveggja <ref idTerm=3270>mengja</ref> $X$ og $Y$, táknað $X \times Y$, er mengi <ref idTerm=5627>para</ref> $(x,y)$ þannig að $x \in X$ og $y \in Y$ &  & Koma með almenna skilgreiningu líka í leiðinni eða setja hana í skýringu? Síðan þarf kannski að minnast á það að það er farið svolítið losaralega að stundum með faldmengi og litið á $(X \times Y) \times Z$ sem $X \times Y \times Z$ \\
\hline
3280 & mengjamunur, mengjamismunur & \textit{Mengjamunur} tveggja <ref idTerm=3270>mengja</ref> $A$ og $B$ er mengi þeirra <ref idTerm=5128>staka</ref> sem eru í $A$ en ekki í $B$. Þetta mengi er táknað með $A \setminus B$ (lesið: $A$ án $B$) og það má rita á forminu \[ A \setminus B = \{x \in A \mid x \notin B\}. \] &  & Setja í skýringu: Óformlega má segja að mengjamunurinn $A \setminus B$ sé sá hluti mengisins $A$ sem er utan mengisins $B$ og á Venn-myndinni hér að neðan er hann litaður rauður. (ásamt mynd)

Fyrir öll mengi $A$ og $B$ má rita mengjamuninn $A \setminus B$ sem sniðmengi $A$ og fyllimengis $B$, þ.e. \[ A \setminus B = A \cap B^c. \] Jafnframt gildir fyrir öll mengi $A$ að $A \setminus \varnothing = A$ og $A \setminus A = \varnothing$. \\
\hline
3281 & mengjaríki &  &  &  \\
\hline
3282 & vörpun & \textit{Vörpun} $f: X \to Y$ frá <ref idTerm=3270>mengi</ref> $X$ yfir í mengið $Y$ er forskrift eða regla sem úthlutar sérhverju <ref idTerm=5128>staki</ref> úr $X$ nákvæmlega einu staki úr $Y$. Mengið $X$ kallast \textit{skilgreiningarmengi} vörpunarinnar og mengið $Y$ kallast \textit{bakmengi} hennar &  & Líka taka fram að ef bakmengið er $\mathbb{R}$ eða $\mathbb{C}$ er talað um fall frekar en vörpun? \\
\hline
3283 & merkingafræði &  &  &  \\
\hline
3284 & Mersenne-frumtala &  &  &  \\
\hline
3285 & meta &  &  &  \\
\hline
3286 & meta að ofan, yfirmeta &  &  &  \\
\hline
3287 & metakerfi, stigveldi &  &  &  \\
\hline
3288 & metill &  &  &  \\
\hline
3289 & mettað mengi &  &  &  \\
\hline
3290 & mettaður &  &  &  \\
\hline
3291 & mettun &  &  &  \\
\hline
3292 & mið (punkts) við kúlu, veldi (punkts) við kúlu &  &  &  \\
\hline
3293 & mið-, miðju-, miðpunkts- &  &  &  \\
\hline
3294 & mið, veldi &  &  &  \\
\hline
3295 & miðalína, jafnveldislína &  &  &  \\
\hline
3296 & miðamiðja, jafnveldismiðja &  &  &  \\
\hline
3297 & miðaslétta, jafnveldisslétta &  &  &  \\
\hline
3298 & miðbaugur &  &  &  \\
\hline
3299 & miðgildi, miðtala &  &  &  \\
\hline
3300 & miðgildispróf &  &  &  \\
\hline
3301 & miðhlutfalla stærð &  &  &  \\
\hline
3302 & miðhlutfallajafna, samfelld hlutfallajafna &  &  &  \\
\hline
3303 & miðja sjónhorfsmyndunar, sjónhorfsmiðja &  &  &  \\
\hline
3304 & miðja, miðdepill, miðpunktur &  &  &  \\
\hline
3305 & miðjufjarlægð &  &  &  \\
\hline
3306 & miðjuhorn, miðpunktshorn &  &  &  \\
\hline
3307 & miðjungur &  &  &  \\
\hline
3308 & miðjuofanvarp, miðjuvarp, keiluofanvarp &  &  &  \\
\hline
3309 & miðjuruna, miðjuturn &  &  &  \\
\hline
3310 & miðjusamhverfni, miðjusamhvarf &  &  &  \\
\hline
3311 & miðjusamhverfur, samhverfur um punkt &  &  &  \\
\hline
3312 & miðjusjónhorf &  &  &  \\
\hline
3313 & miðjustak &  &  &  \\
\hline
3314 & miðjuvíðátta &  &  &  \\
\hline
3315 & miðliður &  &  &  \\
\hline
3316 & miðlína &  &  &  \\
\hline
3317 & miðlína &  &  &  \\
\hline
3318 & miðmunur &  &  &  \\
\hline
3319 & miðpunktur &  &  &  \\
\hline
3320 & miðsett hending, miðsett slembibreyta &  &  &  \\
\hline
3321 & miðsett úrtak &  &  &  \\
\hline
3322 & miðsett vikbil &  &  &  \\
\hline
3323 & miðsettur &  &  &  \\
\hline
3324 & miðslétta &  &  &  \\
\hline
3325 & miðstrengur &  &  &  \\
\hline
3326 & miðtengilína &  &  &  \\
\hline
3328 & miðþverill, miðþverlína &  &  &  \\
\hline
3329 & miðþverslétta &  &  &  \\
\hline
3330 & miklu minni &  &  &  \\
\hline
3331 & miklu stærri &  &  &  \\
\hline
3332 & milli, á milli &  &  &  \\
\hline
3333 & milligildi &  &  &  \\
\hline
3334 & milligildissetning &  &  &  \\
\hline
3335 & milliliðalaus fáguð framlenging &  &  &  \\
\hline
3336 & milliliður, innliður &  &  &  \\
\hline
3337 & millimetrapappír &  &  &  \\
\hline
3338 & millireikningur &  &  &  \\
\hline
3339 & millisvið &  &  &  \\
\hline
3340 & millumleikafrumsenda &  &  &  \\
\hline
3341 & millumleikavensl &  &  &  \\
\hline
3342 & millumleiki &  &  &  \\
\hline
3343 & Minkowski-rúm &  &  &  \\
\hline
3344 & Minkowski-rúmfræði &  &  &  \\
\hline
3345 & minni &  &  &  \\
\hline
3346 & minni eða jafn, minni en eða jafn og, frekar minni &  &  &  \\
\hline
3347 & minni, smærri &  &  &  \\
\hline
3348 & minni, smærri &  &  &  \\
\hline
3349 & minnka, falla, lækka &  &  &  \\
\hline
3350 & minnkandi atburðaruna &  &  &  \\
\hline
3351 & minnkandi, fallandi, lækkandi &  &  &  \\
\hline
3352 & minnsta frávik &  &  &  \\
\hline
3353 & minnsta samfeldi, minnsta sameiginlega margfeldi & <ref idTerm=3347>Minnsta</ref> <ref idTerm=3475>náttúrulega talan</ref> sem tvær <ref idTerm=1896>heiltölur</ref> $n$ og $m$ <ref idTerm=1599>ganga upp í</ref> kallast \textit{minnsta} <ref idTerm=4343>\textit{samfeldi}</ref> $n$ og $m$ &  & Táknað með $\text{msf}(n,m)$ á íslensku? Held að flestir myndu frekar þekkja $\text{lcm}(n,m)$ \\
\hline
3354 & minnsta stig, lágstig &  &  &  \\
\hline
3355 & minnsti samnefnari &  &  &  \\
\hline
3356 & mínus, mínusmerki, frádráttarmerki &  &  &  \\
\hline
3357 & mínus, mínusmerki, mínusformerki, neikvætt formerki &  &  &  \\
\hline
3358 & mínúta &  &  &  \\
\hline
3359 & mínúta, hornmínúta, bogamínúta &  &  &  \\
\hline
3360 & misgeng ósamfelldni &  &  &  \\
\hline
3361 & misgilding, misgildi &  &  &  \\
\hline
3362 & mishverf marglínuleg vörpun, misvíxlin marglínuleg vörpun &  &  &  \\
\hline
3363 & mishverft fall, misvíxlið fall &  &  &  \\
\hline
3364 & mishverft fylki, misvíxlið fylki &  &  &  \\
\hline
3365 & mishverft tvílínulegt form, misvíxlið tvílínulegt form &  &  &  \\
\hline
3366 & mishverfur hnútur &  &  &  \\
\hline
3367 & mishverfur, misvíxlinn &  &  &  \\
\hline
3368 & mislægar línur &  &  &  \\
\hline
3369 & misleit hnit &  &  &  \\
\hline
3370 & misleitur &  &  &  \\
\hline
3371 & mismunaaðferð &  &  &  \\
\hline
3372 & mismunaaðferð &  &  &  \\
\hline
3373 & mismunahlutfall, mismunakvóti &  &  &  \\
\hline
3374 & mismunajafna &  &  &  \\
\hline
3375 & mismunareikningur &  &  &  \\
\hline
3376 & mismunavirki &  &  &  \\
\hline
3377 & missamhverf vensl, mishverf vensl &  &  &  \\
\hline
3378 & missamhverfur, mishverfur &  &  &  \\
\hline
3379 & misstór, misjafn &  &  &  \\
\hline
3380 & mjúkt knippi &  &  &  \\
\hline
3381 & mjúkur &  &  &  \\
\hline
3382 & Möbíusarræma &  &  &  \\
\hline
3383 & Möbíusarrúmfræði, hringarúmfræði &  &  &  \\
\hline
3384 & Möbíusarslétta &  &  &  \\
\hline
3385 & möndulslétta &  &  &  \\
\hline
3386 & Monte-Carlo-aðferð &  &  &  \\
\hline
3387 & Monte-Carlo-útreikningur, hermiútreikningur &  &  &  \\
\hline
3388 & mörkuð reglustika &  &  &  \\
\hline
3389 & mörkun, tölusetning &  &  &  \\
\hline
3390 & möskvahnútur, hnútur &  &  &  \\
\hline
3391 & möskvaskipting &  &  &  \\
\hline
3392 & möskvastærð &  &  &  \\
\hline
3393 & möskvi &  &  &  \\
\hline
3394 & mótanaríki, örvaríki &  &  &  \\
\hline
3395 & mótdæmi, gagndæmi &  &  &  \\
\hline
3396 & mótlæg hlið, öndverð hlið, gagnstæð hlið &  &  &  \\
\hline
3397 & mótlæg horn &  &  &  \\
\hline
3398 & mótlægt horn &  &  &  \\
\hline
3399 & mótlægur hliðarflötur, öndverður hliðarflötur &  &  &  \\
\hline
3400 & mótlægur, mótstæður, gagnstæður, öndverður, andstæður, öfugur &  &  &  \\
\hline
3401 & mótsagnakenndur, í mótsögn &  &  &  \\
\hline
3402 & mótskilyrðing &  &  &  \\
\hline
3403 & mótsögn &  &  &  \\
\hline
3404 & mótulfræði &  &  &  \\
\hline
3405 & mótulgrunnur &  &  &  \\
\hline
3406 & mótulinnanmótun &  &  &  \\
\hline
3407 & mótull &  &  &  \\
\hline
3408 & mótulmótun, línuleg vörpun &  &  &  \\
\hline
3409 & mótun &  &  &  \\
\hline
3410 & mótun, sammótun &  &  &  \\
\hline
3411 & munur, mismunur, afgangur (í frádrætti) &  &  &  \\
\hline
3412 & mynd &  &  &  \\
\hline
3413 & mynd &  &  &  \\
\hline
3414 & mynd, rit &  &  &  \\
\hline
3415 & mynd, rúmmynd &  &  &  \\
\hline
3416 & mynda, setja upp &  &  &  \\
\hline
3417 & myndmengi, mynd & Fyrir <ref idTerm=3282>vörpun</ref> $f: X \to Y$ kallast <ref idTerm=3270>mengi</ref> allra <ref idTerm=1090>gilda</ref> sem vörpunin tekur \textit{mynd} vörpunarinnar $f$. Hana má rita á forminu \[ f(X) = \{ f(x) \in Y \mid x \in X \}. \] &  & Tala um mynd hlutmengja í skilgreiningamenginu í skýringu \\
\hline
3418 & myndmengi, seiling &  &  &  \\
\hline
3419 & myndrit &  &  &  \\
\hline
3420 & myndtala &  &  &  \\
\hline
3421 & myndun, uppsetning &  &  &  \\
\hline
3422 & myndunarregla &  &  &  \\
\hline
3423 & mynstur &  &  &  \\
\hline
3424 & mynsturgrúpa &  &  &  \\
\hline
3425 & mynsturknippi &  &  &  \\
\hline
3426 & nægjanleg hending, tæmandi hending, nægjanleg reiknihending, tæmandi reiknihending &  &  &  \\
\hline
3427 & nægjanlegt skilyrði, fullnægjandi skilyrði, tæmandi skilyrði &  &  &  \\
\hline
3428 & nægjanlegur, fullnægjandi, tæmandi &  &  &  \\
\hline
3429 & nægjanleiki &  &  &  \\
\hline
3430 & nægjanleiki, tæmanleiki &  &  &  \\
\hline
3431 & nærbaugur &  &  &  \\
\hline
3432 & nærfágað deildi &  &  &  \\
\hline
3433 & nærfágað fall &  &  &  \\
\hline
3434 & nærfágaður &  &  &  \\
\hline
3435 & nærfáguð vörpun &  &  &  \\
\hline
3436 & nærfágun &  &  &  \\
\hline
3437 & nærsvið &  &  &  \\
\hline
3438 & nærþjappaður &  &  &  \\
\hline
3439 & nærþjappleiki &  &  &  \\
\hline
3440 & næstum alls staðar, nærri alls staðar &  &  &  \\
\hline
3441 & næstum lotubundið fall &  &  &  \\
\hline
3442 & næstum örugglega &  &  &  \\
\hline
3443 & næstum öruggur, öruggur með líkindum &  &  &  \\
\hline
3444 & næstum sjálfoka &  &  &  \\
\hline
3445 & nafla, naflavirki &  &  &  \\
\hline
3446 & naflapunktur &  &  &  \\
\hline
3447 & nágranni, granni &  &  &  \\
\hline
3448 & nákvæm deildajafna &  &  &  \\
\hline
3449 & nákvæm lausn &  &  &  \\
\hline
3450 & nákvæmni, hittni &  &  &  \\
\hline
3451 & nákvæmt deildaform &  &  &  \\
\hline
3452 & nákvæmt deildi, fleygað deildi &  &  &  \\
\hline
3453 & nákvæmur &  &  &  \\
\hline
3454 & nálægð, nánd &  &  &  \\
\hline
3455 & nálægur &  &  &  \\
\hline
3456 & nálga, nálægja, námunda &  &  &  \\
\hline
3457 & nálgast &  &  &  \\
\hline
3458 & nálgast &  &  &  \\
\hline
3459 & nálgun í jöfnum mæli &  &  &  \\
\hline
3460 & nálgun, námundun &  &  &  \\
\hline
3461 & nálgunar- &  &  &  \\
\hline
3462 & nálgunaraðferð &  &  &  \\
\hline
3463 & nálgunarbrot &  &  &  \\
\hline
3464 & nálgunareining &  &  &  \\
\hline
3465 & nálgunarfall &  &  &  \\
\hline
3466 & nálgunarfræði &  &  &  \\
\hline
3467 & nálgunargildi &  &  &  \\
\hline
3468 & nálgunarlausn &  &  &  \\
\hline
3469 & nálgunarsetning &  &  &  \\
\hline
3470 & nálgunarskekkja &  &  &  \\
\hline
3471 & nálgunarstaðall, nákvæmd &  &  &  \\
\hline
3472 & náttúrleg einsmótun &  &  &  \\
\hline
3473 & náttúrleg færsla, varpafærsla, varpamótun &  &  &  \\
\hline
3474 & náttúrleg mótun &  &  &  \\
\hline
3475 & náttúrleg tala & <ref idTerm=5128>Stak</ref> í <ref idTerm=3270>menginu</ref> $\{0, 1, 2, \dots\}$ (stundum er <ref idTerm=5397>talan</ref> <ref idTerm=3549>0</ref> ekki tekin með) &  & Minnast á Peano frumsemdur í skýringu? \\
\hline
3476 & náttúrleg vörpun &  &  &  \\
\hline
3477 & náttúrlegt grannmynstur &  &  &  \\
\hline
3478 & náttúrlegt ívarp &  &  &  \\
\hline
3479 & náttúrlegt jafngildi, varpajafngildi &  &  &  \\
\hline
3480 & náttúrlegt ofanvarp &  &  &  \\
\hline
3481 & náttúrlegur logri &  &  &  \\
\hline
3482 & náttúrlegur, náttúrulegur, eðlilegur &  &  &  \\
\hline
3483 & náttúrlegur, varpa- &  &  &  \\
\hline
3484 & náttúrlegur, viðtekinn, sjálfgefinn, aðallegur, eðlilegur, kór- &  &  &  \\
\hline
3485 & náttúrleiki &  &  &  \\
\hline
3486 & nauðsyn &  &  &  \\
\hline
3487 & nauðsynlegt og nægjanlegt skilyrði, nauðsynlegt og fullnægjandi skilyrði &  &  &  \\
\hline
3488 & nauðsynlegur &  &  &  \\
\hline
3489 & neðra heildi, neðra tegur &  &  &  \\
\hline
3490 & neðra mark, lágmark &  &  &  \\
\hline
3491 & neðra markgildi, minnsti þéttipunktur, neðsti þéttipunktur &  &  &  \\
\hline
3492 & neðra öryggisbil &  &  &  \\
\hline
3493 & neðra þríhyrningsfylki &  &  &  \\
\hline
3494 & neðri stallagerð, neðri þrepagerð, neðri þrepaskipan &  &  &  \\
\hline
3495 & neðri þríhyrningsgrúpa &  &  &  \\
\hline
3496 & neðsta gildi, lægsta gildi, minnsta gildi, minnsta lággildi, víðfeðmt lággildi &  &  &  \\
\hline
3497 & nefnari & Sjá <ref idTerm=234>almennt brot</ref> &  &  \\
\hline
3498 & neikvæð fylgni &  &  &  \\
\hline
3499 & neikvæð stefna &  &  &  \\
\hline
3500 & neikvæður atburður, óhapp, bilun &  &  &  \\
\hline
3501 & neikvæður hluti &  &  &  \\
\hline
3502 & neikvæður snúningur, snúningur réttsælis &  &  &  \\
\hline
3503 & neikvæður, frádrægur &  &  &  \\
\hline
3504 & neikvætt ákveðinn &  &  &  \\
\hline
3505 & neikvætt áttaður &  &  &  \\
\hline
3506 & neikvætt hálfákveðinn &  &  &  \\
\hline
3507 & neitun & Fyrir sérhverja <ref idTerm=1515>yrðingu</ref> $p$ er $\neg p$ (lesið: „ekki $p$“) sú yrðing sem fæst með því að \textit{neita} yrðingunni $p$. Yrðingin $\neg p$ hefur þess vegna alltaf öfugt <ref idTerm=4560>sanngildi</ref> við yrðinguna $p$ &  & Setja í skýringu: Þessari staðreynd má lýsa með eftirfarandi töflu, sem sýnir sanngildi $\neg p$ eftir ólíkum sanngildum $p$: (plús mynd) \\
\hline
3508 & neitunarháttur &  &  &  \\
\hline
3509 & neitunarmerki &  &  &  \\
\hline
3510 & net &  &  &  \\
\hline
3511 & net &  &  &  \\
\hline
3512 & netafræði, netjafræði, graffræði &  &  &  \\
\hline
3513 & Newton-aðferð &  &  &  \\
\hline
3514 & niðji &  &  &  \\
\hline
3516 & niðurstaða &  &  &  \\
\hline
3517 & niðurstaða, ályktun, lokaályktun &  &  &  \\
\hline
3518 & níhyrningstala, níhyrnd tala &  &  &  \\
\hline
3519 & níhyrningur &  &  &  \\
\hline
3520 & níupróf &  &  &  \\
\hline
3521 & níupunktahringur &  &  &  \\
\hline
3522 & níusetning &  &  &  \\
\hline
3523 & Noether-baugur &  &  &  \\
\hline
3524 & Noether-mótull &  &  &  \\
\hline
3525 & noetherskur, Noether- &  &  &  \\
\hline
3526 & norðurskaut, norðurpóll &  &  &  \\
\hline
3527 & normgera &  &  &  \\
\hline
3528 & normleg hending &  &  &  \\
\hline
3529 & normleg hlutgrúpa, óbreytt hlutgrúpa &  &  &  \\
\hline
3530 & normleg hnit &  &  &  \\
\hline
3531 & normleg jafna &  &  &  \\
\hline
3532 & normleg útvíkkun, normleg víkkun &  &  &  \\
\hline
3533 & normlegt fylki &  &  &  \\
\hline
3534 & normlegt grannmynstur &  &  &  \\
\hline
3535 & normlegt rúm &  &  &  \\
\hline
3536 & normlegt svið &  &  &  \\
\hline
3537 & normlegt tíðnifall &  &  &  \\
\hline
3538 & normlegt víðerni &  &  &  \\
\hline
3539 & normlegt þýði &  &  &  \\
\hline
3540 & normlegur &  &  &  \\
\hline
3541 & normlegur ferill &  &  &  \\
\hline
3542 & normlegur flötur &  &  &  \\
\hline
3543 & normlegur grunnur &  &  &  \\
\hline
3544 & normlegur turn &  &  &  \\
\hline
3545 & normlegur virki &  &  &  \\
\hline
3546 & normleikafrumsenda &  &  &  \\
\hline
3547 & normleikagrúpa &  &  &  \\
\hline
3548 & normleiki &  &  &  \\
\hline
3550 & núllbaugur &  &  &  \\
\hline
3551 & núllbaugur, örbaugur &  &  &  \\
\hline
3552 & núlldeilir &  &  &  \\
\hline
3553 & núllendurkvæmur &  &  &  \\
\hline
3554 & núllendurkvæmur, núllafturkvæmur &  &  &  \\
\hline
3555 & núllfall &  &  &  \\
\hline
3556 & núllfylki &  &  &  \\
\hline
3557 & núllgagnvegur &  &  &  \\
\hline
3558 & núllgild margliða &  &  &  \\
\hline
3559 & núllhlutur &  &  &  \\
\hline
3560 & núllíðal &  &  &  \\
\hline
3561 & núlllausn &  &  &  \\
\hline
3562 & núlllausn &  &  &  \\
\hline
3563 & núllmargliða & <ref idTerm=0>Fastamargliðan</ref> 0 kallast \textit{núllmargliðan}. &  &  \\
\hline
3564 & núllmengi &  &  &  \\
\hline
3565 & núllrót &  &  &  \\
\hline
3566 & núllrót &  &  &  \\
\hline
3567 & núllsækin runa &  &  &  \\
\hline
3568 & núllsamtoga &  &  &  \\
\hline
3569 & núllsamtogun &  &  &  \\
\hline
3570 & núllstætt fall &  &  &  \\
\hline
3571 & núllstöð &  &  &  \\
\hline
3572 & núllstöðvamengi &  &  &  \\
\hline
3573 & núllstöðvasetning &  &  &  \\
\hline
3574 & núllsummuleikur &  &  &  \\
\hline
3575 & núlltilgáta &  &  &  \\
\hline
3576 & núllvalda &  &  &  \\
\hline
3577 & núllvíður &  &  &  \\
\hline
3578 & núllvigur, núllvektor &  &  &  \\
\hline
3579 & nútímarökfræði &  &  &  \\
\hline
3580 & nýgráða &  &  &  \\
\hline
3581 & nykra &  &  &  \\
\hline
3582 & nykrað vigurrúm, nykurrúm &  &  &  \\
\hline
3583 & nykraður, nykur- &  &  &  \\
\hline
3584 & nykruð samlínumótun &  &  &  \\
\hline
3585 & nykrun &  &  &  \\
\hline
3586 & nykrunarsetning &  &  &  \\
\hline
3587 & nykrunarvarpi &  &  &  \\
\hline
3588 & nykurbundin &  &  &  \\
\hline
3589 & nykurfærsla &  &  &  \\
\hline
3590 & nykurferill &  &  &  \\
\hline
3591 & nykurgrind &  &  &  \\
\hline
3592 & nykurgrunnur &  &  &  \\
\hline
3593 & nykurgrúpa, nykurvíxlgrúpa &  &  &  \\
\hline
3595 & nykurlögmál, nykurvensl &  &  &  \\
\hline
3596 & nykurnet &  &  &  \\
\hline
3597 & nykurríki &  &  &  \\
\hline
3598 & nykursetning, nykruð setning &  &  &  \\
\hline
3599 & nykurvarpslétta, nykurslétta &  &  &  \\
\hline
3600 & nykurverkefni &  &  &  \\
\hline
3601 & nykurvirki &  &  &  \\
\hline
3602 & nykurvörpun, nykruð vörpun &  &  &  \\
\hline
3603 & nýmínúta &  &  &  \\
\hline
3604 & nýraferill &  &  &  \\
\hline
3605 & nyrðra hálfhvel, norðurhvel &  &  &  \\
\hline
3606 & nýsekúnda &  &  &  \\
\hline
3607 & nyt &  &  &  \\
\hline
3608 & nytjafall &  &  &  \\
\hline
3609 & nytjafræði &  &  &  \\
\hline
3610 & óabelskur &  &  &  \\
\hline
3611 & óaðgengileg fjöldatala &  &  &  \\
\hline
3612 & óaðgengilegur jaðarpunktur &  &  &  \\
\hline
3613 & óaðlægur, ósamlægur, ósamliggjandi, grunnstæður &  &  &  \\
\hline
3614 & óaðskiljanlegt net &  &  &  \\
\hline
3615 & óaðskiljanlegur &  &  &  \\
\hline
3616 & óaðskiljanleikaröð &  &  &  \\
\hline
3617 & óaðskiljanleikastig &  &  &  \\
\hline
3618 & óaðskiljanleikavísir &  &  &  \\
\hline
3619 & óaðskiljanleiki &  &  &  \\
\hline
3620 & óafmáanlegur &  &  &  \\
\hline
3621 & óafsannanlegur &  &  &  \\
\hline
3622 & óákvarðaður, óákveðinn &  &  &  \\
\hline
3623 & óákveðið form &  &  &  \\
\hline
3624 & óákveðið heildi, óákveðið tegur &  &  &  \\
\hline
3625 & óákveðin stærð &  &  &  \\
\hline
3626 & óákveðinn stuðull &  &  &  \\
\hline
3627 & óákveðinn, óákvarðaður &  &  &  \\
\hline
3628 & óákveðnipunktur &  &  &  \\
\hline
3629 & óandhverfanlegur &  &  &  \\
\hline
3630 & óargeskur &  &  &  \\
\hline
3631 & óarkimedísk rúmfræði &  &  &  \\
\hline
3632 & óarkimedískur &  &  &  \\
\hline
3633 & óáttaður &  &  &  \\
\hline
3634 & óáttaður, óstefndur, óstefnubundinn &  &  &  \\
\hline
3635 & óáttanlegur &  &  &  \\
\hline
3636 & óaugljós &  &  &  \\
\hline
3637 & óbein ályktun &  &  &  \\
\hline
3638 & óbein deildun, óbein diffrun, fólgin deildun, fólgin diffrun &  &  &  \\
\hline
3639 & óbein skilgreining, fólgin skilgreining &  &  &  \\
\hline
3640 & óbein sönnun &  &  &  \\
\hline
3641 & óbein úrtaka &  &  &  \\
\hline
3642 & óbjagað próf &  &  &  \\
\hline
3643 & óbjagað úrtak &  &  &  \\
\hline
3644 & óbjagað vikbil, óbjagað öryggisbil &  &  &  \\
\hline
3645 & óbjagaður metill, meðalgildisréttur metill &  &  &  \\
\hline
3646 & óbjöguð úrtaka &  &  &  \\
\hline
3647 & óblandaður &  &  &  \\
\hline
3648 & óbreyta &  &  &  \\
\hline
3649 & óbreytni &  &  &  \\
\hline
3650 & óbreytt hlutrúm &  &  &  \\
\hline
3651 & óbreytt mengi &  &  &  \\
\hline
3652 & óbreyttur &  &  &  \\
\hline
3653 & óbreyttur eiginleiki &  &  &  \\
\hline
3654 & óbreyttur frá hægri &  &  &  \\
\hline
3655 & óbreyttur frá vinstri &  &  &  \\
\hline
3656 & óbreyttur við hliðrun, óháður hliðrun &  &  &  \\
\hline
3657 & óðal &  &  &  \\
\hline
3658 & odda-, ójafn, hvass &  &  &  \\
\hline
3659 & oddatala, ójöfn tala, hvöss tala & <ref idTerm=1896>Heil tala</ref> sem $2$ <ref idTerm=1599>gengur ekki upp í</ref>, þ.e. heil tala $n$ sem rita má á forminu $n = 2 \cdot k + 1$ þar sem $k$ er heil tala &  &  \\
\hline
3660 & oddklofi &  &  &  \\
\hline
3661 & oddpunktur, topppunktur, broddur & Sjá <ref idTerm=2198>horn</ref> &  &  \\
\hline
3662 & oddstæðni, oddleiki, hvassleiki &  &  &  \\
\hline
3663 & oddstæður & <ref idTerm=1078>Fall</ref> $f: X \to Y$ þar sem $X$ og $Y$ eru <ref idTerm=2098>hlutmengi</ref> í <ref idTerm=3270>mengi</ref> <ref idTerm=4093>rauntalna</ref> kallast \textit{oddstætt} ef fyrir öll $x \in X$ gildi að \[ f(-x) = - f(x). \] &  &  \\
\hline
3664 & oddstæður &  &  &  \\
\hline
3665 & oddsvigi &  &  &  \\
\hline
3666 & ódeilanlegur &  &  &  \\
\hline
3667 & ódeilishluti, hluti en ekki deilir &  &  &  \\
\hline
3668 & óeiginleg afleiða &  &  &  \\
\hline
3669 & óeiginleg deildun, óeiginleg diffrun &  &  &  \\
\hline
3670 & óeiginleg evklíðsk hreyfing &  &  &  \\
\hline
3671 & óeiginleg ósamleitni &  &  &  \\
\hline
3672 & óeiginleg þverstaðalvörpun, snúspeglun &  &  &  \\
\hline
3673 & óeiginlega heildanlegur, óeiginlega tegranlegur &  &  &  \\
\hline
3674 & óeiginlegt brot &  &  &  \\
\hline
3675 & óeiginlegt heildi, óeiginlegt tegur &  &  &  \\
\hline
3676 & óeiginlegt rætt fall &  &  &  \\
\hline
3677 & óeiginlegur &  &  &  \\
\hline
3678 & óeiginlegur flokkur &  &  &  \\
\hline
3679 & óeiginlegur strengur &  &  &  \\
\hline
3680 & óeind &  &  &  \\
\hline
3681 & óendanleg fjöldatala, ofurendanleg fjöldatala &  &  &  \\
\hline
3682 & óendanleg raðtala, ofurendanleg raðtala &  &  &  \\
\hline
3683 & óendanleg röð, röð &  &  &  \\
\hline
3684 & óendanleg runa, runa &  &  &  \\
\hline
3685 & óendanleg slétta &  &  &  \\
\hline
3686 & óendanleg stærð &  &  &  \\
\hline
3687 & óendanleg tala &  &  &  \\
\hline
3688 & óendanleg þáttun, margfeldisliðun, feldisliðun &  &  &  \\
\hline
3689 & óendanlega lítill &  &  &  \\
\hline
3690 & óendanlega margir &  &  &  \\
\hline
3691 & óendanlega oft deildanlegur, óendanlega oft diffranlegur, þjáll &  &  &  \\
\hline
3692 & óendanlega stór &  &  &  \\
\hline
3693 & óendanlega vaxandi &  &  &  \\
\hline
3694 & óendanlega víður &  &  &  \\
\hline
3695 & óendanlegt margfeldi &  &  &  \\
\hline
3696 & óendanlegt mengi &  &  &  \\
\hline
3697 & óendanlegt tugabrot &  &  &  \\
\hline
3698 & óendanlegur &  &  &  \\
\hline
3699 & óendanlegur fjöldi &  &  &  \\
\hline
3700 & óendanlegur, endalaus &  &  &  \\
\hline
3701 & óendanleikafrumsenda, óendanleikafrumsetning, frumsenda um óendanleika, frumsetning um óendanleika &  &  &  \\
\hline
3702 & óendanleikapunktur &  &  &  \\
\hline
3703 & óendanleiki &  &  &  \\
\hline
3704 & óendurkvæmur, óafturkvæmur &  &  &  \\
\hline
3705 & óevklíðsk rúmfræði &  &  &  \\
\hline
3706 & óevklíðskur &  &  &  \\
\hline
3707 & óevklíðskur flutningur, óevklíðsk hreyfing &  &  &  \\
\hline
3708 & ófær, ómögulegur &  &  &  \\
\hline
3709 & ófærusönnun, ómöguleikasönnun &  &  &  \\
\hline
3710 & ofanvarp &  &  &  \\
\hline
3711 & ofanvarpi, ofanvarp &  &  &  \\
\hline
3712 & ofanvarpslína &  &  &  \\
\hline
3713 & ofaukinn, óþarfur &  &  &  \\
\hline
3714 & ofhorn &  &  &  \\
\hline
3715 & öfug röð, gagnstæð röð &  &  &  \\
\hline
3716 & öfug þríliða &  &  &  \\
\hline
3717 & öfugt markgildi &  &  &  \\
\hline
3718 & ófullgerður &  &  &  \\
\hline
3719 & ófullkomin tala &  &  &  \\
\hline
3720 & ófullkominn &  &  &  \\
\hline
3721 & ófullkomleikasetning Gödels &  &  &  \\
\hline
3722 & ofurbreiðger &  &  &  \\
\hline
3723 & ofurendanleg rakning &  &  &  \\
\hline
3724 & ofurendanleg tala &  &  &  \\
\hline
3725 & ofurendanleg þrepun &  &  &  \\
\hline
3726 & ofurendanleg þrepunarsönnun, ofurendanleg þrepasönnun &  &  &  \\
\hline
3727 & ofurendanlegur &  &  &  \\
\hline
3728 & ofurfirð &  &  &  \\
\hline
3729 & ofurfirðarójafna, sterk þríhyrningsójafna &  &  &  \\
\hline
3730 & ofurfirðrúm &  &  &  \\
\hline
3731 & ofurmargfeldi, ofurfeldi &  &  &  \\
\hline
3732 & ofurnet &  &  &  \\
\hline
3733 & ofurrúmkynja fall &  &  &  \\
\hline
3734 & ofurrúmkynja röð &  &  &  \\
\hline
3735 & ofursía &  &  &  \\
\hline
3736 & ofurveldi &  &  &  \\
\hline
3737 & ofurvirðing, óarkimedísk virðing &  &  &  \\
\hline
3738 & ofurþéttipunktur &  &  &  \\
\hline
3739 & og &  &  &  \\
\hline
3740 & og-tákn &  &  &  \\
\hline
3741 & ógegnvirkur &  &  &  \\
\hline
3742 & ógegnvirkur, ekki gegnvirkur &  &  &  \\
\hline
3743 & ögn &  &  &  \\
\hline
3744 & ogun & Fyrir sérhverjar <ref idTerm=1515>yrðingar</ref> $p$ og $q$ er $p \wedge q$ (lesið: „$p$ \textit{og} $q$“) sú yrðing sem segir að $p$ og $q$ séu báðar <ref idTerm=4565>sannar</ref>. Yrðingin $p \wedge q$ er þess vegna sönn þegar yrðingarnar $p$ og $q$ eru báðar sannar, en annars er hún <ref idTerm=3890>ósönn</ref> &  & Setja í skýringu: Þessu má lýsa með eftirfarandi töflu, sem sýnir sanngildi yrðingarinnar $p \wedge q$ eftir ólíkum sanngildum $p$ annars vegar og $q$ hins vegar: (ásamt mynd) \\
\hline
3745 & ogunarhluti, ogunarþáttur &  &  &  \\
\hline
3746 & óháð úrtaka &  &  &  \\
\hline
3747 & óháðar frumsendur, óháðar frumsetningar &  &  &  \\
\hline
3748 & óháðar hendingar &  &  &  \\
\hline
3749 & óháðir atburðir &  &  &  \\
\hline
3750 & óháðir tveir og tveir &  &  &  \\
\hline
3751 & óháður &  &  &  \\
\hline
3752 & óhæði &  &  &  \\
\hline
3753 & óhæðisönnun &  &  &  \\
\hline
3754 & óhliðrað jöfnuhneppi &  &  &  \\
\hline
3755 & óhliðraður &  &  &  \\
\hline
3756 & óhliðruð línuleg jafna &  &  &  \\
\hline
3757 & óhlutbundinn hyrnuklasi &  &  &  \\
\hline
3758 & óhlutvera &  &  &  \\
\hline
3759 & óhnýttur &  &  &  \\
\hline
3760 & ójafn &  &  &  \\
\hline
3761 & ójafn, ólíkur &  &  &  \\
\hline
3762 & ójafna &  &  &  \\
\hline
3763 & ójafnaðarmerki &  &  &  \\
\hline
3764 & ójafnarma þríhyrningur, mishliða þríhyrningur &  &  &  \\
\hline
3765 & ójafnhliða og skakkur samsíðungur &  &  &  \\
\hline
3766 & ójöfnumerki &  &  &  \\
\hline
3767 & ójöfnumerki, ójöfnutákn &  &  &  \\
\hline
3768 & ójöfnuþvingun &  &  &  \\
\hline
3769 & ókleyfi, óþættanleiki, ókljúfanleiki &  &  &  \\
\hline
3770 & ókleyft tilvik &  &  &  \\
\hline
3771 & ókljúfanleg algebrulegt mengi &  &  &  \\
\hline
3772 & ókvíslaður &  &  &  \\
\hline
3773 & óleysanlegur &  &  &  \\
\hline
3774 & óleysanlegur &  &  &  \\
\hline
3775 & ólínuleg hneigð &  &  &  \\
\hline
3776 & ólínulegt aðhvarf &  &  &  \\
\hline
3777 & ólínulegur &  &  &  \\
\hline
3778 & ólínuleiki &  &  &  \\
\hline
3779 & ólokuð segð, ólokuð fullyrðing, ólokuð yrðing &  &  &  \\
\hline
3780 & ólotubundið ástand, lotulaust ástand &  &  &  \\
\hline
3781 & ólotubundinn &  &  &  \\
\hline
3782 & ólotubundinn, lotulaus &  &  &  \\
\hline
3783 & ómælanlegur &  &  &  \\
\hline
3784 & ómagur, ekki magur &  &  &  \\
\hline
3785 & ómagurt mengi, ómagurt hlutmengi &  &  &  \\
\hline
3786 & ómarkverður, fánýtur, klénn &  &  &  \\
\hline
3787 & ómega-fullkominn &  &  &  \\
\hline
3788 & óminnkanleiki &  &  &  \\
\hline
3789 & ómögulegur atburður &  &  &  \\
\hline
3790 & öndverð aðferð, gagnstæð aðgerð &  &  &  \\
\hline
3791 & öndverð grúpa, gagnstæð grúpa &  &  &  \\
\hline
3792 & önnur afleiða &  &  &  \\
\hline
3793 & opið bil & <ref idTerm=457>Bil</ref> sem inniheldur ekki endapunkta sína &  &  \\
\hline
3794 & opið hálfrúm &  &  &  \\
\hline
3795 & opið hlutrúm &  &  &  \\
\hline
3796 & opið mengi, opið hlutmengi &  &  &  \\
\hline
3797 & opin grennd &  &  &  \\
\hline
3798 & opin hringskífa &  &  &  \\
\hline
3799 & opin kúla &  &  &  \\
\hline
3800 & opin segð, opin fullyrðing, opin yrðing & <ref idTerm=1515>Yrðing</ref> sem inniheldur eina eða fleiri óþekktar <ref idTerm=5102>stærðir</ref> &  &  \\
\hline
3801 & opin vörpun &  &  &  \\
\hline
3802 & opin þakning &  &  &  \\
\hline
3803 & opinleiki &  &  &  \\
\hline
3804 & opinn &  &  &  \\
\hline
3805 & opinn ferill &  &  &  \\
\hline
3806 & opinn flötur &  &  &  \\
\hline
3807 & opinn forgrunnur &  &  &  \\
\hline
3808 & opinn og lokaður &  &  &  \\
\hline
3809 & ör, örvarleggur, stefndur leggur &  &  &  \\
\hline
3810 & ör, píla &  &  &  \\
\hline
3811 & óráðið form, óákveðið form &  &  &  \\
\hline
3812 & óráðið stærðtákn &  &  &  \\
\hline
3813 & óræðni &  &  &  \\
\hline
3814 & óræður & <ref idTerm=4093>Rauntala</ref> kallast \textit{óræð} ef hún er ekki <ref idTerm=4039>ræð</ref> &  &  \\
\hline
3815 & órásaður fléttingur, óhringaður fléttingur &  &  &  \\
\hline
3816 & órásaður, óhringaður &  &  &  \\
\hline
3817 & orð &  &  &  \\
\hline
3818 & orðalengd &  &  &  \\
\hline
3819 & orðaverkefni &  &  &  \\
\hline
3820 & óregluleg frumtala &  &  &  \\
\hline
3821 & óreglulegur &  &  &  \\
\hline
3822 & óreglulegur sérstöðupunktur &  &  &  \\
\hline
3823 & óreiða &  &  &  \\
\hline
3824 & óreiknanlegur &  &  &  \\
\hline
3825 & óreiknanleiki &  &  &  \\
\hline
3826 & örfærsla &  &  &  \\
\hline
3827 & örgrúpa &  &  &  \\
\hline
3828 & órífkandi &  &  &  \\
\hline
3829 & óröðuð tvennd &  &  &  \\
\hline
3830 & örsmæð &  &  &  \\
\hline
3831 & örsmæða- &  &  &  \\
\hline
3832 & örsmár &  &  &  \\
\hline
3833 & örspönnuður &  &  &  \\
\hline
3834 & örvaelti &  &  &  \\
\hline
3835 & örvahringur &  &  &  \\
\hline
3836 & örvanet, stefnt net, áttað net &  &  &  \\
\hline
3837 & örvarás &  &  &  \\
\hline
3838 & örvarit &  &  &  \\
\hline
3839 & örvasamanhangandi, örvasamhangandi &  &  &  \\
\hline
3840 & öryggisbelti &  &  &  \\
\hline
3841 & öryggisbil, vikbil &  &  &  \\
\hline
3842 & öryggisbilsmat, bilsmat &  &  &  \\
\hline
3843 & öryggislíkindi &  &  &  \\
\hline
3844 & öryggismark, vikmark &  &  &  \\
\hline
3845 & öryggismat &  &  &  \\
\hline
3846 & öryggismengi &  &  &  \\
\hline
3847 & öryggissporbaugur &  &  &  \\
\hline
3848 & öryggissporvala &  &  &  \\
\hline
3849 & öryggisstig, öryggi &  &  &  \\
\hline
3850 & öryggisstuðull &  &  &  \\
\hline
3851 & öryggissvæði &  &  &  \\
\hline
3852 & öryggissvæðismat &  &  &  \\
\hline
3853 & ós, svelgur &  &  &  \\
\hline
3854 & ósamanhangandi, ósamhangandi &  &  &  \\
\hline
3855 & ósambærilegur &  &  &  \\
\hline
3856 & ósamfelldni &  &  &  \\
\hline
3857 & ósamfelldni frá hægri &  &  &  \\
\hline
3858 & ósamfelldni frá vinstri &  &  &  \\
\hline
3859 & ósamfelldnipunktur &  &  &  \\
\hline
3860 & ósamfelldur &  &  &  \\
\hline
3861 & ósamfelldur frá hægri &  &  &  \\
\hline
3862 & ósamfelldur frá vinstri &  &  &  \\
\hline
3863 & ósamfellt fall &  &  &  \\
\hline
3864 & ósamhverf dreifing &  &  &  \\
\hline
3865 & ósamhverfur &  &  &  \\
\hline
3866 & ósamhverfur &  &  &  \\
\hline
3867 & ósamkvæmni, ósamrýmanleiki &  &  &  \\
\hline
3868 & ósamkvæmur metill &  &  &  \\
\hline
3869 & ósamkvæmur, ósamrýmanlegur &  &  &  \\
\hline
3870 & ósamleifa &  &  &  \\
\hline
3871 & ósamleitinn &  &  &  \\
\hline
3872 & ósamleitni &  &  &  \\
\hline
3873 & ósamleitnipunktur &  &  &  \\
\hline
3874 & ósamlína &  &  &  \\
\hline
3875 & ósammælanlegur &  &  &  \\
\hline
3876 & ósammælanleiki &  &  &  \\
\hline
3877 & ósammiðja, hjámiðja &  &  &  \\
\hline
3878 & ósampunkta &  &  &  \\
\hline
3879 & ósamrýndur bútur &  &  &  \\
\hline
3880 & ósamsnið &  &  &  \\
\hline
3881 & ósamsniða &  &  &  \\
\hline
3882 & ósamsniða &  &  &  \\
\hline
3883 & ósamþátta tveir og tveir, ósamþættir tveir og tveir &  &  &  \\
\hline
3884 & ósamþátta, ósamþættur &  &  &  \\
\hline
3885 & ósamþjappanlegur &  &  &  \\
\hline
3886 & ósannanlegur &  &  &  \\
\hline
3887 & ósannanleiki &  &  &  \\
\hline
3888 & ósanngjarn leikur &  &  &  \\
\hline
3889 & ósannindi, ósannleikur &  &  &  \\
\hline
3891 & ósérstæður algebrulegur ferill, reglulegur algebrulegur ferill &  &  &  \\
\hline
3892 & ósérstæður algebrulegur flötur, reglulegur algebrulegur flötur &  &  &  \\
\hline
3893 & ósérstæður fágaður ferill, reglulegur fágaður ferill &  &  &  \\
\hline
3894 & ósérstæður fágaður flötur, reglulegur fágaður flötur &  &  &  \\
\hline
3895 & ósérstæður, reglulegur &  &  &  \\
\hline
3896 & ósérstæður, reglulegur &  &  &  \\
\hline
3897 & ósérstætt algebrulegt víðerni, reglulegt algebrulegt víðerni &  &  &  \\
\hline
3898 & ósérstætt fágað víðerni, reglulegt fágað víðerni &  &  &  \\
\hline
3899 & óskilgreindur &  &  &  \\
\hline
3900 & óskilvirkur &  &  &  \\
\hline
3901 & óskilvirkur metill &  &  &  \\
\hline
3902 & óskilyrt ójafna &  &  &  \\
\hline
3903 & óslembdur &  &  &  \\
\hline
3904 & óslembinn &  &  &  \\
\hline
3905 & ósléttanlegur &  &  &  \\
\hline
3906 & ósmækkanlegt spannmengi, ósmækkanlegt spannandi mengi, óminnkanlegt spannmengi, óminnkanlegt spannandi mengi &  &  &  \\
\hline
3907 & óspegilvirkur &  &  &  \\
\hline
3908 & óstækkanleiki &  &  &  \\
\hline
3909 & óstefnt net, óáttað net &  &  &  \\
\hline
3910 & óstöðugleiki &  &  &  \\
\hline
3911 & óstöðugur &  &  &  \\
\hline
3912 & óstrjáll &  &  &  \\
\hline
3913 & óstyttanleg frymin þáttun &  &  &  \\
\hline
3914 & ósundurlægur, skaraður &  &  &  \\
\hline
3915 & ótakmarkað bil, óskorðað bil &  &  &  \\
\hline
3916 & ótakmarkað fall, óskorðað fall &  &  &  \\
\hline
3917 & ótakmarkaður virki, óskorðaður virki &  &  &  \\
\hline
3918 & ótakmarkaður, óskorðaður &  &  &  \\
\hline
3919 & óteljanlegur &  &  &  \\
\hline
3920 & ótenginn &  &  &  \\
\hline
3921 & ótiltekið stak &  &  &  \\
\hline
3922 & ótiltekinn &  &  &  \\
\hline
3923 & ótiltekinn, hver sem er &  &  &  \\
\hline
3924 & ótvíræð lausn, ótvírætt ákvörðuð lausn &  &  &  \\
\hline
3925 & ótvíræð skilgreining &  &  &  \\
\hline
3926 & ótvíræð þáttun &  &  &  \\
\hline
3927 & ótvíræðni, einræðni &  &  &  \\
\hline
3928 & ótvíræðnifrumsenda, ótvíræðnifrumsetning &  &  &  \\
\hline
3929 & ótvíræðnisetning &  &  &  \\
\hline
3930 & ótvíræðniskilyrði &  &  &  \\
\hline
3931 & ótvíræður, einræður &  &  &  \\
\hline
3932 & ótvírætt ákvarðaður & <ref idTerm=5102>Stærð</ref> $A$ kallast \textit{ótvírætt ákvörðuð} af tilteknum <ref idTerm=4797>skilyrðum</ref> ef hún

\begin{enumerate}
    \item uppfyllir þau skilyrði ($A$ er \textit{ákvörðuð} af skilyrðunum),
    \item ef önnur stærð $B$ uppfyllir sömu skilyrði þá gildir $A=B$
\end{enumerate} &  &  \\
\hline
3933 & óumhverfanlegur &  &  &  \\
\hline
3934 & óuppbyggjandi &  &  &  \\
\hline
3935 & óuppbyggjandi tilvistarsönnun &  &  &  \\
\hline
3936 & óúrkynjaður, ósérstæður &  &  &  \\
\hline
3937 & óúrskurðanleg fullyrðing, óúrskurðanleg yrðing &  &  &  \\
\hline
3938 & óúrskurðanleg kenning &  &  &  \\
\hline
3939 & óúrskurðanlegur &  &  &  \\
\hline
3940 & óúrskurðanleikasetning &  &  &  \\
\hline
3941 & óúrskurðanleiki &  &  &  \\
\hline
3942 & óvanaleg stærðfræðigreining, óvanaleg greining &  &  &  \\
\hline
3943 & óvanaleg tala &  &  &  \\
\hline
3944 & óvanalegur &  &  &  \\
\hline
3945 & óverulegur &  &  &  \\
\hline
3946 & óverulegur sérstöðupunktur &  &  &  \\
\hline
3947 & óverulegur, hverfandi &  &  &  \\
\hline
3948 & óveruritháttur &  &  &  \\
\hline
3949 & óviðkennd tala &  &  &  \\
\hline
3950 & óviðkomandi rót &  &  &  \\
\hline
3951 & óvissa &  &  &  \\
\hline
3952 & óvíxlanlegur &  &  &  \\
\hline
3953 & óvíxlinn &  &  &  \\
\hline
3954 & óvíxlinn, óabelskur &  &  &  \\
\hline
3957 & óþættanlegur, ókljúfanlegur, ókleyfur &  &  &  \\
\hline
3959 & óþættanleiki, ókljúfanleiki, ókleyfi &  &  &  \\
\hline
3960 & óþekkt stærð &  &  &  \\
\hline
3961 & óþekktur &  &  &  \\
\hline
3962 & óþjappaður &  &  &  \\
\hline
3963 & óþörf frumsenda, óþörf frumsetning &  &  &  \\
\hline
3964 & Papposar-, pappeskur &  &  &  \\
\hline
3965 & Papposarregla &  &  &  \\
\hline
3966 & Papposarsetning, setning Papposar &  &  &  \\
\hline
3967 & Papposarslétta, pappesk slétta &  &  &  \\
\hline
3968 & paraður samanburður &  &  &  \\
\hline
3969 & Pascal-þríhyrningur, þríhyrningur Pascals &  &  &  \\
\hline
3970 & Peano-ferill &  &  &  \\
\hline
3971 & Peano-talnafræði, Peano-reikningur &  &  &  \\
\hline
3972 & Pfaff-ákveða &  &  &  \\
\hline
3973 & Pfaff-ferill &  &  &  \\
\hline
3974 & Pfaff-form, \(1\)-deildaform, \(1\)-diffurform &  &  &  \\
\hline
3975 & Phragmén-Lindelöf-fall, vísifall &  &  &  \\
\hline
3976 & pí (\(\pi\)) &  &  &  \\
\hline
3977 & Picard-setningin meiri, stóra Picard-setningin &  &  &  \\
\hline
3978 & Picard-setningin minni, litla Picard-setningin &  &  &  \\
\hline
3979 & platónskt net &  &  &  \\
\hline
3980 & platónskur &  &  &  \\
\hline
3981 & plús, plúsmerki, samlagningarmerki &  &  &  \\
\hline
3982 & plúsformerki, jákvætt formerki &  &  &  \\
\hline
3983 & Poisson-dreifing &  &  &  \\
\hline
3984 & Poisson-ferli &  &  &  \\
\hline
3985 & pólskt rúm &  &  &  \\
\hline
3986 & próf &  &  &  \\
\hline
3987 & próf &  &  &  \\
\hline
3988 & prófhending &  &  &  \\
\hline
3989 & prófunarfall &  &  &  \\
\hline
3990 & prósent, hundraðshluti, hundraðstala &  &  &  \\
\hline
3991 & prósenta &  &  &  \\
\hline
3992 & prósentureikningur &  &  &  \\
\hline
3993 & prósentustöplarit, hlutfallsstöplarit &  &  &  \\
\hline
3994 & punktað mengi &  &  &  \\
\hline
3995 & punktað rúm &  &  &  \\
\hline
3996 & punktamengi, mengi &  &  &  \\
\hline
3997 & punktaröð, punktavöndur &  &  &  \\
\hline
3998 & punktferli &  &  &  \\
\hline
3999 & punktmat, stakt mat &  &  &  \\
\hline
4000 & punktur sem er ekki hvarfpunktur &  &  &  \\
\hline
4001 & punktur, depill &  &  &  \\
\hline
4002 & Pýþagórasarregla, setning Pýþagórasar & \textit{Setning Pýþagórasar} segir að fyrir <ref idTerm=4173>rétthyrndan þríhyrning</ref> með <ref idTerm=0>hliðarlengdir</ref> $a,b$ og $c$ þar sem $c$ er lengd <ref idTerm=2789>langhliðarinnar</ref> gildi \[c^2 = a^2 + b^2\] &  &  \\
\hline
4003 & Pýþagórasarsvið, pýþagórskt svið &  &  &  \\
\hline
4004 & pýþagórsk aljafna &  &  &  \\
\hline
4005 & pýþagórsk þrennd &  &  &  \\
\hline
4006 & pýþagórskur þríhyrningur &  &  &  \\
\hline
4007 & raða &  &  &  \\
\hline
4008 & raða &  &  &  \\
\hline
4009 & raðað mengi, raðmengi, hlutraðað mengi, hlaðmengi &  &  &  \\
\hline
4010 & raðað úrtak &  &  &  \\
\hline
4011 & raðaður grunnur, raðgrunnur &  &  &  \\
\hline
4012 & raðaður, hlutraðaður &  &  &  \\
\hline
4013 & raðaður, rað- &  &  &  \\
\hline
4014 & raðanlegur &  &  &  \\
\hline
4015 & raðanleiki &  &  &  \\
\hline
4016 & raðarlausn &  &  &  \\
\hline
4017 & raðeinsmótun &  &  &  \\
\hline
4018 & raðfylgni &  &  &  \\
\hline
4019 & raðfylgnistuðull &  &  &  \\
\hline
4020 & raðfylgnistuðull Kendalls &  &  &  \\
\hline
4021 & raðfylgnistuðull Spearmans &  &  &  \\
\hline
4022 & raðgerð &  &  &  \\
\hline
4023 & raðgrúpa, röðuð grúpa &  &  &  \\
\hline
4024 & raðhending, sætishending &  &  &  \\
\hline
4025 & raðhyrna &  &  &  \\
\hline
4026 & ráðinn, ákveðinn &  &  &  \\
\hline
4027 & raðkvarði, raðstigi &  &  &  \\
\hline
4028 & raðmynstur &  &  &  \\
\hline
4029 & raðrækinn &  &  &  \\
\hline
4030 & raðslétta &  &  &  \\
\hline
4031 & raðstig &  &  &  \\
\hline
4032 & raðsvið, raðað svið, raðkroppur, raðaður kroppur &  &  &  \\
\hline
4033 & raðtölu- &  &  &  \\
\hline
4034 & raðtölureikningur &  &  &  \\
\hline
4035 & raðtöluruna &  &  &  \\
\hline
4036 & raðvarpslétta, röðuð varpslétta &  &  &  \\
\hline
4037 & raðvildarslétta, röðuð vildarslétta &  &  &  \\
\hline
4038 & raðþrennd, þrennd, röðuð þrennd &  &  &  \\
\hline
4039 & ræð tala & <ref idTerm=4093>Rauntala</ref> sem er hægt að rita sem <ref idTerm=234>almennt brot</ref> &  &  \\
\hline
4040 & ræður &  &  &  \\
\hline
4041 & ræður ferill &  &  &  \\
\hline
4042 & ræma &  &  &  \\
\hline
4043 & rætt fall &  &  &  \\
\hline
4044 & rakið fall &  &  &  \\
\hline
4045 & rakið mengi &  &  &  \\
\hline
4046 & rakið reiknirit &  &  &  \\
\hline
4047 & rakin nálgunaraðferð &  &  &  \\
\hline
4048 & rakin runa &  &  &  \\
\hline
4049 & rakin umsögn &  &  &  \\
\hline
4050 & rakin vensl &  &  &  \\
\hline
4051 & rakinn &  &  &  \\
\hline
4052 & rakningarformúla &  &  &  \\
\hline
4053 & rakningarfræði &  &  &  \\
\hline
4054 & ramma &  &  &  \\
\hline
4055 & rammanlegur &  &  &  \\
\hline
4056 & rammanleiki &  &  &  \\
\hline
4057 & rammasvið &  &  &  \\
\hline
4058 & rammi &  &  &  \\
\hline
4059 & rangstöðuaðferð &  &  &  \\
\hline
4061 & rás, lokuð keðja &  &  &  \\
\hline
4062 & rás, rásuð uppstokkun, rásuð umröðun &  &  &  \\
\hline
4063 & rásað hlutmengi &  &  &  \\
\hline
4064 & rásað rúm &  &  &  \\
\hline
4065 & rásað svið &  &  &  \\
\hline
4066 & rásaður &  &  &  \\
\hline
4067 & rásaður ferhyrningur, hringferhyrningur &  &  &  \\
\hline
4068 & rásaður turn &  &  &  \\
\hline
4069 & rásaður vigur, rásaður vektor &  &  &  \\
\hline
4070 & rásaframsetning &  &  &  \\
\hline
4071 & rásuð grúpa &  &  &  \\
\hline
4072 & rásuð útvíkkun, rásuð víkkun &  &  &  \\
\hline
4073 & raufasvæði &  &  &  \\
\hline
4074 & raun- &  &  &  \\
\hline
4075 & raunbreyta, raunbreytistærð &  &  &  \\
\hline
4076 & raundreifingarfall, hlutfallslegt safntíðnifall, mælt dreifingarfall &  &  &  \\
\hline
4077 & raunfágað fall &  &  &  \\
\hline
4078 & raunfágað mynstur &  &  &  \\
\hline
4079 & raunfágað rúm &  &  &  \\
\hline
4080 & raunfágaður &  &  &  \\
\hline
4081 & raunfáguð víðátta &  &  &  \\
\hline
4082 & raunfall, raungilt fall, rauntölugilt fall &  &  &  \\
\hline
4083 & raunfallagreining, raungreining &  &  &  \\
\hline
4084 & raunfylki, rauntalnafylki &  &  &  \\
\hline
4085 & raungildur &  &  &  \\
\hline
4086 & raungildur, rauntölugildur &  &  &  \\
\hline
4087 & raunhluti &  &  &  \\
\hline
4088 & raunlína &  &  &  \\
\hline
4089 & raunlokun &  &  &  \\
\hline
4090 & raunpunktur &  &  &  \\
\hline
4091 & raunrúmfræði &  &  &  \\
\hline
4092 & raunslétta &  &  &  \\
\hline
4093 & rauntala &  &  &  \\
\hline
4094 & rauntalnakerfi &  &  &  \\
\hline
4095 & rauntalnarúm &  &  &  \\
\hline
4096 & rauntalnaslétta &  &  &  \\
\hline
4097 & rauntalnasvið &  &  &  \\
\hline
4098 & raunvarplína, varplína yfir rauntölurnar &  &  &  \\
\hline
4099 & raunvarprúm, varprúm yfir rauntölurnar &  &  &  \\
\hline
4100 & raunvarpslétta, varpslétta yfir rauntölurnar &  &  &  \\
\hline
4101 & raunvigurrúm, raunlínulegt rúm, vigurrúm yfir rauntölurnar, línulegt rúm yfir rauntölurnar &  &  &  \\
\hline
4102 & refsifall &  &  &  \\
\hline
4103 & regla &  &  &  \\
\hline
4104 & regla, lögmál &  &  &  \\
\hline
4105 & regluleg skilyrt dreifing &  &  &  \\
\hline
4106 & regluleg strýta, reglulegur píramíti &  &  &  \\
\hline
4107 & reglulegt gildi &  &  &  \\
\hline
4108 & reglulegt keilusnið &  &  &  \\
\hline
4109 & reglulegt mál &  &  &  \\
\hline
4110 & reglulegt net &  &  &  \\
\hline
4111 & reglulegt rúm &  &  &  \\
\hline
4112 & reglulegt víðerni &  &  &  \\
\hline
4113 & reglulegur &  &  &  \\
\hline
4114 & reglulegur að innan &  &  &  \\
\hline
4115 & reglulegur að utan &  &  &  \\
\hline
4116 & reglulegur áttflötungur &  &  &  \\
\hline
4117 & reglulegur ferilbogi &  &  &  \\
\hline
4118 & reglulegur ferill &  &  &  \\
\hline
4119 & reglulegur fjórflötungur &  &  &  \\
\hline
4120 & reglulegur margflötungur, platónskur margflötungur &  &  &  \\
\hline
4121 & reglulegur marghyrningur, reglulegur hyrningur &  &  &  \\
\hline
4122 & reglulegur punktur &  &  &  \\
\hline
4123 & reglulegur sérstöðupunktur &  &  &  \\
\hline
4124 & reglulegur strendingur &  &  &  \\
\hline
4125 & reglulegur tólfflötungur &  &  &  \\
\hline
4126 & reglulegur tvítugflötungur &  &  &  \\
\hline
4127 & regluleiki &  &  &  \\
\hline
4128 & reglustika &  &  &  \\
\hline
4129 & reikna, reikna út &  &  &  \\
\hline
4130 & reikna, reikna út &  &  &  \\
\hline
4131 & reiknanlegur &  &  &  \\
\hline
4133 & reiknanleikafræði &  &  &  \\
\hline
4134 & reiknanleiki &  &  &  \\
\hline
4135 & reikni- &  &  &  \\
\hline
4136 & reikniaðferð &  &  &  \\
\hline
4137 & reikniborð, reiknibretti &  &  &  \\
\hline
4138 & reiknihending, lýsifall &  &  &  \\
\hline
4139 & reiknimeistari, reikningsmaður, reikningsgarpur &  &  &  \\
\hline
4140 & reiknimeistari, reikningsmeistari &  &  &  \\
\hline
4141 & reiknimynd &  &  &  \\
\hline
4142 & reiknimyndafræði &  &  &  \\
\hline
4143 & reiknings-, reikni- &  &  &  \\
\hline
4144 & reikningsaðferð &  &  &  \\
\hline
4145 & reikningsaðgerð, grunnreikningsaðgerð &  &  &  \\
\hline
4146 & reikningsdæmi &  &  &  \\
\hline
4147 & reikningur í tugakerfi &  &  &  \\
\hline
4148 & reikningur í tylftakerfi &  &  &  \\
\hline
4149 & reikningur, talnareikningur, reikningslist, tölvísi &  &  &  \\
\hline
4150 & reikniregla &  &  &  \\
\hline
4151 & reiknirit Evklíðs, keðjudeiling &  &  &  \\
\hline
4152 & reiknirit Gauss, útrýmingaraðferð Gauss, eyðingaraðferð Gauss &  &  &  \\
\hline
4153 & reiknirit, algórismi, reikningsaðferð &  &  &  \\
\hline
4154 & reiknirits-, reikni- &  &  &  \\
\hline
4155 & reikniskekkja &  &  &  \\
\hline
4156 & reiknistokkur &  &  &  \\
\hline
4157 & reiknivél, reiknir &  &  &  \\
\hline
4158 & reiknivilla &  &  &  \\
\hline
4159 & reikul runa, hvarflandi runa &  &  &  \\
\hline
4160 & reisa &  &  &  \\
\hline
4161 & reitafylki &  &  &  \\
\hline
4162 & reitahögun &  &  &  \\
\hline
4163 & reitarit &  &  &  \\
\hline
4164 & reitur &  &  &  \\
\hline
4166 & rekjanlega teljanlegur &  &  &  \\
\hline
4167 & rétt horn & <ref idTerm=2198>Horn</ref> kallast \textit{rétt horn} ef <ref idTerm=334>armar</ref> þess eru <ref idTerm=6469>þverstæðir</ref>. Slíkt horn <ref idTerm=2223>mælist</ref> $\frac{\pi}{2}$ <ref idTerm=478>radíana</ref> eða 90 <ref idTerm=1667>gráður</ref> &  &  \\
\hline
4168 & réttanlegur &  &  &  \\
\hline
4169 & réttanlegur ferill &  &  &  \\
\hline
4170 & rétthyrnd dreifing &  &  &  \\
\hline
4171 & rétthyrndur &  &  &  \\
\hline
4172 & rétthyrndur &  &  &  \\
\hline
4173 & rétthyrndur þríhyrningur & <ref idTerm=6438>Þríhyrningur</ref> með eitt <ref idTerm=4167>rétt horn</ref> &  &  \\
\hline
4174 & rétthyrni &  &  &  \\
\hline
4175 & rétthyrningur & <ref idTerm=1127>Ferhyrningur</ref> með öll <ref idTerm=2198>horn</ref> <ref idTerm=4167>rétt</ref> &  &  \\
\hline
4176 & réttilína &  &  &  \\
\hline
4177 & réttislétta &  &  &  \\
\hline
4178 & réttsælis, sólarsinnis &  &  &  \\
\hline
4179 & réttsköpuð segð, segð &  &  &  \\
\hline
4180 & réttsköpuð stæða &  &  &  \\
\hline
4181 & réttstæð keila &  &  &  \\
\hline
4182 & réttstæð snúðkeila, réttstæð hringkeila &  &  &  \\
\hline
4183 & réttstæð strýta, réttstæður píramíti &  &  &  \\
\hline
4184 & sívalningur, réttstæður sívalningur & <ref idTerm=6532>Rúmmynd</ref> sem afmarkast af tveimur jafnstórum <ref idTerm=6589>hringskífum</ref>, sem liggja í <ref idTerm=4509>samsíða</ref> <ref idTerm=4884>sléttum</ref> og hafa miðjur sem liggja á <ref idTerm=2911>línu</ref> sem er <ref idTerm=6469>hornrétt</ref> á skífurnar, og <ref idTerm=488>sveigðum</ref> <ref idTerm=4171>rétthyrndum</ref> <ref idTerm=4660>hliðarfleti</ref> sem fæst með því að tengja saman <ref idTerm=4001>punkta</ref> <ref idTerm=2295>hringjanna</ref> sem skífurnar afmarkast af með <ref idTerm=2958>línustriki</ref> &  & Dálítið óþjál skilgreining \\
\hline
4185 & réttstæður snúðsívalningur, réttstæður hringsívalningur &  &  &  \\
\hline
4186 & réttsýnishringur &  &  &  \\
\hline
4187 & réttur &  &  &  \\
\hline
4188 & réttur tölustafur &  &  &  \\
\hline
4189 & Riemann-firð &  &  &  \\
\hline
4190 & Riemann-flötur &  &  &  \\
\hline
4191 & Riemann-heildi, Riemann-tegur &  &  &  \\
\hline
4192 & Riemann-krappaþinur, Riemann-sveigjuþinur, krappaþinur Riemanns, sveigjuþinur Riemanns &  &  &  \\
\hline
4193 & Riemann-samgangur &  &  &  \\
\hline
4194 & Riemann-Stieltjes-heildi, Riemann-Stieltjes-tegur &  &  &  \\
\hline
4195 & Riemann-víðátta &  &  &  \\
\hline
4196 & Riesz-framsetningarsetning &  &  &  \\
\hline
4197 & rífkandi &  &  &  \\
\hline
4198 & ríki &  &  &  \\
\hline
4199 & ríki af ríkjum, ríkjaríki &  &  &  \\
\hline
4200 & ríkjafræði &  &  &  \\
\hline
4201 & ríkuleg tala &  &  &  \\
\hline
4202 & rim &  &  &  \\
\hline
4203 & ris &  &  &  \\
\hline
4204 & risaðferð &  &  &  \\
\hline
4205 & risaðferð &  &  &  \\
\hline
4206 & risaðferð mesta bratta &  &  &  \\
\hline
4207 & risskref &  &  &  \\
\hline
4208 & rit &  &  &  \\
\hline
4209 & ritháttur &  &  &  \\
\hline
4210 & röð &  &  &  \\
\hline
4211 & röð &  &  &  \\
\hline
4214 & röðuð skipting &  &  &  \\
\hline
4220 & róf, litróf &  &  &  \\
\hline
4221 & róffræði, litrófsfræði &  &  &  \\
\hline
4222 & rófgeisli &  &  &  \\
\hline
4223 & rófgreining &  &  &  \\
\hline
4224 & rófheildi, róftegur &  &  &  \\
\hline
4225 & rófliðun &  &  &  \\
\hline
4226 & rófmál &  &  &  \\
\hline
4227 & rófsetning &  &  &  \\
\hline
4247 & rómverskar tölur, rómversk tölutákn &  &  &  \\
\hline
4248 & rós, rósarferill &  &  &  \\
\hline
4249 & rót &  &  &  \\
\hline
4250 & rót &  &  &  \\
\hline
4251 & rót &  &  &  \\
\hline
4252 & rót &  &  &  \\
\hline
4253 & rót, lausn & Ef $a$ er <ref idTerm=4093>rauntala</ref> þannig að $a \geq 0$ og $n \geq 1$ er <ref idTerm=3475>náttúruleg tala</ref> þá hefur <ref idTerm=2478>jafnan</ref> $x^n = a$ nákvæmlega eina lausn $t \geq 0$ og þessi lausn er kölluð $n$\textit{-ta rótin} af $a$ og táknuð með $\sqrt[n]{a}$. <ref idTerm=5397>Talan</ref> $a$ kallast \textit{rótarstofn} og $n$ kallast \textit{rótarvísir} &  & Pússa orðalag \\
\hline
4254 & rót, núllstöð & Fyrir <ref idTerm=1948>margliðu</ref> $P$ kallast <ref idTerm=4093>rauntala</ref> $r$ þannig að $P(r)=0$ \textit{rót} eða \textit{núllstöð} margliðunnar &  & Minnast á margliður sem eru ekki raungildar í skýringu \\
\hline
4255 & rótarfall & <ref idTerm=1078>Fall</ref> sem varpar sérhverju <ref idTerm=5128>staki</ref> $x$ úr <ref idTerm=1343>skilgreiningarmenginu</ref> í $n$-tu <ref idTerm=4253>rót</ref> þess $\sqrt[n]{x}$, þar sem $n \geq 2$ er <ref idTerm=3475>náttúruleg tala</ref> &  &  \\
\hline
4256 & rótaríðal &  &  &  \\
\hline
4257 & rótarpróf &  &  &  \\
\hline
4258 & rótarstæða &  &  &  \\
\hline
4259 & rótarstofn &  &  &  \\
\hline
4260 & rótartré &  &  &  \\
\hline
4261 & rótarvísir &  &  &  \\
\hline
4262 & rótlaus, rótvana, hvirflalaus &  &  &  \\
\hline
4263 & rudd algebra &  &  &  \\
\hline
4264 & rudd stallagerð, rudd þrepagerð, rudd þrepaskipan &  &  &  \\
\hline
4265 & ruðningur &  &  &  \\
\hline
4266 & rúm &  &  &  \\
\hline
4267 & rúm sem fullnægir fyrri teljanleikafrumsendu &  &  &  \\
\hline
4268 & rúm sem fullnægir seinni teljanleikafrumsendu &  &  &  \\
\hline
4269 & rúm-, rúmstæður, rúmlægur &  &  &  \\
\hline
4270 & rúmferill, rúmlægur ferill, undinn ferill &  &  &  \\
\hline
4271 & rúmfræði &  &  &  \\
\hline
4272 & rúmfræði fjölda, talningarrúmfræði &  &  &  \\
\hline
4273 & rúmfræði talna &  &  &  \\
\hline
4274 & rúmfræðilegt meðaltal, faldmeðaltal &  &  &  \\
\hline
4275 & rúmfræðilegur, rúmfræði- &  &  &  \\
\hline
4276 & rúmheildi &  &  &  \\
\hline
4277 & rúmhnit &  &  &  \\
\hline
4278 & rúmhorn &  &  &  \\
\hline
4279 & rúmlægur ferhliðungur &  &  &  \\
\hline
4280 & rúmlægur marghyrningur &  &  &  \\
\hline
4281 & rúmmál, rúmtak, rýmd &  &  &  \\
\hline
4282 & rúmmálsform &  &  &  \\
\hline
4283 & rúmmálsfrymi &  &  &  \\
\hline
4284 & rúmpunktur, punktur í rúmi &  &  &  \\
\hline
4285 & rúmstæða, stæða í rúmi &  &  &  \\
\hline
4286 & runa &  &  &  \\
\hline
4287 & runa, samstæða &  &  &  \\
\hline
4288 & runa, töluhlaup &  &  &  \\
\hline
4289 & runufullkominn &  &  &  \\
\hline
4290 & runulokaður &  &  &  \\
\hline
4291 & runurúm &  &  &  \\
\hline
4292 & runusamleitinn &  &  &  \\
\hline
4293 & runuþjappaður &  &  &  \\
\hline
4294 & Russell-þversögn &  &  &  \\
\hline
4295 & rússnesk talnagrind, rússnesk kúlnagrind, rússnesk reiknigrind &  &  &  \\
\hline
4296 & rutt leifakerfi &  &  &  \\
\hline
4297 & ryðja &  &  &  \\
\hline
4298 & sæti &  &  &  \\
\hline
4299 & sæti &  &  &  \\
\hline
4300 & sætiskerfi &  &  &  \\
\hline
4301 & sætisritun, sætistalnaritun &  &  &  \\
\hline
4302 & sætisröðun &  &  &  \\
\hline
4304 & sætistala, sæti &  &  &  \\
\hline
4305 & safn &  &  &  \\
\hline
4306 & safn, mengi &  &  &  \\
\hline
4307 & safn, samsafn, mengi &  &  &  \\
\hline
4308 & safn, samsafn, mengi &  &  &  \\
\hline
4309 & safnast fyrir, hlaðast upp &  &  &  \\
\hline
4310 & safneðlisfræði, safnaflfræði &  &  &  \\
\hline
4311 & safnferill, summuferill &  &  &  \\
\hline
4312 & safntíðni, hlaðin tíðni &  &  &  \\
\hline
4313 & safntíðnifall &  &  &  \\
\hline
4314 & sáldur &  &  &  \\
\hline
4315 & sáldur Eratosþenesar &  &  &  \\
\hline
4316 & samanburðarfall &  &  &  \\
\hline
4317 & samanburðarhópur, viðmiðunarhópur &  &  &  \\
\hline
4318 & samanburðarpróf &  &  &  \\
\hline
4319 & samanburðarsetning &  &  &  \\
\hline
4320 & samanburður &  &  &  \\
\hline
4321 & samanhangandi mengi, samhangandi mengi &  &  &  \\
\hline
4322 & samanhangandi, samhangandi, tengdur &  &  &  \\
\hline
4323 & samantekt, samtak, samtekt &  &  &  \\
\hline
4324 & samása &  &  &  \\
\hline
4325 & sambærilegur &  &  &  \\
\hline
4326 & samblástur &  &  &  \\
\hline
4327 & sambrigðin afleiða &  &  &  \\
\hline
4328 & sambrigðinn varpi, meðbrigðinn varpi &  &  &  \\
\hline
4329 & sambrigðinn þinur, meðbrigðinn þinur &  &  &  \\
\hline
4330 & sambrigðni &  &  &  \\
\hline
4331 & sambrigður, meðbrigður &  &  &  \\
\hline
4332 & samdeilir &  &  &  \\
\hline
4333 & samdrægni, herpanleiki &  &  &  \\
\hline
4334 & samdrægur, herpanlegur &  &  &  \\
\hline
4335 & samdráttur &  &  &  \\
\hline
4336 & samdreginn þinur &  &  &  \\
\hline
4337 & samdreifing &  &  &  \\
\hline
4338 & sameiginleg lóðlína, sameiginleg þverlína &  &  &  \\
\hline
4339 & sameiginlegt mál &  &  &  \\
\hline
4340 & sameiginlegur þáttur &  &  &  \\
\hline
4341 & sameinað úrtak &  &  &  \\
\hline
4342 & samfallandi &  &  &  \\
\hline
4343 & samfeldi, sameiginlegt margfeldi & <ref idTerm=3475>Náttúruleg tala</ref> $n$ kallast \textit{samfeldi} tveggja <ref idTerm=1896>heiltalna</ref> ef þær <ref idTerm=1599>ganga báðar upp í</ref> $n$ &  &  \\
\hline
4344 & samfella &  &  &  \\
\hline
4345 & samfelld hending &  &  &  \\
\hline
4346 & samfelldni &  &  &  \\
\hline
4347 & samfelldni frá hægri &  &  &  \\
\hline
4348 & samfelldni frá vinstri &  &  &  \\
\hline
4349 & samfelldni í jöfnum mæli, jöfn samfelldni &  &  &  \\
\hline
4350 & samfelldnibil &  &  &  \\
\hline
4351 & samfelldnipunktur &  &  &  \\
\hline
4352 & samfelldur &  &  &  \\
\hline
4353 & samfelldur á köflum &  &  &  \\
\hline
4354 & samfelldur deildanleiki, samfelldur diffranleiki &  &  &  \\
\hline
4355 & samfelldur frá hægri &  &  &  \\
\hline
4356 & samfelldur frá vinstri &  &  &  \\
\hline
4357 & samfelldur í jöfnum mæli &  &  &  \\
\hline
4358 & samfellt deildanlegur, samfellt diffranlegur &  &  &  \\
\hline
4359 & samfellt fall &  &  &  \\
\hline
4360 & samfellt ferli &  &  &  \\
\hline
4361 & samfellt róf &  &  &  \\
\hline
4362 & samfellt samleitinn &  &  &  \\
\hline
4363 & samfellt slembiferli &  &  &  \\
\hline
4364 & samfellufrumsenda Cantors, samfellufrumsetning Cantors, samfelldnifrumsenda Cantors, samfelldnifrumsetning Cantors &  &  &  \\
\hline
4365 & samfellufrumsenda Dedekinds, samfellufrumsetning Dedekinds, samfelldnifrumsenda Dedekinds, samfelldnifrumsetning Dedekinds &  &  &  \\
\hline
4366 & samfellufrumsenda, samfellufrumsetning, samfelldnifrumsenda, samfelldnifrumsetning &  &  &  \\
\hline
4367 & samfellutilgáta &  &  &  \\
\hline
4368 & samfellutilgáta Cantors &  &  &  \\
\hline
4369 & samfylgni, samdreifni, samvik &  &  &  \\
\hline
4370 & samfylgnifall, samvikafall &  &  &  \\
\hline
4371 & samfylgnifylki, samvikafylki &  &  &  \\
\hline
4372 & samfylgnigreining, samvikagreining &  &  &  \\
\hline
4373 & samfylgnitala, samfylgnivísir &  &  &  \\
\hline
4374 & samgangsform &  &  &  \\
\hline
4375 & samgangshliðrun &  &  &  \\
\hline
4376 & samgangur &  &  &  \\
\hline
4377 & samheldinn &  &  &  \\
\hline
4378 & samheldni &  &  &  \\
\hline
4379 & samhending &  &  &  \\
\hline
4380 & samhengi &  &  &  \\
\hline
4381 & samhengi &  &  &  \\
\hline
4382 & samhengisþáttur, þáttur &  &  &  \\
\hline
4384 & samhliðungur &  &  &  \\
\hline
4385 & samhliðungur (í \(n\) víddum) &  &  &  \\
\hline
4386 & samhlutfalla, í beinu hlutfalli við, í réttu hlutfalli við &  &  &  \\
\hline
4387 & samhvarfsás, samhverfuás &  &  &  \\
\hline
4388 & samhverf algebra &  &  &  \\
\hline
4389 & samhverf grúpa &  &  &  \\
\hline
4390 & samhverf margliða &  &  &  \\
\hline
4391 & samhverf vensl &  &  &  \\
\hline
4392 & samhverfa, samhvarf &  &  &  \\
\hline
4393 & samhverfing &  &  &  \\
\hline
4394 & samhverfni &  &  &  \\
\hline
4395 & samhverft deildaform, samhverft diffurform &  &  &  \\
\hline
4396 & samhverft fall &  &  &  \\
\hline
4397 & samhverft fylki &  &  &  \\
\hline
4398 & samhverft tvílínulegt form &  &  &  \\
\hline
4399 & samhverfugrúpa &  &  &  \\
\hline
4400 & samhverfugrúpa &  &  &  \\
\hline
4401 & samhverfumiðja, samhvarfsmiðja &  &  &  \\
\hline
4402 & samhverfur &  &  &  \\
\hline
4403 & samhverfur mismunur &  &  &  \\
\hline
4404 & samhverfur virki &  &  &  \\
\hline
4405 & samhverfur þinur &  &  &  \\
\hline
4406 & samhverfuslétta &  &  &  \\
\hline
4407 & samkvæmni &  &  &  \\
\hline
4408 & samkvæmt mat &  &  &  \\
\hline
4409 & samkvæmt próf &  &  &  \\
\hline
4410 & samkvæmur, mótsagnalaus &  &  &  \\
\hline
4411 & samlagning &  &  &  \\
\hline
4412 & samlagning &  &  &  \\
\hline
4413 & samlagning fjöldatalna &  &  &  \\
\hline
4414 & samlagning horna, hornasamlagning &  &  &  \\
\hline
4415 & samlagningar- &  &  &  \\
\hline
4416 & samlagningargrúpa &  &  &  \\
\hline
4417 & samlagningarhlutleysa, núllstak, núll &  &  &  \\
\hline
4418 & samlagningarlokaður, lokaður með tilliti til samlagningar &  &  &  \\
\hline
4419 & samlagningarregla &  &  &  \\
\hline
4420 & samlagningarríki &  &  &  \\
\hline
4421 & samlagningarsetning &  &  &  \\
\hline
4422 & samlagningartafla &  &  &  \\
\hline
4423 & samlagningarumhverfa, samlagningarandhverfa &  &  &  \\
\hline
4424 & samlega &  &  &  \\
\hline
4425 & samleggjandi, liður &  &  &  \\
\hline
4426 & samleggjanlegur &  &  &  \\
\hline
4427 & samleggjanleiki &  &  &  \\
\hline
4428 & samleguaðferð &  &  &  \\
\hline
4429 & samleifa &  &  &  \\
\hline
4430 & samleitið heildi, samleitið tegur &  &  &  \\
\hline
4431 & samleitin en ekki alsamleitin &  &  &  \\
\hline
4432 & samleitin röð &  &  &  \\
\hline
4433 & samleitin runa &  &  &  \\
\hline
4434 & samleitinn &  &  &  \\
\hline
4435 & samleitinn í hverjum punkti &  &  &  \\
\hline
4436 & samleitinn í jöfnum mæli &  &  &  \\
\hline
4437 & samleitni &  &  &  \\
\hline
4438 & samleitni fyrstu stéttar &  &  &  \\
\hline
4439 & samleitni í hverjum punkti &  &  &  \\
\hline
4440 & samleitni í jöfnum mæli, jöfn samleitni &  &  &  \\
\hline
4441 & samleitnibil &  &  &  \\
\hline
4442 & samleitnigeisli &  &  &  \\
\hline
4443 & samleitnihálfslétta &  &  &  \\
\hline
4444 & samleitnihegðun &  &  &  \\
\hline
4445 & samleitnihraði &  &  &  \\
\hline
4447 & samleitnihringskífa &  &  &  \\
\hline
4448 & samleitnihringur &  &  &  \\
\hline
4449 & samleitnilína &  &  &  \\
\hline
4450 & samleitnimark, samleitniláhnit &  &  &  \\
\hline
4451 & samleitnimengi, samleitnisvæði &  &  &  \\
\hline
4452 & samleitnipunktur &  &  &  \\
\hline
4453 & samleitniraðstig, samleitnistétt &  &  &  \\
\hline
4454 & samleitniskilyrði, samleitnipróf &  &  &  \\
\hline
4455 & samleitnisvæði &  &  &  \\
\hline
4456 & samleitnisvæði &  &  &  \\
\hline
4457 & samliggjandi heilar tölur (í vaxandi röð) &  &  &  \\
\hline
4458 & samliggjandi, í röð &  &  &  \\
\hline
4459 & samlína, á sömu línu &  &  &  \\
\hline
4460 & samlínuleiki, línulega ({\it kvk.}) &  &  &  \\
\hline
4461 & samlínumótun, samlínuvörpun &  &  &  \\
\hline
4462 & sammælanlegur &  &  &  \\
\hline
4463 & sammælanleiki &  &  &  \\
\hline
4464 & sammengi & \textit{Sammengi} tveggja <ref idTerm=3270>mengja</ref> $A$ og $B$ er mengi þeirra <ref idTerm=5128>staka</ref> sem eru að minnsta kosti í öðru þeirra. Þetta mengi er táknað með $A \cup B$ (lesið: $A$ sam $B$) og það má rita á forminu \[ A \cup B = \{x \mid x \in A \text{ eða } x \in B\}. \] &  & Setja í skýringu: Á Venn-myndinni hér að neðan er $A \cup B$ litað rautt. (ásamt mynd)

Til að einfalda rithátt er sammengi $n$ mengja $A_1 \cup A_2 \cup \cdots \cup A_n$ oft ritað á forminu $\bigcup_{i=1}^n A_i$.

Sammengi eru tengin og víxlin, þ.e. fyrir öll mengi $A$, $B$ og $C$ gildir að \[ (A \cup B) \cup C = A \cup (B \cup C) \quad \text{og} \quad A \cup B = B \cup A. \] Jafnframt eru sammengi dreifin yfir sniðmengi, þ.e. fyrir öll mengi $A$, $B$ og $C$ gildir að \[ A \cup (B \cap C) = (A \cup B) \cap (A \cup C) \quad \text{og} \quad (B \cap C) \cup A = (B \cup A) \cap (C \cup A). \] Loks gildir fyrir öll mengi $A$ að $A \cup \varnothing = A$ og $A \cup A = A$. \\
\hline
4465 & sammengisfrumsenda, sammengisfrumsetning, frumsenda um sammengi, frumsetning um sammengi &  &  &  \\
\hline
4466 & sammengismerki, sammengistákn &  &  &  \\
\hline
4467 & sammerking &  &  &  \\
\hline
4468 & sammiðja &  &  &  \\
\hline
4469 & samnefnari &  &  &  \\
\hline
4470 & samnefnd brot &  &  &  \\
\hline
4471 & samneitun, sameiginleg neitun &  &  &  \\
\hline
4472 & samok, samoka fall &  &  &  \\
\hline
4473 & samok, samoka fylki &  &  &  \\
\hline
4474 & samok, samoka tala, samoka tvinntala &  &  &  \\
\hline
4475 & samoka &  &  &  \\
\hline
4476 & samoka algebruleg tala &  &  &  \\
\hline
4477 & samoka fylki &  &  &  \\
\hline
4478 & samoka hlutgrúpa &  &  &  \\
\hline
4479 & samoka latneskur ferningur &  &  &  \\
\hline
4480 & samoka miðstrengur &  &  &  \\
\hline
4481 & samoka röðun &  &  &  \\
\hline
4482 & samoka virki &  &  &  \\
\hline
4483 & samoka þýtt fall &  &  &  \\
\hline
4484 & samokun &  &  &  \\
\hline
4485 & samokunarflokkur &  &  &  \\
\hline
4486 & sampunkta, sem liggja gegnum sama punkt &  &  &  \\
\hline
4487 & sampunktleiki &  &  &  \\
\hline
4488 & samruna &  &  &  \\
\hline
4489 & samrýmanlegur &  &  &  \\
\hline
4490 & samrýmanleiki &  &  &  \\
\hline
4491 & samrýnd fjölskylda &  &  &  \\
\hline
4492 & samrýndur bútur &  &  &  \\
\hline
4493 & samsama &  &  &  \\
\hline
4494 & samsemd &  &  &  \\
\hline
4495 & samsemdarmótun &  &  &  \\
\hline
4496 & samsemdarsetning &  &  &  \\
\hline
4497 & samsemdarvensl &  &  &  \\
\hline
4498 & samsemdarvörpun & <ref idTerm=3282>Vörpunin</ref> frá <ref idTerm=3270>mengi</ref> $X$ yfir í $X$ sem varpar sérhverju <ref idTerm=5128>staki</ref> úr $X$ í sjálft sig kallast \textit{samsemdarvörpun} mengisins $X$ og er táknuð með $\mathrm{id}_X$ &  &  \\
\hline
4499 & samsett dreifing &  &  &  \\
\hline
4500 & samsett Markoff-keðja &  &  &  \\
\hline
4501 & samsett tala, þættanleg tala &  &  &  \\
\hline
4502 & samsett tilgáta &  &  &  \\
\hline
4503 & samsett þríliða &  &  &  \\
\hline
4504 & samsettur atburður &  &  &  \\
\hline
4505 & samsettur, þættanlegur &  &  &  \\
\hline
4506 & samsíða leggur &  &  &  \\
\hline
4507 & samsíða lína &  &  &  \\
\hline
4508 & samsíða ofanvarp &  &  &  \\
\hline
4509 & samsíða, samhliða &  &  &  \\
\hline
4510 & samsíðufrumsenda, samsíðufrumsetning, frumsetning um samsíða línur, frumsenda um samsíða línur &  &  &  \\
\hline
4511 & samsíðuhorn &  &  &  \\
\hline
4512 & samsíðungsjafna &  &  &  \\
\hline
4513 & samsíðungsregla &  &  &  \\
\hline
4514 & samsíðungur &  &  &  \\
\hline
4515 & samsíðuskipan &  &  &  \\
\hline
4516 & samskeytanlegur &  &  &  \\
\hline
4517 & samskeyting &  &  &  \\
\hline
4518 & samskeyting & Fyrir <ref idTerm=3282>varpanir</ref> $f: X \to Y$ og $g: Y \to Z$ kallast vörpunin $g \circ f: X \to Z$ sem skilgreind er með forskriftinni $(g \circ f)(x) = g(f(x))$ \textit{samskeyting} varpananna $f$ og $g$ &  &  \\
\hline
4519 & samskeytt fall &  &  &  \\
\hline
4520 & samskeytt vörpun &  &  &  \\
\hline
4521 & samskeyttur &  &  &  \\
\hline
4522 & samskrepping &  &  &  \\
\hline
4523 & samslétta, í sömu sléttu &  &  &  \\
\hline
4524 & samslyngir &  &  &  \\
\hline
4525 & samslyngja, samhlekkja, hlekkja saman &  &  &  \\
\hline
4526 & samsnið, aljöfnuður &  &  &  \\
\hline
4527 & samsniða fylki &  &  &  \\
\hline
4528 & samsniða myndir, aljafnar myndir, eins myndir &  &  &  \\
\hline
4529 & samsniða, aljafn, eins &  &  &  \\
\hline
4530 & samsniðafrumsenda &  &  &  \\
\hline
4531 & samsömun &  &  &  \\
\hline
4532 & samstak, ok &  &  &  \\
\hline
4533 & samstétta &  &  &  \\
\hline
4534 & samstíga &  &  &  \\
\hline
4535 & samstíga hornalínugerningur &  &  &  \\
\hline
4536 & samstríkkaðar myndir &  &  &  \\
\hline
4537 & samsvarandi &  &  &  \\
\hline
4538 & samsvarandi horn &  &  &  \\
\hline
4539 & samsvipa &  &  &  \\
\hline
4540 & samsvipur &  &  &  \\
\hline
4541 & samsvörun &  &  &  \\
\hline
4542 & samtakagrannfræði, samtektagrannfræði, fléttugrannfræði &  &  &  \\
\hline
4543 & samtakarúmfræði, samtektarúmfræði, flétturúmfræði &  &  &  \\
\hline
4544 & samtenging, tenging, tengill &  &  &  \\
\hline
4545 & samtoga &  &  &  \\
\hline
4546 & samtogun, togun &  &  &  \\
\hline
4547 & samtogunarflokkur, togflokkur &  &  &  \\
\hline
4548 & samtogunarlyfting, togunarlyfting &  &  &  \\
\hline
4549 & samur &  &  &  \\
\hline
4550 & samur, alsamur, aljafn &  &  &  \\
\hline
4551 & samþætting &  &  &  \\
\hline
4552 & samþættur &  &  &  \\
\hline
4553 & samþáttur, sameiginlegur þáttur &  &  &  \\
\hline
4554 & sanna &  &  &  \\
\hline
4555 & sannanafræði &  &  &  \\
\hline
4556 & sannanleg segð, sannanleg yrðing, sannanleg fullyrðing &  &  &  \\
\hline
4557 & sannanlegur &  &  &  \\
\hline
4558 & sannanleiki &  &  &  \\
\hline
4559 & sannfall &  &  &  \\
\hline
4560 & sanngildi &  &  &  \\
\hline
4561 & sanngjarn leikur &  &  &  \\
\hline
4562 & sannleikapróf, sannpróf &  &  &  \\
\hline
4563 & sannleiki &  &  &  \\
\hline
4564 & sanntafla &  &  &  \\
\hline
4566 & segð &  &  &  \\
\hline
4567 & seinni teljanleikafrumsendan &  &  &  \\
\hline
4568 & seintekinn &  &  &  \\
\hline
4569 & sekans &  &  &  \\
\hline
4570 & sekúnda &  &  &  \\
\hline
4571 & sekúnda, hornsekúnda, bogasekúnda &  &  &  \\
\hline
4572 & sem gera átti &  &  &  \\
\hline
4573 & sem sanna átti &  &  &  \\
\hline
4574 & sem snertir að innan, sem snertir að innanverðu &  &  &  \\
\hline
4575 & sem snertir að utan, sem snertir utan frá &  &  &  \\
\hline
4576 & sennileikaaðferð &  &  &  \\
\hline
4577 & sennileikahlutfall, líknahlutfall &  &  &  \\
\hline
4578 & sennileikamat &  &  &  \\
\hline
4579 & sennileikamat &  &  &  \\
\hline
4580 & sennileikametill &  &  &  \\
\hline
4581 & sennileiki, líkur &  &  &  \\
\hline
4582 & sérfræðingur í algebru &  &  &  \\
\hline
4583 & sérfræðingur í stærðfræðigreiningu &  &  &  \\
\hline
4584 & sérlausn &  &  &  \\
\hline
4585 & sérlausn, sérstök lausn &  &  &  \\
\hline
4586 & sérlausn, sérstök lausn &  &  &  \\
\hline
4587 & sérstaða, sérstæða &  &  &  \\
\hline
4588 & sérstæð hjásvipfræði &  &  &  \\
\hline
4589 & sérstæð keðja &  &  &  \\
\hline
4590 & sérstæð rás &  &  &  \\
\hline
4591 & sérstæð svipfræði &  &  &  \\
\hline
4592 & sérstæður &  &  &  \\
\hline
4593 & sérstæður ferill &  &  &  \\
\hline
4594 & sérstæður línulegur virki, óandhverfanlegur línulegur virki &  &  &  \\
\hline
4595 & sérstætt fylki, óandhverfanlegt fylki, óumhverfanlegt fylki, óreglulegt fylki &  &  &  \\
\hline
4596 & sérstakt gildi &  &  &  \\
\hline
4597 & sérstöðuleg &  &  &  \\
\hline
4598 & sérstöðupunktur &  &  &  \\
\hline
4599 & sérstök línuleg grúpa &  &  &  \\
\hline
4600 & sértækur, óhlutbundinn, hlutfirrtur &  &  &  \\
\hline
4601 & sértak &  &  &  \\
\hline
4602 & sértekinn &  &  &  \\
\hline
4603 & sértekning, hlutfirring &  &  &  \\
\hline
4604 & setafall Riemanns, zetafall Riemanns, \(\zeta\)-fall Riemanns &  &  &  \\
\hline
4605 & setafall, zetafall, \(\zeta\)-fall &  &  &  \\
\hline
4606 & setja &  &  &  \\
\hline
4607 & setja í staðinn, setja inn &  &  &  \\
\hline
4608 & setja í sviga &  &  &  \\
\hline
4609 & setja inn &  &  &  \\
\hline
4610 & setja jafnt &  &  &  \\
\hline
4611 & setja út horn &  &  &  \\
\hline
4612 & setja út strik &  &  &  \\
\hline
4613 & setning &  &  &  \\
\hline
4614 & setning &  &  &  \\
\hline
4615 & setning um einhalla samleitni &  &  &  \\
\hline
4616 & setning um fólgin föll &  &  &  \\
\hline
4617 & setning um lokað varprit &  &  &  \\
\hline
4618 & setning um opna vörpun &  &  &  \\
\hline
4619 & setning um ótvíræða þáttun &  &  &  \\
\hline
4620 & setning um staðbundið markgildi &  &  &  \\
\hline
4621 & setning, meginsetning, höfuðsetning &  &  &  \\
\hline
4622 & setningin um fjóra hvirfilpunkta &  &  &  \\
\hline
4623 & setningin um níupunktahringinn &  &  &  \\
\hline
4624 & setningin um ytra horn & Setning í <ref idTerm=1030>Evklíðskri rúmfræði</ref> sem segir að <ref idTerm=6295>ytra horn</ref> <ref idTerm=6438>þríhyrnings</ref> sé jafn <ref idTerm=2223>stórt</ref> og samanlögð fjarlægari innri horn &  &  \\
\hline
4625 & sexflötungur & <ref idTerm=3133>Margflötungur</ref> með sex <ref idTerm=1325>hliðar</ref> &  &  \\
\hline
4626 & sexhyrndur &  &  &  \\
\hline
4627 & sexhyrningstala, sexhyrnd tala &  &  &  \\
\hline
4628 & sexhyrningur &  &  &  \\
\hline
4629 & sexstrendingur &  &  &  \\
\hline
4630 & sextándakerfi &  &  &  \\
\hline
4631 & sextugakerfi &  &  &  \\
\hline
4632 & sexund &  &  &  \\
\hline
4633 & sía &  &  &  \\
\hline
4634 & síaður &  &  &  \\
\hline
4635 & síðasta setning Fermats &  &  &  \\
\hline
4636 & sigð &  &  &  \\
\hline
4637 & sigð Hippókratesar &  &  &  \\
\hline
4638 & sigð, hringbogasigð &  &  &  \\
\hline
4639 & sígild aflfræði &  &  &  \\
\hline
4640 & sígild grúpa &  &  &  \\
\hline
4641 & sígild mengjafræði, mengjafræði Cantors &  &  &  \\
\hline
4642 & sígild rökfræði &  &  &  \\
\hline
4643 & sígild yrðingarökfræði &  &  &  \\
\hline
4644 & sigma-algebra &  &  &  \\
\hline
4645 & síkveða &  &  &  \\
\hline
4646 & Simpson-regla &  &  &  \\
\hline
4647 & sínus &  &  &  \\
\hline
4648 & sínusferill &  &  &  \\
\hline
4649 & sínusregla & Um <ref idTerm=6438>þríhyrning</ref> $\Delta ABC$ gildir \textit{sínusreglan} $$ \frac{\sin(A)}{a}=\frac{\sin(B)}{b}= \frac{\sin(C)}{c}. $$ Með öðrum orðum er <ref idTerm=2072>hlutfallið</ref> <ref idTerm=3332>milli</ref> <ref idTerm=4647>sínusar</ref> af <ref idTerm=2198>horni</ref> þríhyrnings og <ref idTerm=3400>mótlægrar</ref> <ref idTerm=2048>hliðar</ref> hornsins það sama fyrir öll hornin. Þetta hlutfall er einnig jafnt $\frac{1}{2 R}$ þar sem $R$ er <ref idTerm=2295>geisli</ref> <ref idTerm=5784>umritaðs hrings</ref> þríhyrningsins. &  &  \\
\hline
4650 & sínusvefja &  &  &  \\
\hline
4651 & síósanna &  &  &  \\
\hline
4652 & sísanna &  &  &  \\
\hline
4653 & sísannur &  &  &  \\
\hline
4654 & sístætt ástand, æstætt ástand &  &  &  \\
\hline
4655 & sístætt ferli, æstætt ferli, jafnvægisferli, ferli í jafnvægi &  &  &  \\
\hline
4656 & sístætt Markov-ferli, æstætt Markov-ferli, Markov-ferli í jafnvægi &  &  &  \\
\hline
4657 & sístætt streymi, æstætt streymi, jafnvægisstreymi &  &  &  \\
\hline
4658 & síugrunnur &  &  &  \\
\hline
4659 & síun &  &  &  \\
\hline
4660 & sívalningsflötur & <ref idTerm=3279>Faldmengi</ref> $\mathcal{C} \times I$ þar sem $\mathcal{C} \in \mathbb{R}^2$ er <ref idTerm=1131>ferill</ref> og $I \subseteq \mathbb{R}$ &  &  \\
\hline
4661 & sívalningshnit &  &  &  \\
\hline
4662 & sívalningsstúfur, stýfður sívalningur &  &  &  \\
\hline
4663 & sívalningur &  &  &  \\
\hline
4664 & sjaldgæfur atburður &  &  &  \\
\hline
4665 & sjálfgefin lausn &  &  &  \\
\hline
4666 & sjálfhverf vensl, spegilvirk vensl &  &  &  \\
\hline
4667 & sjálfhverfa, sjálfhverf vörpun &  &  &  \\
\hline
4668 & sjálfhverfni, spegilvirkni &  &  &  \\
\hline
4669 & sjálfhverft Banach-rúm &  &  &  \\
\hline
4670 & sjálfhverfumiðja &  &  &  \\
\hline
4671 & sjálfhverfur &  &  &  \\
\hline
4672 & sjálfhverfur, spegilvirkur &  &  &  \\
\hline
4673 & sjálfmóta fall &  &  &  \\
\hline
4674 & sjálfmótanagrúpa &  &  &  \\
\hline
4675 & sjálfmótanagrúpa &  &  &  \\
\hline
4676 & sjálfmótun &  &  &  \\
\hline
4677 & sjálfnykraður &  &  &  \\
\hline
4678 & sjálfoka fylki, hermískt fylki &  &  &  \\
\hline
4679 & sjálfoka virki &  &  &  \\
\hline
4680 & sjálfoka, sjálftengdur &  &  &  \\
\hline
4681 & sjálfskautaður þríhyrningur &  &  &  \\
\hline
4682 & sjálfskerandi &  &  &  \\
\hline
4683 & sjálfskerandi ferhliðungur &  &  &  \\
\hline
4684 & sjálfskerandi marghyrningur &  &  &  \\
\hline
4685 & sjálfskurðarpunktur &  &  &  \\
\hline
4686 & sjálfskurður &  &  &  \\
\hline
4687 & sjálfstýring &  &  &  \\
\hline
4688 & sjálfsvipaður &  &  &  \\
\hline
4689 & sjálfsvipur &  &  &  \\
\hline
4690 & sjálfvalda &  &  &  \\
\hline
4691 & sjálfvalda stak &  &  &  \\
\hline
4692 & sjálfvirki &  &  &  \\
\hline
4693 & sjálfvirkjafræði &  &  &  \\
\hline
4694 & sjálfvirkur &  &  &  \\
\hline
4695 & sjálfþver &  &  &  \\
\hline
4696 & sjálfþver vigur &  &  &  \\
\hline
4697 & sjöhyrningstala, sjöhyrnd tala &  &  &  \\
\hline
4698 & sjöhyrningur &  &  &  \\
\hline
4699 & sjóndeildar- &  &  &  \\
\hline
4700 & sjóndeildarjaðar &  &  &  \\
\hline
4701 & sjóndeildarlína &  &  &  \\
\hline
4702 & sjóndeildarpunktur &  &  &  \\
\hline
4703 & sjóndeildarslétta &  &  &  \\
\hline
4704 & sjóndeildarstak &  &  &  \\
\hline
4705 & sjónhorf &  &  &  \\
\hline
4706 & sjónhorfs-, í sjónhorfi &  &  &  \\
\hline
4707 & sjónhorfsás &  &  &  \\
\hline
4708 & sjónhorfsmiðja &  &  &  \\
\hline
4709 & sjónhorfsmótun, sjónhorfsfærsla &  &  &  \\
\hline
4710 & sjónhorfsmyndun &  &  &  \\
\hline
4711 & sjötta stigs &  &  &  \\
\hline
4712 & sjötta stigs ferill &  &  &  \\
\hline
4713 & sjötta stigs flötur &  &  &  \\
\hline
4714 & sjötta stigs jafna &  &  &  \\
\hline
4715 & sjöttungsmark &  &  &  \\
\hline
4716 & skæling &  &  &  \\
\hline
4717 & skælingarhlutfall &  &  &  \\
\hline
4718 & skælingarhorn &  &  &  \\
\hline
4719 & skáhæð &  &  &  \\
\hline
4720 & skáhorn &  &  &  \\
\hline
4721 & skáhyrndur þríhyrningur &  &  &  \\
\hline
4722 & skakkt ofanvarp &  &  &  \\
\hline
4723 & skakkur samhliðungur, skásettur samhliðungur &  &  &  \\
\hline
4724 & skammás &  &  &  \\
\hline
4725 & skammhlið &  &  &  \\
\hline
4726 & skammleiðarverkefni, verkefni skemmstu leiðar &  &  &  \\
\hline
4727 & skammsveifarferill &  &  &  \\
\hline
4728 & skammtímaferill &  &  &  \\
\hline
4729 & skarast &  &  &  \\
\hline
4730 & skásett hnit, skökk hnit &  &  &  \\
\hline
4731 & skásett hnitakerfi, skakkt hnitakerfi &  &  &  \\
\hline
4732 & skásett hringkeila, skökk hringkeila &  &  &  \\
\hline
4733 & skásettur hringsívalningur, skakkur hringsívalningur &  &  &  \\
\hline
4734 & skásettur, skakkur &  &  &  \\
\hline
4735 & skástrendingur, skakkur strendingur &  &  &  \\
\hline
4736 & skástýfð keila &  &  &  \\
\hline
4737 & skástýfður sívalningur &  &  &  \\
\hline
4738 & skaut &  &  &  \\
\hline
4739 & skaut, póll &  &  &  \\
\hline
4740 & skaut, póll, skauthnitamiðja, pólhnitamiðja &  &  &  \\
\hline
4741 & skautframsetning, skauthnitaframsetning, framsetning í skauthnitum (pólhnitum) &  &  &  \\
\hline
4742 & skauthluti &  &  &  \\
\hline
4743 & skauthnit, pólhnit & <ref idTerm=2145>Hnitakerfi</ref> á <ref idTerm=1031>evklíðska sléttu</ref> þar sem <ref idTerm=4001>punktar</ref> eru staðsettir með því að tilgreina <ref idTerm=1204>fjarlægð</ref> þeirra og <ref idTerm=5135>stefnu</ref> séð frá gefnum punkti $O$ sem er kallaður \textit{skaut} eða \textit{póll} &  & Tala um hvernig maður breytir frá kassahnitum í pólhnit í skýringu \\
\hline
4744 & skauthnitaferill, pólhnitaferill, ferill í skauthnitum (pólhnitum) &  &  &  \\
\hline
4745 & skauthnitajafna, pólhnitajafna &  &  &  \\
\hline
4746 & skauthnitakerfi, pólhnitakerfi &  &  &  \\
\hline
4747 & skauthnitás, skauthnitaás, pólhnitás, pólhnitaás &  &  &  \\
\hline
4748 & skauthringur &  &  &  \\
\hline
4749 & skautlína &  &  &  \\
\hline
4750 & skautun &  &  &  \\
\hline
4751 & skautunarjafna &  &  &  \\
\hline
4752 & skautunarmiðja &  &  &  \\
\hline
4753 & skautþáttun, pólþáttun &  &  &  \\
\hline
4754 & skautþríhyrningur &  &  &  \\
\hline
4755 & skeif dreifing &  &  &  \\
\hline
4756 & skeifni &  &  &  \\
\hline
4757 & skekking &  &  &  \\
\hline
4758 & skekkja &  &  &  \\
\hline
4759 & skekkjubil &  &  &  \\
\hline
4760 & skekkjufall, Gauss-skekkjufall &  &  &  \\
\hline
4761 & skekkjugreining &  &  &  \\
\hline
4762 & skekkjuliður &  &  &  \\
\hline
4763 & skekkjuliður &  &  &  \\
\hline
4764 & skekkjulíkindi &  &  &  \\
\hline
4765 & skekkjumark, skekkjuskorða &  &  &  \\
\hline
4766 & skekkjumat &  &  &  \\
\hline
4767 & skekkjuútbreiðsla &  &  &  \\
\hline
4768 & skelferill &  &  &  \\
\hline
4769 & skema &  &  &  \\
\hline
4770 & skera, sníða &  &  &  \\
\hline
4771 & skerast, sníðast, skarast &  &  &  \\
\hline
4772 & skerðing &  &  &  \\
\hline
4773 & skert dreifing &  &  &  \\
\hline
4774 & skeyta saman &  &  &  \\
\hline
4775 & skífurit &  &  &  \\
\hline
4776 & skila &  &  &  \\
\hline
4777 & skilapunktur &  &  &  \\
\hline
4778 & skilgreina með þrepun &  &  &  \\
\hline
4779 & skilgreina, skýrgreina &  &  &  \\
\hline
4780 & skilgreinandi &  &  &  \\
\hline
4781 & skilgreinandi jafna &  &  &  \\
\hline
4782 & skilgreinandi, skilgreiningarorð &  &  &  \\
\hline
4783 & skilgreinanlegur &  &  &  \\
\hline
4784 & skilgreinanleiki &  &  &  \\
\hline
4785 & skilgreining, skýrgreining &  &  &  \\
\hline
4786 & skilgreiningarbil &  &  &  \\
\hline
4787 & skilgreiningarefni &  &  &  \\
\hline
4788 & skilgreiningarmengi virkja &  &  &  \\
\hline
4789 & skilorðslaus samleitni &  &  &  \\
\hline
4790 & skilorðslaust samleitinn &  &  &  \\
\hline
4791 & skilun &  &  &  \\
\hline
4792 & skilvirkasti metill &  &  &  \\
\hline
4793 & skilvirkni &  &  &  \\
\hline
4794 & skilvirkt mat &  &  &  \\
\hline
4795 & skilvirkt próf &  &  &  \\
\hline
4796 & skilvirkur metill &  &  &  \\
\hline
4797 & skilyrði &  &  &  \\
\hline
4798 & skilyrði um fallandi keðjur, skilyrði um minnkandi keðjur &  &  &  \\
\hline
4799 & skilyrði um vaxandi keðjur, skilyrði um stígandi keðjur &  &  &  \\
\hline
4800 & skilyrðing &  &  &  \\
\hline
4801 & skilyrt dreifing &  &  &  \\
\hline
4802 & skilyrt líkindi &  &  &  \\
\hline
4803 & skilyrt mál &  &  &  \\
\hline
4804 & skilyrt meðalgildi &  &  &  \\
\hline
4805 & skilyrt meðalgildisfall &  &  &  \\
\hline
4806 & skilyrt óreiða &  &  &  \\
\hline
4807 & skilyrt próf &  &  &  \\
\hline
4808 & skilyrt samleitinn, skilorðssamleitinn &  &  &  \\
\hline
4809 & skilyrt samleitni, skilorðsbundin samleitni &  &  &  \\
\hline
4810 & skilyrt þéttleikafall &  &  &  \\
\hline
4811 & skilyrtur atburður &  &  &  \\
\hline
4812 & skilyrtur þéttleiki &  &  &  \\
\hline
4813 & skilyrtur, skilorðsbundinn &  &  &  \\
\hline
4814 & skipta í þríhyrninga &  &  &  \\
\hline
4815 & skipta, hluta &  &  &  \\
\hline
4816 & skiptanlegur í þríhyrninga &  &  &  \\
\hline
4817 & skiptanleiki í þríhyrninga &  &  &  \\
\hline
4818 & skiptihlutfall &  &  &  \\
\hline
4819 & skipting &  &  &  \\
\hline
4820 & skipting í tilfelli, skipting í tilvik &  &  &  \\
\hline
4821 & skipting, deildaskipting, flokkaskipting, hlutun &  &  &  \\
\hline
4822 & skipting, hlutun &  &  &  \\
\hline
4824 & skiptingarfall, hlutunarfall &  &  &  \\
\hline
4825 & skiptipunktur, hlutunarpunktur &  &  &  \\
\hline
4826 & skipun &  &  &  \\
\hline
4827 & skógur &  &  &  \\
\hline
4828 & skorða, takmörkun &  &  &  \\
\hline
4829 & skorðumynstur &  &  &  \\
\hline
4830 & skorðun &  &  &  \\
\hline
4831 & skorðurúm &  &  &  \\
\hline
4832 & skörun &  &  &  \\
\hline
4833 & skrá &  &  &  \\
\hline
4834 & skráarvinnsla &  &  &  \\
\hline
4835 & skref &  &  &  \\
\hline
4836 & skrefabestun &  &  &  \\
\hline
4837 & skrefaður, skrefa- &  &  &  \\
\hline
4838 & skrefagreining &  &  &  \\
\hline
4839 & skrefamat &  &  &  \\
\hline
4840 & skrefapróf &  &  &  \\
\hline
4841 & skrúfflötur &  &  &  \\
\hline
4842 & skrúflína, gormur &  &  &  \\
\hline
4843 & skrúfuflutningur, skrúfuhreyfing &  &  &  \\
\hline
4844 & skúffuregla Dirichlets, skúffuregla &  &  &  \\
\hline
4845 & skurðpunktur, sniðpunktur &  &  &  \\
\hline
4846 & skurður &  &  &  \\
\hline
4847 & skýjakljúfsknippi &  &  &  \\
\hline
4848 & skylt horn &  &  &  \\
\hline
4849 & skyndiúrtak &  &  &  \\
\hline
4850 & skýribreyta &  &  &  \\
\hline
4851 & skýringabreyta &  &  &  \\
\hline
4852 & skýringarmynd &  &  &  \\
\hline
4853 & slakbreyta &  &  &  \\
\hline
4854 & slanga &  &  &  \\
\hline
4855 & slanga &  &  &  \\
\hline
4856 & slemba &  &  &  \\
\hline
4857 & slembdur &  &  &  \\
\hline
4858 & slembi-, slembinn, tilviljanakenndur &  &  &  \\
\hline
4859 & slembiáhrif &  &  &  \\
\hline
4860 & slembibestun &  &  &  \\
\hline
4861 & slembidreifing &  &  &  \\
\hline
4862 & slembiferli &  &  &  \\
\hline
4863 & slembiferli með óháðum viðbótum &  &  &  \\
\hline
4864 & slembifræði &  &  &  \\
\hline
4865 & slembifylki &  &  &  \\
\hline
4866 & slembigangur, slembiganga &  &  &  \\
\hline
4867 & slembiheildi, slembitegur &  &  &  \\
\hline
4868 & slembileikur &  &  &  \\
\hline
4869 & slembilíkan, líkindalíkan &  &  &  \\
\hline
4870 & slembing &  &  &  \\
\hline
4871 & slembingur, slembileiki &  &  &  \\
\hline
4872 & slembióháður, óháður &  &  &  \\
\hline
4873 & slembióhæði, óhæði &  &  &  \\
\hline
4874 & slembiruna &  &  &  \\
\hline
4875 & slembiskekkja &  &  &  \\
\hline
4876 & slembistak &  &  &  \\
\hline
4877 & slembistund &  &  &  \\
\hline
4878 & slembitala &  &  &  \\
\hline
4879 & slembitalnatafla &  &  &  \\
\hline
4880 & slembiúrtak &  &  &  \\
\hline
4881 & slembiúrtaka &  &  &  \\
\hline
4882 & sleppa &  &  &  \\
\hline
4883 & sleppa &  &  &  \\
\hline
4884 & slétta, sléttur flötur, plan &  &  &  \\
\hline
4885 & sléttanlegur &  &  &  \\
\hline
4886 & sléttanlegur flötur &  &  &  \\
\hline
4887 & sléttu- &  &  &  \\
\hline
4888 & sléttuferill, ferill í sléttu &  &  &  \\
\hline
4889 & sléttuhnit &  &  &  \\
\hline
4890 & sléttuhornafræði &  &  &  \\
\hline
4891 & sléttumynd, flatarmynd &  &  &  \\
\hline
4892 & sléttunet &  &  &  \\
\hline
4893 & slétturúmfræði, flatarmyndafræði &  &  &  \\
\hline
4894 & sléttuskiki &  &  &  \\
\hline
4895 & sléttusnið &  &  &  \\
\hline
4896 & sléttustæða &  &  &  \\
\hline
4897 & sléttuþríbaugur &  &  &  \\
\hline
4898 & slönguferill &  &  &  \\
\hline
4899 & slöngugrennd &  &  &  \\
\hline
4900 & smáhringur &  &  &  \\
\hline
4901 & smárahnútur, einfaldur hnútur &  &  &  \\
\hline
4902 & smituð dreifing &  &  &  \\
\hline
4903 & snara &  &  &  \\
\hline
4904 & sneið, kúlusneið, kúlukollur &  &  &  \\
\hline
4905 & snemmtekinn &  &  &  \\
\hline
4906 & snerti-, snertil- &  &  &  \\
\hline
4907 & snertibundin, snertlabundin &  &  &  \\
\hline
4908 & snertiferill &  &  &  \\
\hline
4909 & snertihringur &  &  &  \\
\hline
4910 & snertihvel, snertikúla &  &  &  \\
\hline
4911 & snertikeila, snertlakeila &  &  &  \\
\hline
4912 & snertill, snertilína & <ref idTerm=2911>Lína</ref> sem sker <ref idTerm=2295>hring</ref> í nákvæmlega einum <ref idTerm=4001>punkti</ref> &  &  \\
\hline
4913 & snerting &  &  &  \\
\hline
4914 & snertingarslétta &  &  &  \\
\hline
4915 & snertipunktur &  &  &  \\
\hline
4916 & snertirúm, snertlarúm &  &  &  \\
\hline
4917 & snertislétta &  &  &  \\
\hline
4918 & snertistig &  &  &  \\
\hline
4919 & snertistrengur &  &  &  \\
\hline
4920 & snertivigur, snertill &  &  &  \\
\hline
4921 & snertiþáttur, snertiliður &  &  &  \\
\hline
4922 & snertlahorn &  &  &  \\
\hline
4923 & snertlasvið &  &  &  \\
\hline
4924 & snertlasvið &  &  &  \\
\hline
4925 & snið &  &  &  \\
\hline
4926 & snið &  &  &  \\
\hline
4927 & snið, mót &  &  &  \\
\hline
4928 & snið, sniðmengi, skurðmengi, skarmengi & \textit{Sniðmengi} tveggja <ref idTerm=3270>mengja</ref> $A$ og $B$ er mengi þeirra <ref idTerm=5128>staka</ref> sem eru í báðum þeirra. Þetta mengi er táknað með $A \cap B$ (lesið: $A$ snið $B$) og það má rita á forminu \[ A \cap B = \{x \mid x \in A \text{ og } x \in B\}. \] &  & Á Venn-myndinni hér að neðan er $A \cap B$ litað rautt. Til að einfalda rithátt er sniðmengi $n$ mengja $A_1 \cap A_2 \cap \cdots \cap A_n$ oft ritað á forminu $\bigcap_{i=1}^n A_i$. (ásamt mynd)

Ef mengin $A$ og $B$ hafa ekkert sameiginlegt stak, þ.e. ef $A \cap B = \varnothing$, er sagt að þau séu sundurlæg.

Sniðmengi eru tengin og víxlin, þ.e. fyrir öll mengi $A$, $B$ og $C$ gildir að \[ (A \cap B) \cap C = A \cap (B \cap C) \quad \text{og} \quad A \cap B = B \cap A. \] Jafnframt eru sniðmengi dreifin yfir sammengi, þ.e. fyrir öll mengi $A$, $B$ og $C$ gildir að \[ A \cap (B \cup C) = (A \cap B) \cup (A \cap C) \quad \text{og} \quad (B \cup C) \cap A = (B \cap A) \cup (C \cap A). \] Loks gildir fyrir öll mengi $A$ að $A \cap \varnothing = \varnothing$ og $A \cap A = A$. \\
\hline
4929 & sniðill & <ref idTerm=2911>Lína</ref> sem sker <ref idTerm=2295>hring</ref> í tveimur ólíkum <ref idTerm=4001>punktum</ref> &  &  \\
\hline
4930 & sniðkrappi, sniðsveigja, Riemann-krappi, Riemann-sveigja &  &  &  \\
\hline
4931 & sniðlína, skurðlína &  &  &  \\
\hline
4932 & sniðmargfeldni &  &  &  \\
\hline
4933 & sniðmengismerki, sniðmengistákn &  &  &  \\
\hline
4934 & sniðstak, mót &  &  &  \\
\hline
4935 & snigill, snigilferill, snigillína &  &  &  \\
\hline
4936 & snúðflötur &  &  &  \\
\hline
4937 & snúðhjólflötur &  &  &  \\
\hline
4938 & snúðkeila &  &  &  \\
\hline
4939 & keila, snúðkeila, hringkeila & <ref idTerm=6532>Rúmmynd</ref> sem afmarkast af <ref idTerm=6589>hringskífu</ref> og <ref idTerm=2579>fletinum</ref> sem fæst með því að draga <ref idTerm=2958>línustrik</ref> milli <ref idTerm=4001>punkta</ref> <ref idTerm=2295>hringsins</ref> sem skífan afmarkast af og eins punkts sem liggur á <ref idTerm=2911>línu</ref> sem er <ref idTerm=6469>hornrétt</ref> á skífuna og fer í gegnum miðju hennar &  & Þarf punkturinn að liggja fyrir utan grunnflötinn? - Rosalega óþjál skilgreining \\
\hline
4940 & snúðkeiluflötur, hringkeiluflötur &  &  &  \\
\hline
4941 & snúðsívalningsflötur, hringsívalningsflötur &  &  &  \\
\hline
4942 & snúðsívalningur, hringsívalningur &  &  &  \\
\hline
4943 & snúðsívalningur, hringsívalningur &  &  &  \\
\hline
4944 & snúður, snúðþykkvi &  &  &  \\
\hline
4945 & snúningagrúpa &  &  &  \\
\hline
4946 & snúningsás, möndull &  &  &  \\
\hline
4947 & snúningshorn &  &  &  \\
\hline
4948 & snúningsóbreyta &  &  &  \\
\hline
4949 & snúningsóbreyttur &  &  &  \\
\hline
4950 & snúningssamhverfni, snúsamhverfni &  &  &  \\
\hline
4951 & snúningssamhverfur, snúsamhverfur &  &  &  \\
\hline
4952 & snúningsskaut, snúningspóll, snúningsmiðja &  &  &  \\
\hline
4953 & snúningur &  &  &  \\
\hline
4954 & snúningur & Ef $\rho$ er <ref idTerm=5779>færsla</ref> á tiltekinni <ref idTerm=4884>sléttu</ref> og $M$ er <ref idTerm=4001>punktur</ref> í sléttunni þá kallast $\rho$ \textit{snúningur} með <ref idTerm=3304>miðju</ref> í punktinum $M$ ef hún hefur eftirfarandi eiginleika:

\begin{enumerate}
\item punktarnir $A$ og $\rho(A)$ eru í sömu <ref idTerm=1204>fjarlægð</ref> frá $M$ fyrir alla punkta $A$ í sléttunni,

\item <ref idTerm=2198>hornin</ref> $\angle AM\rho(A)$ eru jafn stór fyrir alla punkta $A\neq M$.
\end{enumerate} &  &  \\
\hline
4955 & snústríkkun &  &  &  \\
\hline
4956 & snyrta, jafna &  &  &  \\
\hline
4957 & snyrting, jöfnun &  &  &  \\
\hline
4958 & snyrtiskekkja, jöfnunarskekkja &  &  &  \\
\hline
4959 & söðulflötur &  &  &  \\
\hline
4960 & söðulpunktsaðferð &  &  &  \\
\hline
4961 & söðulpunktur &  &  &  \\
\hline
4962 & sólsproti &  &  &  \\
\hline
4963 & sönn fullyrðing, sönn yrðing &  &  &  \\
\hline
4964 & sönnun &  &  &  \\
\hline
4965 & sóphorn &  &  &  \\
\hline
4966 & spá, forsögn &  &  &  \\
\hline
4967 & spáfræði &  &  &  \\
\hline
4968 & spágildi, reiknað gildi &  &  &  \\
\hline
4969 & spanna &  &  &  \\
\hline
4970 & spanna &  &  &  \\
\hline
4971 & spanna, framleiða &  &  &  \\
\hline
4972 & spannandi hlutnet &  &  &  \\
\hline
4973 & spannandi skógur, spannskógur &  &  &  \\
\hline
4974 & spannandi tré, spanntré, hátré &  &  &  \\
\hline
4975 & spannandi vensl &  &  &  \\
\hline
4976 & spannarmiðja &  &  &  \\
\hline
4977 & spannfjölskylda, spannandi fjölskylda &  &  &  \\
\hline
4979 & spannmengi (íðals), íðalspannmengi, framleiðandi mengi, spannandi mengi &  &  &  \\
\hline
4980 & spáskekkja &  &  &  \\
\hline
4981 & spegilmynd &  &  &  \\
\hline
4982 & spegilsamhvarf, spegilsamhverfni, sléttusamhverfni, tvíhvarf, tvíhverfni, tvíhliða samhvarf, tvíhliða samhverfni &  &  &  \\
\hline
4983 & spegilsamhverfur &  &  &  \\
\hline
4984 & spegilsamsniða, óeiginlega samsniða &  &  &  \\
\hline
4985 & spegiltala &  &  &  \\
\hline
4986 & speglun &  &  &  \\
\hline
4987 & speglunarsetning, speglunarlögmál &  &  &  \\
\hline
4988 & spenna &  &  &  \\
\hline
4989 & spennuþinur &  &  &  \\
\hline
4990 & spil, leikur, keppni &  &  &  \\
\hline
4991 & spilfræði, spilafræði, leikjafræði, keppnisfræði &  &  &  \\
\hline
4992 & spinalgebra &  &  &  \\
\hline
4993 & spingreining, spinagreining &  &  &  \\
\hline
4994 & spingrúpa &  &  &  \\
\hline
4995 & spinsvið, spinasvið &  &  &  \\
\hline
4996 & spinur, spunaþinur &  &  &  \\
\hline
4997 & splæsifall &  &  &  \\
\hline
4998 & spönn &  &  &  \\
\hline
4999 & spönnuður &  &  &  \\
\hline
5000 & spor &  &  &  \\
\hline
5001 & sporbaugsfari &  &  &  \\
\hline
5002 & sporbaugur &  &  &  \\
\hline
5003 & sporger &  &  &  \\
\hline
5004 & sporger ferill &  &  &  \\
\hline
5005 & sporger fleygbogaflötur &  &  &  \\
\hline
5006 & sporger hlutafleiðujafna &  &  &  \\
\hline
5007 & sporger hreyfing, hreyfing á sporbaug &  &  &  \\
\hline
5008 & sporger rúmfræði &  &  &  \\
\hline
5009 & sporger sjálfmótun, sporger færsla &  &  &  \\
\hline
5010 & sporger virki &  &  &  \\
\hline
5011 & sporgerð &  &  &  \\
\hline
5012 & sporgerð &  &  &  \\
\hline
5013 & sporgert fall &  &  &  \\
\hline
5014 & sporgert fall Jacobis &  &  &  \\
\hline
5015 & sporgert heildi, sporgert tegur &  &  &  \\
\hline
5016 & sporlína &  &  &  \\
\hline
5017 & sporvala &  &  &  \\
\hline
5018 & sporvöluflötur, sporbaugsflötur &  &  &  \\
\hline
5019 & sporvöluhnit &  &  &  \\
\hline
5020 & sporvölusnúðflötur &  &  &  \\
\hline
5021 & sporvölusnúður &  &  &  \\
\hline
5022 & spunabundin &  &  &  \\
\hline
5023 & spunaferill &  &  &  \\
\hline
5024 & spunafylki &  &  &  \\
\hline
5025 & spunagrúpa &  &  &  \\
\hline
5026 & spunaútsetning &  &  &  \\
\hline
5027 & spuni &  &  &  \\
\hline
5028 & spyrðing &  &  &  \\
\hline
5029 & spyrðing af hámarksstærð &  &  &  \\
\hline
5030 & spyrðingarsetning &  &  &  \\
\hline
5031 & spyrður &  &  &  \\
\hline
5032 & spyrtur hnútur &  &  &  \\
\hline
5033 & stað- &  &  &  \\
\hline
5034 & staða &  &  &  \\
\hline
5035 & staða, stelling &  &  &  \\
\hline
5036 & staðal- &  &  &  \\
\hline
5037 & staðal- &  &  &  \\
\hline
5038 & staðalfrávik &  &  &  \\
\hline
5039 & staðalfrávik úrtaks &  &  &  \\
\hline
5040 & staðalgebra &  &  &  \\
\hline
5041 & staðalgerð, normleg gerð, staðalframsetning, normleg framsetning &  &  &  \\
\hline
5042 & staðalgrannmynstur &  &  &  \\
\hline
5043 & staðalgrannmynstur &  &  &  \\
\hline
5044 & staðaljafna &  &  &  \\
\hline
5045 & staðall, lengd &  &  &  \\
\hline
5046 & staðalrækinn &  &  &  \\
\hline
5047 & staðalsamleitinn &  &  &  \\
\hline
5048 & staðalsamleitni &  &  &  \\
\hline
5049 & staðalsamleitni &  &  &  \\
\hline
5050 & staðalskekkja &  &  &  \\
\hline
5051 & staðalskekkja meðaltals &  &  &  \\
\hline
5052 & staðalþverlægni &  &  &  \\
\hline
5053 & staðarfall &  &  &  \\
\hline
5054 & staðarvigur, stöðuvigur &  &  &  \\
\hline
5055 & staðaugljós &  &  &  \\
\hline
5056 & staðbaugað rúm &  &  &  \\
\hline
5057 & staðbaugur &  &  &  \\
\hline
5058 & staðbinda &  &  &  \\
\hline
5059 & staðbinding &  &  &  \\
\hline
5060 & staðbundið hágildi, hágildi &  &  &  \\
\hline
5061 & staðbundið lággildi, lággildi &  &  &  \\
\hline
5062 & staðbundið útgildi &  &  &  \\
\hline
5063 & staðbundin Lie-grúpa, Lie-grúpukím &  &  &  \\
\hline
5064 & staðbundinn eiginleiki &  &  &  \\
\hline
5065 & staðbundinn, stað- &  &  &  \\
\hline
5066 & staðendanlega spannaður &  &  &  \\
\hline
5067 & staðendanlegur &  &  &  \\
\hline
5068 & staðfastur &  &  &  \\
\hline
5069 & staðfrjáls &  &  &  \\
\hline
5070 & staðgengill &  &  &  \\
\hline
5071 & staðgrannmóta &  &  &  \\
\hline
5072 & staðhæfing &  &  &  \\
\hline
5073 & staðhnit, staðhnitafall &  &  &  \\
\hline
5074 & staðhnitakerfi, staðhnit, kort &  &  &  \\
\hline
5075 & staðkúptur &  &  &  \\
\hline
5076 & staðla &  &  &  \\
\hline
5077 & staðlað línulegt rúm, staðlað vigurrúm, staðlað vektorrúm &  &  &  \\
\hline
5078 & staðlaður &  &  &  \\
\hline
5079 & staðlaður &  &  &  \\
\hline
5080 & staðlaður &  &  &  \\
\hline
5081 & staðlaður grunnur &  &  &  \\
\hline
5082 & staðlanlegur &  &  &  \\
\hline
5083 & staðlanlegur &  &  &  \\
\hline
5084 & staðlítill, staðsmár &  &  &  \\
\hline
5085 & staðlokaður &  &  &  \\
\hline
5086 & staðsamanhangandi, staðsamhangandi &  &  &  \\
\hline
5087 & staðstiki, staðbundinn stiki &  &  &  \\
\hline
5088 & staðstikun, staðbundin einstikun &  &  &  \\
\hline
5089 & staðsvið, staðkroppur &  &  &  \\
\hline
5090 & staðtakmarkaður, staðskorðaður &  &  &  \\
\hline
5091 & staðtala, lýsitala &  &  &  \\
\hline
5092 & staður &  &  &  \\
\hline
5093 & staðvegsamanhangandi, staðvegsamhangandi &  &  &  \\
\hline
5094 & staðþjappaður &  &  &  \\
\hline
5095 & staðþjappleiki, staðþjöppun &  &  &  \\
\hline
5096 & stæða &  &  &  \\
\hline
5097 & stæða, stærðtákn &  &  &  \\
\hline
5098 & stæðafjöldi &  &  &  \\
\hline
5099 & stæði &  &  &  \\
\hline
5100 & stækkun, viðbót, auki &  &  &  \\
\hline
5102 & stærð &  &  &  \\
\hline
5103 & stærð, fjöldatala &  &  &  \\
\hline
5104 & stærð, stórleiki &  &  &  \\
\hline
5105 & stærð, þrep &  &  &  \\
\hline
5106 & stærðarþrep, tugaþrep &  &  &  \\
\hline
5107 & stærðfræði &  &  &  \\
\hline
5113 & stærðfræðilegt líkan &  &  &  \\
\hline
5114 & stærðfræðilegur hlutur &  &  &  \\
\hline
5116 & stærðfræðimynstur, stærðfræðilegt mynstur &  &  &  \\
\hline
5117 & stærðfræðingur &  &  &  \\
\hline
5118 & stærri &  &  &  \\
\hline
5119 & stærri eða jafn, stærri eða jafn og, frekar stærri &  &  &  \\
\hline
5120 & stærsta stak, efsta stak, hæsta stak, síðasta stak &  &  &  \\
\hline
5121 & stærsti samdeilir, stærsti sameiginlegi deilir &  &  &  \\
\hline
5122 & stafræn tölva &  &  &  \\
\hline
5123 & stafrænn &  &  &  \\
\hline
5124 & stafrófsröð, stafrófsröðun &  &  &  \\
\hline
5125 & stafur, bókstafur, tákn &  &  &  \\
\hline
5126 & stak fylkis, stuðull fylkis &  &  &  \\
\hline
5127 & stak utan hornalínu &  &  &  \\
\hline
5128 & stak, íbúi &  &  &  \\
\hline
5129 & stak, stuðull &  &  &  \\
\hline
5130 & stakbreyta &  &  &  \\
\hline
5131 & stallagerð, þrepagerð, þrepaskipan &  &  &  \\
\hline
5132 & standandi bylgja &  &  &  \\
\hline
5133 & stefna &  &  &  \\
\hline
5134 & stefna á &  &  &  \\
\hline
5135 & stefna, átt &  &  &  \\
\hline
5136 & stefnd lína, áttuð lína &  &  &  \\
\hline
5137 & stefndur til hægri &  &  &  \\
\hline
5138 & stefndur til vinstri &  &  &  \\
\hline
5139 & stefnt línustrik &  &  &  \\
\hline
5140 & stefnt mengi &  &  &  \\
\hline
5141 & stefnu- &  &  &  \\
\hline
5142 & stefnuafleiða &  &  &  \\
\hline
5143 & stefnubundin fjarlægð &  &  &  \\
\hline
5144 & stefnudeildun, stefnudiffrun &  &  &  \\
\hline
5145 & stefnuhorn &  &  &  \\
\hline
5146 & stefnuskrá Hilberts &  &  &  \\
\hline
5147 & stefnustuðull, stefnukósínus &  &  &  \\
\hline
5148 & stefnusvið, einingarvigursvið &  &  &  \\
\hline
5149 & stefnuvigur &  &  &  \\
\hline
5150 & sterk fylgni &  &  &  \\
\hline
5151 & sterk lausn &  &  &  \\
\hline
5152 & sterk samleitni &  &  &  \\
\hline
5153 & sterkari kenning &  &  &  \\
\hline
5154 & sterklega &  &  &  \\
\hline
5155 & sterklega samleitinn &  &  &  \\
\hline
5156 & sterkt grannmynstur &  &  &  \\
\hline
5157 & sterkur &  &  &  \\
\hline
5158 & stétt, myndvídd, metorð, tign &  &  &  \\
\hline
5159 & stétt, tign &  &  &  \\
\hline
5160 & stig &  &  &  \\
\hline
5161 & stig & Sjá <ref idTerm=1948>margliða</ref> &  & Setja i skýringu: Stig núllmargliðunnar er ekki unnt að skilgreina með skynsamlegum hætti sem heila tölu og er því stig hennar sagt vera óskilgreint. \\
\hline
5162 & stig &  &  &  \\
\hline
5164 & stíga, vaxa, hækka &  &  &  \\
\hline
5165 & stigaður &  &  &  \\
\hline
5166 & stígandi keðja, vaxandi keðja &  &  &  \\
\hline
5167 & stígandi, vaxandi &  &  &  \\
\hline
5168 & stígandi, vaxandi, hækkandi &  &  &  \\
\hline
5169 & stiglækkun &  &  &  \\
\hline
5170 & stigskipt úrtaka &  &  &  \\
\hline
5171 & stigulaðferð, aðferð mesta bratta &  &  &  \\
\hline
5172 & stigull &  &  &  \\
\hline
5173 & stigun &  &  &  \\
\hline
5174 & stika, hlaupa &  &  &  \\
\hline
5175 & stikabundið próf &  &  &  \\
\hline
5176 & stikabundin tilgáta &  &  &  \\
\hline
5177 & stikaður ferill, hlaupinn ferill &  &  &  \\
\hline
5178 & stikaður, hlaupinn &  &  &  \\
\hline
5179 & stikaður, stika- &  &  &  \\
\hline
5180 & stikaframsetning, hlaupaframsetning &  &  &  \\
\hline
5181 & stikajafna, hlaupajafna &  &  &  \\
\hline
5182 & stikalaus &  &  &  \\
\hline
5183 & stikalaus &  &  &  \\
\hline
5184 & stikalaust próf &  &  &  \\
\hline
5185 & stikarúm, hlauparúm &  &  &  \\
\hline
5186 & stiki, hlaupi, hlaupabreyta, aukabreyta, hlaupastærð, færibreyta &  &  &  \\
\hline
5187 & stikun &  &  &  \\
\hline
5188 & stilkur &  &  &  \\
\hline
5189 & Stirling-tala &  &  &  \\
\hline
5190 & stjarfleiki &  &  &  \\
\hline
5191 & stjarfur &  &  &  \\
\hline
5192 & stjarfur hlutur &  &  &  \\
\hline
5193 & stjarna &  &  &  \\
\hline
5194 & stjarna &  &  &  \\
\hline
5195 & stjökun &  &  &  \\
\hline
5196 & stjökunarvirki &  &  &  \\
\hline
5197 & stjörf hreyfing, stjarfur flutningur &  &  &  \\
\hline
5198 & stjórn, stýring &  &  &  \\
\hline
5199 & stjórna, stýra &  &  &  \\
\hline
5200 & stjórnbreyta, stýribreyta, ákvörðunarbreyta &  &  &  \\
\hline
5201 & stjörnualgebra &  &  &  \\
\hline
5202 & stjörnuferill &  &  &  \\
\hline
5203 & stjörnumargflötungur, stirndur margflötungur &  &  &  \\
\hline
5204 & stjörnusvæði, stjörnulaga svæði &  &  &  \\
\hline
5205 & stjörnuþykkvi, stjörnulaga þykkvi, stjörnulaga rúmmynd &  &  &  \\
\hline
5206 & stoð &  &  &  \\
\hline
5207 & stoðafjölskylda &  &  &  \\
\hline
5208 & stoðfall &  &  &  \\
\hline
5209 & stoðhending, hjálparhending &  &  &  \\
\hline
5210 & stoðlína &  &  &  \\
\hline
5211 & stöðluð algebra &  &  &  \\
\hline
5212 & stöðluð dreifing &  &  &  \\
\hline
5213 & stöðluð Gauss-dreifing, stöðluð normaldreifing, stöðluð normleg dreifing &  &  &  \\
\hline
5214 & stöðluð hending &  &  &  \\
\hline
5215 & stöðluð margliða & <ref idTerm=1948>Margliða</ref> þar sem gildi forystustuðulsins er 1 &  &  \\
\hline
5216 & stöðlun &  &  &  \\
\hline
5217 & stoðslétta &  &  &  \\
\hline
5218 & stöðug lausn &  &  &  \\
\hline
5219 & stöðugleikagrúpa &  &  &  \\
\hline
5220 & stöðugleiki &  &  &  \\
\hline
5221 & stöðugt ástand &  &  &  \\
\hline
5222 & stöðugt ferli &  &  &  \\
\hline
5223 & stöðugt hlutmengi &  &  &  \\
\hline
5224 & stöðugur &  &  &  \\
\hline
5225 & stöðugur, stöðu- &  &  &  \\
\hline
5226 & stöðurúm, svigrúm, stellingarúm &  &  &  \\
\hline
5227 & stofnbrot & <ref idTerm=4039>Ræð tala</ref> skrifuð sem <ref idTerm=234>almennt brot</ref> þar sem <ref idTerm=5462>teljarinn</ref> er 1 og <ref idTerm=3497>nefnarinn</ref> er <ref idTerm=2535>jákvæð</ref> <ref idTerm=1896>heiltala</ref> &  &  \\
\hline
5228 & stofnbrot, hlutbrot &  &  &  \\
\hline
5229 & stofnbrotaliðun, hlutbrotaliðun, liðun í stofnbrot, liðun í hlutbrot &  &  &  \\
\hline
5230 & stofnfall &  &  &  \\
\hline
5231 & stofnfrumsenda, stofnfrumsetning, frumsenda um stofn, frumsetning um stofn, regluleikafrumsenda, regluleikafrumsetning &  &  &  \\
\hline
5232 & stökk, misgengi &  &  &  \\
\hline
5233 & stökkfall &  &  &  \\
\hline
5234 & stöng hlutvera &  &  &  \\
\hline
5235 & stöplarit &  &  &  \\
\hline
5236 & stórgeisli, umgeisli, geisli umritaðs hrings, geisli umritaðrar kúlu &  &  &  \\
\hline
5237 & stórhringur &  &  &  \\
\hline
5238 & stórt ríki &  &  &  \\
\hline
5239 & stórt ríki &  &  &  \\
\hline
5240 & stranglega &  &  &  \\
\hline
5241 & stranglega einhalla &  &  &  \\
\hline
5242 & stranglega jákvæður, stranglega viðlægur &  &  &  \\
\hline
5243 & stranglega margfeldinn &  &  &  \\
\hline
5244 & stranglega neikvæður, stranglega frádrægur &  &  &  \\
\hline
5245 & strangleiki &  &  &  \\
\hline
5246 & strangminnkandi, stranglega minnkandi, strangfallandi, stranglega fallandi, stranglækkandi, stranglega lækkandi, síminnkandi, sífallandi, sílækkandi &  &  &  \\
\hline
5247 & strangt þríhyrningsfylki &  &  &  \\
\hline
5248 & strangur &  &  &  \\
\hline
5249 & strangur, rökfastur &  &  &  \\
\hline
5250 & strangvaxandi, stranglega vaxandi, strangstígandi, stranglega stígandi, stranghækkandi, stranglega hækkandi, sívaxandi, sístígandi, síhækkandi &  &  &  \\
\hline
5251 & straumfall &  &  &  \\
\hline
5252 & straumlína &  &  &  \\
\hline
5253 & straumur &  &  &  \\
\hline
5254 & strendingsbróðir &  &  &  \\
\hline
5255 & strendingsflötur &  &  &  \\
\hline
5256 & strendingsfrændi &  &  &  \\
\hline
5257 & strendingsstúfur, stýfður strendingur &  &  &  \\
\hline
5258 & strendingur &  &  &  \\
\hline
5259 & streng- &  &  &  \\
\hline
5260 & strengfirð &  &  &  \\
\hline
5261 & strengur & <ref idTerm=2958>Línustrik</ref> sem hefur báða <ref idTerm=0>endapunkta</ref> sína á <ref idTerm=2295>hring</ref> kallast \textit{strengur} í hringnum &  &  \\
\hline
5262 & strengur &  &  &  \\
\hline
5263 & streymi &  &  &  \\
\hline
5264 & strik, merki &  &  &  \\
\hline
5265 & stríkkun & Ef $\lambda$ er <ref idTerm=5779>færsla</ref> á tiltekinni <ref idTerm=4884>sléttu</ref>, $k$ er <ref idTerm=4093>rauntala</ref> og $M$ er <ref idTerm=4001>punktur</ref> í sléttunni þá kallast $\lambda$ \textit{stríkkun} með <ref idTerm=3304>miðju</ref> í punktinum $M$ ef fyrir alla punkta $P$ í sléttunni gildi að <ref idTerm=2072>hlutfallið</ref> milli <ref idTerm=2860>lengd</ref> <ref idTerm=2958>strikanna</ref> $MP$ og $M\lambda(P)$ sé $k$, sem kallast þá \textit{kvarði} stríkkunarinnar &  &  \\
\hline
5266 & stríkkunarhlutfall &  &  &  \\
\hline
5267 & stríkkunarstuðull &  &  &  \\
\hline
5268 & strjál bestun &  &  &  \\
\hline
5269 & strjál breyta &  &  &  \\
\hline
5270 & strjál dreifing &  &  &  \\
\hline
5271 & strjál hending &  &  &  \\
\hline
5272 & strjál stærðfræði &  &  &  \\
\hline
5273 & strjála &  &  &  \\
\hline
5274 & strjáll tími &  &  &  \\
\hline
5275 & strjált dreifingarfall &  &  &  \\
\hline
5276 & strjált ferli &  &  &  \\
\hline
5277 & strjált róf &  &  &  \\
\hline
5278 & strjált streymi &  &  &  \\
\hline
5279 & strjálun &  &  &  \\
\hline
5280 & strjálunarskekkja &  &  &  \\
\hline
5281 & ströng leiðing &  &  &  \\
\hline
5282 & ströng ójafna &  &  &  \\
\hline
5283 & ströng röðun &  &  &  \\
\hline
5284 & ströng sönnun &  &  &  \\
\hline
5285 & strútur &  &  &  \\
\hline
5286 & strýta, píramíti &  &  &  \\
\hline
5287 & strýtustúfur, píramítastúfur, stýfð strýta, stýfður píramíti &  &  &  \\
\hline
5288 & stuðlafylki &  &  &  \\
\hline
5289 & stuðlasamanburður, stuðlaákvörðun &  &  &  \\
\hline
5290 & stuðull &  &  &  \\
\hline
5291 & stúfabundin &  &  &  \\
\hline
5292 & stúfur &  &  &  \\
\hline
5293 & stúfur &  &  &  \\
\hline
5294 & stuttur, skammur &  &  &  \\
\hline
5295 & stýfa &  &  &  \\
\hline
5296 & stýfð röð &  &  &  \\
\hline
5297 & stýfður &  &  &  \\
\hline
5298 & stýfing &  &  &  \\
\hline
5299 & stýfingarskekkja &  &  &  \\
\hline
5300 & stýrð breyta &  &  &  \\
\hline
5301 & stýriferill &  &  &  \\
\hline
5302 & stýrifræði &  &  &  \\
\hline
5303 & stýrifræði &  &  &  \\
\hline
5304 & stýrilína &  &  &  \\
\hline
5305 & styrkur, styrkleiki &  &  &  \\
\hline
5306 & stytt margföldun, stuttamargföldun &  &  &  \\
\hline
5307 & stytta &  &  &  \\
\hline
5308 & stytta, stytta út, stytta burt &  &  &  \\
\hline
5309 & styttanlegt stak &  &  &  \\
\hline
5310 & styttanlegur &  &  &  \\
\hline
5311 & styttanleiki &  &  &  \\
\hline
5312 & styttin hálfgrúpa &  &  &  \\
\hline
5313 & stytting &  &  &  \\
\hline
5314 & stytting &  &  &  \\
\hline
5315 & styttingarskekkja &  &  &  \\
\hline
5316 & styttinn &  &  &  \\
\hline
5317 & styttiregla &  &  &  \\
\hline
5318 & suðurskaut, suðurpóll &  &  &  \\
\hline
5319 & súlurit, stuðlarit &  &  &  \\
\hline
5320 & summa, samtala &  &  &  \\
\hline
5321 & summumyndun &  &  &  \\
\hline
5322 & summumyndunaraðferð, summuaðferð &  &  &  \\
\hline
5323 & summutákn &  &  &  \\
\hline
5324 & sundurgreining, greining, krufning &  &  &  \\
\hline
5325 & sundurlæg skipting &  &  &  \\
\hline
5326 & sundurlægir atburðir &  &  &  \\
\hline
5327 & sundurlægir tveir og tveir &  &  &  \\
\hline
5328 & sundurlægt sammengi &  &  &  \\
\hline
5329 & sundurlægur, óskaraður &  &  &  \\
\hline
5330 & sundurlaus, strjáll &  &  &  \\
\hline
5331 & sundurlaust grannmynstur, strjált grannmynstur &  &  &  \\
\hline
5332 & sundurlaust rúm, strjált rúm &  &  &  \\
\hline
5333 & sundurleitni &  &  &  \\
\hline
5334 & sundurliðun &  &  &  \\
\hline
5335 & svæði &  &  &  \\
\hline
5336 & svæðisheildi, svæðistegur &  &  &  \\
\hline
5337 & svæðismat &  &  &  \\
\hline
5338 & svar, svörun &  &  &  \\
\hline
5339 & svarferill &  &  &  \\
\hline
5340 & svarflötur &  &  &  \\
\hline
5341 & sveifarferill &  &  &  \\
\hline
5342 & sveifla &  &  &  \\
\hline
5343 & sveifla &  &  &  \\
\hline
5344 & sveigð hnit &  &  &  \\
\hline
5345 & sveigfeldi &  &  &  \\
\hline
5346 & sveimferli &  &  &  \\
\hline
5347 & sveimstuðull &  &  &  \\
\hline
5348 & sveipferill &  &  &  \\
\hline
5349 & svið &  &  &  \\
\hline
5350 & svið ræðra talna &  &  &  \\
\hline
5351 & svið, kroppur &  &  &  \\
\hline
5352 & sviðafræði, kroppafræði &  &  &  \\
\hline
5353 & sviðsfræði &  &  &  \\
\hline
5354 & sviðsútvíkkun, sviðsvíkkun &  &  &  \\
\hline
5355 & svigalaus ritháttur &  &  &  \\
\hline
5356 & svigi &  &  &  \\
\hline
5357 & svigrúmsvídd &  &  &  \\
\hline
5358 & svipalgebra &  &  &  \\
\hline
5359 & svipflokkur &  &  &  \\
\hline
5360 & svipfræði &  &  &  \\
\hline
5361 & svipgrúpa &  &  &  \\
\hline
5362 & sviphjávídd, hjávídd, dýpt &  &  &  \\
\hline
5363 & svipjafngildi &  &  &  \\
\hline
5364 & svipjafngildur &  &  &  \\
\hline
5365 & svipkenning &  &  &  \\
\hline
5366 & sviplest &  &  &  \\
\hline
5367 & svipvídd &  &  &  \\
\hline
5368 & syðra hálfhvel, suðurhvel &  &  &  \\
\hline
5369 & Sylow-hlutgrúpa, Sylow-grúpa &  &  &  \\
\hline
5370 & sýndaróendanleiki &  &  &  \\
\hline
5371 & tæming &  &  &  \\
\hline
5372 & tæmingaraðferð &  &  &  \\
\hline
5373 & tafla &  &  &  \\
\hline
5374 & taka &  &  &  \\
\hline
5375 & taka (gildi) &  &  &  \\
\hline
5376 & taka enda &  &  &  \\
\hline
5377 & taka saman &  &  &  \\
\hline
5378 & taka til láns &  &  &  \\
\hline
5379 & taka út fyrir sviga &  &  &  \\
\hline
5380 & takmarkað fall, skorðað fall & <ref idTerm=1078>Fall</ref> $f: X \to Y$ kallast \textit{takmarkað} ef til er <ref idTerm=4093>rauntala</ref> $C$ þannig að $|f(x)| \leq C$ fyrir öll $x \in X$ &  & Hægt að endurnýta skilgreininguna um takmörkuð mengi með því að segja að fall sé takmarkað ef myndmengi þess er takmarkað mengi \\
\hline
5381 & takmarkað mengi, skorðað mengi & <ref idTerm=4093>Rauntalna</ref><ref idTerm=3270>mengi</ref> $A$ kallast \textit{takmarkað} ef til er rauntala $C$ þannig að $|x| \leq C$ fyrir öll $x$ úr $A$ &  & Minnast á tvinntalnamengi í skýringu \\
\hline
5382 & takmarkað svæði, skorðað svæði &  &  &  \\
\hline
5383 & takmarkað vik &  &  &  \\
\hline
5384 & takmarkað vik, skorðað vik &  &  &  \\
\hline
5385 & takmarkaður að mestu, næstum takmarkaður, skorðaður að mestu, næstum skorðaður &  &  &  \\
\hline
5386 & takmarkaður að neðan, skorðaður að neðan & <ref idTerm=4093>rauntalna</ref><ref idTerm=3270>mengi</ref> kallast \textit{takmarkað að neðan} ef það hefur <ref idTerm=5818>undirtölu</ref> &  &  \\
\hline
5387 & takmarkaður að ofan, skorðaður að ofan & <ref idTerm=4093>rauntalna</ref><ref idTerm=3270>mengi</ref> kallast \textit{takmarkað að ofan} ef það hefur <ref idTerm=6280>yfirtölu</ref> &  &  \\
\hline
5388 & takmarkaður í hverjum punkti, skorðaður í hverjum punkti &  &  &  \\
\hline
5389 & takmarkaður í jöfnum mæli, skorðaður í jöfnum mæli &  &  &  \\
\hline
5390 & takmarkaður, skorðaður &  &  &  \\
\hline
5391 & takmörkuð afstæðiskenning &  &  &  \\
\hline
5392 & takmörkuð keðja, skorðuð keðja &  &  &  \\
\hline
5393 & takmörkun í jöfnum mæli, skorðun í jöfnum mæli, jafnmælisskorðun &  &  &  \\
\hline
5395 & táknmál &  &  &  \\
\hline
5396 & táknrænn &  &  &  \\
\hline
5397 & tala &  &  &  \\
\hline
5398 & tala &  &  &  \\
\hline
5399 & tala með formerki &  &  &  \\
\hline
5400 & tálmafall &  &  &  \\
\hline
5401 & tálmi &  &  &  \\
\hline
5402 & talnabil &  &  &  \\
\hline
5403 & talnabrot &  &  &  \\
\hline
5404 & talnaeining &  &  &  \\
\hline
5405 & talnafasti &  &  &  \\
\hline
5406 & talnafræði margföldunar &  &  &  \\
\hline
5407 & talnafræði samlagningar &  &  &  \\
\hline
5408 & talnafræði, talnalist &  &  &  \\
\hline
5409 & talnafræðifall, talnafræðilegt fall &  &  &  \\
\hline
5410 & talnafræðilegur &  &  &  \\
\hline
5411 & talnafræðingur &  &  &  \\
\hline
5412 & talnagögn &  &  &  \\
\hline
5413 & talnagrind, kúlnagrind, reiknigrind &  &  &  \\
\hline
5414 & talnahvel Riemanns &  &  &  \\
\hline
5415 & talnakerfi &  &  &  \\
\hline
5416 & talnalína, talnaás, rauntalnalína, rauntalnaás, raunás & <ref idTerm=2911>Lína</ref> með tvo aðgreinda <ref idTerm=4001>punkta</ref> $O$ (\textit{núllpunktur}) og $E$ (\textit{einingarpunktur}) sem svara til <ref idTerm=3475>náttúrulegu talnanna</ref> 0 og 1. <ref idTerm=1204>Fjarlægðin</ref> milli þessa punkta er notuð sem <ref idTerm=2860>lengdareining</ref> og <ref idTerm=1896>heiltölunum</ref> er raðað á línuna þannig að <ref idTerm=4458>samliggjandi</ref> tölur eru í þessari fjarlægð frá hvor annarri &  & Þarf að minnast á hvernig ræðar og óræðar tölur dreifast inn á línuna líka? \\
\hline
5417 & talnalykill &  &  &  \\
\hline
5418 & talnareikningur &  &  &  \\
\hline
5419 & talnaröð &  &  &  \\
\hline
5420 & talnarúm &  &  &  \\
\hline
5421 & talnaruna &  &  &  \\
\hline
5422 & talnaslétta &  &  &  \\
\hline
5423 & talnastærð &  &  &  \\
\hline
5424 & talnastuðull &  &  &  \\
\hline
5425 & talnatákn, tölutákn &  &  &  \\
\hline
5426 & talnatvennd &  &  &  \\
\hline
5427 & talning, upptalning &  &  &  \\
\hline
5428 & talningarferli &  &  &  \\
\hline
5429 & talningarfræði &  &  &  \\
\hline
5430 & talningarmál &  &  &  \\
\hline
5431 & tangens &  &  &  \\
\hline
5432 & tangensregla &  &  &  \\
\hline
5433 & tapfall &  &  &  \\
\hline
5434 & Tauber-setning &  &  &  \\
\hline
5435 & Taylor-liðun &  &  &  \\
\hline
5436 & Taylor-röð &  &  &  \\
\hline
5437 & tegund &  &  &  \\
\hline
5438 & tegund &  &  &  \\
\hline
5439 & teikna &  &  &  \\
\hline
5440 & teikna, draga &  &  &  \\
\hline
5441 & teikna, draga &  &  &  \\
\hline
5442 & teiknifræði, lýsandi rúmfræði &  &  &  \\
\hline
5443 & teikniheildun, teiknitegrun &  &  &  \\
\hline
5444 & teiknihorn &  &  &  \\
\hline
5445 & teiknilausn &  &  &  \\
\hline
5446 & teikning &  &  &  \\
\hline
5447 & teikning með hringfara &  &  &  \\
\hline
5448 & teikning með hringfara og reglustiku &  &  &  \\
\hline
5449 & teikning, dráttur, uppdráttur, rúmfræðileg teikning &  &  &  \\
\hline
5450 & teikning, uppdráttur, dráttur &  &  &  \\
\hline
5451 & teikniverkefni, dráttardæmi &  &  &  \\
\hline
5452 & telja &  &  &  \\
\hline
5453 & teljanlega óendanlegur &  &  &  \\
\hline
5454 & teljanlega viðbætinn, sigma-viðbætinn, teljanlega samleggjandi, sigma-samleggjandi &  &  &  \\
\hline
5455 & teljanlega þjappaður, sigma-þjappaður &  &  &  \\
\hline
5456 & teljanlegt mengi &  &  &  \\
\hline
5457 & teljanlegur &  &  &  \\
\hline
5458 & teljanlegur óendanleiki &  &  &  \\
\hline
5459 & teljanleikafrumsenda, teljanleikafrumsetning &  &  &  \\
\hline
5460 & teljanleikaskilyrði, teljanleikaforsenda, teljanleikafrumsetning, teljanleikafrumsenda &  &  &  \\
\hline
5461 & teljanleiki &  &  &  \\
\hline
5462 & teljari & Sjá <ref idTerm=234>almennt brot</ref> &  &  \\
\hline
5463 & tempruð dreif, temprað dreififelli &  &  &  \\
\hline
5464 & tengd ástönd &  &  &  \\
\hline
5465 & tengi- &  &  &  \\
\hline
5466 & tengilína &  &  &  \\
\hline
5467 & tengill &  &  &  \\
\hline
5468 & tengimótun &  &  &  \\
\hline
5469 & tengin aðgerð, tengin reikniaðgerð &  &  &  \\
\hline
5470 & tengin algebra &  &  &  \\
\hline
5471 & tenging &  &  &  \\
\hline
5472 & tenginn &  &  &  \\
\hline
5473 & tengiregla & <ref idTerm=5702>Tvístæð reikniaðgerð</ref> $\odot : X \times X \to X$ á <ref idTerm=3270>mengi</ref> $X$ kallast \textit{tengin} eða að hún fullnægi \textit{tengireglunni} ef fyrir sérhver <ref idTerm=5128>stök</ref> $x,y,z \in X$ gildi að $(x \odot y) \odot z = x \odot (y \odot z)$ &  &  \\
\hline
5474 & tengivegur &  &  &  \\
\hline
5475 & tengni &  &  &  \\
\hline
5476 & tengslatafla, slembitafla &  &  &  \\
\hline
5477 & tengt frumíðal &  &  &  \\
\hline
5478 & teningsgrúpa, áttflötungsgrúpa &  &  &  \\
\hline
5479 & teningstala, verpiltala &  &  &  \\
\hline
5480 & teningur, verpill, reglulegur sexflötungur & <ref idTerm=2559>Kassi</ref> þar sem allar <ref idTerm=1325>hliðar</ref> eru <ref idTerm=1163>ferningar</ref> &  &  \\
\hline
5481 & tíðni &  &  &  \\
\hline
5482 & tíðnidreifing &  &  &  \\
\hline
5483 & tíðnifall &  &  &  \\
\hline
5484 & tíðnirit &  &  &  \\
\hline
5485 & tígulflötungur &  &  &  \\
\hline
5486 & tígull, tigull &  &  &  \\
\hline
5487 & tíhyrningur &  &  &  \\
\hline
5488 & tilfærsla &  &  &  \\
\hline
5489 & tilfærsla á stuðlum, stuðlatilfærsla &  &  &  \\
\hline
5490 & tilfelli, tilvik &  &  &  \\
\hline
5491 & tilgáta &  &  &  \\
\hline
5492 & tilgáta &  &  &  \\
\hline
5493 & tilgáta Riemanns &  &  &  \\
\hline
5494 & tilheyra, vera í, vera stak í &  &  &  \\
\hline
5495 & tilraun &  &  &  \\
\hline
5496 & tilraunafall &  &  &  \\
\hline
5497 & tilraunahögun, högun &  &  &  \\
\hline
5498 & tilraunaskekkja &  &  &  \\
\hline
5499 & tilsvarandi óhliðrað hneppi &  &  &  \\
\hline
5500 & tiltölulega lokað mengi &  &  &  \\
\hline
5501 & tiltölulega opið mengi &  &  &  \\
\hline
5502 & tiltölulega þjappað mengi &  &  &  \\
\hline
5503 & tilvikasönnun &  &  &  \\
\hline
5504 & tilviljun &  &  &  \\
\hline
5505 & tilviljunarákvörðun, tilviljanakennd ákvörðun &  &  &  \\
\hline
5506 & tilvist &  &  &  \\
\hline
5507 & tilvistar- og ótvíræðnisetning &  &  &  \\
\hline
5508 & tilvistarleysi &  &  &  \\
\hline
5509 & tilvistarlokun &  &  &  \\
\hline
5510 & tilvistarmagnari, tilvistarvirki &  &  &  \\
\hline
5511 & tilvistarsetning &  &  &  \\
\hline
5512 & tilvistarsönnun &  &  &  \\
\hline
5513 & tilvistartákn &  &  &  \\
\hline
5514 & tímabreyta &  &  &  \\
\hline
5515 & tímaóháður, sjálfstýrandi &  &  &  \\
\hline
5516 & tímaraðagreining &  &  &  \\
\hline
5517 & tímaröð, tímaruna &  &  &  \\
\hline
5518 & tímarúm &  &  &  \\
\hline
5519 & töfraferningur &  &  &  \\
\hline
5520 & togferill &  &  &  \\
\hline
5521 & togfræði &  &  &  \\
\hline
5522 & toggerð &  &  &  \\
\hline
5523 & toggrúpa &  &  &  \\
\hline
5524 & togjafngildi &  &  &  \\
\hline
5525 & togjafngildur &  &  &  \\
\hline
5526 & toglest &  &  &  \\
\hline
5527 & togunarinndragi &  &  &  \\
\hline
5528 & togunarinndráttur &  &  &  \\
\hline
5529 & tólfflata &  &  &  \\
\hline
5530 & tólfflötungsgrúpa, tvítugflötungsgrúpa &  &  &  \\
\hline
5531 & tólfflötungur &  &  &  \\
\hline
5532 & tólfhyrningur &  &  &  \\
\hline
5533 & tölfræði &  &  &  \\
\hline
5534 & tölfræðileg ályktun &  &  &  \\
\hline
5535 & tölfræðileg marktekt &  &  &  \\
\hline
5536 & tölfræðileg tilgáta &  &  &  \\
\hline
5537 & tölfræðilega marktækur &  &  &  \\
\hline
5538 & tölfræðilegt eftirlit &  &  &  \\
\hline
5539 & tölfræðilegt próf, marktektarpróf &  &  &  \\
\hline
5540 & tölfræðingur &  &  &  \\
\hline
5541 & tölugildur &  &  &  \\
\hline
5542 & tölukrappi, krappatala &  &  &  \\
\hline
5543 & töluleg aðferð &  &  &  \\
\hline
5544 & töluleg deildun, töluleg diffrun &  &  &  \\
\hline
5545 & töluleg gögn, staðtölur, lýsitölur &  &  &  \\
\hline
5546 & töluleg greining &  &  &  \\
\hline
5547 & töluleg heildun, töluleg tegrun &  &  &  \\
\hline
5548 & töluleg lausn, talnalausn &  &  &  \\
\hline
5549 & töluleg línuleg algebra &  &  &  \\
\hline
5550 & töluleg nálgunarfræði &  &  &  \\
\hline
5551 & tölulegur, talnalegur &  &  &  \\
\hline
5552 & tölumargfeldi (vigurs) &  &  &  \\
\hline
5553 & tölumargföldun &  &  &  \\
\hline
5554 & töluóbreyta &  &  &  \\
\hline
5556 & tölusett net, markað net, merkt net &  &  &  \\
\hline
5557 & tölusettur &  &  &  \\
\hline
5558 & tölustafur & <ref idTerm=5394>Tákn</ref> sem stendur fyrir ákveðna <ref idTerm=3475>náttúrulega tölu</ref> &  &  \\
\hline
5559 & tölustafur, fingur &  &  &  \\
\hline
5560 & tölustafur, fingur &  &  &  \\
\hline
5561 & tölusvið &  &  &  \\
\hline
5562 & tölva, reiknivél &  &  &  \\
\hline
5563 & tölvufræði, tölvunarfræði &  &  &  \\
\hline
5564 & tómt mengi, tómamengi & <ref idTerm=3270>Mengið</ref> með engin <ref idTerm=5128>stök</ref>, ýmist táknað með $\emptyset, \varnothing$ eða $\{\}$ &  &  \\
\hline
5565 & tómur &  &  &  \\
\hline
5566 & tómur flokkur &  &  &  \\
\hline
5567 & topphorn &  &  &  \\
\hline
5568 & topphorn, gagnstæð horn & <ref idTerm=3400>Gagnstæðar</ref> <ref idTerm=1840>hálflínur</ref> við <ref idTerm=334>arma</ref> <ref idTerm=2198>horns</ref> mynda <ref idTerm=3398>mótlægt horn</ref> upphaflega hornsins og sagt er að þessi tvö horn séu \textit{gagnstæð horn} &  &  \\
\hline
5569 & topppunktur &  &  &  \\
\hline
5570 & topppunktur, tindur, toppur, broddur &  &  &  \\
\hline
5571 & torræð tala &  &  &  \\
\hline
5572 & torræðni &  &  &  \\
\hline
5573 & torræðnigrunnur &  &  &  \\
\hline
5574 & torræðnistig &  &  &  \\
\hline
5575 & torræður &  &  &  \\
\hline
5576 & torrætt fall &  &  &  \\
\hline
5577 & trapisa, hálfsamsíðungur &  &  &  \\
\hline
5578 & tré &  &  &  \\
\hline
5579 & trefja & <ref idTerm=1366>Frummynd</ref> <ref idTerm=3282>vörpunnar</ref> $f: X \to Y$ af einstökungnum $\{y\}$ þar sem $y \in Y$ kallast \textit{trefja} vörpunarinnar $f$ af stakinu $y$ &  &  \\
\hline
5580 & trefjabundin &  &  &  \\
\hline
5581 & trefjamargfeldi &  &  &  \\
\hline
5582 & trefjarækin vörpun &  &  &  \\
\hline
5583 & trefjarúm &  &  &  \\
\hline
5584 & trefjun &  &  &  \\
\hline
5585 & tregða &  &  &  \\
\hline
5586 & tregðuregla, tregðuregla Sylvesters &  &  &  \\
\hline
5587 & tregðuvísir &  &  &  \\
\hline
5588 & troðsla &  &  &  \\
\hline
5589 & troðsluverkefni &  &  &  \\
\hline
5590 & troðsluþéttleiki &  &  &  \\
\hline
5591 & truflaður &  &  &  \\
\hline
5592 & truflanafræði &  &  &  \\
\hline
5593 & truflanaliður &  &  &  \\
\hline
5594 & truflanastiki &  &  &  \\
\hline
5595 & truflun &  &  &  \\
\hline
5596 & trygg túlkun &  &  &  \\
\hline
5597 & trygg útsetning &  &  &  \\
\hline
5598 & trygg verkun &  &  &  \\
\hline
5599 & tryggilega flatur &  &  &  \\
\hline
5600 & tryggingastærðfræði &  &  &  \\
\hline
5601 & tryggur &  &  &  \\
\hline
5602 & tryggur mótull &  &  &  \\
\hline
5603 & tryggur varpi &  &  &  \\
\hline
5604 & Tsjebysjev-margliða &  &  &  \\
\hline
5606 & tugabrot &  &  &  \\
\hline
5607 & tugabrot &  &  &  \\
\hline
5608 & tugabrotshluti &  &  &  \\
\hline
5609 & tugabrotsliðun, tugabrotsrakning, tugabrotsteygð &  &  &  \\
\hline
5611 & tugakerfi &  &  &  \\
\hline
5612 & tugalogri, venjulegur logri, logri með grunntölu, logri með stofntölu, Briggs-logri &  &  &  \\
\hline
5613 & tugasæti &  &  &  \\
\hline
5614 & tugastafur &  &  &  \\
\hline
5615 & tugastafur &  &  &  \\
\hline
5616 & tugatala, tala í tugakerfi &  &  &  \\
\hline
5617 & túlkun &  &  &  \\
\hline
5618 & tunna &  &  &  \\
\hline
5619 & tunnurúm &  &  &  \\
\hline
5620 & Turing-reiknanlegur &  &  &  \\
\hline
5621 & Turing-vél &  &  &  \\
\hline
5622 & turn &  &  &  \\
\hline
5623 & turn af sviðum &  &  &  \\
\hline
5624 & tveggja manna leikur &  &  &  \\
\hline
5625 & tveggjaferningasetning &  &  &  \\
\hline
5626 & tveir og tveir &  &  &  \\
\hline
5627 & tvennd, raðtvennd, röðuð tvennd & Par <ref idTerm=2123>hluta</ref> þar sem röð skiptir máli &  &  \\
\hline
5628 & tvíaðfelldur þríhyrningur &  &  &  \\
\hline
5629 & tvíblaða &  &  &  \\
\hline
5630 & tvíblaða breiðbogaflötur, tvískála breiðbogaflötur &  &  &  \\
\hline
5631 & tvíblaða þekjurúm &  &  &  \\
\hline
5632 & tvídeildanlegur, tvídiffranlegur &  &  &  \\
\hline
5633 & tvífléttingur &  &  &  \\
\hline
5634 & tvíflokkun &  &  &  \\
\hline
5635 & tvíflötungsgrúpa &  &  &  \\
\hline
5636 & tvíflötungur &  &  &  \\
\hline
5637 & tvígild rökfræði &  &  &  \\
\hline
5638 & tvígildur &  &  &  \\
\hline
5639 & tvígildur &  &  &  \\
\hline
5640 & tvíhliða &  &  &  \\
\hline
5641 & tvíhliða flötur &  &  &  \\
\hline
5642 & tvíhliða íðal &  &  &  \\
\hline
5643 & tvíhliða mótull &  &  &  \\
\hline
5644 & tvíhliða öryggisbil &  &  &  \\
\hline
5645 & tvíhliða próf &  &  &  \\
\hline
5646 & tvíhliða umhverfa &  &  &  \\
\hline
5647 & tvíhliða, frá báðum hliðum &  &  &  \\
\hline
5648 & tvíhlutanet &  &  &  \\
\hline
5649 & tvíhornalínufylki &  &  &  \\
\hline
5650 & tvíhyrningur &  &  &  \\
\hline
5651 & tvíhyrningur, kúlutvíhyrningur, kúlusigð &  &  &  \\
\hline
5652 & tvíkeila, tvöföld keila &  &  &  \\
\hline
5653 & tvíkostur &  &  &  \\
\hline
5654 & tvíkostur Fredholms, tilvik Fredholms &  &  &  \\
\hline
5655 & tvíkryppudreifing, tvítindadreifing &  &  &  \\
\hline
5656 & tvíliða jafna, tvíliðujafna &  &  &  \\
\hline
5657 & tvíliða, tvíliða stærð &  &  &  \\
\hline
5658 & tvíliðu-, tvíliða &  &  &  \\
\hline
5659 & tvíliðudreifing, tvíkostadreifing, Bernoulli-dreifing &  &  &  \\
\hline
5660 & tvíliðuregla &  &  &  \\
\hline
5661 & tvíliðuröð &  &  &  \\
\hline
5662 & tvíliðusetning &  &  &  \\
\hline
5663 & tvíliðustuðull &  &  &  \\
\hline
5664 & tvílíkindafylki &  &  &  \\
\hline
5665 & tvílínuleg vörpun &  &  &  \\
\hline
5666 & tvílínulegt form &  &  &  \\
\hline
5667 & tvílínulegur &  &  &  \\
\hline
5668 & tvílitun &  &  &  \\
\hline
5669 & tvílotubundinn, tvílotulegur &  &  &  \\
\hline
5670 & tvílotuleiki &  &  &  \\
\hline
5671 & tvinn-, tvinntalna- &  &  &  \\
\hline
5672 & tvinnbreyta, tvinnbreytistærð &  &  &  \\
\hline
5673 & tvinnfall, tvinngilt fall, tvinntölugilt fall &  &  &  \\
\hline
5674 & tvinnfallagreining, tvinnfallafræði, tvinngreining &  &  &  \\
\hline
5675 & tvinnfylki, tvinntalnafylki &  &  &  \\
\hline
5676 & tvinngildur, tvinntölugildur &  &  &  \\
\hline
5677 & tvinnlína &  &  &  \\
\hline
5678 & tvinnlínubundin &  &  &  \\
\hline
5679 & tvinnpunktur &  &  &  \\
\hline
5680 & tvinnrúmfræði &  &  &  \\
\hline
5681 & tvinnslétta &  &  &  \\
\hline
5682 & tvinntala &  &  &  \\
\hline
5683 & tvinntalnakerfi &  &  &  \\
\hline
5684 & tvinntalnarúm &  &  &  \\
\hline
5685 & tvinntalnaslétta, tvinnslétta &  &  &  \\
\hline
5686 & tvinntalnasvið &  &  &  \\
\hline
5687 & tvinnun &  &  &  \\
\hline
5688 & tvinnvarplína &  &  &  \\
\hline
5689 & tvinnvarprúm, varprúm yfir tvinntölurnar &  &  &  \\
\hline
5690 & tvinnvarpslétta, varpslétta yfir tvinntölurnar &  &  &  \\
\hline
5691 & tvinnvigurrúm, tvinnlínulegt rúm, vigurrúm yfir tvinntölurnar, línulegt rúm yfir tvinntölurnar &  &  &  \\
\hline
5692 & tvinnþverstaðalgrúpa, þverstaðalgrúpa yfir tvinntölurnar &  &  &  \\
\hline
5693 & tvípunktur, tvöfaldur punktur &  &  &  \\
\hline
5694 & tvíraðafylgni, tvírunufylgni &  &  &  \\
\hline
5695 & tvírétthyrndur &  &  &  \\
\hline
5696 & tvískífa &  &  &  \\
\hline
5697 & tvískilyrðing, gagnkvæm skilyrðing &  &  &  \\
\hline
5698 & tvískipting &  &  &  \\
\hline
5699 & tvískipting, tvíhlutun &  &  &  \\
\hline
5700 & tvískiptitala &  &  &  \\
\hline
5701 & tvísniðill &  &  &  \\
\hline
5702 & tvístæð aðgerð & <ref idTerm=46>Aðgerð</ref> á <ref idTerm=3270>mengi</ref> $X$ af gerðinni $X \times X \to X$ kallast \textit{tvístæð aðgerð} &  &  \\
\hline
5703 & tvístæð samtenging, tvístæður tengill &  &  &  \\
\hline
5704 & tvístæð umsögn &  &  &  \\
\hline
5705 & tvístæð vensl &  &  &  \\
\hline
5706 & tvístæður &  &  &  \\
\hline
5707 & tvístakaruna &  &  &  \\
\hline
5708 & tvístigaður &  &  &  \\
\hline
5709 & tvístika fjölskylda &  &  &  \\
\hline
5710 & tvístökungur, tveggjastakamengi &  &  &  \\
\hline
5711 & tvístrun &  &  &  \\
\hline
5712 & tvisvar, tvisvar sinnum, tví- &  &  &  \\
\hline
5713 & tvítengt net &  &  &  \\
\hline
5714 & tvítugakerfi &  &  &  \\
\hline
5715 & tvítugflata &  &  &  \\
\hline
5716 & tvítugflötungur &  &  &  \\
\hline
5717 & tvíunda- &  &  &  \\
\hline
5718 & tvíundaframsetning &  &  &  \\
\hline
5719 & tvíundakerfi &  &  &  \\
\hline
5720 & tvíundasæti &  &  &  \\
\hline
5721 & tvíundatala &  &  &  \\
\hline
5722 & tvíundatalnaritun &  &  &  \\
\hline
5723 & tvíundatré, tvískiptitré &  &  &  \\
\hline
5724 & tvíveldasumma, ferningasumma &  &  &  \\
\hline
5725 & tvíveldi, annað veldi, ferningur &  &  &  \\
\hline
5726 & tvívíð dreifing &  &  &  \\
\hline
5727 & tvívíð Gaussdreifing, tvívíð normleg dreifing, tvívíð normaldreifing &  &  &  \\
\hline
5728 & tvíviðbætinn, tvísamleggjandi &  &  &  \\
\hline
5729 & tvívíður &  &  &  \\
\hline
5730 & tvívíður &  &  &  \\
\hline
5731 & tvívigur &  &  &  \\
\hline
5732 & tvívíxill &  &  &  \\
\hline
5733 & tvíþverill &  &  &  \\
\hline
5734 & tvíþýð deildajafna, tvíþýð hlutafleiðujafna &  &  &  \\
\hline
5735 & tvíþýtt fall &  &  &  \\
\hline
5736 & tvöfaldlega samanhangandi, tvöfaldlega samhangandi &  &  &  \\
\hline
5737 & tvöfaldur &  &  &  \\
\hline
5738 & tvöfaldur beygjuskilakrosspunktur &  &  &  \\
\hline
5739 & tvöfalt heildi, tvöfalt tegur &  &  &  \\
\hline
5740 & tvöföld neitun &  &  &  \\
\hline
5741 & tvöföld röð &  &  &  \\
\hline
5742 & tvöföld rót &  &  &  \\
\hline
5743 & tvöföld runa &  &  &  \\
\hline
5744 & tvöföldun tenings &  &  &  \\
\hline
5745 & tvöföldunarregla &  &  &  \\
\hline
5746 & tylfta-, í tylftakerfi &  &  &  \\
\hline
5747 & tylftabrot &  &  &  \\
\hline
5748 & tylftakerfi &  &  &  \\
\hline
5749 & tylftaritháttur &  &  &  \\
\hline
5750 & úffeldi &  &  &  \\
\hline
5751 & úfur &  &  &  \\
\hline
5752 & umáttun &  &  &  \\
\hline
5753 & umbunarfall &  &  &  \\
\hline
5754 & umfangsfrumsenda, umfangsfrumsetning, frumsenda um umfang, frumsetning um umfang &  &  &  \\
\hline
5755 & umferðarheildi, umferðartegur &  &  &  \\
\hline
5756 & umhverf tala, umhverfa, margföldunarumhverfa, andhverf tala, andhverfa, margföldunarandhverfa, gagnkvæm tala &  &  &  \\
\hline
5757 & umhverfa, andhverfa & \textit{Andhverfa} <ref idTerm=5128>staks</ref> $x$ í <ref idTerm=3270>mengi</ref> $X$ með tilliti til <ref idTerm=46>reikniaðgerðarinnar</ref> $\odot$ með <ref idTerm=0>hlutleysu</ref> $e$ er stak $y \in X$ sem uppfyllir $x \odot y = y \odot x = e$ &  & Tala um $-x$ og $x^{-1}$ í skýringu, tengni neyðir andhverfu til að vera otvírætt ákvarðaða - Ekki alveg skýrt út frá skilgreiningunni að andhverfa þarf ekki endilega að vera til? \\
\hline
5758 & umhverfanlegt fylki, andhverfanlegt fylki, reglulegt fylki, ósérstætt fylki &  &  &  \\
\hline
5759 & umhverfanlegur, andhverfanlegur &  &  &  \\
\hline
5760 & umhverfanlegur, andhverfanlegur, reglulegur, ósérstæður &  &  &  \\
\hline
5761 & umhverfður veldisvísisferill &  &  &  \\
\hline
5762 & umhverfing &  &  &  \\
\hline
5763 & umhverfing &  &  &  \\
\hline
5764 & umhverfingarhringur, spegilhringur &  &  &  \\
\hline
5765 & umhverfingarmiðja, speglunarmiðja &  &  &  \\
\hline
5766 & umhverft brot &  &  &  \\
\hline
5767 & umhverft fall, marföldunarumhverfa falls, margföldunarandhverfa falls &  &  &  \\
\hline
5768 & umhverfur &  &  &  \\
\hline
5769 & umhverfur, margföldunarumhverfur, andhverfur, margföldunarandhverfur &  &  &  \\
\hline
5770 & umhverfuröð, þýð röð &  &  &  \\
\hline
5771 & umkvörðun, kvarðafærsla, kvarðabreyting &  &  &  \\
\hline
5772 & umlykja, fela í sér, innifela &  &  &  \\
\hline
5773 & umlykjandi &  &  &  \\
\hline
5774 & ummál &  &  &  \\
\hline
5775 & ummálsjafnaðarójafna &  &  &  \\
\hline
5776 & ummálsjafnaðarverkefni &  &  &  \\
\hline
5777 & ummálsregla, hálfhornsregla &  &  &  \\
\hline
5778 & ummiðja &  &  &  \\
\hline
5779 & ummyndun, myndun, færsla & <ref idTerm=3282>Vörpun</ref> $\phi : \mathbb{R}^2 \to \mathbb{R}^2$ sem uppfyllir

\begin{enumerate}
    \item punktar í tiltekinni röð á sömu línu færast á punkta í sömu röð á einhverri línu,
    \item strik sem eru jafn löng færast á strik sem eru líka jafn löng
\end{enumerate} &  & Er ekki bara verið að lýsa línulegri vörpun? \\
\hline
5780 & ummyndun, umskrift, umritun, breyting &  &  &  \\
\hline
5781 & umreikningur &  &  &  \\
\hline
5782 & umrita &  &  &  \\
\hline
5783 & umritaður &  &  &  \\
\hline
5784 & umritaður hringur & Ef <ref idTerm=2235>hornpunktar</ref> <ref idTerm=3161>marghyrnings</ref> liggja allir á gefnum <ref idTerm=2295>hring</ref>, þá kallast hringurinn \textit{umritaður hringur} marghyrningsins &  &  \\
\hline
5785 & umrituð kúla &  &  &  \\
\hline
5786 & umritun &  &  &  \\
\hline
5787 & umritun &  &  &  \\
\hline
5788 & umröðun &  &  &  \\
\hline
5789 & umröðunarsetning Riemanns &  &  &  \\
\hline
5790 & umsagnarbreyta &  &  &  \\
\hline
5791 & umsagnareikningur &  &  &  \\
\hline
5792 & umsagnarökfræði &  &  &  \\
\hline
5793 & umsagnartákn &  &  &  \\
\hline
5794 & umskiptanleiki &  &  &  \\
\hline
5795 & umskipting, tvírás &  &  &  \\
\hline
5796 & umsögn &  &  &  \\
\hline
5797 & umstikun, endurstikun, stikabreyting, stikaskipti &  &  &  \\
\hline
5798 & umtak, yfirgrip, umfang &  &  &  \\
\hline
5799 & undanfari &  &  &  \\
\hline
5800 & undantekningar- &  &  &  \\
\hline
5801 & undinn &  &  &  \\
\hline
5802 & undinn þriðja stigs ferill &  &  &  \\
\hline
5803 & undirbúningssetning &  &  &  \\
\hline
5804 & undirfall &  &  &  \\
\hline
5805 & undirferðarpunktur &  &  &  \\
\hline
5806 & undirforrit &  &  &  \\
\hline
5807 & undirhornalína &  &  &  \\
\hline
5808 & undirliður &  &  &  \\
\hline
5809 & undirmat, mat að neðan &  &  &  \\
\hline
5810 & undirmeta, meta að neðan &  &  &  \\
\hline
5811 & undirmetanlegur &  &  &  \\
\hline
5812 & undirröð &  &  &  \\
\hline
5813 & undirsetning, undirforsenda &  &  &  \\
\hline
5814 & undirskipað mengi &  &  &  \\
\hline
5815 & undirskipaður &  &  &  \\
\hline
5816 & undirsporger &  &  &  \\
\hline
5817 & undirstaða &  &  &  \\
\hline
5818 & undirstak, undirskorða, undirtala & Fyrir <ref idTerm=4093>rauntalna</ref><ref idTerm=3270>mengi</ref> $A$ kallast rauntalan $L$ \textit{undirtala} mengisins $A$ ef fyrir öll $x \in A$ gildir að $x \geq L$ &  &  \\
\hline
5819 & undirstöðuflokkur &  &  &  \\
\hline
5820 & undirstöðugrúpa &  &  &  \\
\hline
5821 & undirstöðugrýpi &  &  &  \\
\hline
5822 & undirstöðusetning algebrunnar, grunnsetning algebrunnar & \textit{Undirstöðusetning algebrunnar} segir að sérhver <ref idTerm=1948>margliða</ref> yfir <ref idTerm=5682>tvinntölurnar</ref> $\mathbb{C}$ með stig $\geq 1$ eigi sér núllstöð í $\mathbb{C}$ &  &  \\
\hline
5823 & undirstöðusetning örsmæðareikningsins, grunnsetning reiknivísinnar &  &  &  \\
\hline
5824 & undirstöðusetning reikningslistarinnar, grunnsetning reikningslistarinnar, frumþáttunarsetning & \textit{Undirstöðusetning reikningslistarinnar} segir að sérhverja <ref idTerm=3475>náttúrulega tölu</ref> $n$ megi skrifa á nákvæmlega einn hátt sem <ref idTerm=3118>margfeldi</ref> af <ref idTerm=1468>frumtölum</ref>. Nánar tiltekið gildir fyrir sérhverja náttúrulega tölu $n$ að til eru <ref idTerm=3932>ótvírætt ákvarðaðar</ref> frumtölur $p_{1}, \ldots ,p_{k}$ og náttúrulegar tölur $a_{1}, \ldots ,a_{k}$ þannig að:

\[n = {p_{1}}^{a_{1}} \cdot {p_{2}}^{a_{2}} \cdot \ldots \cdot {p_{k-1}}^{a_{k-1}} \cdot {p_{k}}^{a_{k}} \quad \text{þar sem $p_{i} < p_{j}$ þegar $i < j$}. \]

Að skrifa töluna $n$ á þessu formi kallast að \textit{frumþátta} hana. <ref idTerm=5397>Tölurnar</ref> $p_{1}, \ldots, p_{k}$ kallast þá \textit{frumþættir} tölunnar $n$ og $a_{i}$ kallast \textit{margfeldni} frumþáttarins $p_{i}$ &  &  \\
\hline
5825 & undirsumma &  &  &  \\
\hline
5826 & undirviðbætinn, undirsamleggjandi &  &  &  \\
\hline
5827 & undirþýður &  &  &  \\
\hline
5828 & uppblástur &  &  &  \\
\hline
5829 & uppbygging, smíð &  &  &  \\
\hline
5830 & uppbyggingaraðferð &  &  &  \\
\hline
5831 & uppbyggingarhyggja &  &  &  \\
\hline
5832 & uppbyggjandi &  &  &  \\
\hline
5833 & uppbyggjandi greining &  &  &  \\
\hline
5834 & uppbyggjandi tilvistarsetning &  &  &  \\
\hline
5835 & uppfylling fernings &  &  &  \\
\hline
5836 & uppfyllingarsetning &  &  &  \\
\hline
5837 & uppfyllingarsetning Steinitz &  &  &  \\
\hline
5838 & upphæð &  &  &  \\
\hline
5839 & upphaf, upphafspunktur, hnitmiðja, hnitamiðja &  &  &  \\
\hline
5840 & upphafs- &  &  &  \\
\hline
5841 & upphafsástand &  &  &  \\
\hline
5842 & upphafsgildi, byrjunargildi &  &  &  \\
\hline
5843 & upphafsgildisverkefni, byrjunargildisverkefni, Cauchy-verkefni &  &  &  \\
\hline
5844 & upphafsgrannmynstur &  &  &  \\
\hline
5845 & upphafshlutur &  &  &  \\
\hline
5846 & upphafshnútur, upphafsendi &  &  &  \\
\hline
5847 & upphafsliður &  &  &  \\
\hline
5848 & upphafspunktur &  &  &  \\
\hline
5849 & upphafspunktur &  &  &  \\
\hline
5850 & upphafspunktur &  &  &  \\
\hline
5851 & upphafsskilyrði, byrjunarskilyrði &  &  &  \\
\hline
5852 & upphafsstúfur &  &  &  \\
\hline
5853 & upphenging, hengirúm &  &  &  \\
\hline
5854 & uppleysing &  &  &  \\
\hline
5855 & upplýsingafræði &  &  &  \\
\hline
5856 & upplýsingamagn &  &  &  \\
\hline
5857 & upplýsingar &  &  &  \\
\hline
5858 & uppröðun &  &  &  \\
\hline
5859 & uppsafnaður, hlaðinn &  &  &  \\
\hline
5860 & uppsöfnuð áhrif &  &  &  \\
\hline
5861 & uppsöfnuð skekkja, heildarskekkja &  &  &  \\
\hline
5862 & uppspretta, lind &  &  &  \\
\hline
5863 & uppsprettulaus &  &  &  \\
\hline
5864 & uppsprettulaust svið &  &  &  \\
\hline
5865 & uppsprettustyrkur, lindarstyrkur &  &  &  \\
\hline
5866 & uppstokka, stokka upp, umraða &  &  &  \\
\hline
5867 & uppstokkanagrúpa, umraðanagrúpa &  &  &  \\
\hline
5868 & uppstokkanaútsetning, umraðanaútsetning &  &  &  \\
\hline
5869 & uppstokkun með endurtekningu, umröðun með endurtekningu &  &  &  \\
\hline
5870 & uppstokkun, umröðun &  &  &  \\
\hline
5871 & úrkynjaður &  &  &  \\
\hline
5872 & úrkynjun &  &  &  \\
\hline
5873 & úrskurðanleg fullyrðing, úrskurðanleg yrðing &  &  &  \\
\hline
5874 & úrskurðanleg kenning &  &  &  \\
\hline
5875 & úrtak &  &  &  \\
\hline
5876 & úrtaka án skila &  &  &  \\
\hline
5877 & úrtaka með skilum &  &  &  \\
\hline
5878 & úrtaka, úrtakskönnun, könnun (með úrtaki) &  &  &  \\
\hline
5879 & úrtakseinkenni &  &  &  \\
\hline
5880 & úrtakshögun &  &  &  \\
\hline
5881 & úrtaksmeðaltal, meðaltal &  &  &  \\
\hline
5882 & úrtaksmiðtala &  &  &  \\
\hline
5883 & úrtakssamfylgni, úrtakssamvik &  &  &  \\
\hline
5884 & úrtaksskekkja &  &  &  \\
\hline
5885 & úrtaksspönn &  &  &  \\
\hline
5886 & úrtaksstærð, stærð úrtaks &  &  &  \\
\hline
5887 & út-, útgildis- &  &  &  \\
\hline
5888 & utanverð víxlhorn, ytri víxlhorn &  &  &  \\
\hline
5889 & útbreiddur &  &  &  \\
\hline
5890 & útbreiðsla &  &  &  \\
\hline
5891 & útbreiðsla &  &  &  \\
\hline
5892 & útbreiðslutala &  &  &  \\
\hline
5893 & útbreitt rúm &  &  &  \\
\hline
5894 & útflötur, hlið &  &  &  \\
\hline
5895 & útgeisli &  &  &  \\
\hline
5896 & útgildispunktur &  &  &  \\
\hline
5897 & útgildisverkefni &  &  &  \\
\hline
5898 & úthjólferill &  &  &  \\
\hline
5899 & úthluta &  &  &  \\
\hline
5900 & úthlutun &  &  &  \\
\hline
5901 & úthlutunarverkefni &  &  &  \\
\hline
5902 & úthorn, fyllihorn &  &  &  \\
\hline
5903 & útilokun &  &  &  \\
\hline
5904 & útinnritaður &  &  &  \\
\hline
5905 & útinnritaður hringur, úthringur, utanverður snertihringur &  &  &  \\
\hline
5906 & útinnrituð kúla, útkúla, utanverð snertikúla &  &  &  \\
\hline
5907 & útkoma, niðurstaða &  &  &  \\
\hline
5908 & útkomandi &  &  &  \\
\hline
5909 & útkomurúm &  &  &  \\
\hline
5910 & útliður &  &  &  \\
\hline
5911 & útlína, myndjaðar &  &  &  \\
\hline
5912 & útmengi, útsvæði &  &  &  \\
\hline
5913 & útmiðja, úthringsmiðja &  &  &  \\
\hline
5914 & útör &  &  &  \\
\hline
5915 & útpunktur &  &  &  \\
\hline
5916 & útpunktur &  &  &  \\
\hline
5917 & útreikningur &  &  &  \\
\hline
5918 & útreikningur, reikningur &  &  &  \\
\hline
5919 & útrýma, eyða &  &  &  \\
\hline
5920 & útrýming, eyðing &  &  &  \\
\hline
5921 & útrýmingaraðferð, eyðingaraðferð &  &  &  \\
\hline
5922 & útsetning &  &  &  \\
\hline
5923 & útsetningafræði &  &  &  \\
\hline
5924 & útsetningarrúm &  &  &  \\
\hline
5925 & útslag, sveifluhæð, ölduhæð &  &  &  \\
\hline
5926 & útstæður þvervigur, útstæður þverill, ytri þvervigur, ytri þverill &  &  &  \\
\hline
5927 & útstætt horn &  &  &  \\
\hline
5928 & útstig &  &  &  \\
\hline
5929 & útsveifarferill &  &  &  \\
\hline
5930 & úttak &  &  &  \\
\hline
5931 & útvíkka, framlengja &  &  &  \\
\hline
5932 & útvíkkanlegur, framlengjanlegur &  &  &  \\
\hline
5933 & útvíkkuð rauntalnalína &  &  &  \\
\hline
5934 & útvíkkuð tvinntalnaslétta &  &  &  \\
\hline
5935 & útvíkkun &  &  &  \\
\hline
5936 & útvíkkun grunnbaugsins, víkkun grunnbaugsins &  &  &  \\
\hline
5937 & útvíkkun, víkkun &  &  &  \\
\hline
5938 & vægi &  &  &  \\
\hline
5939 & vægi &  &  &  \\
\hline
5940 & vægi um meðalgildi &  &  &  \\
\hline
5941 & vægisfall &  &  &  \\
\hline
5942 & vægisframleiðandi fall &  &  &  \\
\hline
5943 & vægisstuðull &  &  &  \\
\hline
5944 & vægisverkefni &  &  &  \\
\hline
5945 & vægjaaðferð &  &  &  \\
\hline
5946 & vægjafylki &  &  &  \\
\hline
5947 & vægjahlutfall &  &  &  \\
\hline
5948 & væntanlegt gildi, væntigildi, meðalgildi &  &  &  \\
\hline
5949 & vafningstala &  &  &  \\
\hline
5950 & val &  &  &  \\
\hline
5951 & val &  &  &  \\
\hline
5952 & val, valfrelsi &  &  &  \\
\hline
5953 & valfall, valvörpun &  &  &  \\
\hline
5954 & valfrumsenda, valfrumsetning, frumsenda um val, frumsetning um val &  &  &  \\
\hline
5955 & valstund &  &  &  \\
\hline
5956 & valtari &  &  &  \\
\hline
5957 & valúrtaka &  &  &  \\
\hline
5958 & vanákvarðað kerfi &  &  &  \\
\hline
5959 & vanaleg tala &  &  &  \\
\hline
5960 & vanalegur &  &  &  \\
\hline
5961 & vanhorn &  &  &  \\
\hline
5962 & varðveisla, rækni &  &  &  \\
\hline
5963 & varðveita, rækja &  &  &  \\
\hline
5964 & varfærið próf &  &  &  \\
\hline
5965 & varp- &  &  &  \\
\hline
5966 & varpanarúm &  &  &  \\
\hline
5967 & varparíki &  &  &  \\
\hline
5968 & varpgildi &  &  &  \\
\hline
5969 & varphnit &  &  &  \\
\hline
5970 & varpi &  &  &  \\
\hline
5971 & varpkeilusnið &  &  &  \\
\hline
5972 & varplína &  &  &  \\
\hline
5973 & varpmótull &  &  &  \\
\hline
5974 & varpmótun, línuleg varpmótun &  &  &  \\
\hline
5975 & varpmyndun &  &  &  \\
\hline
5976 & varpóháður &  &  &  \\
\hline
5977 & varprúm &  &  &  \\
\hline
5978 & varprúmfræði &  &  &  \\
\hline
5979 & varpslétta &  &  &  \\
\hline
5980 & vaxtarfall &  &  &  \\
\hline
5981 & vaxtarhraði &  &  &  \\
\hline
5982 & vaxtarregla &  &  &  \\
\hline
5983 & vaxtarröð, vaxtarraðstig &  &  &  \\
\hline
5984 & vaxtavextir &  &  &  \\
\hline
5985 & veðfé &  &  &  \\
\hline
5986 & veðmál &  &  &  \\
\hline
5987 & vefja Arkimedesar, kuðungur Arkimedesar, kuðungslína Arkimedesar, snigill Arkimedesar &  &  &  \\
\hline
5988 & vefja, kuðungslína, kuðungsferill, kuðungur, snigill &  &  &  \\
\hline
5989 & vefur &  &  &  \\
\hline
5990 & vegið meðaltal &  &  &  \\
\hline
5991 & veik lausn &  &  &  \\
\hline
5992 & veik lokun &  &  &  \\
\hline
5993 & veik ójafna &  &  &  \\
\hline
5994 & veik samleitni &  &  &  \\
\hline
5995 & veikt grannmynstur, gróft grannmynstur &  &  &  \\
\hline
5996 & veikt markgildi &  &  &  \\
\hline
5997 & veikur &  &  &  \\
\hline
5998 & vel framsett verkefni &  &  &  \\
\hline
5999 & vel skilgreindur &  &  &  \\
\hline
6000 & vélarmál &  &  &  \\
\hline
6001 & vélbúnaður &  &  &  \\
\hline
6002 & veldaraðaaðferð &  &  &  \\
\hline
6003 & veldaraðabaugur &  &  &  \\
\hline
6004 & veldaraðarframsetning &  &  &  \\
\hline
6005 & veldaraðarliðun &  &  &  \\
\hline
6006 & veldaröð &  &  &  \\
\hline
6007 & veldi &  &  &  \\
\hline
6008 & veldi af frumtölu &  &  &  \\
\hline
6009 & veldisfall & <ref idTerm=1078>Fall</ref> sem varpar sérhverju <ref idTerm=5128>staki</ref> $x$ í $n$-ta <ref idTerm=6007>veldi</ref> þess $x^n$, þar sem $n \geq 2$ er <ref idTerm=3475>náttúruleg tala</ref> &  &  \\
\hline
6010 & veldisfrumsenda, veldisfrumsetning, frumsenda um veldismengi, frumsetning um veldismengi &  &  &  \\
\hline
6011 & veldisleif &  &  &  \\
\hline
6012 & veldisleifatákn &  &  &  \\
\hline
6013 & veldismengi & <ref idTerm=3270>Mengið</ref> sem samanstendur af öllum <ref idTerm=2098>hlutmengjum</ref> í tilteknu mengi $X$ kallast \textit{veldismengi} mengisins $X$ og er táknað með $\mathcal{P}(X)$ &  & Setja í skýringu: Almennt gildir að ef fjöldi staka í menginu $A$ er $n$, þar sem $n$ er náttúruleg tala, er fjöldi staka í veldismengi þess $2^n$. \\
\hline
6014 & veldisstofn &  &  &  \\
\hline
6015 & veldisvísajafna &  &  &  \\
\hline
6016 & veldisvísaregla &  &  &  \\
\hline
6017 & veldisvísir &  &  &  \\
\hline
6018 & veldisvísis- &  &  &  \\
\hline
6019 & veldisvísisdreifing, vísisdreifing &  &  &  \\
\hline
6020 & veldisvísisfall, vísisfall &  &  &  \\
\hline
6021 & veldisvísisferill, vísisferill &  &  &  \\
\hline
6022 & veldisvísisgerð &  &  &  \\
\hline
6023 & veldisvísisvöxtur, vísisvöxtur &  &  &  \\
\hline
6024 & vélheildun, véltegrun &  &  &  \\
\hline
6025 & vellauðugur &  &  &  \\
\hline
6026 & velraða &  &  &  \\
\hline
6027 & velraðað mengi &  &  &  \\
\hline
6028 & velraðaður &  &  &  \\
\hline
6029 & velröðun &  &  &  \\
\hline
6030 & velröðunarlögmál, velröðunarsetning &  &  &  \\
\hline
6031 & veltibraut &  &  &  \\
\hline
6032 & veltiferill &  &  &  \\
\hline
6033 & veltilína, veltibraut &  &  &  \\
\hline
6034 & vélvit, tölvuvit, gervigreind &  &  &  \\
\hline
6035 & vendi- &  &  &  \\
\hline
6036 & vendiaðferð &  &  &  \\
\hline
6037 & vendidálkur &  &  &  \\
\hline
6038 & vendifylki &  &  &  \\
\hline
6039 & vendilína &  &  &  \\
\hline
6040 & vendistak &  &  &  \\
\hline
6041 & vendistuðull &  &  &  \\
\hline
6042 & venjuleg deildajafna, venjuleg diffurjafna, venjuleg afleiðujafna &  &  &  \\
\hline
6043 & venjuleg firð &  &  &  \\
\hline
6044 & venjulegt grannmynstur &  &  &  \\
\hline
6045 & venjulegt innfeldi &  &  &  \\
\hline
6046 & venjulegt róf &  &  &  \\
\hline
6047 & venjulegur &  &  &  \\
\hline
6048 & venjulegur &  &  &  \\
\hline
6049 & venjulegur punktur &  &  &  \\
\hline
6050 & Venn-mynd & Myndræn <ref idTerm=1416>framsetning</ref> á innbyrðis afstöðu ólíkra <ref idTerm=3270>mengja</ref> þar sem <ref idTerm=3016>lokaðir ferlar</ref> (eins og <ref idTerm=2295>hringir</ref>, <ref idTerm=5002>sporbaugar</ref>, <ref idTerm=1127>ferhyrningar</ref> o.s.frv.) eru notaðir til að tákna mengi og <ref idTerm=5335>svæðið</ref> sem er innan ferlanna táknar <ref idTerm=5128>stökin</ref> í menginu &  & Setja myndir í skýringu \\
\hline
6051 & vensl &  &  &  \\
\hline
6052 & venslarit, vensl &  &  &  \\
\hline
6053 & venslastæða &  &  &  \\
\hline
6054 & venslatenging &  &  &  \\
\hline
6055 & vera á eftir &  &  &  \\
\hline
6056 & vera jafn &  &  &  \\
\hline
6057 & vera ósamleitinn &  &  &  \\
\hline
6058 & vera samleitinn &  &  &  \\
\hline
6059 & vera samur &  &  &  \\
\hline
6060 & vera til &  &  &  \\
\hline
6061 & vera víxlinn &  &  &  \\
\hline
6062 & verka &  &  &  \\
\hline
6063 & verka frá hægri &  &  &  \\
\hline
6064 & verka frá vinstri &  &  &  \\
\hline
6065 & verka frjálslega, verka án kyrrapunkta, verka án fastapunkta &  &  &  \\
\hline
6066 & verka samfellt, verka með samfelldum hætti &  &  &  \\
\hline
6067 & verka tryggilega &  &  &  \\
\hline
6068 & verkefni &  &  &  \\
\hline
6069 & verkröðun &  &  &  \\
\hline
6070 & verkröðunarverkefni &  &  &  \\
\hline
6071 & verkun &  &  &  \\
\hline
6072 & verkun frá hægri &  &  &  \\
\hline
6073 & verkun frá vinstri &  &  &  \\
\hline
6074 & verulegur sérstöðupunktur &  &  &  \\
\hline
6075 & vextir &  &  &  \\
\hline
6076 & víðátta &  &  &  \\
\hline
6077 & víðátta án jaðars &  &  &  \\
\hline
6078 & víðátta með jaðri &  &  &  \\
\hline
6079 & viðbætið fall, samleggjandi fall &  &  &  \\
\hline
6080 & viðbætið mengjafall, samleggjandi mengjafall &  &  &  \\
\hline
6081 & viðbætinn, samleggjandi &  &  &  \\
\hline
6082 & viðbót &  &  &  \\
\hline
6083 & viðbótarskilyrði, viðbótarkrafa, viðbótarforsenda &  &  &  \\
\hline
6084 & viðbótarupplýsingar, fylgiupplýsingar &  &  &  \\
\hline
6085 & vídd, stétt &  &  &  \\
\hline
6086 & vídd, stig &  &  &  \\
\hline
6087 & vídd, víddartala, rúmvídd &  &  &  \\
\hline
6088 & víddarfræði &  &  &  \\
\hline
6089 & víðerni &  &  &  \\
\hline
6090 & viðfangskenning &  &  &  \\
\hline
6091 & viðfangsmál &  &  &  \\
\hline
6092 & viðfangssetning &  &  &  \\
\hline
6093 & víðfeðm einstikun &  &  &  \\
\hline
6094 & víðfeðm greining, víðgreining &  &  &  \\
\hline
6095 & víðfeðm tilvist &  &  &  \\
\hline
6096 & víðfeðmt sterkasta próf &  &  &  \\
\hline
6097 & víðfeðmt útgildi &  &  &  \\
\hline
6098 & víðfeðmur &  &  &  \\
\hline
6099 & víðfeðmur hágildispunktur &  &  &  \\
\hline
6100 & víðfeðmur lággildispunktur &  &  &  \\
\hline
6101 & víðfeðmur útgildispunktur &  &  &  \\
\hline
6102 & víðgildi &  &  &  \\
\hline
6103 & viðkennd tala &  &  &  \\
\hline
6104 & viðmið &  &  &  \\
\hline
6105 & viðmiðun &  &  &  \\
\hline
6106 & víðtækur, almennur, víður &  &  &  \\
\hline
6107 & viðtaka, viðtekt &  &  &  \\
\hline
6108 & viðtökulíkindi &  &  &  \\
\hline
6109 & viðtökumark &  &  &  \\
\hline
6110 & viðtökupróf &  &  &  \\
\hline
6111 & viðunandi &  &  &  \\
\hline
6112 & viðunandi gæðamörk &  &  &  \\
\hline
6113 & viðureign &  &  &  \\
\hline
6114 & viðvörunarmörk &  &  &  \\
\hline
6115 & vigur- &  &  &  \\
\hline
6116 & vigur, vektor &  &  &  \\
\hline
6117 & viguralgebra &  &  &  \\
\hline
6118 & vigurbundin &  &  &  \\
\hline
6119 & vigurfall &  &  &  \\
\hline
6120 & vigurgildur &  &  &  \\
\hline
6121 & vigurgilt fall, vigurfall &  &  &  \\
\hline
6122 & vigurgreining &  &  &  \\
\hline
6123 & vigursamlagning, vigrasamlagning &  &  &  \\
\hline
6124 & vigursumma &  &  &  \\
\hline
6125 & vigursumma &  &  &  \\
\hline
6126 & vigursvið &  &  &  \\
\hline
6127 & vik &  &  &  \\
\hline
6128 & vild samlínumótun &  &  &  \\
\hline
6129 & vildarás &  &  &  \\
\hline
6130 & vildarfall, línufall &  &  &  \\
\hline
6131 & vildargrúpa &  &  &  \\
\hline
6132 & vildarhjúpur &  &  &  \\
\hline
6133 & vildarhnit, samsíða hnit &  &  &  \\
\hline
6134 & vildarlína &  &  &  \\
\hline
6135 & vildarlokun, vildarhjúpur &  &  &  \\
\hline
6136 & vildarrúm &  &  &  \\
\hline
6137 & vildarrúmfræði &  &  &  \\
\hline
6138 & vildarsamgangur &  &  &  \\
\hline
6139 & vildarslétta &  &  &  \\
\hline
6140 & vildarvíðerni &  &  &  \\
\hline
6141 & vildarvörpun, vildarmótun &  &  &  \\
\hline
6142 & vildur, vildar- &  &  &  \\
\hline
6143 & villa, mistök &  &  &  \\
\hline
6144 & vinatölur &  &  &  \\
\hline
6145 & vindingalaus grúpa, vindingslaus grúpa &  &  &  \\
\hline
6146 & vindingsfeldi &  &  &  \\
\hline
6147 & vindingsgeisli &  &  &  \\
\hline
6148 & vindingsgrúpa &  &  &  \\
\hline
6150 & vindingslaus, vindingalaus &  &  &  \\
\hline
6151 & vindingsmótull &  &  &  \\
\hline
6152 & vindingsstak &  &  &  \\
\hline
6153 & vindingsþinur &  &  &  \\
\hline
6154 & vindingur &  &  &  \\
\hline
6155 & vindingur &  &  &  \\
\hline
6156 & vinningsbandalag &  &  &  \\
\hline
6157 & vinningsfall &  &  &  \\
\hline
6158 & vinningsfylki &  &  &  \\
\hline
6159 & vinningur &  &  &  \\
\hline
6160 & vinnsla &  &  &  \\
\hline
6161 & vinstra einingarstak &  &  &  \\
\hline
6162 & vinstra hjámengi &  &  &  \\
\hline
6163 & vinstra íðal &  &  &  \\
\hline
6166 & vinstri deildanleiki, deildanleiki frá vinstri, vinstri diffranleiki, diffranleiki frá vinstri &  &  &  \\
\hline
6167 & vinstri dreifiregla & Sjá <ref idTerm=723>dreifiregla</ref> &  &  \\
\hline
6168 & vinstri endapunktur &  &  &  \\
\hline
6169 & vinstri grennd &  &  &  \\
\hline
6170 & vinstri handar hnitakerfi &  &  &  \\
\hline
6171 & vinstri hávísir &  &  &  \\
\hline
6172 & vinstri hlið &  &  &  \\
\hline
6173 & vinstri hornklofi &  &  &  \\
\hline
6174 & vinstri lágvísir &  &  &  \\
\hline
6175 & vinstri mótull &  &  &  \\
\hline
6176 & vinstri oddklofi &  &  &  \\
\hline
6177 & vinstri oddsvigi &  &  &  \\
\hline
6178 & vinstri samsemdarmótun &  &  &  \\
\hline
6179 & vinstri svigi &  &  &  \\
\hline
6180 & vinstri umhverfa, umhverfa frá vinstri, vinstri andhverfa, andhverfa frá vinstri &  &  &  \\
\hline
6181 & vinstrisnúin skrúflína &  &  &  \\
\hline
6182 & vinstrisnúinn &  &  &  \\
\hline
6183 & vinstrisnúinn ferill &  &  &  \\
\hline
6184 & virðing &  &  &  \\
\hline
6185 & virðingarbaugur &  &  &  \\
\hline
6186 & virðingarfræði &  &  &  \\
\hline
6187 & virðingargrannmynstur &  &  &  \\
\hline
6188 & virki &  &  &  \\
\hline
6189 & virkjagildur &  &  &  \\
\hline
6190 & virkjajafna &  &  &  \\
\hline
6191 & virkjareikningur &  &  &  \\
\hline
6192 & virkjastaðall &  &  &  \\
\hline
6193 & virkni &  &  &  \\
\hline
6194 & virkt hliðarskilyrði &  &  &  \\
\hline
6195 & virkt mengi &  &  &  \\
\hline
6196 & virkt stak &  &  &  \\
\hline
6197 & vísajafna, veldisvísajafna &  &  &  \\
\hline
6198 & vísalækkun &  &  &  \\
\hline
6199 & vísamengi, vísitölumengi, skilgreiningarmengi &  &  &  \\
\hline
6200 & vísasetja &  &  &  \\
\hline
6201 & vísasetning &  &  &  \\
\hline
6202 & vísasettur &  &  &  \\
\hline
6203 & vísir, vísitala &  &  &  \\
\hline
6204 & vísis- &  &  &  \\
\hline
6205 & vísisfallaröð &  &  &  \\
\hline
6206 & vísissetning &  &  &  \\
\hline
6207 & vísisvirðing &  &  &  \\
\hline
6208 & vísitala &  &  &  \\
\hline
6209 & vísitala &  &  &  \\
\hline
6210 & vísitala &  &  &  \\
\hline
6211 & vísitala &  &  &  \\
\hline
6212 & vísitala &  &  &  \\
\hline
6213 & viss atburður, öruggur atburður &  &  &  \\
\hline
6214 & viss, öruggur &  &  &  \\
\hline
6215 & vissa, fullvissa &  &  &  \\
\hline
6216 & vítahringur, hringrök &  &  &  \\
\hline
6218 & víxill &  &  &  \\
\hline
6219 & víxla &  &  &  \\
\hline
6220 & víxlagrúpa, afleidd grúpa &  &  &  \\
\hline
6221 & víxlanlegur &  &  &  \\
\hline
6222 & víxlast, vera víxlanlegur &  &  &  \\
\hline
6223 & víxlbaugur, víxlinn baugur &  &  &  \\
\hline
6224 & víxlhorn &  &  &  \\
\hline
6225 & víxlið örvarit &  &  &  \\
\hline
6226 & víxlimengi, víxlialgebra &  &  &  \\
\hline
6227 & víxlin aðgerð, víxlin reikniaðgerð &  &  &  \\
\hline
6228 & víxlinn &  &  &  \\
\hline
6229 & víxlmerkjamengi &  &  &  \\
\hline
6230 & víxlmerkjaröð &  &  &  \\
\hline
6231 & víxlmerkjasumma &  &  &  \\
\hline
6232 & víxlni &  &  &  \\
\hline
6233 & víxlregla & <ref idTerm=5702>Tvístæð reikniaðgerð</ref> $\odot : X \times X \to X$ á <ref idTerm=3270>mengi</ref> $X$ kallast \textit{víxlin} eða að hún fullnægi \textit{víxlreglunni} ef fyrir sérhver <ref idTerm=5128>stök</ref> $x,y \in X$ gildi að $x \odot y = y \odot x$ &  &  \\
\hline
6234 & víxlstefnuaðferð &  &  &  \\
\hline
6235 & víxlunarjafna &  &  &  \\
\hline
6236 & vogarskálaaðferð &  &  &  \\
\hline
6237 & vökvaaflfræði, straumfræði &  &  &  \\
\hline
6238 & vökvastöðufræði, kyrrðarfræði &  &  &  \\
\hline
6239 & vöndur af hringum, hringavöndur, hneppi af hringum, hringahneppi &  &  &  \\
\hline
6240 & vöndur af línum, línuvöndur, hneppi af línum, línuhneppi &  &  &  \\
\hline
6241 & vöndur af samsíða línum, samsíðuvöndur, hneppi af samsíða línum, samsíðuhneppi &  &  &  \\
\hline
6242 & vöndur af sléttum, sléttuvöndur, hneppi af sléttum, sléttuhneppi &  &  &  \\
\hline
6243 & vöndur, hneppi &  &  &  \\
\hline
6244 & vörpunarkeila &  &  &  \\
\hline
6245 & vörpunarsetning Riemanns &  &  &  \\
\hline
6246 & vörpunarsívalningur &  &  &  \\
\hline
6247 & vörpunarstig &  &  &  \\
\hline
6248 & vöxtur &  &  &  \\
\hline
6249 & Weierstrass-próf &  &  &  \\
\hline
6250 & Wiener-ferli &  &  &  \\
\hline
6251 & Wronski-ákveða &  &  &  \\
\hline
6252 & x-ás &  &  &  \\
\hline
6253 & y-ás &  &  &  \\
\hline
6254 & yfirákvarðaður, ofákvarðaður &  &  &  \\
\hline
6255 & yfirbaugur &  &  &  \\
\hline
6256 & yfirborð &  &  &  \\
\hline
6257 & yfirbreyta &  &  &  \\
\hline
6258 & yfirferðarpunktur &  &  &  \\
\hline
6259 & yfirgnæfa &  &  &  \\
\hline
6260 & yfirgnæfa &  &  &  \\
\hline
6261 & yfirgnæfð röð &  &  &  \\
\hline
6262 & yfirgnæfð samleitni &  &  &  \\
\hline
6263 & yfirgripsregla &  &  &  \\
\hline
6264 & yfirhjólferill &  &  &  \\
\hline
6265 & yfirhornalína &  &  &  \\
\hline
6266 & yfirkenning &  &  &  \\
\hline
6267 & yfirliður &  &  &  \\
\hline
6268 & yfirmál, yfirtungumál &  &  &  \\
\hline
6269 & yfirmengi &  &  &  \\
\hline
6270 & yfirmetanlegur &  &  &  \\
\hline
6271 & yfirmótull &  &  &  \\
\hline
6272 & yfirráðasvæði &  &  &  \\
\hline
6273 & yfirrásuð grúpa &  &  &  \\
\hline
6274 & yfirregla &  &  &  \\
\hline
6275 & yfirsetning &  &  &  \\
\hline
6276 & yfirsetning, yfirforsenda &  &  &  \\
\hline
6277 & yfirsporger &  &  &  \\
\hline
6278 & yfirstærðfræði &  &  &  \\
\hline
6279 & yfirstærðfræðilegur &  &  &  \\
\hline
6280 & yfirstak, yfirskorða, yfirtala & Fyrir <ref idTerm=4093>rauntalna</ref><ref idTerm=3270>mengi</ref> $A$ kallast rauntalan $U$ \textit{yfirtala} mengisins $A$ ef fyrir öll $x \in A$ gildir að $x \leq U$ &  &  \\
\hline
6281 & yfirstrik &  &  &  \\
\hline
6282 & yfirsumma &  &  &  \\
\hline
6283 & yfirsvið &  &  &  \\
\hline
6284 & yfirviðbætinn, yfirsamleggjandi &  &  &  \\
\hline
6285 & yfirvíxlin grúpa &  &  &  \\
\hline
6286 & yfirþýður &  &  &  \\
\hline
6287 & yrðingabreyta, fullyrðingabreyta &  &  &  \\
\hline
6288 & yrðingafall &  &  &  \\
\hline
6289 & yrðingamót &  &  &  \\
\hline
6290 & yrðingareikningur, fullyrðingareikningur &  &  &  \\
\hline
6291 & yrðingarökfræði, fullyrðingarökfræði &  &  &  \\
\hline
6292 & ytra deildaform, ytra diffurform, deildaform, diffurform &  &  &  \\
\hline
6293 & ytra deildi, ytra diffur &  &  &  \\
\hline
6294 & ytra flatarmál &  &  &  \\
\hline
6295 & ytra horn, utanvert horn & <ref idTerm=1677>Grannhorn</ref> tiltekins <ref idTerm=2198>horns</ref> í <ref idTerm=3161>marghyrningi</ref>. Horn marghyrningsins kallast þá \textit{innri horn} &  &  \\
\hline
6296 & ytra horn, utanvert horn &  &  &  \\
\hline
6297 & ytra innihald, ytra Jordan-innihald &  &  &  \\
\hline
6298 & ytra mál, utanmál &  &  &  \\
\hline
6299 & ytra mál, utanmál &  &  &  \\
\hline
6300 & ytra meðalfervik, ytri dreifni &  &  &  \\
\hline
6301 & ytra veldi &  &  &  \\
\hline
6302 & ytri aðgerð, ytri reikniaðgerð &  &  &  \\
\hline
6303 & ytri afleiða &  &  &  \\
\hline
6304 & ytri algebra, Grassmann-algebra &  &  &  \\
\hline
6305 & ytri einslögunarmiðja &  &  &  \\
\hline
6306 & ytri fylgni &  &  &  \\
\hline
6307 & ytri gagnstæð horn, utanverð gagnstæð horn &  &  &  \\
\hline
6308 & ytri helmingalína &  &  &  \\
\hline
6309 & ytri margföldun &  &  &  \\
\hline
6310 & ytri punktur &  &  &  \\
\hline
6311 & ytri sjálfmótun &  &  &  \\
\hline
6312 & ytri, utanverður &  &  &  \\
\hline
6313 & z-ás &  &  &  \\
\hline
6314 & Zermelo-Fraenkel-mengjafræði &  &  &  \\
\hline
6315 & Zermelo-Fraenkel-mengjafræði með valfrumsendu &  &  &  \\
\hline
6316 & þætta, kljúfa, sundurliða &  &  &  \\
\hline
6317 & þætta, leysa upp í þætti &  &  &  \\
\hline
6318 & þættaður &  &  &  \\
\hline
6319 & þættanleg margliða &  &  &  \\
\hline
6320 & þættanlegur &  &  &  \\
\hline
6321 & þættanlegur, kljúfanlegur, kleyfur &  &  &  \\
\hline
6322 & þættanleiki, kljúfanleiki, kleyfi &  &  &  \\
\hline
6323 & þak &  &  &  \\
\hline
6324 & þakning &  &  &  \\
\hline
6325 & þáttabaugur &  &  &  \\
\hline
6326 & þáttagreining &  &  &  \\
\hline
6327 & þáttalíkan &  &  &  \\
\hline
6328 & þáttatilraun &  &  &  \\
\hline
6330 & þáttun &  &  &  \\
\hline
6331 & þáttun, klofningur, sundurliðun, skipting &  &  &  \\
\hline
6332 & þáttur &  &  &  \\
\hline
6333 & þáttur &  &  &  \\
\hline
6334 & þekjufærsla &  &  &  \\
\hline
6335 & þekjugrúpa &  &  &  \\
\hline
6336 & þekjurúm &  &  &  \\
\hline
6337 & þekjuvídd &  &  &  \\
\hline
6338 & þetafall &  &  &  \\
\hline
6339 & þéttast &  &  &  \\
\hline
6340 & þéttibil &  &  &  \\
\hline
6341 & þéttipunktur &  &  &  \\
\hline
6342 & þéttipunktur &  &  &  \\
\hline
6343 & þéttipunktur &  &  &  \\
\hline
6344 & þéttisporbaugur &  &  &  \\
\hline
6345 & þéttleikafall, líkindaþéttleikafall, dreifingarþéttleikafall &  &  &  \\
\hline
6346 & þéttleiki &  &  &  \\
\hline
6347 & þéttraðað mengi &  &  &  \\
\hline
6348 & þéttskipaður &  &  &  \\
\hline
6349 & þéttteljanlegt firðrúm &  &  &  \\
\hline
6350 & þéttteljanlegt grannmynstur &  &  &  \\
\hline
6351 & þéttteljanlegt grannrúm &  &  &  \\
\hline
6352 & þéttteljanlegur &  &  &  \\
\hline
6353 & þéttteljanleiki &  &  &  \\
\hline
6354 & þéttur &  &  &  \\
\hline
6355 & þinalgebra &  &  &  \\
\hline
6356 & þinbundin &  &  &  \\
\hline
6357 & þinfeldi &  &  &  \\
\hline
6359 & þinföldun &  &  &  \\
\hline
6360 & þinföldun &  &  &  \\
\hline
6361 & þingreining &  &  &  \\
\hline
6362 & þinreikningur &  &  &  \\
\hline
6363 & þinsvið &  &  &  \\
\hline
6364 & þinur &  &  &  \\
\hline
6365 & þjál víðátta, óendanlega oft deildanleg víðátta &  &  &  \\
\hline
6366 & þjál vörpun &  &  &  \\
\hline
6367 & þjálga &  &  &  \\
\hline
6368 & þjálgun &  &  &  \\
\hline
6369 & þjáll á köflum &  &  &  \\
\hline
6370 & þjáll ferill &  &  &  \\
\hline
6371 & þjáll flötur &  &  &  \\
\hline
6372 & þjált fall &  &  &  \\
\hline
6373 & þjappað mengi &  &  &  \\
\hline
6374 & þjappað rúm &  &  &  \\
\hline
6375 & þjappað rúm, þjappað mengi &  &  &  \\
\hline
6376 & þjappaður &  &  &  \\
\hline
6377 & þjappaður virki &  &  &  \\
\hline
6378 & þjappleiki &  &  &  \\
\hline
6379 & þjöppuð eyðing &  &  &  \\
\hline
6380 & þjöppuð samleitni &  &  &  \\
\hline
6381 & þjöppuð víðátta, lokuð víðátta &  &  &  \\
\hline
6382 & þjöppun &  &  &  \\
\hline
6383 & þol, þolstig &  &  &  \\
\hline
6384 & þolandi &  &  &  \\
\hline
6385 & þolbil &  &  &  \\
\hline
6386 & þoldreifing &  &  &  \\
\hline
6387 & þolmark &  &  &  \\
\hline
6388 & þolmat &  &  &  \\
\hline
6389 & þolsvæði &  &  &  \\
\hline
6390 & þolþáttur &  &  &  \\
\hline
6391 & þraut &  &  &  \\
\hline
6392 & þrefaldur &  &  &  \\
\hline
6393 & þrefalt heildi, þrefalt tegur &  &  &  \\
\hline
6394 & þreföld rót &  &  &  \\
\hline
6395 & þrenging, takmörkun, takmarkandi skilyrði &  &  &  \\
\hline
6396 & þrengja &  &  &  \\
\hline
6397 & þrennd &  &  &  \\
\hline
6398 & þrep &  &  &  \\
\hline
6399 & þrep, kvíslitala, kvíslunartala &  &  &  \\
\hline
6400 & þrep, stig &  &  &  \\
\hline
6401 & þrepafall &  &  &  \\
\hline
6402 & þrepasönnun, þrepunarsönnun, rakin sönnun &  &  &  \\
\hline
6403 & þrepun niður á við &  &  &  \\
\hline
6404 & þrepun upp á við &  &  &  \\
\hline
6405 & þrepun, stærðfræðileg þrepun &  &  &  \\
\hline
6406 & þrepunarályktun, þrepunarniðurstaða &  &  &  \\
\hline
6407 & þrepunarforsenda &  &  &  \\
\hline
6408 & þrepunarfrumsenda, þrepunarfrumsetning, frumsenda um þrepun, frumsetning um þrepun &  &  &  \\
\hline
6409 & þrepunarnálgun &  &  &  \\
\hline
6410 & þrepunarskilgreining, rakin skilgreining &  &  &  \\
\hline
6411 & þrepunarskref &  &  &  \\
\hline
6412 & þrepunarupphaf &  &  &  \\
\hline
6413 & þríaðfelldur þríhyrningur &  &  &  \\
\hline
6414 & þríbaugur &  &  &  \\
\hline
6415 & þríblaða rós, smáraferill &  &  &  \\
\hline
6416 & þriðja rót, teningsrót, verpilrót &  &  &  \\
\hline
6417 & þriðja stigs ferill &  &  &  \\
\hline
6418 & þriðja stigs flötur &  &  &  \\
\hline
6419 & þriðja stigs, tenings-, verpils- &  &  &  \\
\hline
6420 & þrífeldi &  &  &  \\
\hline
6421 & þríflokkun &  &  &  \\
\hline
6422 & þríforkur Newtons &  &  &  \\
\hline
6423 & þrígild rökfræði &  &  &  \\
\hline
6424 & þríhliðungur &  &  &  \\
\hline
6425 & þríhnattaverkefni, þríagnaverkefni &  &  &  \\
\hline
6426 & þríhornalínufylki &  &  &  \\
\hline
6427 & þríhyrningar í sjónhorfi &  &  &  \\
\hline
6428 & þríhyrningaregla &  &  &  \\
\hline
6429 & þríhyrningaskipting &  &  &  \\
\hline
6430 & þríhyrningaskipting &  &  &  \\
\hline
6431 & þríhyrningaskipting &  &  &  \\
\hline
6432 & þríhyrningsfylki &  &  &  \\
\hline
6433 & þríhyrningsgerlegur &  &  &  \\
\hline
6434 & þríhyrningsgerningur &  &  &  \\
\hline
6435 & þríhyrningshnit &  &  &  \\
\hline
6436 & þríhyrningsójafna &  &  &  \\
\hline
6437 & þríhyrningstala, þríhyrnd tala &  &  &  \\
\hline
6438 & þríhyrningur & <ref idTerm=3161>Marghyrningur</ref> með þrjár <ref idTerm=2048>hliðar</ref> &  &  \\
\hline
6439 & þríhyrnt reitafylki &  &  &  \\
\hline
6440 & þríliða &  &  &  \\
\hline
6441 & þrílínulegur &  &  &  \\
\hline
6442 & þrírétthyrndur &  &  &  \\
\hline
6443 & þrískiptiferill &  &  &  \\
\hline
6444 & þrískiptilína &  &  &  \\
\hline
6445 & þrískipting &  &  &  \\
\hline
6446 & þrístæð vensl &  &  &  \\
\hline
6447 & þrístæður &  &  &  \\
\hline
6448 & þrístrendingur &  &  &  \\
\hline
6449 & þríunda- &  &  &  \\
\hline
6450 & þríundakerfi &  &  &  \\
\hline
6451 & þríundamengi &  &  &  \\
\hline
6452 & þríundatala &  &  &  \\
\hline
6453 & þríveldi, þriðja veldi, teningur, verpill &  &  &  \\
\hline
6454 & þrívíður &  &  &  \\
\hline
6455 & þrívítt rúm &  &  &  \\
\hline
6456 & þröngur &  &  &  \\
\hline
6457 & þröngur, í þröngum skilningi &  &  &  \\
\hline
6458 & þröngvaður &  &  &  \\
\hline
6459 & þröngvun &  &  &  \\
\hline
6460 & þröngvunarfall &  &  &  \\
\hline
6461 & þröskuldsfall &  &  &  \\
\hline
6462 & þungahnit, þungamiðjuhnit, voghnit &  &  &  \\
\hline
6463 & þungamiðja, þyngdarpunktur &  &  &  \\
\hline
6464 & þungamiðja, þyngdarpunktur &  &  &  \\
\hline
6465 & þungamiðjuhlutun &  &  &  \\
\hline
6466 & þunnskipað fylki &  &  &  \\
\hline
6467 & þvælistiki, amastiki &  &  &  \\
\hline
6468 & þver-, þverkyns, ímyndaður &  &  &  \\
\hline
6469 & þver, þverstæður, hornréttur, þverlægur &  &  &  \\
\hline
6470 & þvera &  &  &  \\
\hline
6471 & þverafleiða &  &  &  \\
\hline
6472 & þverás &  &  &  \\
\hline
6473 & þverás &  &  &  \\
\hline
6474 & þverbrennistrengur &  &  &  \\
\hline
6475 & þverbundin, þverlabundin &  &  &  \\
\hline
6476 & þvereining, þvertalan \(i\) &  &  &  \\
\hline
6477 & þverferill, þverbraut &  &  &  \\
\hline
6478 & þverfyllirúm, hornrétt fyllirúm &  &  &  \\
\hline
6479 & þvergrunnur, þverstæður grunnur, hornréttur grunnur &  &  &  \\
\hline
6480 & þverhluti &  &  &  \\
\hline
6481 & þverhnit, þverstæð hnit, hornrétt hnit, rétthyrnd hnit &  &  &  \\
\hline
6482 & þverhnitakerfi, þverstætt hnitakerfi, hornrétt hnitakerfi, rétthyrnt hnitakerfi &  &  &  \\
\hline
6483 & þverill, þverlína &  &  &  \\
\hline
6484 & þverjafna &  &  &  \\
\hline
6485 & þverkyns rót &  &  &  \\
\hline
6486 & þverlæg fjölskylda, þverstæð fjölskylda &  &  &  \\
\hline
6487 & þverlægni, þverstaða &  &  &  \\
\hline
6488 & þverlarammi, þverrammi &  &  &  \\
\hline
6489 & þvermál & <ref idTerm=2295>Hringur</ref> með geisla $r$ hefur \textit{þvermál} $2r$ &  &  \\
\hline
6490 & þversagnakennt mengi &  &  &  \\
\hline
6491 & þverslétta &  &  &  \\
\hline
6492 & þverslétta &  &  &  \\
\hline
6493 & þversnið &  &  &  \\
\hline
6494 & þversnið, snið &  &  &  \\
\hline
6495 & þversögn Bertrands, Bertrand-þversögn &  &  &  \\
\hline
6496 & þversögn, þverstæða &  &  &  \\
\hline
6497 & þverstaðalámóta &  &  &  \\
\hline
6498 & þverstaðalgrúpa &  &  &  \\
\hline
6499 & þverstaðalvörpun, þverstaðalfærsla &  &  &  \\
\hline
6500 & þverstaðlað fylki, þverfylki &  &  &  \\
\hline
6501 & þverstaðlað hnitakerfi &  &  &  \\
\hline
6502 & þverstaðlaður &  &  &  \\
\hline
6503 & þverstaðlaður grunnur, staðlaður þvergrunnur, hornréttur einingargrunnur &  &  &  \\
\hline
6504 & þverstæð röð &  &  &  \\
\hline
6505 & þverstæð runa &  &  &  \\
\hline
6506 & þverstæð þáttun, þverþáttun, hornrétt þáttun &  &  &  \\
\hline
6507 & þverstætt rúm, þverlægt rúm, hornrétt rúm &  &  &  \\
\hline
6508 & þverstöðluð hnit &  &  &  \\
\hline
6509 & þverstöðlun &  &  &  \\
\hline
6510 & þverstöðlunaraðgerð &  &  &  \\
\hline
6511 & þversumma &  &  &  \\
\hline
6512 & þvertala &  &  &  \\
\hline
6513 & þverun &  &  &  \\
\hline
6514 & þverunaraðferð, þverunaraðgerð &  &  &  \\
\hline
6515 & þvervigur, þvervektor &  &  &  \\
\hline
6516 & þverþáttur, þverliður &  &  &  \\
\hline
6517 & þvingaður, með hliðarskilyrðum &  &  &  \\
\hline
6518 & þvingandi &  &  &  \\
\hline
6519 & þvingun, aukaskilyrði, hliðarskilyrði &  &  &  \\
\hline
6520 & þýð fernd (af punktum) &  &  &  \\
\hline
6521 & þýð skipting &  &  &  \\
\hline
6522 & þýða &  &  &  \\
\hline
6523 & þýðandi &  &  &  \\
\hline
6524 & þýði &  &  &  \\
\hline
6525 & þýðing &  &  &  \\
\hline
6526 & þýðismeðalgildi, raunverulegt meðalgildi &  &  &  \\
\hline
6527 & þýðisstiki &  &  &  \\
\hline
6528 & þýður &  &  &  \\
\hline
6529 & þýður ferhliðungur &  &  &  \\
\hline
6530 & þykkt &  &  &  \\
\hline
6531 & þykkvafræði, þrívíð rúmfræði &  &  &  \\
\hline
6532 & þykkvi, rúmskiki, líkami, hlutur, rúmmynd, gripur &  &  &  \\
\hline
6533 & þyrping &  &  &  \\
\hline
6534 & þýtt fall &  &  &  \\
\hline
6535 & þýtt hlutfall &  &  &  \\
\hline
6536 & þýtt meðaltal &  &  &  \\
\hline
6537 & þýtt samok &  &  &  \\
\hline
6539 & aðfeldisfall &  &  &  \\
\hline
6540 & aðfeldi, hrópmerkt tala &  &  &  \\
\hline
6541 & aðgreinandi, aðgreinir, aðskilja ({\it kvk.}) &  &  &  \\
\hline
6542 & aðoka fylki &  &  &  \\
\hline
6543 & aðoki, aðoka varpi &  &  &  \\
\hline
6544 & aðoki, aðoka virki &  &  &  \\
\hline
6545 & alhæfður &  &  &  \\
\hline
6546 & alhæfð runa, nót, net &  &  &  \\
\hline
6547 & alhæfing &  &  &  \\
\hline
6548 & alhæfingarregla &  &  &  \\
\hline
6549 & alhæfi, altæki &  &  &  \\
\hline
6550 & alhæft meðaltal &  &  &  \\
\hline
6551 & alhæfur, altækur &  &  &  \\
\hline
6552 & andhverf vensl &  &  &  \\
\hline
6553 & andsamhverfur, oddhverfur &  &  &  \\
\hline
6554 & andsamhverf marglínuleg vörpun, oddhverf marglínuleg vörpun &  &  &  \\
\hline
6555 & andsamhverf vensl &  &  &  \\
\hline
6556 & angle}: þríhyrningur ákvarðast af  tveimur hliðum og horninu á &  &  &  \\
\hline
6557 & annars stigs jafna, ferningsjafna &  &  &  \\
\hline
6558 & annars stigs jafna, ferningsjafna, annars stigs flötur &  &  &  \\
\hline
6559 & biðröð &  &  &  \\
\hline
6560 & biðröð &  &  &  \\
\hline
6561 & biðtími &  &  &  \\
\hline
6562 & bil &  &  &  \\
\hline
6563 & bilanatíðni, óhappatíðni &  &  &  \\
\hline
6564 & bilanet &  &  &  \\
\hline
6565 & binda &  &  &  \\
\hline
6566 & biti, tvíundatölustafur &  &  &  \\
\hline
6567 & deildastuðull, deildahlutfall, diffurkvóti, afleiða (í punkti) &  &  &  \\
\hline
6568 & eining, einingarstak, margföldunarhlutleysa &  &  &  \\
\hline
6569 & einmóta, ímóta &  &  &  \\
\hline
6570 & endapunktur, heildunarmark, heildismark, tegrunarmark, tegurmark &  &  &  \\
\hline
6571 & eyðir, eyðingaríðal &  &  &  \\
\hline
6572 & eyðir, eyðingarvirki &  &  &  \\
\hline
6573 & formerkislaus, án formerkis &  &  &  \\
\hline
6574 & formleg kenning, formgerð kenning &  &  &  \\
\hline
6575 & frumstak &  &  &  \\
\hline
6576 & frumstak &  &  &  \\
\hline
6577 & fullt horn, \(360^\circ\)  horn &  &  &  \\
\hline
6578 & fylliatburður, andstæða, andstæður atburður &  &  &  \\
\hline
6579 & færslulíkindi &  &  &  \\
\hline
6580 & gagnvegur &  &  &  \\
\hline
6581 & Galois-grúpa &  &  &  \\
\hline
6582 & Galois-kenning &  &  &  \\
\hline
6583 & Galois-útvíkkun, Galois-víkkun &  &  &  \\
\hline
6584 & grunn-, undirstöðu-, meg\hbox{in{-}.}\ekkipunktur &  &  &  \\
\hline
6585 & hjáfylki, bylt hjáþáttafylki, bylt fylgiþáttafylki &  &  &  \\
\hline
6586 & hlutleysa, hlutlaust stak &  &  &  \\
\hline
6587 & hlutmengjafrumsenda, hlutmengjafrumsetning, frumsenda um hlutmengi, frumsetning um hlutmengi &  &  &  \\
\hline
6588 & hornið á milli &  &  &  \\
\hline
6589 & hringskífa, hringflötur &  &  &  \\
\hline
6590 & hægri hlutleysa, hlutleysa frá hægri &  &  &  \\
\hline
6591 & ívarp &  &  &  \\
\hline
6592 & jafngildisvensl &  &  &  \\
\hline
6593 & karteskur, Descartes-, Kartes\hbox{ar-.}\ekkipunktur &  &  &  \\
\hline
6594 & kartesk hnit, Kartesarhnit &  &  &  \\
\hline
6595 & kartesk slétta, Kartesarslétta &  &  &  \\
\hline
6596 & keilusnið &  &  &  \\
\hline
6597 & kennimargliða &  &  &  \\
\hline
6598 & leggjanet, aðoka net, afleitt net &  &  &  \\
\hline
6599 & leysikjarni &  &  &  \\
\hline
6600 & líkindaþéttleiki, dreifingarþéttleiki, þéttleiki &  &  &  \\
\hline
6601 & lyfting &  &  &  \\
\hline
6602 & margföldunar- &  &  &  \\
\hline
6603 & margföldun fjöldatalna &  &  &  \\
\hline
6604 & margföldun frá hægri, hægri margföldun &  &  &  \\
\hline
6605 & margföldun frá vinstri, vinstri margföldun &  &  &  \\
\hline
6606 & margföldun í kross &  &  &  \\
\hline
6607 & margföldun, reikniaðgerð rituð sem margföldun &  &  &  \\
\hline
6608 & margföld flokkun &  &  &  \\
\hline
6609 & margföld núllstöð &  &  &  \\
\hline
6610 & margföld rakning &  &  &  \\
\hline
6611 & margföld rót &  &  &  \\
\hline
6612 & margföld runa &  &  &  \\
\hline
6613 & margföld röð &  &  &  \\
\hline
6614 & margföld ör &  &  &  \\
\hline
6615 & milli þeirra &  &  &  \\
\hline
6616 & millumleikafrumsenda, millumleikafrumsetning &  &  &  \\
\hline
6617 & niðurstigningaraðferð &  &  &  \\
\hline
6618 & niðurstigning, niðurstigningaraðferð &  &  &  \\
\hline
6619 & nykurrúm &  &  &  \\
\hline
6620 & ofanvarpsmiðja, varpmiðja &  &  &  \\
\hline
6621 & óhliðruð línuleg deildajafna, óhliðruð línuleg diffurjafna &  &  &  \\
\hline
6622 & óræð tala &  &  &  \\
\hline
6623 & ósamrýmanlegir tveir og tveir, sundurlægir tveir og tveir &  &  &  \\
\hline
6624 & ósannur &  &  &  \\
\hline
6625 & óþættanleg margliða &  &  &  \\
\hline
6626 & óþættanleg Markov-keðja &  &  &  \\
\hline
6627 & raðtala &  &  &  \\
\hline
6628 & rakning &  &  &  \\
\hline
6629 & reitur (í töflu) &  &  &  \\
\hline
6630 & rótarmerki &  &  &  \\
\hline
6631 & röðuð \(n\)-d, röðuð \(n\)-und, \(n\)-d, \(n\)-und &  &  &  \\
\hline
6632 & röðunareinsmótun, raðeinsmótun &  &  &  \\
\hline
6633 & röðunarfrumsenda, raðfrumsenda, röðunarfrumsetning, raðfrumsetning &  &  &  \\
\hline
6634 & röðunargrannmynstur &  &  &  \\
\hline
6635 & röðunarsjálfmótun, raðsjálfmótun &  &  &  \\
\hline
6636 & röðun, hlutröðun, raðvensl &  &  &  \\
\hline
6637 & röð, þrep &  &  &  \\
\hline
6638 & röð, þrep &  &  &  \\
\hline
6639 & rökaðgerð &  &  &  \\
\hline
6640 & rökalgebra &  &  &  \\
\hline
6641 & rökbreyta &  &  &  \\
\hline
6642 & rökfasti &  &  &  \\
\hline
6643 & rökfesta, stærðfræðileg rökfesta &  &  &  \\
\hline
6644 & rökfrumsenda, rökfrumsetning &  &  &  \\
\hline
6645 & rökfræði &  &  &  \\
\hline
6646 & rökfræðilegur, rök-, röklegur, rökrænn &  &  &  \\
\hline
6647 & rökfræði annarrar stéttar &  &  &  \\
\hline
6648 & rökfræði fyrstu stéttar &  &  &  \\
\hline
6649 & rökhenda &  &  &  \\
\hline
6650 & rökhyggja &  &  &  \\
\hline
6651 & rökleiðsla, röksemdafærsla &  &  &  \\
\hline
6652 & rökleiðsluregla, rökregla, ályktunarregla, afleiðsluregla &  &  &  \\
\hline
6653 & rökreikningur &  &  &  \\
\hline
6654 & röktákn &  &  &  \\
\hline
6655 & röktenging, röktengill, tengill, tenging &  &  &  \\
\hline
6656 & rökvilla &  &  &  \\
\hline
6657 & rökvirki &  &  &  \\
\hline
6658 & samoka stak &  &  &  \\
\hline
6659 & sannur &  &  &  \\
\hline
6660 & spannmengi, íðalspannmengi, framleiðandi mengi, spannandi mengi &  &  &  \\
\hline
6661 & stýrihringur &  &  &  \\
\hline
6662 & stærðfræðilegur, stærð\hbox{fræði-.}\ekkipunktur &  &  &  \\
\hline
6663 & stærðfræðileg bestun, bestun &  &  &  \\
\hline
6664 & stærðfræðileg frumsenda, stærðfræðileg frumsetning &  &  &  \\
\hline
6665 & stærðfræðileg rökfræði &  &  &  \\
\hline
6666 & stærðfræðileg tölfræði &  &  &  \\
\hline
6667 & stærðfræðileg þversögn &  &  &  \\
\hline
6668 & talnaritunarkerfi &  &  &  \\
\hline
6669 & tugafylling, tugafyllitala &  &  &  \\
\hline
6670 & tugatala, tala í tugakerfi &  &  &  \\
\hline
6671 & tuga- &  &  &  \\
\hline
6672 & umhverfa, andhverfa &  &  &  \\
\hline
6673 & útgildi &  &  &  \\
\hline
6674 & vinstri aðoki &  &  &  \\
\hline
6675 & vinstri afleiða, afleiða frá vinstri &  &  &  \\
\hline
6676 & vinstri andhverfa, andhverfa frá vinstri &  &  &  \\
\hline
6677 & vinstri hlutleysa, hlutleysa frá vinstri &  &  &  \\
\hline
6678 & vörpun, færsla, myndun, ummyndun &  &  &  \\
\hline
6679 & yfirfall &  &  &  \\
\hline
6680 & yfirröð &  &  &  \\
\hline
6681 & þekjuvörpun &  &  &  \\
\hline
6682 & þriðja stigs jafna &  &  &  \\
\hline
6683 & þýð fernd &  &  &  \\
\hline
6684 & -gildur &  &  &  \\
\hline
\end{longtable}
